% =====================================================================
%  Harald Høffding, MINDRE ARBEJDER [I].  København: Det Nordiske Forlag (Bogforlaget
%  Ernst Bojesen), 1899 (Thieles Bogtrykkeri).  VIII, 243 pp.
%
%  TRANSCRIPTION -- GENERATED; DO NOT EDIT BY HAND.  Rebuilt by (ebook-method/mindre-arbejder/, see the docstrings):
%    python3 ebook-method/mindre-arbejder/draft.py --emit --force && python3 ebook-method/collate/check.py mindre-arbejder --apply
%  draft.py: Danish Wikisource's page-by-page transcription (Index "Høffding - Mindre Arbejder.djvu", another scan of
%  the same edition; all text pages at level "Valideret"; cached by wikisource.py), templates expanded, page markers at
%  Wikisource's page turns, then PATCHES (readings at the image a word decision cannot make).  check.py --apply: the
%  decisions in decisions.tsv, taken against the Royal Danish Library's scan of a copy of the same edition
%  (~/bibliotek/Høffding, Harald/1899-1913-mindre-arbejder-bd1.pdf; SCANS.tsv; PDF = printed + 15).
%  Contents: Forord pp. V-VIII (dated 22 May 1899), Indholdsfortegnelse, twelve essays pp. 1-243.
%
%  WORD-ACCURACY (2026-10-05).  The draft was collated word by word, with punctuation, against the KB scan's own text
%  layer (check.py): 1,292 disagreements.  348 settled by check.py's rules (glyph faults, noise, line order); 493 by
%  rules.py with tesseract (dan) on the KB images as a THIRD witness (392 where tesseract read as Wikisource against
%  the layer; footnote marks; page furniture run into the layer's lines).  The remaining 454 were read at the image:
%  294 as BLIND crops (the reader saw only the printed lines) and 160 on crops whose labels showed two readings, the
%  readers transcribing in both cases; reads.py decided 317 by comparison, 137 were decided by hand, 22 of them re-read
%  at the image.  OPEN 0; no decision left unapplied.  Wikisource is not diplomatic: it lower-cases nouns (standpunkt,
%  samvittighed), has at for af, Fersken for Forsken, 1899 for 1892, Breslan, ger for gør, an extra "i" (p. 92), etc.,
%  all corrected here; where it silently corrected a misprint ({{rettelse}}) the print's reading is kept.
%  Italics: Wikisource's italics checked against the KB images by stroke slant (check.py --italics); 115 mismatches
%  read at the image (sheets i01-i12): 7 corrected (Virkninger p. 52, Glæden p. 53, Fællesskab p. 154 italic; p. 169
%  only "de første" italic; p. 225 the note to ET TILBAGEBLIK roman), the rest Wikisource right.  OPEN 0.
%  Footnotes: every note read in full at the image (17 sheets, 46 pages); the readings compared with the text
%  (notescmp.py) and every difference re-read at the image: 3 corrections (p. 107 spaced points, p. 181 the dash, p. 37
%  misprint marked); the readers' own normalisations (Göttingen, Königsberg for the print's ø; Fremskrid) rejected.
%  Unattested-word sweep (sweep.py: words in no other transcription of the corpus and one edit from one that is):
%  26 suspects read at the image; misprints marked.  Front matter (title page, verso, Forord, Indhold) proofread in
%  full at the image: 2 commas of Wikisource's removed (p. VII), a stray colon removed (p. VIII), the rule under the
%  Indhold added.  Spot-read in full, pp. 71 and 140 and the Forord, with two planted control errors each (6 of 6
%  caught): no real fault on pp. 71, 140.
%  SECOND READING BY TRANSLATION (2026-10-05, TRANSLATION-PLAYBOOK §6c): the translators' 138 odd spots were
%  screened and 59 read at the image (crop sheets o01-o06, z01): 5 corrections (p. 24 „paavist“ one word; p. 177 the
%  stop after „Tælleren“; p. 188 „fælles“ for Wikisource's and the layer's „tælles“; pp. 108 and 153 n. commas for
%  Wikisource's stops) and 6 misprints marked (pp. 20, 75, 97, 124, 183, 234); the rest read as transcribed.
%  Not established: letterspacing (Wikisource marks it only in titles; not measured); misprints that every witness
%  reads alike as a real word; the run-over parts of three notes (pp. 107-108, 136-137, 152-153) are set whole at the
%  note's mark.
%
%  CONVENTIONS.  "% --- p. N ---" + \opage{N} at the first word of printed page N (roman for the Forord); a word broken
%  over the turn is joined round the marker.  Italics and letterspacing \emph.  Footnotes as \footnote at the mark (the
%  print marks them *) or ¹) per page).  Guillemets as printed, » « (Unicode).  Folios, running heads and signature
%  lines are not transcribed.  The print governs; misprints are kept and marked "% PRINTED AS IS" (18): »Social
%  Pessimisme» (p. VI), Forestilinger (54), Vekselvirking (50), uathængigt (92), Athængighed (93, 136), Amercian (130 n.),
%  orklaredes (131), til without its stop (174), Punker (37 n.), Religionen. da (20), at »Gabrielis« (75), en stort (97),
%  han mødrene (124), personsonligt (183), paavirket at Kant (234), and those Wikisource corrected (pp. 103, 108).
%  Headings as printed in the text; where the Indhold differs (it has »Hedenske Sandhedssøgere« without the quotation
%  marks the heading has, and »Kampen mellem gamle og nye Tanker. Et Tilbageblik«), both readings stand.
%  Earlier transcriptions of single essays (texts/hoeffding/realisme, forholdet-tro-viden) are superseded by this one:
%  realisme was found to read „Kants“ ... „Spencers“ (p. 6) and „Udgangspunkter og“ (p. 7) where the print has
%  „Kant's“ ... „Spencer's“ and „Udgangs- og Støttepunkter“.
% =====================================================================
\documentclass[12pt,a4paper,openany]{book}
\usepackage[utf8]{inputenc}
\usepackage[T1]{fontenc}
\usepackage{libertinus}
\usepackage{libertinust1math}
\usepackage[danish]{babel}
\usepackage[protrusion=true,expansion=true]{microtype}
\usepackage{setspace}
\usepackage[margin=1.2in]{geometry}
\usepackage{marginnote}
\usepackage{hyperref}
\hypersetup{pdftitle={Mindre Arbejder}, pdfauthor={Harald Høffding}, hidelinks}
\setstretch{1.3}
\newcommand{\opage}[1]{\marginnote{\footnotesize\textit{[#1]}}}
\newcommand{\essay}[2][]{\chapter*{#2}\addcontentsline{toc}{chapter}{#1}}   % [toc title]{title as printed, with its note}
\newcommand{\streg}{\begin{center}\rule{3em}{0.4pt}\end{center}}
\begin{document}
\frontmatter
\begin{titlepage}\centering\vspace*{2cm}
{\large\emph{HARALD HØFFDING}}\\[1ex]\rule{3em}{0.4pt}\\[3ex]{\Huge MINDRE ARBEJDER}\par\vfill
\emph{KØBENHAVN}\\[1ex]\emph{DET NORDISKE FORLAG}\\[0.5ex]{\small BOGFORLAGET ERNST BOJESEN}\\[2ex]1899
\end{titlepage}
\thispagestyle{empty}\vspace*{\fill}\begin{center}\emph{Oplag: 1200 Ekspl.}\\[3ex]{\small THIELES BOGTRYKKERI}\end{center}\clearpage

% --- p. V ---
\opage{V}
\chapter*{FORORD.}

Denne Bog indeholder en Række af Afhandlinger, der
i sin Tid ere fremkomne dels som Foredrag, dels som
undersøgende Anmeldelser. Den ydre Anledning til, at jeg har
samlet dem, var, at jeg --- under Forberedelsen til et
paatænkt Arbejde --- følte Trang til at se, hvor langt tilbage
jeg kunde følge den Opfattelse af videnskabelige, sociale og
religiøse Grundspørgsmaal, som nu i adskillige  Aar har dannet
Grundlaget for mit filosofiske Arbejde. Der afgrænsede sig
da for mig en Periode fra 1882, da jeg paa
Studentersamfundets første Mødeaften holdt Foredrag om »Realisme i
Videnskab og Tro«, til 1898, da jeg i Biologisk Selskab
holdt Foredrag om Vitalisme. Disse Udtalelser, og alle de
mellemliggende Smaaarbejder vidne om en og samme
Grundopfattelse, som det nu er mit Ønske --- hvis Tid og Evne
gøre det muligt --- at gennemarbejde og udføre. Det var da
et naturligt Ønske at have disse forskellige Arbejder samlede
i Trykken, dels fordi de kunne tjene som Indledning til det
paatænkte Arbejde, dels fordi de ogsaa ellers indeholde
Adskilligt, som endnu kunde være af Interesse, og jeg er det
højtærede Forlag taknemmelig, fordi det gik ind paa min
Tanke.

% --- p. VI ---
\opage{VI}Disse Arbejder udgaa nu til en anden Slægt end den,
til hvilken de første af dem henvendtes. Jeg følte mig, da
Rækken af disse Arbejder begyndte, endnu som hørende til
de Yngre. Nu hører jeg snart til Veteranerne; men jeg
har ingen Grund til at tro, at jeg ikke for det, jeg her har
at sige, skulde finde Modtagelighed hos En og Anden af den
Slægt, der nu er den unge.

En særlig Glæde har det været for mig at optrykke
Afhandlingen om F. C. Sibbern. Jeg vilde gerne, at hans
Værk og hans Personlighed skulde bevares i ærefuldt Minde,
og jeg vilde være glad, dersom min Virksomhed i nogen
Grad kunde fortjene at blive betragtet som en Fortsættelse
af hans. ---

I og for sig kunde flere ældre Smaaarbejder meget vel
være medtagne. Saaledes Foredraget »Darwinismen og
Filosofien« (Nær og Fjærn 1874) og min udførlige Kritik af
Martensens Etik (Nær og Fjærn 1878). Men der maatte
sættes en Grænse, og det sidstnævnte Arbejde var aldeles
overvejende af polemisk Natur. Af denne sidste Grund er
heller ikke en Række Artikler i »Tilskueren« fra November
1889 til Maj 1890 medtagne; den polemiserede mod den i
mine Øjne overdrevne Betydning, der var bleven tillagt
Nietzsches Ideer og Skrifter. ---

Om et enkelt af de optrykte Arbejder ønsker jeg at
gøre et Par Bemærkninger, nemlig om det, jeg har kaldet
»Social Pessimisme»% PRINTED AS IS (Wikisource corrects to "«")
, og som fremkaldtes ved Ferdinand
Tönnies’ Bog »Gemeinschaft und Gesellschaft«. Forfatteren
til dette Skrift, som jeg nu har den Glæde at kunne kalde
min personlige Ven, gjorde lige straks, da Afhandlingen var
bleven trykt (1890), en Indvending mod Titelen. »Hvis social
% --- p. VII ---
\opage{VII}Pessimisme,« skrev han til mig (26. Juli 1890), »skal betyde
en praktisk Retning, saa hylder jeg den ikke; tvertimod er
jeg tilbøjelig til at betragte alle Arter af Reform med
Forhaabninger.« Og denne Betænkelighed ved Titelen har han
udtrykt nyligt igen, da han hørte, at Afhandlingen skulde
optrykkes. Hans Hensigt med Bogen var en rent sociologisk
Undersøgelse og desuden en Paralleliseren af de hinanden
afløsende sociale Former med psykologiske Forhold. Skønt
denne Undersøgelse i Virkeligheden fører til en skarp Kritik
af de nu bestaaende Samfundsformer, saa ser han dog med
Haab og Tillid en ny Kultur i Møde og mener, at Arbejdet
paa den er i fuld Gang. I hvert Tilfælde er hans Blik paa
Forholdene blevet lysere siden Bogens Udgivelse. «Jeg
indrømmer Dem gerne,« skriver han (14. Maj 1899), »at jeg, da
Bogen blev skreven, saa Fremtiden i et mørkere Lys, hvad
en mulig Genfødelse angaar, og at der findes Spor deraf i
Bogen; men praktisk stod jeg allerede den Gang
Socialdemokratiets Standpunkt meget nær.« --- Jeg maa altsaa bede
Læseren tage Ordet »Pessimisme« som Karakteristik af
Tönnies’ Anskuelse med alt muligt Forbehold. Det var nu
egentligt ikke heller saa meget min Mening at ville udtrykke, at
Tönnies saa Fremtiden i et mørkt Lys, som at han saa
Fortidens Udvikling --- saaledes som den er gaaet for sig siden
Middelalderens Slutning --- i et mørkt Lys. Naturligvis
hænge begge Dele temmeligt nøje sammen. Det er
vanskeligt at forstaa, hvorledes der kan være lyst for Taaen, naar
der er aldeles mørkt for Hælen. Der er her en
Vanskelighed, som Tönnies’ Opfattelse deler med Marxismen, hvilken
Anskuelse han i det Hele staar nær. Jeg maa tilføje, at
det kun er mod Titelen paa Afhandlingen, min ærede Ven
% --- p. VIII ---
\opage{VIII}har gjort Indsigelse. »Jeg har,« skrev han i det sidst
anførte Brev, »paany læst Deres Essay, og jeg anerkender
paany den udmærkede Grundighed og Omhu, som De har
anvendt paa min Bog, hvem jeg ikke kunde ønske en bedre
filosofisk Fortolker, hvorfor ogsaa den nye Udgave for mig
er meget glædelig.« Jeg tør derfor stole paa, at mit Essay
vil være en god Forberedelse for den, der kunde ønske at
studere Tönnies’ Bog, som langsomt, men sikkert, har
vundet en fremragende Plads i den sociologiske Literatur.

\emph{Den 22. Maj 1899.}

\textbf{Harald Høffding.}


\chapter*{INDHOLDSFORTEGNELSE.}
\streg{}
\begin{flushleft}
Om Realisme i Videnskab og Tro (1884) \dotfill\ 1\\
Tvivlens Historie i den nyere Tid (1891) \dotfill\ 15\\
Om Vitalisme (1898) \dotfill\ 40\\
Sindsbevægelsernes Fysiologi (1886) \dotfill\ 51\\
Frederik Christian Sibbern (1885) \dotfill\ 61\\
Alexander Hamilton og den nordamerikanske Unionsforfatning (1889) \dotfill\ 113\\
Social Pessimisme (1890) \dotfill\ 142\\
Hedenske Sandhedssøgere (1892) \dotfill\ 158\\
Filosofi som Kunst (1893) \dotfill\ 180\\
Forholdet mellem Tro og Viden i dets historiske Udvikling (1885) \dotfill\ 195\\
Det religiøse Problem (1897) \dotfill\ 216\\
Kampen mellem gamle og nye Tanker. Et Tilbageblik (1895) \dotfill\ 225\\
\end{flushleft}
\begin{center}\rule{10em}{0.4pt}\end{center}

\mainmatter
% --- p. 1 ---
\opage{1}
\essay[OM REALISME I VIDENSKAB OG TRO]{OM REALISME I VIDENSKAB OG TRO\footnote{I det Væsentlige et Foredrag, holdt i „Studentersamfundet“ den
2. September 1882.}.}

\streg{}

Ordet Realisme bruges i vore Dage hyppigst som
Betegnelse for en Retning i den poetiske Literatur. Den
anvendes snart i dadlende, snart i rosende Betydning, men
sjeldent med klar Bevidsthed om hvad det egentligt udtrykker
og om Grænserne for Gyldigheden af det Begreb, det
angiver. En klar Redegørelse herfor vil maaske først komme,
naar den moderne æstetiske Retning fuldstændigere og
tydeligere har udfoldet sit Væsen og sine Konsekvenser. Det er
ikke den æstetiske Realisme, jeg her vil drøfte. Jeg vil
derimod søge at fremhæve det Ejendommelige ved den
realistiske Tendens, der i vor Tid ikke mindre tydeligt ytrer sig
paa Videnskabens og paa Livsanskuelsens Omraade. Det
følger af sig selv, at der maa være en indre Forbindelse
mellem Tendenserne paa de forskellige Omraader, thi lige saa
lidt i den aandelige som i den fysiske Verden gives der
noget Tilfældigt og Isoleret. Den menneskelige Aand kan
ikke udstykkes saaledes, at den i sin Fantasi følger Veje,
som aldeles ikke have noget Analogt og Parallelt med de
Veje, den slaar ind paa i sin Tænkning og i sin Tro. Det
er kun de mest doktrinære og forstokkede af den moderne
æstetiske Realismes Modstandere, som hævde Kunst og Poesi
en saa afsondret Plads i Aandens Verden, at Strømningerne
fra andre Omraader ikke kunne eller ikke bør kunne naa
derhen. Ved at drøfte Realismens Natur og Berettigelse paa
ét Omraade, give vi indirekte Bidrag til Forstaaelse af dens
Fremtræden paa andre Omraader.

% --- p. 2 ---
\opage{2}
\section*{I.}

I den mest omfattende Betydning kunne vi definere
Realismen som \emph{de naturlige Aarsagers Princip}. Naar Realismen
har Modstandere, og naar den først langsomt har arbejdet
sig frem til klar Bevidsthed, maa det altsaa komme af, at
man har søgt og endnu søger Forklaringen af Fænomener og
Begivenheder ad andre Veje end ved Hjælp af de naturlige
Aarsager.

Menneskets Trang til at finde Tingenes Aarsager har
ytret sig længe før Grundlæggelsen af en realistisk
Videnskab. Der kan ikke paavises nogen Periode i den
menneskelige Aands Udviklingshistorie, i hvilken den har staaet
aldeles passivt over for Sanseindtrykkene, uden det mindste
Forsøg paa at forklare sig dem paa en eller anden Maade.
De lavest staaende Vilde have deres Naturforklaring, der,
sét fra deres Standpunkt, kan have en fuldkommen rationel
Karakter. Og dersom vi kunde sætte os tilstrækkeligt ind
i Dyrets Bevidsthed, vilde vi sikkert finde, at det fortolker
de Indtryk, det modtager, efter visse Synspunkter. Det
følger allerede af, at Lovene for Forestillingernes Forbindelse
(Idéassociationen) vise sig at gælde lige saa vel for den
dyriske som for den menneskelige Bevidsthed. Erkendelsens
Historie begynder altsaa ikke med en Periode, i hvilken der
blot indsamles enkelte Indtryk, som saa i en følgende Periode
forbindes og bearbejdes. Paa ethvert Trin af Erkendelsens
Udvikling ytrer der sig en mere eller mindre udpræget
Stræben efter at forene det Mangfoldige og sammensmelte
det Beslægtede.

Modsætningen mellem de forskellige Perioder i
Erkendelsens Historie er derimod karakteriseret ved de forskellige
\emph{Forklarings- eller Fortolkningsprinciper}, der anvendes paa
Erfaringerne. Og her staa især to store Principer i skarp
Modsætning til hinanden.

Det ene fortolker Naturfænomenerne efter Regler, der
hentes \emph{uden for} Naturen; det andet fortolker Naturen ved
\emph{Naturen selv}.

Under den første Art af Naturforklaring hører al Myto%
% --- p. 3 ---
\opage{3}logi. Den mytologiske Naturopfattelse forklarer Fænomenerne
som Udslag af personlige Væseners Vilje. For den er
Naturen befolket af mere eller mindre menneskelignende
Aander, hvis Indgriben især spores i mærkelige og
betydningsfulde Begivenheder. Fra Fantasiens og Drømmenes Land
befolkes den virkelige Verden. Paa det primitive Aandstrin,
hvor Drømmelandets Beboere ere lige saa gode Realiteter
som den virkelige Verdens Fænomener, er det i sig selv
fuldstændigt konsekvent og berettiget at lægge dem til Grund
ved Naturforklaringen. I de højere Folkereligioners Teologi
er det sidste Forklaringsprincip for Alt et overnaturligt Væsens
Vilje. Ligeledes aldeles konsekvent, naar man én Gang har
slaaet Realiteten af et saadant overnaturligt Væsen fast.
Men den Naturforklaring, som herved kan naas, betragter
ikke Naturen i dennes eget Lys. Lige som Arkimedes
ønskede et Punkt uden for Jorden for at kunne bevæge
denne, saaledes har Teologien søgt sig et Punkt \emph{uden for}
Naturen for at kunne forklare denne.

Noget Lignende gælder ogsaa om mange Former for
filosofisk Idealisme. Naar man finder den egentlige
Forklaring af et Naturfænomen i det Formaal, i hvis Tjeneste det
antages at staa, eller den Idé, man mener, det udtrykker,
saa forklarer man heller ikke Naturen ved Naturen selv.
Ti hvorfra kende vi Formaalet eller Ideen? Selve Erfaringen
lærer os dem ikke; Kundskaben om dem maa altsaa hentes
fra en anden Kilde end den, gennem hvilken
Naturfænomenerne blive os tilgængelige. Vi maa altsaa have en særegen
Evne til at erkende »Ideerne«, og vor Erkendelsesevne
spaltes i to Grene, hvis Sammenhæng ikke kan paavises.
Vi have da paa den ene Side en Evne til at opfatte det i
Erfaringen Givne, paa den anden Side en Evne til at træde i
Forhold til en højere Verden og fra denne at hente
Principerne til Forstaaelse af Erfaringsverdenen. For saa vidt
de naturlige Aarsagers Princip anerkendes, betragtes det
blot som hørende under den lavere Erkendelse. Den sande
Erkendelse faa vi først, naar vi genkende Ideerne i
Erfaringsverdenen. Denne af \emph{Platon}, den spekulative Idealismes Stam%
% --- p. 4 ---
\opage{4}fader, grundlagte Opfattelse gaar under forskellige Former
igen lige til den nyeste Tid.

Realismen med de naturlige Aarsagers Princip udleder
derimod stedse Naturfænomenerne af andre Naturfænomener, de
mere sammensatte Fænomener af de simplere. Naar det
hedder, at Naturen skal forklares \emph{ved sig selv}, maa det
erindres, at Naturen er en uoverskuelig Verden af Fænomener,
der udfolde sig for os i Tiden og i Rummet. Ethvert Punkt
inden for denne Verden bestemme vi ved dets Forhold til
andre Punkter. Vi gaa ikke udenfor Naturen, da vi have
nok at gøre med at forbinde den uendelige Mængde af
Elementer inden for den. Eller rigtigere sagt: Udtrykkene
»inden for« og »uden for« tabe nu deres Betydning. Hvor
langt vi end gaa fremad fra Aarsag til Virkning, eller hvor
langt vi gaa tilbage fra Virkning til Aarsag, --- en Grænse
øjne vi ikke. Vor Forsken kan staa stille, fordi den mangler
Midler til at komme videre, saaledes som f. Ex. Kants og
Laplaces Hypotese ikke fører os længere tilbage end til den
roterende Urtaage. Men hvor som helst vi begynde eller
ende, staar der stedse et nyt Spørgsmaal. Det Princip, at
Naturen skal forklares ved sig selv, stiller derfor en uendelig
Opgave. Naturen som Helhed, som den Totalitet, der
indbefatter alle Naturfænomener i indbyrdes Sammenhæng, er
et Ideal for vor Erkendelse, til hvilket den Skridt for
Skridt nærmer sig, men som den aldrig fuldt vil kunne
virkeliggøre.

\section*{II.}

Naar vi gennemtænke Realismens Begreb paa denne
Maade, falder den absolut fjendtlige Modsætning til Idealismen
bort. Det er en af Idealismens Grundtanker, at vi ikke
kunne blive staaende ved de enkelte, isolerede Kendsgerninger,
men maa søge et Baand, der kan bringe Sammenhæng og
Enhed til Veje. Idealismen har fejlet ved den Antagelse, at
vi kunne naa denne Enhed og Sammenhæng ad anden Vej
end netop ved fortsat Bearbejdelse af Erfaringens
Kendsgerninger. I sin Utaalmodighed foregreb den et Ideal, som
% --- p. 5 ---
\opage{5}kun langsomt og kun tilnærmelsesvis kan naas; og da den
ikke fandt dette Ideal fuldt begrundet ved Erfaringen, udledte
den det af en højere Sfære.

Kun da blive Realisme og Idealisme fuldstændige
Modsætninger, naar Realismen helt opgiver den Grundsætning,
at sand Erkendelse ikke bestaar i ophobede Erfaringer, men
er Indsigt i den indre Sammenhæng mellem Erfaringerne.
Men denne Grundsætning \emph{kan} Realismen ikke opgive, da
den er al Videnskabs første og vigtigste Grundsætning. Naar
Realismen vender sig polemisk og tvivlende mod den
idealistiske Spekulation, kan det kun være, fordi den mener, at
Idealismen har misforstaaet og misbrugt sin Grundsætning.
Dens Tvivl er Udtryk for en dybere Tro: alte dubitat qvi
altius credit! Tvivl og Negation ere ikke Videnskabens sidste
Ord. Der er naturligvis mange populære Fordomme, der for
den videnskabelige Kritik opløses i Intet, mange Problemer
og Spørgsmaal, der ved nærmere Prøvelse vise sig at skylde
illusoriske Forudsætninger deres Opstaaen. Den
menneskelige Aand er oprindeligt sangvinsk. Den venter ikke
taalmodig paa Erfaringens Belæring, men haster forud med sine
Dogmer og Hypoteser. Med fremskridende Erfaring følger
derfor ogsaa fremskridende Kritik. Erfaringsvidenskabens
Udvikling medfører en Inddæmmen af
Forestillingsforbindelser, der hidtil bredte sig trygt i alle Retninger som en Flod,
der er gaaet over sine Bredder. Men denne negative,
hæmmende Virksomhed er kun den ene Side af Sagen.

Kampen mellem de forskellige Naturopfattelser afgøres
paa samme Maade, som Kampen mellem Drøm og
Virkelighed. Jeg tillægger det vaagne Liv og dets Indhold Realitet,
fordi det er mere omfattende og mere sammenhængende end
Drømmen, der altid er mere eller mindre fragmentarisk. Den
større og fastere Sammenhæng gør den mindre og løsere
Sammenhæng begribelig, ikke omvendt. Kun derved, at
Realismen giver os en mere sammenhængende
Verdensopfattelse end den spekulative Idealisme og Teologien formaaede,
kun derved vil den kunne bestaa i Kampen. Om
Enkelthederne vil der altid kunne strides. Fordi en Række
Erfaringer maa forklares efter Realismens Princip, kunde andre
% --- p. 6 ---
\opage{6}Erfaringer i sig selv godt være at forklare idealistisk eller
teologisk. Realismen bliver derfor først en ejendommelig
Verdensopfattelse, naar den vover sig ud over det blot Givne
eller det hidtil Givne og aabner Udsigten til en total og
gennemført Anvendelse af de naturlige Aarsagers Princip.
Det er dette, som er sket i de store naturvidenskabelige
Hypoteser, hvis Række aabnes af Kant’s og Laplace’s Teori
og slutter med Darwin’s og Spencer’s. Disse Hypoteser ere
voxede naturligt frem som Konsekvenser af Videnskabens
Udviklingsgang. I dem træder Realismen tydeligt frem som
almindeligt videnskabeligt Princip. Det er ikke altid sagt,
at Detailforskerne blive sig dette Princip bevidst. Arbejdets
Deling paa Videnskabens Omraade har ført det med sig, at
den enkelte Forsker med Lethed og tilsyneladende ogsaa med
god Samvittighed kan grave sig ned paa et enkelt Punkt
uden at bekymre sig om, hvor vidt de Principer og Metoder,
han følger, have almindeligere Betydning og indeholde
Konsekvenser af indgribende Betydning for hele Verdensopfattelsen.
Mange Naturforskeres Stilling til Darwinismen finder sin
Forklaring herved. Den, som derimod fæster sit Blik paa
Videnskabernes almindelige Udviklingsgang og sér, hvorledes
ethvert Fremskridt i Forstaaelse er vundet ved Anvendelse
af de naturlige Aarsagers Princip, betragter dette Princip
som det, der er kaldet til at beherske hele vor Naturopfattelse.
Det maa udfolde alle sine Konsekvenser; vi maa tilegne os
det; uden om det fører ingen Vej.

Er Realismen da bevist? Ingenlunde. Og den kan lige
saa lidt i egentlig Forstand bevises, som noget Princip i det
Hele kan det. Vi naa til et Princip, naar vi ved at følge en
Udviklingsgang opdage en herskende Regel eller en Tendens,
der træder des tydeligere for Dagen, jo mere Udviklingen
skrider frem. Ethvert Princip er for saa vidt en Hypotese,
som vi prøve at anvende i Virkeligheden. Vi prøve det ene
Princip efter det andet, indtil vi finde det, der bedst formaar
at forstaaeliggøre det Givne. Mytologiens Guder og
Platonismens Ideer vilde endnu herske i Videnskaberne, hvis den
fremskridende Erfaring forklaredes lige saa godt ved dem
som ved de naturlige, for selve Erfaringen tilgængelige Aar%
% --- p. 7 ---
\opage{7}sager. Den videnskabelige Realismes Gyldighed og
Berettigelse vil bero paa, om den virkelig angiver det eneste
Princip, som kan bringe Enhed og Harmoni i vor Tænkning,
og om den vinder fortsat Bekræftelse under Videnskabens
Fremskriden. Foreløbigt angiver den den eneste Maade, paa
hvilken vor almindelige Verdensanskuelse kan komme i Harmoni
med den Aand, i hvilken Enkeltundersøgelserne ledes, hvilket
kommer af, at den selv ikke er andet end den moderne
Videnskabs egen Aand.

\section*{III.}

Udtrykket: naturlige Aarsager behøver en nærmere
Forklaring. Dersom vi ved det Naturlige forstaa det, der kan
blive Genstand for sikker menneskelig Erfaring, saa høre de
aandelige Fænomener lige saa vel til Naturen som de
materielle. Naar den levende Følelse sætter Fantasien i Bevægelse,
eller naar vi ved en Viljesanstrængelse genkalde Noget i
vor Erindring, da have vi her lige saa vel et naturligt
Aarsagsforhold, som naar Legemerne udvides ved Tilførsel
af Varme eller presses sammen ved stærkt Tryk. Naturen
indbefatter altsaa to Grupper eller Rækker af Fænomener,
som vi kun ad Hypotesens eller Spekulationens Vej kunne
bringe i Forbindelse med hinanden: de materielle Fænomener,
der udfolde sig i Rummet, og Bevidsthedslivets Fænomener,
som kun ved vor Selviagttagelse blive os umiddelbart
tilgængelige, og som kun paa symbolsk Maade kunne
fremstilles i Rummets Form. For begge Grupper gælder de
naturlige Aarsagers Princip. Den videnskabelige Psykologi
søger de sjælelige Fænomeners indbyrdes Sammenhæng, lige
som Fysiken og Kemien søge de materielle Fænomeners
indbyrdes Sammenhæng. Begge Steder forklare vi Naturen ved
Naturen selv.

Vanskeligheden ligger, som allerede bemærket, i at finde
Forbindelsen mellem de to Erfaringsomraader. Materialismen
hugger Knuden over, idet den oversér deres Forskellighed
og uden videre gør det Aandelige til en Form for eller et
Produkt af det Materielle. Men det maa vel fastholdes, at
% --- p. 8 ---
\opage{8}Spørgsmaalet om Forholdet mellem det Aandelige og det
Materielle er et aabent Spørgsmaal inden for Realismen.
Realismen kan meget vel anerkende, at der er en saadan
Forskel mellem de to Erfaringsgrupper, at den ene ikke kan
føres tilbage til den anden, skønt de vise sig at staa i
inderlig Forbindelse med hinanden. Vi kende ingen aandelige
Fænomener, der ikke ere knyttede til materielle Fænomener,
men vi have intet fælles Maal, der kunde tydeliggøre os,
hvorledes en materiel Bevægelse kan frembringe en Tanke
eller omvendt.

I tidligere Tider traadte denne Vanskelighed ikke altid
saa klart frem. Man lod ubekymret Hjernen frembringe
Tanker og Følelser og omvendt Viljen sætte Lemmerne i
Bevægelse, saa længe man hverken havde nogen klar
Forestilling om de materielle Aarsagslove eller om de sjælelige
Fænomeners Grundejendommelighed. Men lige som de
naturlige Aarsagers Princip er trængt igennem i Videnskaben i
det Hele, saaledes er særligt i Naturvidenskaben det Princip
mere og mere blevet anvendt, at der til ethvert materielt
Fænomen ogsaa skal paavises en materiel Aarsag, og at en
materiel Aarsag kun kan have en materiel Virkning. Ogsaa
Fysiologien, Videnskaben om det organiske Liv, arbejder
efter dette Princip og opstiller den Fordring, at alle
organiske Fænomener --- altsaa ogsaa alt det, der foregaar i vor
Hjerne, naar vi tænke, føle og ville --- skal forklares efter
de materielle Aarsagers Princip. Fysiologien indrømmer, at
der endnu er saare langt tilbage, inden denne Fordring kan
ske Fyldest, men den afviser enhver Forklaring, som ikke
hviler paa dette Princip. Selve de sjælelige Fænomener
bryder Fysiologien sig egentligt kun om, for saa vidt de ere
knyttede til visse materielle Funktioner. Det er disse
materielle Funktioner, Fysiologien skal forklare os; at de
ledsages af Tanker og Følelser, er den strengt taget
ligegyldigt.

Men dersom vi ikke længere kunne lade Tanker og
Følelser gribe ind i de materielle Aarsagers Række, hvilken
Plads og Betydning skulle vi da tillægge dem? Deres Realitet
kan Ingen fornægte; den er os umiddelbar vis, vissere end
% --- p. 9 ---
\opage{9}alt Andet, ti det er ved Tænkningen, vi naa til al begrundet
Vished. Snarere vilde vi kunne drives over til den Antagelse,
at Tankens og Følelsens Verden er den eneste reale, og at
hvad vi kalde den materielle Verden kun er et Tankeprodukt,
der os selv ubevidst og i Kraft af en uundgaaelig Illusion
danner sig for os. Men selv da vilde Spørgsmaalet komme
igen: hvorledes skulle vi tænke os Forholdet mellem vore
Bevidsthedsfænomener og de materielle Fænomener, vi nu én
Gang forestille os nøje forbundne med dem?

Der bliver neppe nogen anden Udvej tilbage end at
opfatte Bevidsthedsfænomenerne som indre, sjælelige Udtryk for
det samme, der for de ydre Sanser fremtræder som materielt
Fænomen. Den virkelige Tilværelse er altsaa kun én, men
vi opfatte den (paa visse Punkter i det Mindste) under en
dobbelt Form. Hvilken af disse Former der er den
oprindelige, den konstituerende --- og om der maaske skulde
gives dybere liggende Principer, af hvilke de begge ere
afledede, --- det er Spørgsmaal, vi her ikke skulle indlade os
paa. Kun det skal fremhæves, at Realismen her ikke afskærer
videre gaaende Spekulationer, naar de kun have faste
Udgangs- og Støttepunkter i virkelige Erfaringer.

Realismen behøver derfor ikke at bestride det, der er det
egentlige Motiv og den dybeste Interesse i den idealistiske
Verdensanskuelse. Den er meget vel forenelig med en
Bestræbelse for at hævde det aandelige Livs Værdi og
Betydning i Tilværelsen. Vel kan den ikke gaa ind paa, at det
Aandelige er Noget, der fra en Idéverden sænkes ned i den
materielle Verden som i et Fængsel; men den fornegter eller
forringer ikke dets Realitet, naar den opfatter det som den
indre og for os betydningsfuldeste Aabenbaring af den samme
Kraft, som virker i den ydre Natur. Den kan ikke --- men
heller ikke Idealismen har hidtil kunnet det --- lære os,
\emph{hvorfor} og \emph{hvorledes} de aandelige Tilstande og Virksomheder
ere knyttede til de materielle; den indskrænker sig til den
beskedne Antagelse, at da der synes at herske
Aarsagssammenhæng og Lovbundethed i Verden, er det formodentligt
ingen Tilfældighed, at Bevidsthed og Hjernevirksomhed ere
knyttede til hinanden, saaledes som Erfaringen viser os det.
% --- p. 10 ---
\opage{10}Den véd ikke, om der er et Maal, henimod hvilket hele
Tilværelsen styrer, men den finder i sig selv ikke nogen
Umulighed i den Antagelse, at hele denne store
Natursammenhæng med sine ubrødelige Love og sin fast
sammenkædede Række af Aarsager og Virkninger er som en stor
Husholdning, en Verdensorden, hvor Udvikling og Opløsning
skifte, men hvor der dog gennem al Bølgegangen sker en
Forøgelse af det Værdifulde gennem højere Livsformers
Opstaaen.

\section*{IV.}

Allerede ved de sidste Betragtninger ere vi førte fra
Videnskabens Omraade over paa Troens. Ti det mest
almindelige Skelnemærke mellem \emph{Videnskab} og \emph{Tro} er, at
hin forsker efter Tilværelsens Love og indre Sammenhæng,
medens denne spørger om Tilværelsens etiske Betydning og
Værdi.

Naar vi definere Begrebet Tro paa denne Maade, existerer
der ingen nødvendig Strid mellem Videnskab og Tro.
Videnskaben begrunder og forklarer, Troen vurderer. Begges
fælles Grundlag er den for vor Erfaring tilgængelige
Virkelighed. Dennes Love søger Videnskaben, støttet til de
naturlige Aarsagers Princip, og Troen udtaler den Værdi, som
denne lovbestemte Virkelighed maa have for os i Kraft af
vore Idealer. Forklaring og Vurdering behøve ikke at
udelukke hinanden. Hvad der staar for os som højt og herligt,
som skønt og godt, kan dog derfor have haft sin naturlige
Udviklingshistorie. Kun de falske Værdier forsvinde, naar
Forklaringen er vunden. Naar vi tro paa, at
Verdensudviklingen er Mere end et Spil af Atomer, at der --- om det
end sker gennem megen Kamp --- arbejder sig noget Godt
og Værdifuldt frem gennem den, saa udelukker denne Tro
ikke en lige saa fast Overbevisning om, at dette Resultat
kun naas gennem en fast sammenhængende Aarsagsrække.

Grunden til, at vi alligevel skelne mellem Tro og Viden,
er dels, at den ideale Vurdering i og for sig ikke udsiger
% --- p. 11 ---
\opage{11}Noget om de virkelige Betingelser for det Værdifuldes
Opstaaen og Bestaaen, dels, at Overbevisningen om det
Værdifuldes Bestaaen og Gyldighed mere har sin Rod i Følelse og
Vilje end i Erkendelse. Kun da har jeg Anledning til at
opfatte Naturudviklingen som fremskridende Virkeliggørelse
af et idealt Indhold, naar jeg føler ideale Opgaver og
Formaal som det egentlige Tyngdepunkt i min egen Livsførelse.
Det etiske Ideal, jeg anerkender, kaster sit Lys ud over
Naturen. Som jeg er, tænker, føler og vil, med mine Idealer
og med mine Evner, er jeg fremgaaet af den store
Naturudvikling. Hvad der staar fast for mig som det ubetinget
Sande og Gode, er ganske vist betinget ved min endelige
Natur og min forsvindende Stilling i det store Hele; men
derved gøres det ikke til noget Betydningsløst og Tilfældigt.
Det er af Frugterne, man skal kende Træet, og af de
Frugter, Naturen har ladet modne i os og for os, kunne vi
lære, at der i den arbejde Kræfter, der ikke blot virke i
samme Retning som vore højeste Bestræbelser, men der
endog selv fra først af have ledet os ind paa disse
Bestræbelser. Paa et saadant Grundlag kan der dog ikke,
som den spekulative Filosofi antog, opbygges en teoretisk
Lærebygning. Ved den videre Udførelse vilde poetiske
Symboler snart træde i de videnskabelige Begrebers Sted.

Ved \emph{Realisme i Tro} mener jeg nu en Livsanskuelse, der
fastholder Troen paa Verdensudviklingens ideale Værdi, men
samtidigt lige saa inderligt er overbevist om, at denne ideale
Værdi realiseres gennem de naturlige Aarsagers Virken.
Naar denne Tro én Gang er vunden, er den uafhængig af
teoretiske Omskiftelser, ti den er fra først af i Harmoni med
selve Videnskabens Aand.

Den gængse Betydning af Ordet Tro udelukker
derimod en saadan Harmoni. Den fordrer et mere eller mindre
overnaturligt Apparat til sine Formaals Virkeliggørelse og
opstiller faste Former og uangribelige Symboler. Herpaa
beror Forskellen mellem human eller filosofisk Tro og
teologisk Tro. Den teologiske Tro maa i Følge sit Princip
opstille teoretiske Lærdomme, der komme i Konflikt med
Videnskaben. Den har til Genstand Noget, som skal være op%
% --- p. 12 ---
\opage{12}højet over teoretisk Forklaring, og som skal trodse de
naturlige Aarsagers Princip. Under disse Forhold kan Striden
mellem Tro og Viden aldrig faa Ende. Saa længe disse
Forudsætninger hævdes, vil man ikke slippe for Gentagelser
af det uhyggelige Billede, Videnskabens Historie viser os af
Troens Angst og Flugt for Videnskaben. Hver Gang
Videnskaben i Experiment, Bevis eller Hypotese rykker frem til
et nyt Omraade, gøres der Anskrig fra den teologiske Tros
Side. Kopernikus, Galilei og Darwin, Giordano Bruno,
Spinoza og J. G. Fichte ere Exempler paa, hvorledes uhildet
Iagttagelse og logisk Konsekvens føre til Modstrid med en
Tro, som endnu ikke har lært, at Vurdering og Forklaring
ikke ophæve hinanden. Kun en Tro, der kan vise den
Resignation at bøje sig for de naturlige Aarsagers Princip,
fører ud over denne Strid. Ogsaa her gælder det: alte
dubitat qvi altius credit. Ti det er dog vel den dybeste
Tro, der ikke fordrer overnaturlige Indgreb og Aabenbaringer
for at fastholde Overbevisningen om Betydningen og
Gyldigheden af det Ideale i Verden. ---

I en helt anden Betydning af Ordet vilde man ogsaa
inden for den teologiske Tro kunne tale om en realistisk
Tendens i vore Dage. Realismen er overalt betegnet ved
Mistillid til Spekulationen; den søger at have fast Grund
under Fødderne og holder sig derfor til det positivt Givne,
til det Faktiske --- eller det, som hver Art af Realisme anser
for faktisk givet. Naar nu den teologiske Tro paa sin Maade
gribes af Realismens Aand og ser sig om efter sine Fakta
for at holde sig til dem, saa støder den først og fremmest
paa sine overleverede Formler og Bekendelsesskrifter, i hvilke
den har det Indhold givet, som den absolute Autoritet har
sanktioneret. Det vil derfor være klart, at naar en realistisk
Tendens behersker Tiden, vil den inden for den teologiske
Tros Omraade naturligt komme til at fremtræde som potenseret
Autoritetstro. Den samme Tendens, der frembringer
Materialisme uden for den teologiske Tros Omraade, frembringer
Bogstavtro inden for dens Omraade. Teologien har derfor
ogsaa i vore Dage opgivet enhver spekulativ Begrundelse og
trukket sig tilbage til de kirkeligt sanktionerede Dogmer.
% --- p. 13 ---
\opage{13}Derved kommer der en stor Afstand mellem en teologisk
Tænker som Biskop Martensen og den yngre teologiske
Generation. I Slutningen af hans Bog om Jakob Bøhme
kan man læse hans advarende Ord til de unge teologiske
»Realister«. Han har sikkert Ret, ti denne sneverhjærtede
Klamren ved de overleverede Dogmer er ikke holdbar i
Længden.

Intet staar stille i Verden. Selv hvor der paa Overfladen
synes at være Ro og Uforanderlighed, kan der i Dybet foregaa
de betydningsfuldeste Processer og aflejres Lag, der til sin Tid
komme frem for Lyset. De almenmenneskelige Erfaringer øve
deres stille Indflydelse. Gennem den teologisk-religiøse
Udvikling siden Middelalderens Dage kan der spores en stadig
Tendens efter at naa Maalet ad saa kort og direkte en Vej
og med saa simple Midler som muligt, den samme Tendens,
der paa videnskabelig Grund udtrykker sig i den Regel, at
vi ikke uden tvingende Nødvendighed maa forøge vore Principer
og Hypoteser. Det er i Følge den teologiske Tro saare
mange og store Ting, man skal vide Besked om, for at kunne
leve som et Menneske og ikke som den, der æder og drikker,
fordi han véd, at han skal dø i Morgen. Den fremskridende
Udvikling fører her langsomt, men stadigt til større
Sparsommelighed og Beskedenhed. Protestantismen betegner
allerede en vis Reduktion i Forhold til Katolicismen, og i
Forhold til den gammellutherske Ortodoksi ere baade Pietismen,
Rationalismen og Grundtvigianismen afgjorte Fremskridt i
Koncentration og Simplifikation. Der er i denne Henseende
en Parallel mellem Teologien og Medicinen. Den ældre Medicin
udmærkede sig ved en uhyre Masse Medikamenter og et
stort Antal heroiske Kure. Men man hjælper sig efterhaanden
baade med færre Medikamenter for sin legemlige Sundhed og
med færre Dogmer for sin aandelige Sundhed. \emph{Hvor}
Simplifikationen vil ende, er ikke godt at sige. Den vil sikkert
ikke ende paa samme Punkt for alle Individer. Men i
Principet vil Maalet sikkert være naaet, saa snart Harmonien
mellem Vurdering og Forklaring ret er gaaet op for den
almindelige Bevidsthed. Da vil ogsaa den betydningsfulde
Kærne i de højere Former for teologisk Tro kunne bevares,
% --- p. 14 ---
\opage{14}medens Skallerne smuldre bort. Ti hvad der besjæler de
højere Folkereligioner og har givet dem deres vældige Magt
over Tusinder af Slægter, det er tilsidst den ideale Drift,
der vil se Tilværelsen som Andet og Mere end en Sum af
Kendsgerninger, og denne Drift er langt ædlere og vigtigere
end nogen som helst af de Former og Symboler, i hvilke
den har givet sig Udtryk. Den humane Tro kan og vil
bevare denne Drift, ja dens Stræben gaar netop ud paa at lutre
den og at frigøre den for alle hæmmende Skranker.

Kun Blasertheden kaster Kærnen bort med Skallerne.
En saadan Blaserthed er vist den farligste Fjende for den
aandelige Udvikling i vore Dage, farligere end al Intolerance
og Fanatisme, ti i dem er dog Liv, medens Blasertheden
baade selv er forstenet og forstener Alt omkring sig. Men
Blasertheden, for hvem Alt er lige godt, naar det blot er
faktisk, ender naturligt med ikke længere at bryde sig saa
meget om at skælne mellem virkelige og illusoriske Fakta;
den kvæler altsaa tilsidst selve den Sans for Virkelighed, af
hvilken den selv udleder sin Opstaaen, --- forstener sin egen
Kilde. Den opstaar let i stærke Brydningstider; men da
den beror paa en Forvexling af det Tilfældige med det
Væsentlige, vil den forsvinde lige som det Skum, en stærk
Bølgebrydning frembringer.

\streg{}

% --- p. 15 ---
\opage{15}
\essay[TVIVLENS HISTORIE I DEN NYERE TID]{TVIVLENS HISTORIE I DEN NYERE TID.}

\begin{center}\emph{C. N. Starcke:} Skepticismen som Led i de aandelige Bevægelser siden
Reformationen. Kbhvn. 1890.\end{center}

\streg{}

Det er ofte blevet paavist, hvor stor en Plads de
vaklende og tvivlende Naturer indtage i den moderne Literatur.
Hamletskikkelserne gaa stadigt igen under utallige Former
og have maaske næppe nogen Sinde floreret i den Grad som
i vore Dage, skønt deres Dimensioner blive mindre og deres
Indhold ubetydeligere, jo mere moderne Skildringen er. En
saadan staaende Type peger paa en væsentlig
Ejendommelighed ved den moderne Udvikling, den samme
Ejendommelighed, som gør, at det, naar den nyere Tid skal karakteriseres,
i Regelen er negative Udtryk som Frigørelse, Emancipation
o. lign., der først indfinde sig. Det mest Iøjnefaldende er
Løsrivelsen fra det Gamle, ligesom det, der foregaar ved
Fødselen, først og fremmest er Barnets Løsrivelse fra
Moderorganismen; først derefter kommer Spørgsmaalet om det
Nyfødtes Levedygtighed og Livsindhold. Om det, der hidtil
næredes i Moders Liv, nu kan staa paa sine egne Ben og
føre sit eget Liv, maa naturligvis afgøres ved praktisk Prøve.
At man kan leve, beviser man bedst ved at leve. Og dog
vil Mennesket, der, som Homér siger, er et Væsen, som »sér
frem og tilbage«, saa snart Bevidstheden er udviklet,
naturligt komme til at tænke over de Midler, der staa til dets
Raadighed i Livskampen. Saaledes gaar det ogsaa med den
nyere Tid. Alt imedens dens Liv udfolder sig i de
forskellige Retninger, gør Eftertanken sig gældende, og
modstridende Tanker rejse sig, naar det Spørgsmaal opkastes,
hvad der bærer vort Liv, og om det, der bærer det, har
Styrke nok. Der er dem, der mene, at kun i det Gamle,
% --- p. 16 ---
\opage{16}hvis ubestridte Herredømme ophørte med Renaissancen og
Revolutionen, er der Styrke og Fred. I Modsætning hertil
er der dem, der mene, at Fejlen netop er, at Resterne af det
Gamle tynge paa os, og at »vi sejle med et Lig i Lasten«.
Atter Andre mene, at Bruddet var nødvendigt, men at det
har slaget Saar, som kun langsomt kunne læges, og hæmmet
Sider af det aandelige Liv, som først efterhaanden kunne
udvikle sig fyldigt paany. Endeligt er der dem --- og til
deres Tal hører Forfatteren af det ovenfor nævnte Skrift ---
der mene, at Fejlen ligger i Mangel paa Selvtillid og
Konsekvens inden for den nye Livsanskuelse, idet den Tvivl, der
ganske naturligt maatte opstaa ved Bruddet med den
middelalderlige Livsanskuelse og Livsordning, altfor længe er blevet
anset for berettiget og ubetvingelig, og at Hovedopgaven
derfor maa være at fordrive Tvivlen fra ethvert Skjul, i hvilket
den kan forstikke sig.

Dr. Starcke er vel kendt i den filosofiske og sociologiske
Literatur. Foruden sin Doktordisputats om Ludwig
Feuerbach har han udgivet en Undersøgelse om »den primitive
Familie«, der nu foreligger baade paa Tysk, Engelsk og
Fransk, og som jeg i sin Tid har anmeldt i »Tilskueren«,
og en Afhandling om »Etikens teoretiske Grundlag«, der er
optagen i 6. Række af Videnskabernes Selskabs Skrifter. I
det foreliggende Skrift har Forfatteren gjort store Fremskridt
i Fremstillingsevne. Man faar af Bogen et velgørende
Indtryk af selvstændig Tænken og af en ædruelig, og dog
varmhjærtet Opfattelse af det menneskelige Liv og dets
Vilkaar. --- Forfatteren har vist mig den Ære at sætte mit
Navn paa Bogens første Blad sammen med en finsk
Videnskabsmands. Jeg ser deri ingen Hindring for, at jeg her kan
udtale, baade hvad jeg har at anerkende og hvad jeg har
at bestride i det foreliggende Skrift.

Det er en god Tanke, Dr. Starcke har haft, at betragte
den skeptiske Retning i Filosofien som en Slags Termometer
for den aandelige Udvikling. Han undersøger, hvad der
var Aarsag til, at der snart efter Reformationen opstod en
Række skeptiske Forfattere, og hvilke Grundtanker og
Synsmaader der kom frem hos dem. Ved da at vise, hvorledes
% --- p. 17 ---
\opage{17}hos de senere Forfattere, der kunne siges at falde ind under
Begrebet Skepticisme, Motiver og Grundtanker ændres, faar
han en Maalestok for, hvilken Kraft der er vunden, og
hvilken Udsigt en selvstændig Livsanskuelse har til at gøre
sig gældende. Dr. Starcke kommer ad denne Vej ind paa
en Række erkendelsesteoretiske, psykologiske og etiske
Problemer, hvis Behandling han mener er bleven hæmmet ved
den skeptiske Tendens, der har holdt sig længere end
nødvendigt, og han giver selv dygtige Bidrag til deres
Behandling.

Arbejdet falder, som det heraf vil ses, i to Hoveddele,
en overvejende historisk Del og en overvejende teoretisk.
Forfatteren har dog, synes det mig, ikke holdt dem
tilstrækkeligt ude fra hinanden. Navnligt teoretiserer han for
meget allerede i den første Del, skønt denne jo skulde
udgøre en væsentlig Del af hans Bevismateriale. Hans
Fremstilling vilde sikkert have vundet i Klarhed og Interesse,
dersom den historiske Behandling af Skepticismen var holdt
mere for sig selv, og dersom der saa tillige var givet en
mere indgaaende psykologisk og historisk Karakteristik af
de enkelte Forfattere. Derved vilde det sikkert tillige have
vist sig nødvendigt at medtage flere Forfattere; man savner
i det mindste nogle, som ere af Betydning for de Spørgsmaal,
der drøftes i den sidste Del af Bogen. Forfatteren giver saa
meget Godt, at vi naturligt ønske mere, især da det vilde
stemme med den vel udtænkte Plan.

Til Karakteristik af Bogen skal jeg dvæle ved begge
Hovedafsnit, altsaa først betragte Skepticismen som historisk
Fænomen, og dernæst drøfte Dr. Starckes Anskuelse, at
ethvert Spor af Skepticisme kan og skal forsvinde fra
Tænkningen og Livet.

\section*{I.}

Skepticismen bestaar, ifølge Dr. Starcke, paa det
teoretiske Omraade deri, at der mistvivles om Menneskets Evne
til at danne sig en Viden om Verdens og Menneskets
inderste virkelige Beskaffenhed, og paa det praktiske Omraade
% --- p. 18 ---
\opage{18}deri, at der rejser sig Tvivl om Menneskets Evne til paa
egen Haand at lære det Gode at kende og til at elske det.
En saadan Tvivl kan kun opstaa, naar hverken den teoretiske
Verdensopfattelse eller den praktiske Livsanskuelse og
Livsførelse længere helt igennem bestemmes ved Tro paa hellige
Autoriteter. I Renaissancens og Reformationens Aarhundrede
var der derfor en frugtbar Jordbund for den. Den faar en
meget forskellig Karakter og Retning, alt eftersom den især
motiveres af Kærlighed til det Gamle og bruges som Vaaben
mod det Nye, som man af Magelighed eller Frygt vil holde
sig fra Livet, eller den udspringer af Kærligheden til det
Nye, forbunden med Frygt for, at dette Nye ikke kan
gennemføres og være sig selv nok. I Forgrunden kommer dog
stedse til at staa det Spørgsmaal, om der gives en selvstændig
teoretisk Videnskab og en selvstændig Etik.

Efter Dr. Starckes Opfattelse har Skepticismen i den
Form, i hvilken den fremtraadte i det 16. Aarhundrede og i
den følgende Tid, en væsentligt anden Karakter end den
skeptiske Retning inden for den græske Filosofi. Oldtidens
Skepticisme havde, mener han, en mere teoretisk Karakter
og kan paa ingen Maade sammenstilles med den moderne,
hos hvilken der viser sig en stærk Fornemmelse af at være
kommen til kort i sin Kamp for Tilværelsen. Han finder
Forskellen udtrykt i den forskellige Anvendelse, der gøres
af Argumenterne for Erkendelsens Relativitet: hos de græske
Skeptikere bruges disse Argumenter, efter hans Opfattelse,
til at paavise en sikker Erkendelses Umulighed, medens de
i nyere Tid bruges til at vise, at den praktiske Usikkerhed
over for Naturen ikke kan afhjælpes. --- Jeg nærer Tvivl om,
hvor vidt Forfatteren har Ret heri. Man har for meget
vænnet sig til at betragte de græske Skeptikere (hvis egne
Skrifter vi ikke have) fra den rent negative Side. Adskilligt
tyder paa, at deres afvisende Holdning overfor Dogmer og
Spekulationer hang sammen med deres Interesse for
Erfaringsmetoden. Der er en historisk Sammenhæng mellem de
skeptiske Filosofer og den empiriske Lægeskole. Galenos bruger
herom det Udtryk, at de empiriske Læger vare i Medicinen,
hvad Skeptikerne vare i Filosofien, og han fremstiller Pyrrhon,
% --- p. 19 ---
\opage{19}den skeptiske Skoles Grundlægger, som Mønster paa den
sande Empiriker\footnote{Jeg kan her ikke gaa nærmere ind paa dette, men maa henvise til
Kapitlet: „Die Erfahrungslehre der Skeptiker“ i \emph{P. Natorp:}
„Forschungen zur Geschichte des Erkenntnissproblemes im Alterthum“.
(Berlin 1884).}. At Erfaringsvidenskaben i Oldtiden
tidligt hæmmedes i Væksten er sandt; men den gjorde sine
Tilløb, og en Sammenligning mellem de græske og de moderne
Skeptikere vilde derfor ikke mangle Berettigelse og
Interesse. ---

I første Række omtaler Dr. Starcke \emph{Michel de Montaigne}
(1533---1584), Hovedrepræsentanten for den nyere Tids
Skepticisme, for saa vidt denne betegner »Individets Oprør mod
Autoriteten«. Dette Udtryk synes mig noget for stærkt,
naar Talen er om Montaigne. Der er hos ham ikke Tale
om Oprør. Sine mangfoldige Kundskaber, sine sindrige
Sammenstillinger og ofte virkeligt dybt gaaende Tanker benytter
han til paa en lemfældig Maade at skaffe Plads for Udfoldelsen
af Individualiteterne efter deres forskellige Natur. Selve
denne Stræben efter individuel Frigørelse, af hvilken der
iøvrigt i Bogen gives en interessant Skildring, fremtræder
--- hvad jeg tror, man i Regelen overser, naar Talen er om
Montaigne --- ikke paa rent negativ Maade, men hænger
sammen med Montaignes Appel til Naturen, denne »vor store
og mægtige Moder«, som Menneskene have forladt for at
underkaste sig tilfældigt opstaaede og ved Vanen sanktionerede
Dogmer og Love. Montaigne forherliger Uvidenhed, naar
den er forbunden med et godt Hoved, og priser det naturlige
Instinkt, som hjælper bedre end alle kløgtige Paafund. I
hver Enkelt ytrer Naturen sig paa forskellig Maade: Enhver
har sin herskende Evne (forme maistresse). Det gælder da
om at bringe Naturen til dens Ret, hvilket ikke sker ved at
lægge Hænderne i Skødet, men fordrer Eftertanke og
Selvbeherskelse. --- Der er en hel stor positiv Baggrund hos
Montaigne, der endog paa et vist Punkt minder om hans
Samtidige Giordano Brunos begejstrede Naturpanteisme.
Navnet Skeptiker er ikke udtømmende her. Han staar i den
% --- p. 20 ---
\opage{20}franske Literatur som Forløber for Rousseau og Diderot ved
sin Appel til Naturen. Skeptiker er han dog, for saa vidt
hans passive Vedhængen ved det Gamle skærper hans Kritik
over for det Nye, og fordi han føler en af Magelighed motiveret
Skyhed over for den Uro, de nye Retninger vakte baade
paa Religionens Omraade (Luther) og paa Videnskabens
Omraade (Kopernikus, Paracelsus). Det laa ikke i hans
Natur pludseligt at vende op og ned paa alle sine
Forestillinger.

Ogsaa Montaigne’s Efterfølger \emph{Pierre Charron} maa, tror
jeg, opfattes fra en mere positiv Side end det sker i Dr.
Starckes Fremstilling. Charron er i Etikens Historie
mærkelig ved, at vi hos ham, den katolske Gejstlige, der benytter
den skeptiske Metode til at vise, at Mennesket ikke kan
maale Alt efter sin naturlige Fornuft, men maa holde sig til
Kirkens gamle Lære, --- at vi hos ham dog finde den første
udtrykkelige Hævdelse af Etikens Uafhængighed af Religionen.
Religion og Sædelighed, lærer han, bero hver paa sit Motiv
(eller Drivfjeder, \emph{ressort}) og maa derfor ikke sammenblandes.
Fremfor alt maa Godheden ikke gøres afhængig af Religionen. % PRINTED AS IS (a stop for a comma before „da“)
da den derved bliver tilfældig og ikke udspringer af \emph{le bon
ressort de la nature}. »Jeg vil, at man skal være et godt
Menneske, selv om der ikke gives Paradis eller Helvede.
Det er for mig en frygtelig Tale: »»Hvis jeg ikke var
Kristen, hvis jeg ikke frygtede for at blive dømt, vilde jeg
gøre eller undlade det eller det.«« Jeg vil, at du skal være
et godt Menneske, fordi Naturen og Fornuften vil det saa,
fordi det fordres af den almindelige Tingenes Orden, hvoraf
du kun er en Del, --- fordi du ikke kan handle mod dig
selv, dit Væsen, dit Formaal. Lad der saa komme ud deraf
hvad der vil.« (Charron: De la sagesse, Bordeaux. 1600. p.
342). --- Der er hos Charron, ligesom hos Montaigne,
afgjort en positiv Baggrund bag den skeptiske Form og
Metode.

Det Samme gælder for to Tænkeres Vedkommende, hvis
skeptiske Metode især finder Anvendelse paa det teoretiske
Omraade. --- \emph{Sanchez} opstillede ganske vist den Sætning, at
vi Intet vide, endog i Titelen paa sit filosofiske Hovedskrift
% --- p. 21 ---
\opage{21}(Qvod nihil scitur) Men man maa ikke lægge for megen
Vægt paa en paradox Udtryksmaade. Hovedtanken hos
Sanchez er Modsætningen mellem Idealet af en Viden i
Betydning af fuldkommen Indsigt i Tingenes Natur (perfecta
cognitio rei) og vor virkelige Viden. Han har ikke de store
Forhaabninger om hvad forbedrede Metoder ville kunne bringe,
som hans Samtidige Bacon af Verulam; men Mistrøstning var
ikke den eneste Stemning hos ham. Han opfordrer energisk
til at »gaa til selve Tingene« og ikke hilde sig i skolastiske
Slutningsrækker. Og selve det Værk, i hvilket han udvikler
sin Skepsis, er en Indledning til hans lægevidenskabelige
Arbejder. Han har desuden selv fremstillet en Naturfilosofi,
der minder om den tidligere af Italieneren Telesio opstillede,
og som ganske vist ikke har blivende Betydning, men dog
vidner om, at Erkendelsesmodet ikke var stækket hos ham.

Endnu mere gælder dette om \emph{Joseph Glanvil}. Han har
ligesom de foregaaende Forfattere et klart Blik for Alt hvad
der hæmmer vor Erkendelses Udvikling, og giver en bred
Fremstilling deraf. Men han erklærer udtrykkeligt, at
Skepticismen for ham er en Metode, ikke et System. Han vil
Dogmatismen, baade i dens skolastiske og i dens
materialistiske Form, til Livs og bruger den skeptiske Metode til at
undergrave den og til at udvide og frigøre Sindet; ti al
Dogmatisme er for ham Aandssneverhed. Mod Slutningen
af sin Fremstilling siger han (Scepsis scientifica. Kap. XXVI):
»Jeg har kæmpet mod Dogmatismen under Skepticismens
Mærke, men jeg ønsker ikke, at man skal betragte Hensigten
med hele min Undersøgelse fra dette Synspunkt.« Det er,
siger han et andet Sted (Kap. X), kun den dogmatiske
Uvidenhed, der holder den forsigtige Tilbageholdenhed fra
enhver uprøvet Paastand for Skepticisme. Hvad Glanvil især
indskærper, er Nødvendigheden --- og tillige Vanskeligheden
--- af en Erfaringsbekræftelse af alle vore Hypoteser. I sin
Dedikation til »Royal Society« siger han, at uden
Erfaringskendsgerninger »ere vore Hypoteser vilde Drømme og
Romaner, og vor Videnskab blot Gisninger og Meninger. Ti
naar vi danne Systemer uden at raadspørge Fænomenerne,
bygge vi kun i Luften og beskrive en selvskabt og indbildt
% --- p. 22 ---
\opage{22}Verden, som kun er lidet beslægtet med den virkelige ....
Og det er muligt, at alle de hidtil udtænkte Hypoteser byggedes
paa en for snever Betragtning af Tingene«. En saadan
Advarsel trænger enhver Tid til, ikke blot den Tid, der
laa forud for Newton. --- Og den positive Side fremtræder
endog meget bestemt hos Glanvil, idet han slutter sig til den
kartesianske Filosofi og Fysik, til hvilken hele hans Skrift
kan betragtes som en Indledning. Han har egentligt kun
givet en bred Udvikling af den skeptiske Metode, som
Descartes selv havde benyttet i sin »Discours sur la méthode«. ---

Ved disse Exempler har jeg villet vise, at man let kommer
til at lægge for Meget i Ordet Skeptiker, og at adskillige
af de ældre Forfattere, der give sig selv dette Navn
eller figurere under det i Filosofiens Historie, dog staa i
et positivt Forhold til Erkendelsen og dens Opgaver. Jeg
tror, at Dr. Starckes Opfattelse her er noget ensidig, og at
dette har faaet Indflydelse paa hans Tankegang hele Bogen
igennem. Ti ligesom han opfatter de ældre Skeptikere,
f. Eks. Mænd som de nys nævnte, rent fra den negative
Side, saaledes er i det Hele hans Opfattelse af »Skepticismen«
den, at den er Noget, der opstod under Brydningsforhold af
forbigaaende Natur, og at den nu kan og skal forsvinde
fuldstændigt. Hans Bog er en Jagt efter Skepticismen for
at drive den ud af enhver Krog, i hvilken den kan skjule
sig. I Modsætning hertil ser jeg hine »Skeptikeres« Betydning
i, at de hævde den videnskabelige Metode paa alle
Omraader, hvilket ikke kunde ske, uden at der blev lagt
stærkt Vægt paa Afstanden fra Erkendelsens Ideal. Og jeg
kan ikke se, at al Skepticisme er af det Onde, selv den
Dag i Dag, hvor vi dog have set saa Meget, der vilde have
staaet som en vild Drøm for det 17. Aarhundredes Forfattere,
blive fastslaaet som sikker Erkendelse. Dog herom
Mere siden. ---

Fortræffelig og, saa vidt jeg ved, ny er den Belysning, i
hvilken \emph{Pascal}’s Skepticisme stilles. Dr. Starcke skelner
mellem Tænkere som de fleste ovenfor omtalte, hvis Skepticisme
nærmest beror paa manglende Tillid til, at Mennesket
ved egen Hjælp kan bygge sin Verdensanskuelse, og en
% --- p. 23 ---
\opage{23}Pascal, hos hvem det Afgørende er »den Rædsel, der griber
Mennesket, idet han ser Muligheden af, at han paa egen
Haand kan finde sig selv til Rette, men tillige ser denne
Mulighed knyttet til en Betingelse, Verdensmekaniken, som
synes at angribe Livets Værdi i dets inderste Rod«. Det
er nu »allerede selve Forskningens Metode og det Skema
for Verdenshelet, denne forudsætter, som vækker Bekymring.
Tanken om en Verden, der ligesom et Urværk gaar efter
evige mekaniske Love, synes disse Mennesker at maatte
fylde os med Gru.« Ad denne Vej forklarer Dr. Starcke
Lidenskabeligheden i Pascals Polemik mod den naturlige
Fornuft, idet han i ham ser »en Levendegørelse af den store
Tragedie, som spilledes ved Sammenstødet mellem det Gamle
og det Nye.« --- Jeg føjer her kun til, at det er en Tragedie,
der senere ofte har gentaget sig under Tænkningens rytmiske
Udvikling. I vor egen Literatur have vi en original
»Levendegørelse« af denne Art i S. Kierkegaard. Og han bliver
maaske ikke den sidste.

I et følgende Afsnit omhandles den engelske \emph{Deisme} paa
en interessant Maade. --- Deismen var et Kompromis: man
kastede en Del Dogmer over Bord, som man mente havde
en rent ydre, historisk Oprindelse, og beholdt tilbage en lille
udvalgt Kreds af religiøse Forestillinger (især om en personlig
Gud og Udødeligheden), som man mente udsprang af
»Fornuften« og var forenelig med en rent videnskabelig Verdensopfattelse. Det var det Standpunkt, hvorpaa det 18.
Aarhundredes »Oplysning« stod. Dr. Starcke søger at vise, at
Deismen opløses, da den Konsekvens drages, at Mennesket i
sin Livsførelse maa kunne undvære disse resterende Dogmer
lige saa vel som de forkastede, og at Etiken lige saa vel
er uafhængig af den »naturlige« Religion som af den positive
Religion. Det kunde i denne Sammenhæng være blevet paavist, at Deismen ogsaa blev angreben fra en helt anden Side.
I sit mærkelige Skrift »Analogy of Religion, natural and
revealed, to the constitution and course of nature« (London
1736) paaviste \emph{Joseph Butler}, at de vigtigste Indvendinger,
der kunne rettes mod den aabenbarede Religion, ogsaa kunne
rettes mod Troen paa en Forsynsstyrelse i Naturen, idet man
% --- p. 24 ---
\opage{24}ad Erfaringens Vej ikke kan bevise, at Naturens Ophav
besidder Visdom, Retfærdighed og Godhed. Naar man forarges
over den aabenbarede Religion, saa er der, lærer Butler,
nok saa Meget at forarges over i Naturen. I Naturen lider
f. Eks. den Uskyldige ofte for den Skyldige, lige saa vel som i
den kirkelige Forsoningslære. Butler’s Skrift, der --- i
Modsætning til den herskende Optimisme --- fremhæver den
mørke, disharmoniske Side ved Naturen, især saadanne Træk,
som man nu henfører under Synspunktet af den blinde Kamp
for Tilværelsen, spillede en indgribende Rolle i den religiøse
Diskussion i det følgende Aarhundrede og har bl. A. stærkt
paavirket James Mill’s og Stuart Mill’s religiøse Anskuelser.
--- Ogsaa et andet Savn føler man her ved den historiske
Fremstilling. Det vilde have været naturligt her at omtale
\emph{Mandeville}’s Bifabel med dens Grundsætning, at Samfundet
kun kan trives og Civilisationen blomstre paa Moralens
Bekostning, et Indlæg af ikke uvæsentlig Betydning i den
moderne Skepticismes Historie, og som særligt vilde have
været af Interesse for Dr. Starcke, der i det Følgende saa
stærkt fremhæver det Problem, hvorledes Individets
Lykkestræben kan forenes med Samfundets Krav. --- Endeligt vilde
det have været interessant, om Forf., der dvæler ved flere
mindre betydelige Mænd fra den literære Kreds, der omgav
\emph{Frederik den Store}, i Stedet derfor havde belyst den
filosoferende Konges eget Forhold til de i Bogen omhandlede
Spørgsmaal. Kongens Overgang fra Wolffs Dogmatisme
gennem Bayles Skepticisme til Lockes Empirisme vilde være
af Interesse til Belysning af »Skepticismen som Led i de
aandelige Bevægelser siden Reformationen«, ligesom hans
definitive Anskuelser betegne en ejendommelig Modifikation
af den i Oplysningstiden herskende »naturlige Religion«. ---

Men alle disse Bemærkninger rokke ikke ved Værdien
af det historiske Parti i Bogen. Forfatteren har --- ved at
føre sin Fremstilling igennem til Hume, Rousseau, Kant og
Stuart Mill --- klart bevist sin Hovedsætning, at
Skepticismen ender med at have en helt anden Karakter og indtage
en anden Stilling end ved sin Begyndelse. Den begyndte
som Udtryk for Mistillid, Beskedenhed eller Forsigtighed.
% --- p. 25 ---
\opage{25}Den ender med Overbevisningen om, at en hel Række
spekulative og teologiske Spørgsmaal trygt kunne skydes til Side,
idet baade den positive Videnskab og det praktiske
Menneskeliv gaa deres Gang og naa deres Maal, hvad enten hine
Spørgsmaal løses eller ikke. Den gennem mange aandelige
Bevægelser og Kampe naaede Overbevisning, som Forf. især
finder repræsenteret af Rousseau og Kant, udtrykker han
træffende i følgende Ord: »I Livet selv maa Midlerne søges
til at løse Livets Opgaver; findes de ikke dér, findes de
ingen Steder … Kun paa menneskelig Grund kan Vejen
lægges … Efter Kant og Rousseau vil Skepticismen ikke
mere træde frem i saa klare Former som hidtil, fordi
der overhovedet ikke mere stilles noget Valg om, hvor vidt
denne Vej paa menneskelig Grund skal følges eller ikke …
Vi træffe nu ikke længere paa rene, gennemførte skeptiske
Standpunkter; der er i de fleste Standpunkter et skeptisk
Stænk, og det bliver derfor vor Opgave at samle alle disse
isolerede Træk til tydelige Billeder.« Disse sidste Ord
betegne Dr. Starckes Opfattelse af Skepticismen i vore Dage.
Vi ere, mener han, »brudte udover Skepticismens
Landemærker«, og det gælder nu blot om »at trække den sidste
Brod ud.«

Dette Resultat fremtræder ikke blot i
Erfaringsvidenskaberne og den kritiske Filosofis rent teoretiske Resultater.
Det viser sig særligt i den forandrede Stilling, Religionen
indtager i praktisk Henseende. Ved Udviklingens
Begyndelse hentedes baade Verdensanskuelsens og Livsførelsens
Principer fra Religionen. Der har nu dannet sig et ad rent
videnskabelig Vej dannet Verdensbillede. Religionens
Betydning indskrænkes derfor mere og mere til det praktiske,
etiske Omraade. Men ogsaa her foregaar der en Forandring.
I Middelalderen blev Moralen Mennesket kær gennem
Religionen. I den nyere Tid ændredes Forholdet først derhen,
at Religionen hævdedes som Kilde til Moralen. Senere, da
Etikens Selvstændighed og Uafhængighed gennemførtes, blev
Religionen ikke længer dens Kilde, men blot dens Garanti
og Sanktion. Religionen forsvares nu som nødvendig, ikke
fordi man gennem den lærer Moralen at kende, men fordi
% --- p. 26 ---
\opage{26}det menes, at uden Troen paa guddommelig Gengældelse vinder
Moralen ingen Magt i den menneskelige Verden. --- Denne
væsentlige Forandring er et af de vigtigste Resultater,
Skepticismens Historie kan opvise.

\section*{II.}

Den sidste Del af Dr. Starckes Bog søger at vise
Nødvendigheden og Muligheden af, at »den skeptiske Brod
trækkes ud.« Han behandler i den Anledning en Række
erkendelsesteoretiske og etiske Problemer paa en skarpsindig
Maade, men saaledes i det Enkelte, at jeg paa dette Sted
ikke kan følge ham, men maa holde mig til nogle
Hovedpunkter. Og jeg vælger da her, ligesom i det Foregaaende,
især saadanne, hvor jeg afviger fra ham.

Først og fremmest rejser sig det Spørgsmaal, om han
har Ret i, at den »skeptiske Brod«, som er tilbage, kan og
skal trækkes ud. Afgørelsen heraf vil naturligvis for en
Del bero paa, hvor meget man lægger ind i Ordet
Skepticisme. Jeg har i det Foregaaende søgt at vise, at Dr. Starcke
tager noget for tungt paa den ældre Tids »skeptiske«
Forfattere, idet han gør den skeptiske Metode til Et med absolut
Skepticisme og undervurderer de positive Grundtanker, for
hvilke man ved Tvivlen vilde skaffe Plads. Dersom man i
Begrebet Skepticisme lægger Uro og Mistrøstning ind som
nødvendige Træk, saa kan der neppe være Strid om, at
Skepticismen er et underordnet og forbigaaende Fænomen.
Den skyldes da væsentligt en Slags Kontrastvirkning.
Dogmatismen lovede Guld og grønne Skove, Svar paa alle
Spørgsmaal, Afslutning af Tankens Arbejde. Den nyere
Erfaringsvidenskab og Filosofi kan ikke give saadanne Forjættelser.
Den lover os Fremskridt i det Enkelte, nye Kendsgerninger
og større Klarhed, hvad Principerne angaar, og den peger
henimod en stor, indre Sammenhæng i Alt; men den rejser
tillige stedse nye Problemer, opbyder stedse nyt Arbejde,
fordi der stedse frembyder sig nye Skranker at flytte. Som
der ingen absolut Hvile findes i den ydre Natur, men Alt er
i Bevægelse og Virksomhed, saaledes gælder dette særligt
% --- p. 27 ---
\opage{27}for Tanken. Den faar ikke Lov at hvile i kontemplativ
Beskuen; de foreløbigt afsluttende Forestillinger, den stanser
ved, vise sig snart som Symboler, der antyde og signalisere,
men ikke umiddelbart gengive Tilværelsens Virkelighed. Den,
som pludseligt bringes over for denne Opfattelse af vor
Erkendelse, efter forud at have hvilet trygt i Dogmatismens
Favn, maa ganske vist føle sig urolig og mistrøstig. Men
kun saa længe Bruddet er friskt, og Modsætningen til den
dogmatiske Tro føles i hele sin Skarphed, vil denne
Misstemning vare. Den vil forsvinde, jo mere han lever sig ind
i de nye Vilkaar. Resignation vil stedse være nødvendig.
En saadan kan ingen som helst Verdensanskuelse undvære,
da det jo overalt gælder om Indordning under forudsatte
faste Vilkaar. Der ligger aldeles ikke noget Skeptisk i
Anerkendelsen af Resignationens Nødvendighed.

Livet maa stedse begynde med et Overskud af Kraft og
Mod, for hvis Anvendelse Erfaringen successivt sætter sine
Grænser. Ethvert Liv er en Strøm, der dels gennembryder
hvad der staar imod, dels inddæmmes af hvad der ikke kan
brydes. Strømmens endelige Leje bestemmes ved Forholdet
mellem disse to Sider. At lære Livet at kende er at erfare,
hvor langt man med de Kræfter og Drifter, man raader over,
kan udvide de givne Skranker, og i hvilken Grad man
derimod nødes til at indordne sig under dem. Resignationens
Nødvendighed ligger i, at Strømmens endelige Leje ikke \emph{blot}
bestemmes ved Udspringets Kraft, men ogsaa ved de
modstaaende, inddæmmende Kræfter. Dette gælder baade for
vor teoretiske og vor praktiske Stræben.

Den videnskabelige Metode bestaar ikke i passiv
Ophoben af Iagttagelser, af hvilke der senere uddrages
Resultater. Den første Betingelse for frugtbart videnskabeligt
Arbejde er en ledende Ide, en foreløbig Hypotese, der bliver
Motiv til Opsøgen af nye Erfaringer. Disse nye Erfaringer
medføre ofte Nødvendigheden af, at Ideen ændres eller endog
ombyttes med andre Ideer, som man fra først af maaske ikke
havde tænkt sig Muligheden af at komme til at anvende.
Den teoretiske Resignation betinges ved Nødvendigheden af
at skaffe Erfaringsbekræftelse paa vore Ideer og af dog at
% --- p. 28 ---
\opage{28}begynde med Ideer. Her er et Grundforhold, som ikke kan
falde bort, men som maaske endog for Fremtiden vil
fremtræde i skarpere Form end hidtil.

Ogsaa i praktisk Henseende maatte der begyndes med
et Overskud, med Antecipationer, der paa mange Maader
berigtiges. Uden kraftige Forhaabninger og Længsler og uden
store Forbilleder kan intet Stort udrettes. Men
Virkeliggørelsen naas kun gennem Brydning og Vekselvirkning med
de virkelige Forhold. Meget maa opgives, som fra først af
ansaas for værdifuldt, og Meget maa optages og antages, som
fra først af var Genstand for Uvilje.

I Teorien Modsætningen mellem Hypotese og
Erfaringsbekræftelse, i Praksis Modsætningen mellem Ideal og
Virkeliggørelse: disse Modsætninger ville stedse gøre sig gældende
med mere eller mindre Styrke, alt efter Opgavernes Natur.
De ligge i selve Livets stadige Vilkaar. Og de maa ikke
mindst gøre sig gældende i Filosofien, som jo kun er den
klare Besindelse over Principerne for vor teoretiske og
praktiske Stræben. »Skeptiske Stænk« eller en »skeptisk Brod«
ville derfor ikke komme til at mangle. Forhaabentligt da
ikke: ti de ere de fremaddrivende Kræfter, uden hvilke
Udviklingen vilde gaa i Staa. ---

Paa det \emph{erkendelsesteoretiske} Omraade gælder det, at de
sidste Forudsætninger for al vor Erkenden maa være af
hypotetisk Natur. En saadan Forudsætning er f. Eks.
Aarsagssætningen, den Antagelse, at ethvert Fænomen har en
Aarsag. Den ligger til Grund for al induktiv Forsken, men
kan ikke selv direkte bevises. Vi antage den, fordi den gør
real Videnskab mulig, men et Bevis kan der kun føres, for
saa vidt som det successivt kan paavises, at Fænomenerne
faktisk ordne sig i saadanne regelmæssige Rækkefølger, som
Aarsagssætningen fordrer. Et egentligt Bevis er det ikke,
da selve Aarsagsforholdet ikke kan iagttages; hvad vi
iagttage er kun Successionsforhold. Det er og bliver da en
Ejendommelighed for vor Erkendelse, at den fra først til
sidst arbejder med Hypoteser. --- Dette vilde kun være
Skepticisme, dersom man ved Hypotese forstod en vilkaarlig
Paastand eller Indbildning. Men i de Hypoteser, som ligge
% --- p. 29 ---
\opage{29}til Grund for vor Erkendelse, udtrykker sig selve vor Aands
Natur, dens inderste Trang, Trangen efter Enhed og
Sammenhæng, dens Higen efter ikke blot at fyldes med et rigt og
mangfoldigt Indhold, men at kunne ordne og beherske det.
Spørgsmaalet er kun, i hvilken Grad denne Trang kan ske
Fyldest.

Dr. Starcke synes at mene, at selve det, at man endnu
i Filosofien taler om et Aarsagsproblem og om
Vanskeligheden ved at begrunde Aarsagssætningen, vidner om, at der
endnu sidder os for megen Skepticisme i Kroppen. »Lad os
antage«, siger han, »at dette Spørgsmaal om Aarsagslovens
Gyldighed ikke var kommet til os fra en Tid, da man endnu
ikke troede paa den; hvorledes vilde vel saa Mennesket stille
sig over for hele dette Problem? Vilde det saa ikke snarere
være den, som rettede en Tvivl mod Aarsagssætningen, der
skulde bevise at have Ret til at tvivle, end den, som troede
paa den, der skulde begrunde sin Tro?« Dertil maa svares,
at naar det er Filosofiens Opgave at fremdrage de skjulte
Forudsætninger, som betinge vor Erkendelses Gyldighed, saa
er det jo ligegyldigt, hvem Bevisbyrden paahviler. Problemet
er der og maa behandles, hvad enten den enkelte Tænker
nærmer sig det mere fra den negative eller mere fra den
positive Side. Og Begrundelsen af Tvivlen er jo en Kritik
af Troen, ligesom omvendt, saa det gør ikke stort til Sagen,
hvorledes man gaar frem. Flere af de ældre »Skeptikere«
staa som Forløbere for den moderne Erkendelsesteori ved
den Maade, paa hvilken de have angrebet Troen paa
Erkendelsens Principer.

Naar der tales om en Tid, der ikke endnu troede paa
Aarsagssætningen, er dette en unøjagtig Udtryksmaade.
Der har næppe været nogen saadan Tid. Kampen for
Tilværelsen driver til at søge efter Aarsager for at faa Magt
over Tingene, og en Trang i den Retning ligger desuden i
den menneskelige Bevidstheds Natur. Selv i den vildeste
Overtro ytrer denne Trang sig. Derimod har der været
Tider, hvor man ikke til alle Fænomener antog \emph{naturlige}
Aarsager. De forskellige Udviklingstrin i intellektuel
Henseende karakteriseres ved, hvilke Aarsager der anses for gyldige.
% --- p. 30 ---
\opage{30}Miraklet strider ikke mod Aarsagssætningen, da det jo
udtrykkeligt opfattes som Virkning af en overnaturlig Aarsag.
Hvad Kampen stod om i den nyere Videnskabs
Begyndelsesperiode, var, om der er \emph{naturlige} Aarsager til Alt hvad der
sker, det vil sige Aarsager, der selv kunne paavises som
virkelige Fænomener. Og Bevisbyrden tilfaldt her saa meget
mere dem, der hævdede de naturlige Aarsager, som
Paavisningen af disse Aarsager jo netop er Videnskabens positive
Opgave. Det gaar med de naturlige Aarsagers Princip, som
det gaar med Aarsagsprincipet i det Hele: kun fremskridende
Verifikation er muligt.

Dr. Starcke hævder, at fordi vor Erkendelse er uafsluttet,
er den derfor ikke usikker, og han mener, at »den skeptiske
Holdning« i den moderne Erkendelsesteori hænger sammen med
en Forveksling af disse to Begreber. Og dog kommer han selv
i samme Aandedræt (p. 236) til at udtale, at Erkendelsens
Afsluttethed kun vilde have den Værdi at gøre dens \emph{Sikkerhed}
afsluttet. Nu er det dog klart, at en Sikkerhed, der ikke er
»afsluttet«, ikke er nogen fuld Sikkerhed. Et er det, at vi
ere sikre paa, at der kun er én Vej at gaa for at vinde
Erkendelse, et Andet er det, hvor langt denne Vej kan føre
os. Dette vide vi ikke, saa længe vi ikke have gaaet Vejen
til Ende. Der maa stedse voves og haabes; Afstanden mellem
Ideal og Virkelighed paa Erkendelsens Omraade vil ikke
forsvinde. Deri have de gamle Skeptikere Ret. ---

Foruden Problemet om Aarsagssætningen drøfter
Forfatteren ogsaa Problemet om Erkendelsens Forhold til Tingene
og om Forholdet mellem Sjæl og Legeme, for ogsaa her at
ramme »det skeptiske Drag, der gaar gennem vor Tids
Erkendelsesteori«. Jeg kan her ikke gaa nærmere ind paa
disse Punkter, skønt det vilde have sin Interesse. I det
Hele kan jeg desuden samstemme med hvad Forfatteren udvikler
angaaende de nævnte Spørgsmaal. Kun tror jeg, at det er et
»Stænk« af Spiritualisme, der gør, at han slaar fast, at
enhver Bevidsthedstilstand maa være et Atoms Tilstand (skønt,
som han søger at vise, alle vore Bevidsthedstilstande derfor
dog ikke behøve at være samme Atoms Tilstande). Atomet
er et Begreb, som man har konstrueret for at anskueliggøre
% --- p. 31 ---
\opage{31}sig hvad der foregaar ved de kemiske og fysiske Processer.
Der gives dog fremragende Fysikere, som helst ville være
fri for det og holde sig til at beskrive selve Bevægelserne
eller Virksomhederne. I ethvert Tilfælde har
Naturvidenskaben Intet at gøre med og véd Intet om, hvad der
foregaar i Atomets Indre. Det er Forholdet mellem Atomerne,
disses Omlejringer og indbyrdes Tiltrækken og Frastøden,
der spiller en Rolle ved Naturforklaringen. Og det er disse
samme Processer, til hvilke Erfaringen viser os
Bevidsthedsfænomenerne knyttede. Bevidsthedsfænomenerne hænge da
sandsynligvis sammen med hvad der foregaar mellem Atomerne,
med disses Omlejringer og Svingninger, ikke med mystiske
Tilstande i »selve« Atomerne. --- Dog dette er Noget, som kun
en længere Udvikling kunde give fuld Klarhed. ---

Ikke blot paa Erkendelsesteoriens, ogsaa paa \emph{Etikens}
Omraade har ifølge Dr. Starcke Skepticismen endnu for meget
Raaderum.

Rousseau og Kant havde, som allerede omtalt, hævdet
Muligheden af en Etik paa rent menneskelig Grund. Men
de overvandt, efter Dr. Starckes Opfattelse, dog ikke
Skepticismen. Denne ytrer sig hos dem --- og senere i den hele
rationalistiske Retning --- i en Disharmoni mellem Individets
Natur og den etiske Fordring, en Disharmoni, der bliver saa
meget des stærkere, som ifølge Kant den etiske Fordring
udspringer af selve Menneskets Natur. Troen paa, at det
Etiske vil bestaa og sejre i Verden, kan da kun fastholdes,
naar man tager Teologien til Hjælp og i den religiøse Tros
Form slaar det fast, man ikke kan begrunde ad Erkendelsens
Vej. Kant opstillede som bekendt i Form af religiøse
Postulater af etiske Grunde Antagelser, han i sin Fornuftkritik
havde berøvet ethvert teoretisk Grundlag. »Man tror,« som
Dr. Starcke udtrykker det, »paa en verdensstyrende Gud,
fordi man alene saaledes kan værne om sin Tillid til Moralens
Magt.« Det er Manglen af et teoretisk Grundlag, som giver
denne Anskuelse en skeptisk Karakter.

Men ogsaa i andre moderne etiske Retninger finder
Forfatteren en lignende Disharmoni og Skepticisme. Baade
Udviklingslærens Moral og Velfærdsmoralen have efter hans Op%
% --- p. 32 ---
\opage{32}fattelse Skepticismen i sig endnu. De begrunde Moralen fra
Slægtens Synspunkt. Men hvorfor skal den Enkelte ofre sig
for Slægten? Hvilket Motiv skal afholde ham fra saa vidt
muligt at udrydde eller modarbejde de nedarvede Anlæg i
Retning af Sympati og Samvittighed, han maatte besidde?
Med hvilken Ret fordres det af ham, at han skal opgive hvad
han anser for sin egen Lykke, for at en Gang i Fremtiden
andre Individer kunne nyde deres Lykke? --- Det nytter
ikke her at henvise til Sympatien. »Det er jo netop de
sympatiske Tilskyndelser, man opfatter som det store Bedrag,
Naturen øver over for Mennesket, og hvorigennem den faar
det narret til at tjene sine Formaal, idet Mennesket under
Indflydelse af disse Tilskyndelser tror at leve for sin Lykke
og saa dog kun lever for Slægten med Opofrelse af sig selv.«
Der kræves da et Bevis, som endnu ikke er ført, for, at »i
det Store og Hele det moralske Liv, som leves i Samfundet
og Familien, er det, som muliggør det for Individet at bringe
sine Lidelser og Forsagelser ned til et Minimum og for den
billigste Pris at faa den størst mulige Sikkerhed for
personlig Vinding.« Og --- hvad der er Hovedsagen --- denne
Bevisførelse maa foretages ud fra det enkelte, \emph{isolerede} Individs
Standpunkt: »Samfundslivets Fordringer maa formuleres
præcist og udtømmende, og disse Fordringer maa kunne tilegnes
af det isolerede Individ som dets egen Lykkefordring«.
Begrundelsen maa altsaa være egoistisk i den Forstand, at det
skal vises, at naar Individet vil følge sin velforstaaede
Interesse (sin Egoisme i videre Betydning af Ordet) staar
det sig ikke ved at være Egoist i Ordets snævrere
Betydning, det vil sige ved at tilfredsstille sig selv paa Andres
Bekostning.

Ingen er fjernere end Dr. Starcke fra at lege med Ord
eller jage efter »aandrige« Vendinger. Men jeg undrer mig
over, at han i sin Tankeudvikling bruger det tvetydige Ord
Egoisme, saaledes som han gør, uagtet han udtrykkeligt gør
opmærksom paa dets dobbelte Betydning. At dette skulde gøre
fuld Klarhed mulig, formaar jeg ikke at indse. Men om
Ord vilde vi i hvert Tilfælde ikke stride. --- Dog er der, af%
% --- p. 33 ---
\opage{33}set fra Sprogbrugen, Adskilligt, som gør, at jeg ikke kan
tilegne mig Forfatterens Tankegang i Et og Alt. For det
Første maa jeg bemærke, at det havde været ønskeligt, om
han havde givet sin Kritik af Evolutionsmoralen og
Velfærdsmoralen en bestemt Adresse og belagt den med bestemte
Citater. I en Bog af særligt historisk Karakter vilde det
være naturligt. I det Foregaaende har Forf. omhyggeligt
nævnt endog adskillige meget ubetydelige literære Fænomener,
som syntes ham at falde ind under Begrebet Skepticisme;
men her faa vi ikke et Navn, ikke en Henvisning. Derfor
bliver der noget Vagt ved den hele Udvikling. Den faar
Karakteren af et Skøn, og skal jeg mod dette Skøn stille
mit eget, da maa jeg sige, at hans Anke ikke forekommer
mig at være historisk berettiget. Man kan ikke sige, at
Problemet om Forholdet mellem Individ og Samfund er blevet
overset af de etiske Retninger, Forfatteren her kritiserer.
Det er tvert imod lige til den nyeste Tid blevet drøftet
med stor Iver, bl. a. i \emph{Henry Sidgwick}’s »Methods of Ethics«
og de Diskussioner, som dette udmærkede Skrift har givet
Anledning til.

Men Hovedsagen er, hvorledes det forholder sig med
selve dette Problem, saaledes som Forfatteren opstiller det.

For mit Vedkommende vedkender jeg mig at være
»Skeptiker« i Etiken, efter Dr. Starckes Maalestok, idet jeg
ikke tror, at Regnskabet kan gøres op mellem Individ og
Samfund paa den Maade, han fordrer. Der er noget
Irrationelt ved Etikens første Principer, for saa vidt det Grundlag,
ud fra hvilket disse fastsættes og anerkendes, ikke er det,
hvorpaa alle Individer, paa hvilke disse Principer blive at
anvende, allerede i Virkeligheden staa. De etiske Domme,
de Domme om Godt og Ondt, som det enkelte Individ bliver
Genstand for, kunne ikke altid fældes fra hans eget
nuværende Standpunkt. De Goder, om hvilke Etiken (den
praktiske Etik) værner, ere ikke altid Goder i det enkelte
Individs Øjne, og han kan ofte kun gennem en lang
Opdragelse (maaske Optugtelse), gennem en lang Karakterudvikling
lære at betragte dem som Goder. Et Bevis kan ikke altid
føres for ham, som han staar og gaar. Over for Børn, Sinds%
% --- p. 34 ---
\opage{34}syge og Forbrydere er dette især indlysende; dem underkaste
vi en praktisk Behandling, før vi stole paa Ræsonnementernes
Kraft. Men ogsaa hos andre Individer end saadanne, der
maa regnes til de nævnte Klasser, maa Grunden ofte
præpareres, før den rationelle Etik kan begynde sit Værk. Hvad
Etiken fastsætter, er da her de Midler og Veje, det pædagogiske
System (i videste Forstand), som bliver at anvende, for at en
egentligt etisk Behandling kan træde i Kraft.

Den rene, konsekvente Egoist vilde i det Højeste kunne
bringes til at betragte andre Menneskers Vel som et \emph{Middel}
for sit eget. Men han vilde da vedblive at staa i et rent
udvortes og mekanisk Forhold til Samfundslivet, og til alt
Liv uden for hans eget. Og skulde der kunne føres et Bevis
for, at Samfundslivet, som det i det givne Øjeblik er,
frembyder Harmoni mellem alle Individers egoistiske (»isolerede«)
Interesser? Selv om der kunde, hvem vilde egentligt et
saadant Bevis interessere? Ikke den egentlige Egoist; ham er
det nok, naar Samfundet blot er Middel for \emph{hans} Interesser,
--- ligegyldigt, hvorledes det saa gaar med de andre Egoister;
ja, maaske staar han sig bedst ved, at Andre ikke kunne
gøre Regnskabet op paa samme Maade.

Kan der i det Hele taget gives \emph{isolerede} Individer?
Derpaa kommer i denne Sammenhæng Alt an. Hos ethvert
Menneske rører der sig forskellige Følelser, der drive til
Vurdering af Livet og dets Forhold. Hvor mange af dem
skal Individet afføre sig for at blive isoleret? Er det i og
for sig sagt, at Sympatien og de med den beslægtede Følelser
skulle aflægges, naar Regnskabet skal gøres op? --- Mennesket
eksisterer kun som historisk givet Væsen, af bestemt Slægt
og Race, med bestemte Dispositioner og Interesser, og kan
ikke vilkaarligt drage sig ud af dette, heller ikke undlade
at faa sin Vurdering bestemt derved. Det vilde ved
Isolationen, om den var mulig, netop berøve sig nogle af de mest
nødvendige Poster paa Regnskabet. Lad mig udvikle dette
lidt nærmere.

Der er Interesseharmoni mellem to Mennesker, hvor de
kunne betragte hinanden som Middel, og hver for sig staa
som Formaal. Saaledes f. Eks. mellem den Handlende og
% --- p. 35 ---
\opage{35}hans Kunde. Hin trænger til dennes Penge, denne til Hins
Varer. Der var Interesseharmoni mellem Skræderen Mustafa
i Ispahan og hans Ven af samme Fag, ti som Morgiane
siger:

\begin{verse}
„Den Anden roste tit din Faders Kapper; \\
Til Gengæld blev hans Bukser rost igen. \\
Ak, det var Skræderlavets bedste Tid!“
\end{verse}

For en stor Del bestaar Samfundet ved Hjælp af
en saadan Interesseharmoni. At den ikke er til Stede i
fuldt Maal og i den ønskelige Grad, ses af de stadige
Klager over Samfundet og af, at der gives et »socialt
Spørgsmaal«. Man kan altsaa ikke beraabe sig paa den
i det givne, nærværende Øjeblik. Der kræves et Arbejde
for at bringe den til Veje, og der maa føres et særligt
Bevis for, at dette Arbejde kan udføres blot i Kraft af
de Motiver, der kunne findes hos det »isolerede« Individ,
og i Kraft af det, et saadant opfatter som sin »velforstaaede
Interesse«.

Dr. Starcke antager selv, at Regnskabet ikke \emph{endnu} kan
bringes til at stemme: »Det er en Fordel for Samfundet,«
siger han, »at der findes Mennesker, som sætte deres Dyd
over deres Fordel; men det er Samfundets Fejl, at der kan
findes et saadant Misforhold mellem Dyd og Lykke, og dets
Opgave maa være at se at rette den Fejl.« Det hele
Spørgsmaal maa jo saa staa aabent, indtil Fejlen er rettet. Men hvilke
Motiver skal man da henvise til, for at Individerne kunne
bringes til at arbejde paa at faa Fejlen rettet? Det kunde
jo være, at den fulde Interesseharmoni (Harmonien mellem
»isolerede« Individers Interesser) først sent hen i Fremtiden
blev etableret, og da vilde Dr. Starckes Teori være Genstand
for den samme Kritik, han selv retter mod Evolutionsmoralen,
idet den vilde fordre Opofrelse og Arbejde for en Lykke, der
tilfaldt en fremtidig Slægt, men ikke os selv. De »isolerede«
Individer vilde da maaske svare, at da deres
Efterkommere i fjernere Led ikke have gjort noget for dem, ville de
heller ikke gøre noget for disse Efterkommere. Après nous
le déluge! ---

% --- p. 36 ---
\opage{36}Men selv om en saadan Interesseharmoni en Gang skulde
kunne bringes fuldstændigt til Veje, saa at Enhver kunde
betragte alle Andre som Midler for sig selv, saa vilde derved
dog ikke et idealt Samfund, ikke den inderligste og stærkeste
Art af Samfundsliv være naaet. Et kraftigt Samfund
bestaar kun, naar der under Samlivet og Samarbejdet udvikler
sig Fællesfølelse og virkelig Sympati, saa den Ene ikke
skiller sin Sag fra den Andens, ikke blot betragter den
Anden som Middel, men ogsaa som Formaal. Kun under denne
Forudsætning vil det omtalte Regnskab kunne bringes til at
stemme. Det gøres da ikke længere op paa udvortes,
mekanisk Vis. Ti i selve Fællesfølelsen og Sympatien haves da
et af de største Goder, om hvilket der ikke blot er at spørge,
hvad det nytter til, men hvis Nytte er aldeles umiddelbar
(immanent), idet det fylder Sindet med en af de inderligste
Lystfølelser. Individets eget Væsen styrkes og udvides
gennem Hengivelsen; Horizonten udvides, og nye Muligheder
for Virken, Nyden og Liden vise sig; hele Livet føles paa
en rigere og dybere Maade. I ethvert Kærlighedsforhold gør
dette sig gældende, stærkest i de snevreste, mest
umiddelbare og personlige, i aftagende Grader i større Kredse, men
stedse med en Afglans fra hine. Der aabnes et nyt
Synspunkt for alle andre Følelser og Evner: alle de Følelser og
Evner, der føre i samme Retning som Fællesskabsfølelsen og
Sympatien, faa en Sanktion, de før manglede, en Accent, de
ikke kunde have, saa længe de blot vurderedes fra det
»isolerede« Individs ensomme Drejning om sin egen Akse. --- Denne
Akseomdrejning hører iøvrigt ikke op, fordi den hele Klode
nu viser sig at staa under Indflydelse af dybere liggende
Verdenskræfter.

Alt dette vil falde bort, dersom Individet blot
betragter sig selv som sidste Formaal. Joseph Butler brugte
i sin Tid det Udtryk, at »naar vi sætte os hen i en kølig
Stund« (when we sit down in a cool hour), kunne vi ikke
retfærdiggøre vor Stræben efter det Rette og Gode, naar
vi ikke ere overbeviste om, at det vil tjene til vor Lykke
eller i det Mindste ikke hindre denne. Han bruger denne
Paastand til at godtgøre Nødvendigheden af, at Etiken
% --- p. 37 ---
\opage{37}tilsidst støtter sig til Religionen og dens Forjættelser om
det Hinsidige (ligesom senere Kant og i vore Dage
Sidgwick). Men hans Standpunkt er i Virkeligheden det samme,
som det, der kommer frem i vor Forfatters Fordring om
en Begrundelse ud fra det isolerede Individs Standpunkt.
Over for Butler maa man opkaste det Spørgsmaal, hvorfor
den kølige Time her skulde have Ret og Begejstringens
Time, hvor Hengivelsens og Selvforglemmelsens Varme raader,
nødvendigvis Uret? Det maa erindres, at det her ikke
gælder Logik og Matematik, ikke Filosofi over Livet, men
Oplevelse og Konstatering af selve Livets virkelige Elementer.
Dersom »Køligheden« fører til, at nogle af de væsentligste
Elementer ikke kunne gøre sig gældende, saa kan Regnskabet
umuligt blive rigtigt.

Men jeg indrømmer, og det er endog en væsentlig Del
af min etiske Teori, at her kunne Vejene skilles, uden at
der foreløbigt ad teoretisk Vej kan opnaas Enighed\footnote{I 3. Kapitel af min \emph{Etik} (1887) havde jeg nærmere motiveret
dette Standpunkt. I første Kapitel af mine \emph{Etiske Undersøgelser}
(1891) tog jeg Spørgsmaalet op til ny Behandling. I anden
Udgave af \emph{Etik} (1897) er der igen paa flere Punker % PRINTED AS IS (for Punkter)
 søgt at opnaa
fornøden Tydelighed med Hensyn til dette Punkt. (Senere Anm.)}. Det
er muligt at koncentrere sig i sig selv, betragte denne
Selvkoncentration som Et og Alt, og ud fra dette
Synspunkt at gennemføre et konsekvent etisk System (som jeg
dog hellere vilde kalde Individualisme end Egoisme). Til
Fordel for den Retning, der gør Sympatien til Grundlag for
det etiske System, kan vel anføres, at derved føres man ad
mere direkte Vej til Interessernes Harmoni, end naar man
gaar ud fra de sondrede, isolerede Individer. Men et tvingende
Bevis er dette ikke.

Heller ikke nytter det her at fremhæve, hvad der
oven for blev vist, at Individets Liv bliver rigere og dybere
gennem Fællesskab og Sympati, altsaa netop ved, at det
ikke gør sig selv til eneste Formaal. Derom vil den Enkelte
kun kunne overbevise sig ad den praktiske Erfarings Vej.
Og Erfaring betyder her intet Mindre end en Karakterfor%
% --- p. 38 ---
\opage{38}andring. Den, som hengiver sig til de sympatiske
Tilskyndelser og lader sin Egoisme dæmpes, faar derved et helt nyt
Synspunkt for sin Vurdering, saa at han kommer til at drage
Konsekvenserne af sit Eksperiment fra et helt andet
Grundlag, end han begyndte fra.

Der vil derfor stedse være Plads for Skepticisme med
Hensyn til Muligheden for en rationel Etik. Da det
psykologiske og historiske Grundlag er forskelligt, maa de
Systemer, der bygges derpaa, ogsaa blive forskellige, selv om de
alle tage deres Udgangspunkt i den virkelige Menneskenatur.
Denne er i Vorden og Udvikling og omfatter mangfoldige
og ofte modstridende Elementer, af hvilke snart et, snart et
andet skyder sig frem med særlig Kraft. Om den historiske
Udvikling en Gang vil medføre fuldstændig Enhed i Etikens
Grundlag, skal jeg ikke sige Noget om; men før kan i hvert
Tilfælde »den skeptiske Brod« ikke helt drages ud. Den
etiske Tænkning taber derfor ikke sin Betydning. Den viser
os det Lys, i hvilket Handlinger og Livsformer fremtræde,
sete ud fra et bestemt Udviklingstrin, et bestemt i den
virkelige Menneskenatur givet Grundmotiv. Derved oplyser den
de Muligheder, der ligge i den menneskelige Natur, og virker
tilbage paa selve det Grundlag, der bygges paa; gennem
Konsekvenserne lære vi først vor Natur ret at kende.
Ogsaa en Fremstilling paa det af Dr. Starcke hævdede
individualistiske Synspunkt vil have sin store Interesse. En saadan
er i Grunden ikke for Alvor forsøgt endnu, hvorimod vi
have analoge Forsøg i den teologiske Etik,
Nationaløkonomien og Retslæren. Det vilde være overordentligt interessant
at se, i hvilke Punkter den mødes med Fremstillinger ud fra
andre Udgangspunkter. En alvorlig Tænkning bringer altid
sin Frugt. Af den psykologiske og historiske Erfarings rige
Skat øse vi Alle; men netop fordi den er saa rig, kan man
faa Adgang til den fra forskellige Sider.

Den enkelte Forsker maa betragte sig som en Bjærgmand,
der baner sig sin Skakt ind i Bjærget. Naar han føler sig
ensom under dette Arbejde, glæder han sig ved at høre
Kammeratens Hammer i Naboskakten. Han véd, at den
Anden arbejder mod samme Maal, skønt han har valgt sig
% --- p. 39 ---
\opage{39}et andet Udgangspunkt. Under min Beskæftigelse med Dr.
Starckes Bog har denne Kammeratskabsfølelse ikke forladt
mig, saa megen Kritik jeg end har maattet lade komme til
Orde. Den er det bedste af det meget Gode, jeg medtager
fra hans Bog\footnote{I sin Bog „Samvittighedslivet“ (1894---1897) har Dr. Starcke senere
givet en udførlig Fremstilling af sin Etik. (Senere Anm.)}.

\streg{}

% --- p. 40 ---
\opage{40}
\essay[\emph{OM VITALISME}]{\emph{OM VITALISME}\footnote{Et Foredrag, holdt i Biologisk Forening 27. Januar 1898.}.}

\streg{}

Det er mere end et billedligt Udtryk, naar man taler
om, at Tanker leve og do. Ligesom alle organiske Væsener
have Tanker deres bestemte Betingelser, hvoraf deres
Opstaaen, Udvikling eller Forsvinden afhænger. Tankerne
maa, ligesom Organismerne, kæmpe for at leve. Kun de
Tanker leve i Længden, der kunne bestaa med og være
Udtryk for store Kredse af vel prøvede og sammenhængende
Erfaringer. Og ligesom der sker en naturlig Udvælgelse
blandt de organiske Spirer, Tilløb og Variationer, som paa
ethvert Punkt af Udviklingen ere til Stede, saaledes sker
der ogsaa stadigt et Valg mellem de Tanker og Synspunkter,
som til den givne Tid foreligge. Der er et Vekselforhold
mellem Tænkning og Erfaring. Ti de uvilkaarligt dannede
Tanker og Synspunkter, mellem hvilke Erfaringen vælger,
ere naturligvis ikke selv opstaaede uden Indflydelse af
Erfaringen. Men ligesom der i Naturen opstaa langt flere
Spirer og Variationer, end der udvælges til fortsat Udvikling,
saaledes begynder ogsaa Tænkningen med energiske og frodige
Indskud og Antecipationer, der efterhaanden sigtes, begrænses,
udvikles eller vrages.

Dette gælder særligt hvad Forstaaelsen af det
organiske Liv angaar. Kampen mellem Spekulation og Verifikation
bevæger sig paa rytmisk Maade gennem Biologiens Historie,
og i den nyeste Tid synes der at røre sig en Bølgegang, der
fører til at genoptage Synspunkter, hvis Tid afgjort syntes
at være forbi. Man har ment at finde Liv i gamle Tanker,
% --- p. 41 ---
\opage{41}hvis Dødsfald allerede var bekendtgjort. I hvilken Forstand
denne Mening er berettiget, skulle vi nærmere undersøge.

Der er to Grundopfattelser, som her staa over for
hinanden. Den ene støtter sig til, at den klareste Forstaaelse,
vi kunne have af et Fænomen, have vi der, hvor vi kunne
opløse det i dets Elementer og igen se det som Produkt eller
Resultat af disse Elementers Samvirken, --- altsaa naar vi
gaa fra Delene til Helheden. Den saakaldte mekaniske
Naturopfattelse hævder, at denne Forklaringsmaade ikke blot
gælder for de fysiske og kemiske Fænomener, men for alle
Naturfænomener. Herimod gøres der fra den anden Side
gældende, at det Ejendommelige for de organiske Fænomener
er, at Elementernes, Delenes Beskaffenhed, Form og
Virkemaade ikke kan forstaas uden ved at lægge Tanken om den
Helhed, de danne eller skulle danne, til Grund. Organismen
kan ikke forklares som blot Produkt af Organer eller Celler,
og Cellen ikke som blot Produkt af Molekuler. Vi maa
da her antage en særegen Kraft, der virker efter andre Love
end de mekaniske Love. Dette er den vitalistiske Opfattelse.
--- Da den eneste Maade, hvorpaa vi kunne forestille os en
Helhed som bestemmende for Delenes Beskaffenhed og
Samvirken, er at forestille os den som et Formaal, der tilstræbes,
saa vil en vitalistisk Opfattelse stedse blive mere eller mindre
teleologisk.

\streg{}

Naar Historien viser os en rytmisk Skiften af disse
modstaaende Synsmaader, udelukkes derved ingenlunde
Erkendelsens Fremskriden. Det er de nye Erfaringer, som
lade Fænomenerne fremtræde snart mere fra den ene, snart
mere fra den anden Side, og som tillige efterhaanden
omdanner de tidligere Former for de to Grundopfattelser.

En mekanisk Naturopfattelse kommer først frem paa et
forholdsvis fremskredent Standpunkt. Naturmenneskets
Opfattelse er vitalistisk eller animistisk, forklarer alle
Fænomener, der vække Opmærksomhed, ved personlige Væseners
Indgriben. Al senere Vitalisme har et vist Slægtskab med
denne primitive Opfattelse. --- I den græske Oldtid gøres
% --- p. 42 ---
\opage{42}Mekanismens Synspunkt gældende af \emph{Atomisterne} (hvis Lære
senere fremstilledes i Lucretius’ berømte Digt de rerum natura).
\emph{Aristoteles}, der selv var en genial Naturforsker paa det
organiske Livs Omraade, hævdede i Modsætning til
Atomistiken en Opfattelse, der betragtede Naturen som en stor
Trinrække, hvor ethvert Led er Stof og Middel i Forhold
til et højere, medens det er Form og Formaal i Forhold til
et Lavere. Intetsteds i Naturen kunde han tænke sig Helheden
opstaa som Resultat af Elementerne, Formen som Virkning
af Stoffet. Han udstrækker altsaa Vitalismen til hele Naturen.

Ejendommeligt for Livsproblemets Stilling i Oldtiden er
det saaledes, at der ikke sattes nogen skarp Grænse mellem
uorganiske og organiske Fænomener. Baade for den
aristoteliske og for den atomistiske Opfattelse er der i Naturen kun
Gradsforskelle.

Foreløbigt sejrede den aristoteliske Opfattelse, og den
blev den herskende i Middelalderen. Et Vendepunkt kommer
med \emph{Grundlæggelsen af den moderne Naturvidenskab} ved
Leonardo da Vinci, Kepler og Galilei. Man opstillede nu
den Sætning: al Forandring er Bevægelse, og fordrede den
gennemført paa alle Omraader, saa at alle Naturlove kunde
udledes af den mekaniske Bevægelseslæres Principer. \emph{Den
kartesianske Filosofi} generaliserede denne Opfattelse og førte
den ud i store Kredse. Opdagelsen af Blodets Kredsløb og
af Refleksbevægelsen var et erfaringsmæssigt Støttepunkt for
den Overbevisning, at det organiske Livs Omraade ikke
kunde danne nogen Undtagelse.

Som Reaktion mod hele denne Retning fremstaar den
Opfattelse, som man nu nærmest tænker paa, naar man taler
om Vitalisme. \emph{G. E. Stahl} bekæmpede den mekaniske
Fysiologi og hævdede (især i Theoria medica vera 1708), at uden
en Sjæls Indgriben var den organiske Vækst og Organismens
Selvopholdelse uforklarlig. Han gaar altsaa længere tilbage
end til Aristoteles, idet han fornyer den primitive Animisme.
Han mente, at naar den bevidste Sjæl kan være Aarsag til
vilkaarlige Bevægelser, maa den ogsaa kunne være Aarsag
til de uvilkaarlige Processer i Organismen, selv om den her
virker paa mere instinktmæssig Maade, og selv om den ikke
% --- p. 43 ---
\opage{43}altid bærer sig fornuftigt ad, hvad den jo heller ikke altid
gør i sine vilkaarlige Handlinger. --- Den bekendte
Lægeskole i Montpellier ændrede Teorien ved at skelne mellem
Sjæl og Liv, saaledes at der mellem Sjæl og Legeme
indskødes et særeget Livsprincip. I denne Form er Vitalismen
fremstillet af \emph{Barthez} (Nouveaux Elements de la science de
l’homme. Montpellier 1778). Han bekæmper to Sekter,
Mekanisternes og Animisternes, og repræsenterer Vitalismen
i snevreste Betydning af Ordet. Han definerer Livsprincipet
som »den Aarsag, der frembringer alle Livsfænomenerne i
det menneskelige Legeme«. Det følger Love, der ligge over
Fysikens Love (d’un ordre transscendant par rapport aux lois
de la Physique). Dets egentlige Natur, særligt dets Forhold
til den egentlige Sjæl, er ubekendt. Men Barthez opfatter
det tydeligt nok i Analogi med det egentlige Sjæleliv;
tillægger det Instinkt, Temperament og Sanseevne. Det griber
ind i Livsprocessen og leder denne i en bestemt Retning.
Det lægger den Kraft i Musklerne, som behøves til
Overvindelse af en bestemt Modstand, og viser særligt sin Magt
i konvulsiviske Anfald. Det udvider Muskelfibrene i Mælkens
og Sædens Udførselsgange, og det virker paa Safterne, saa
at der dannes de Stoffer, som behøves til Reparation eller
Forplantning. Det regulerer Aandedræt og Varmeudvikling
og betinger en »sympatisk« Forbindelse mellem Organismens
forskellige Dele. --- Forskellig baade fra Barthez’s og fra
Stahls Vitalisme er den Teori, den berømte Anatom \emph{Bichat}
opstillede (i Recherches physiologiques sur la vie et la mort.
Paris 1800). Han opfatter Livskraften som en Egenskab ved
de organiske Væv. Det var en begyndende Decentralisation
af Vitalismen.

\emph{Den klassiske Vitalisme} (saaledes betegner jeg den
Opfattelse, der i forskellige Afskygninger repræsenteres af Stahl,
Barthez og Bichat) har det fælles med den aristoteliske, at
den ved Antagelsen af en Livskraft mener at have givet en
tilstrækkelig Forklaring af Livsfænomenernes
Ejendommelighed. Den adskiller sig fra den aristoteliske Vitalisme ved,
at den skarpt fremhæver Modsætningen mellem det Uorganiske
og det Organiske og endog antager et fjendtligt Forhold
% --- p. 44 ---
\opage{44}mellem dem. Sjælen eller Livskraften maa kæmpe med de
mekaniske og kemiske Forhold for at faa dem til at virke i
en bestemt Retning, ja Bichat definerer endog Livet som
Indbegrebet af de Kræfter, der kæmpe mod Døden.

Et nyt og storartet Opsving af den mekaniske
Naturopfattelse i dens Anvendelse paa Biologien fremkom paa
Grundlag af italienske, franske og schweizerske Naturforskeres
Arbejder, ved Slutningen af det 18. Aarhundrede, som dog
først ret kom til Anerkendelse henimod Midten af vort
Aarhundrede. Ved Aar 1840 reformerede \emph{Dumas} og \emph{Liebig}
Plantefysiologien og paaviste Stoffets Kredsløb i Naturen, og
dermed Sammenhængen mellem den organiske og den
uorganiske Natur. Snart efter opdagede \emph{Robert Mayer} Loven om
Energiens Bestaaen og anvendte den straks paa det organiske
Omraade (i sit Skrift: Die organische Bewegung 1845). En
Sætning af dette Skrift kan staa som Motto for den hele
nye biologiske Opfattelse: »Under Livsprocessen foregaar der
kun en Forvandling af Materie og af Kraft, men aldrig en
Skabelse af denne eller hin.« En hel Række af fremragende,
især tyske Fysiologer (navnligt Johannes Müllers Disciple ---
uagtet Mesteren endnu var Vitalist) viste de nye Synspunkters
Frugtbarhed. Fysiologien syntes nu at kunne opfattes som
anvendt Fysik og Kemi. Antagelsen af en Livskraft
stempledes som Mysticisme og Overtro. \emph{Helmholtz} bruger (i sit
Skrift Das Denken in der Medicin 1878) Vitalisme som
Modsætning til Naturforsken og betragter Spørgsmaalet om
Livskraften som Et med Spørgsmaalet om perpetuum mobile. Og
i sin Mindetale over Helmholtz (1897) udtalte \emph{Dubois
Reymond}, at Læren om Energiens Bestaaen gav Gøglebilledet af
en Livskraft det sidste Stød. Den Periode i Biologien, som
begyndte ved 1840, er en af de frugtbareste i hele dens
Historie, særligt ved mange Anvendelser af fysiske og
kemiske Love paa det organiske Omraade. En stor og
sammenhængende Naturopfattelse kunde herved en Gang for alle synes
at være vunden.

I de seneste Aar er man saa igen begyndt at tale om
Vitalisme som et nødvendigt Standpunkt. I en vis
Sammenhæng hermed staar det, at man taler om, at Naturvidenskaben
% --- p. 45 ---
\opage{45}har skuffet de Forventninger, man havde stillet til den, ja,
at man endog har talt om »Videnskabens Bankerot«. De
Fakta, til hvilke den saakaldte \emph{Nyvitalisme} støtter sig, vare
i Hovedsagen allerede kendte tidligere og fremhævede af
kritiske Forskere. Mænd som Lotze og Virchow, Littré og
Claude Bernard, der alle hævdede, at kun Fysikens og Kemiens
Metoder kunde føre til Indsigt i Livsfænomenernes Love,
mente dog ingenlunde, at Livets inderste Natur fuldstændigt
kunde forstaas ud fra disse Love alene. De betragtede Livet
som et Faktum og saa det som Biologiens Opgave at
undersøge de enkelte Led i den faktisk forløbende Livsproces.
Anledningen til Nyvitalismens Fremtræden kan kun søges i,
at mange Dele af Livsprocessen have vist sig at være \emph{endnu}
mere indviklede, end man paa Forhaand havde trot. Der er
især tre Punkter, der træde i Forgrunden for Nyvitalismen.

1) Selv om der ved hver enkelt Overgang i Livsprocessen
kun skete saadanne Stof- og Kraftomsætninger, som kunde
forklares efter kemiske og fysiske Love, saa vilde dog den
Kendsgerning staa tilbage, at de kemiske og fysiske
Processer under Væksten og Ernæringen forløbe i bestemte
\emph{Retninger}. Dette er saa meget des mere iøjnefaldende, som
Væksten foregaar sporadisk, ud fra forskellige Punkter, saa
at Sammenslutningen til en Helhedsform først sker
efterhaanden. Skulde det da ikke være berettiget at tale om en
Formdrift (nisus formativus)?

2) Ved \emph{Optagelsen af Næringsstof} i Cellerne øves der et
Valg. Hverken Kvantiteten eller Kvaliteten af det, der
optages kan forklares ad rent kemisk eller fysisk Vej. Cellerne
synes her at optræde som smaa Væsener, hvis Trang og Ønske
bestemme deres Forhold til Omverdenen. Flere Nyvitalister
personificere derfor Cellerne, ligesom de klassiske Vitalister
personificerede hele Organismen. Da vi kun fra os selv
kende, hvad Trang og Ønske er, har en af de ivrigste
Nyvitalister, \emph{Bunge}, endog udtalt, at Psykologien egentligt her
burde komme Fysiologien til Hjælp.

3) Det organiske Væv adskiller sig fra Krystallen ved,
at dets Molekuler \emph{aldrig} komme i \emph{fuldstændig Ligevægt}, saa%
% --- p. 46 ---
\opage{46}længe Livet varer. Det at kunne dø er derfor et Kendemærke
paa det Levende. ---

Utvivlsomt er det, at Opmærksomheden overfor disse
Punkter er bleven skærpet i de senere Aar. Tvivlsomt er
det derimod, om der derved virkeligt er kommet et nyt
Stadium frem i den gamle rytmiske Proces, saa at Navnet
Nyvitalisme kunde være berettiget. Hvor vidt man vil
strække Navnet, kan i hvert Tilfælde være tvivlsomt. Forskere
som \emph{Giulio Fano} (La fisiologia come scienza autonoma. Genova
1884) og \emph{Hans Driesch} (Die Biologie als selbständige
Grundwissenschaft. Leipzig 1893), der kun ville hævde Biologiens
Selvstændighed over for Kemien og Fysiken, kunne kun
uegentligt kaldes Nyvitalister. Derimod passer Navnet paa
\emph{G. Bunge} (Kapitlet »Vitalismus und Mechanismus« i
»Lehrbuch der physiologischen und pathologischen Chemie«. Leipzig
1887), der næsten fornyer Stahls Animisme. Fra indbyrdes
meget forskellige Standpunkter have \emph{Haeckel}, \emph{Verworn} og
\emph{Rindfleisch} udtalt sig i lignende Retning\footnote{Paa den betydningsfulde Forskel mellem de Opfattelser, der ofte
under Et betegnes som Nyvitalisme, har allerede \emph{Felice Tocco} gjort
opmærksom i sin Afhandling „Le disfatte della scienza“ (Nuova
Anthologia. 1. Marzo 1896).}.

\streg{}

Det vil have sin Interesse at sammenligne Nyvitalismen
med den klassiske Vitalisme for at se, om den, som Nogle
(f. Eks. \emph{Mosso} i sin Afhandling »Materialismo e Misticismo«
i Nuova Anthologia 1. Decbr. 1895) mene, er Tegn paa en
Tilbagegang, eller om der i Forhold til tidligere, mere eller
mindre beslægtede Retninger kan paavises Forskelle, der
tyde paa Fremgang. En saadan Sammenligning kan føres
tilbage til tre Hovedpunkter.

1) Den klassiske Vitalisme mener, at der ved Antagelsen
af en »Livskraft« er naaet en afsluttende, positiv Forklaring.
Men for de mest kritiske Forskere, som Nyvitalismen tager
% --- p. 47 ---
\opage{47}til Indtægt, gælder det blot at hævde, at Biologien har sit
selvstændige Erfaringsgrundlag og er Andet og Mere end et
Indbegreb af fysiske og kemiske Deduktioner. For saa vidt
der tales om en Livskraft, betyder dette Ord et x, en uopløst
Rest, og staar som Udtryk for de sammensatte Forhold, under
hvilke de almindelige Naturlove finde Anvendelse paa det
organiske Omraade.

2) Den klassiske Vitalisme blev af Claude Bernard kaldt
»en Lære, der gør doven« (une doctrine paresseuse). I sin
Tillid til Livskraften ringeagtede den eksakte Undersøgelser
af de enkelte Forhold. Barthez mente f. Eks. at det var
unødvendigt og urimeligt at søge et anatomisk Grundlag for
Sammenhængen mellem sensoriske og motoriske Funktioner;
Livskraftens Enhed var tilstrækkelig Forklaring af denne
Sammenhæng. Han afviste altsaa paa Forhaand Undersøgelser
som dem, der senere førte Charles Bell til hans Opdagelse af
Forskellen mellem sensoriske og motoriske Nerver. (Se
Barthez: Nouveaux Éléments. p. 142---153; 207 f.)\footnote{Dette Punkt har allerede \emph{Louis Peisse} fremdraget i sin Kritik af
Montpellierskolen. (La Médicine et les Medecins. Paris 1857. I.
p. 291---293).}. Nyvitalismen
derimod fastholder --- trods sin Tilbøjelighed til psykologiske
Analogier --- at den eneste Metode, der faktisk kan være
Tale om at anvende i Biologien, er den, som den mekaniske
Opfattelse har proklameret. Den anser Livsproblemet for
mere indviklet, end den klassiske Vitalisme gjorde.

3) Den mest iøjnefaldende Forskel mellem Nyvitalismen
og den klassiske Vitalisme bestaar i, at Livskraften ikke
eller ikke blot tillægges Organismen som Helhed, men Cellerne,
Organismens Elementer. I Stedet for en Genius, der leder
Organismens Liv som Helhed, er man nu tilbøjelig til at
antage en Art Cellegenier, som afgør hvad og hvor meget
der skal optages fra det indre Medium. \emph{Bunge} udtaler
ligefrem: »Wir übertragen das aus dem eigenen Bewusstsein
geschöpfte auf die Objecte unserer Sinneswahrnehmung, auf \emph{die
Organe}, auf \emph{die Gewebselemente}, auf \emph{jede kleine Zelle}.« Det
betegner utvivlsomt et Fremskridt, at man bruger Kraft%
% --- p. 48 ---
\opage{48}begrebet om de enkelte samvirkende Elementer, i Stedet for
en bloc om hele Organismen. Problemet er blevet mere
lokaliseret; Gaaderne ere lagte længere tilbage. Ganske vist
staar saa det store Problem tilbage, hvorledes de mange
Cellers Liv forbindes til Organismens samlede Liv, og det
var især dette, der vakte den Forundring, hvis teoretiske
Udtryk den gamle Vitalisme var. Nyvitalismen er ogsaa
Udtryk for en Forundring, men især for den, der vækkes
ved det Hemmelighedsfulde i de Processer, der gaa for sig i
Organismens mindste Dele. Man havde ikke ventet, at Livets
Problem skulde komme igen ved de Elementer, i hvilke man
mente at have opløst Livshelheden. ---

Forskellen mellem Nyvitalismen og den klassiske Vitalisme
er saa stor, at det egentligt er uberettiget at sige, at det er
en og samme Opfattelse, der er kommen igen. Med Hensyn
til det, der videnskabeligt paastaas, og til den Metode, der
faktisk anvendes, er der egentligt i Forhold til den nærmest
forudgaaende Periode i Biologiens Historie ingen Forandring
sket. Al Tale om Tilbageskridt (Mosso) eller Bankerot
(Brunetière) er derfor uberettiget, og selve Navnet
Nyvitalisme burde helst undgaas.

\streg{}

Til Slutning vil det være af Interesse at fremhæve
Analogien mellem det Problem, der gør Biologien til Andet og
Mere end anvendt Kemi og Fysik, og en hel Række andre
Problemer, der alle bero paa Modsætningen mellem det
Rationelle og det Empiriske, eller --- som man ogsaa kan udtrykke
det --- mellem det Formale og det Reale i vor Erkendelse. Jo
mere rationel og formal vor Viden er, des mere kan den
bevæge sig frem gennem Definitioner og Deduktioner, saaledes
at enhver Overgang foregaar med anskuelig Nødvendighed.
Den empiriske eller reale Viden maa derimod ofte stanse ved
Fakta, der vel kunne beskrives og analyseres, men ikke defineres
og udledes af andre Fakta. Et saadant Faktum er Livet, og
derfor hører Biologien, og vil maaske altid komme til at høre,
% --- p. 49 ---
\opage{49}til den empiriske Viden. Det var ikke blot Stahl, der i sin
Tid savnede en Definition af hvad Livet er (et Savn, han dog
selv troede at kunne afhjælpe); det er ogsaa i vore Dage (af
Lotze og af Littré, af Bernard og af Fano) blevet udtalt, at en
saadan Definition ikke er funden. Men Livet staar i denne
Henseende ikke alene. Paa en hel Række Punkter afløses
den rationelle Kontinuitet i vor Viden af empiriske Data, der
foreløbigt staa som uopløselige. Det er Kurrer, eller maaske
endog Knuder paa Traaden.

1) Lad os tænke os, at den Sætning, at al Forandring
er Bevægelse, fuldstændigt kunde gennemføres, og at den
mekaniske Naturopfattelses Ideal for saa vidt kunde
virkeliggøres. De empiriske Kurrer eller Knuder paa Traaden vilde
derfor ikke mangle. Paa ethvert Punkt --- hvor langt vi end
gaa tilbage --- vilde Bevægelsen dog optræde med en vis given
Retning og en vis given Hastighed, og de ydre Aarsager, som
ændre denne Retning og denne Hastighed, vilde virke fra visse
bestemte Steder og i visse bestemte Grader og Retninger. ---
Den Sum af Energi, som anvendes ved Bevægelsen og ved
de Aarsager, der ændre den, er desuden af en bestemt given
Størrelse, og dens Fordeling paa de enkelte virkende Elementer
kan ogsaa kun kendes empirisk.

2) Nu optræder tillige Energien i kvalitativt forskellige
Former, og ligeledes have de virkende Elementer deres
kvalitative Ejendommeligheder. Selv om der i Følge Loven om
Energiens og Stoffets Bestaaen stedse kan paavises Ækvivalenter
ved enhver Omsætning fra en Energiform til en anden og ved
ethvert Stofskifte, saa opstaar der tillige ved disse Omsætninger
nye Egenskaber, der ikke kunne udledes af de forudgaaende
Betingelser, men rent empirisk knyttes til dem.

3) Paa det organiske Omraade bliver Indbegrebet af
empiriske Elementer endnu større og mere udpræget. Det er
Opmærksomheden for disse Elementer, der ofte medfører en
fornyet Tilbøjelighed til vitalistiske Anskuelser. Som saadanne
Elementer have vi i det Foregaaende omtalt Vækstretning og
Struktur, kvalitativt og kvantitativt Valg ved Cellernes
Stofskifte og Bevarelsen af den ustadige Ligevægt.

4) Hvis vi tage det psykologiske Omraade med, staa
% --- p. 50 ---
\opage{50}Bevidsthedsfænomenerne --- lige fra den mest elementære
Fornemmelse til de højeste Tanker, Følelser og Viljesakter ---
som et empirisk Givet, der ikke kan udledes af de materielle
Forudsætninger. Psykologien er i endnu højere Grad end de
andre Dele af Biologien en rent deskriptiv Videnskab. Og
indenfor Psykologien kommer der saa igen nye kvalitative
Forskelle og Omsætninger frem, hvor det samme Problem
paany stiller sig.

Det er klart, at hvis man paa alle de Punkter, hvor det
Empiriske, Kvalitative og Individuelle stiller den rationelle
Kontinuitet, vor Erkendelse tilstræber, paa Prøve, vilde
antage nye selvstændige og oprindelige »Kræfter«, vilde man
bedrage sig selv. De Grænser, vi saaledes stadigt støde paa,
ere stedse tillige at betragte som nyt Stof, der stiller nye
Opgaver. Det er Tilværelsens Rigdom og Dybde, der
frembringer disse Kurrer og Knuder paa Traaden og derved
nødvendiggør en stadig Vekselvirking % PRINTED AS IS (for Vekselvirkning)
 mellem Erfaring og
Tænkning, Induktion og Deduktion, Analyse og Syntese, ---
denne Dobbelthed i Metoden, som Goethe har sammenlignet
med Forholdet mellem Indaanding og Udaanding. Kampen
mellem Mekanisme og »Vitalisme« er, ret forstaaet, en Kamp
mellem to Tendenser, mellem hvilke der i Grunden ingen
Modsætning kan være. Vi vilde ikke være fuldkomnere
Væsener, om vi kunde aande ud uden at aande ind. Hin
Dobbelthed gør baade Livet og Tanken rigere.

\streg{}

% --- p. 51 ---
\opage{51}
\essay[SINDSBEVÆGELSERNES FYSIOLOGI]{SINDSBEVÆGELSERNES FYSIOLOGI.}

\begin{center}\emph{C. Lange:} Om Sindsbevægelser. En psykofysiologisk Studie. 1885.\end{center}

\streg{}

Det er en gammel Fordom, at hvad der staar som smukt,
herligt eller ophøjet for vor Følelse, skulde miste sin Glans
og sin Værdi derved, at dets Tilbliven, Udvikling og
Bestaaen forklaredes efter simple naturlige Love. Det var en
Forargelse for Oldtidens Grækere, da Filosofer og
Naturforskere begyndte at tale om, at Sol og Maane vare
materielle Legemer ligesom Jorden; man havde hidtil betragtet
dem som Guder, der majestætisk og ved overnaturlige Kræfter
bevæge sig over Himlen til Beundring og Nytte for Jordens
Beboere. De naturvidenskabelige Forklaringsforsøg have saare
ofte mødt den samme Fordom og den samme Modstand igen.
Men hvor stærk maa Modstanden da ikke blive, naar det er
selve Følelseslivet, man vil forklare paa naturvidenskabelig
Vis? Det vil for Mange staa som om det var at dræbe
Følelsen med det Samme, --- som om det at forklare den var
Et med at bortforklare den. Det foreliggende Skrift, som
tilstræber en rent fysiologisk Undersøgelse og Forklaring af
Sindsbevægelserne, vil næppe undgaa at mødes med saadanne
Fordomme. Og dog er det sikkert, at det kun nærmere
anvender og gennemfører en Opfattelse, som længe har været
gjort gældende i den videnskabelige Psykologi og Fysiologi.

Ved Sindsbevægelser forstaar Forf. \emph{de simpleste, mest
usammensatte Følelsesfænomener}, som danne de Bestanddele,
af hvilke de egentlige Følelser og Lidenskaber bestaa.
Saadanne Sindsbevægelser ere Sorg, Glæde, Frygt, Vrede,
Spænding o. s. v. De komme som Elementer igen i saadanne mere
sammensatte Følelser som Kærlighed, Had, Beundring, For%
% --- p. 52 ---
\opage{52}agt o. s. v. Denne Inddeling stemmer i det Hele med den,
som den nyere Psykologi lægger til Grund, som det jo ogsaa
er den naturligste Fremgangsmaade at gaa fra det
Elementære til det mere Sammensatte. Man kunde dog maaske gøre
Indvendinger mod den Maade, hvorpaa Forf. i det enkelte
drager Grænsen mellem de to Grupper af Følelsesfænomener.
Til de simple Følelser, der benævnes Sindsbevægelser, burde
ogsaa Forundringen være regnet. Der er et Element af
Forundring, hver Gang et Sanseindtryk kommer til Bevidsthed,
idet det er en Betingelse for enhver Fornemmelses Opstaaen,
at den Tilstand, Indtrykket fremkalder, staar i en vis
Modsætning til den forudgaaende Tilstand. Jo mere et Indtryk
er forskelligt fra det foregaaende, og jo pludseligere det
kommer, des stærkere bliver det Element af Forundring, som
gør sig gældende, og som kan danne Indledningen til meget
forskellige Sindstilstande (Frygt, Skuffelse, Foragt og
Kærlighed o. s. v.). Der er saa meget des mere Grund til at
fremhæve dette, som den Grænse, Forf. vil hævde mellem de to
Grupper af Følelsesfænomener, næppe vil kunne fastholdes,
uden naar man til de af ham anførte Skelnemærker føjer
endnu et, nemlig at Sindsbevægelserne mere have Karakteren
af pludselige Opbrusninger, medens de egentlige Følelser eller
Lidenskaberne ere fastere, varigere, mere indgroede Tilstande.

Disse Sindsbevægelser vil Forf. nu studere ved Hjælp
af \emph{den fysiologiske Metode}. Han havde oprindeligt stillet sin
Opgave saaledes, at han spurgte om Sindsbevægelsernes
Virkninger paa de legemlige Livsytringer; men han kom til det
Resultat, at det var en helt urigtig Maade at stille
Spørgsmaalet paa. Sindsbevægelsen er ikke Ét, dens legemlige
Virkning et Andet. Vi have, fysiologisk set, her ikke et
Hovedfænomen, »selve« Sindsbevægelsen, og et Bifænomen,
dens Virkninger paa Organismen. Fysiologen kan ikke
indrømme, at disse Virkninger ere noget Underordnet; det er
for ham netop dem, der udgøre hele Sindsbevægelsen. Man
gaar ud fra en forudfattet og ubegrundet Distinktion mellem
Sjæl og Legeme, naar man henlægger Sindsbevægelsen til
Sjælen og først efter at den er opstaaet dér lader den virke
paa Legemet. Fysiologisk set kunne vi kun studere Sinds%
% --- p. 53 ---
\opage{53}bevægelsen gennem det, der foregaar i Legemet; en
Paastand om, at dette er foraarsaget af Noget, der foregaar i
Sjælen, er uvidenskabelig og ubevislig. Vi maa i hvert
Tilfælde se, hvor langt vi kunne komme ved Hjælp af den blotte
Fysiologi.

Forf. viser dette ved en fysiologisk Skildring af nogle
af de vigtigste Sindsbevægelser. Dette Parti af hans Bog
er skrevet med sandt Mesterskab. Man mærker den øvede
kliniske Iagttager, og man kommer til at tænke paa, at den
indtrængende psykologiske Sans --- som ytrer sig midt under
den fysiologiske Karakteristik --- er et Arvegods i Forfatterens
Slægt, nu allerede i tredje Generation. Her kan jeg kun
meddele nogle Grundtræk af den fyldige Fremstilling.

Ved \emph{Sorgen} svækkes Evnen til vilkaarlig Bevægelse.
Deraf den Sørgendes trætte Udseende, bøjede Holdning og
vaklende Gang. Derimod trække de uvilkaarlige Muskler
sig sammen, især de, der findes i Blodkarrenes Vægge og
som sættes i Virksomhed ved Nerver, der udgaa fra det
»karbevægende« (vasomotoriske) Centrum i Hjernen. Blodet
presses derved ud af de fine Kar, saa at Legemets Væv og
Organer blive fattige paa Blod; deraf følger Bleghed,
Kuldefornemmelser og Slaphed i Trækkene. Taareafsondringen
synes at kunne forklares som Følge af en Slappelse af
Karmusklerne efter den første Sammentrækning; derfor er Graaden
en Lettelse og udebliver ved den voldsomste Sorg. Paa lignende
Maade synes Sukket at være den naturlige Reaktion mod
den ved Sammentrækningen af Lungernes fine Kar fremkaldte
Aandenød. --- Ved Glæden forhøjes derimod Evnen til
vilkaarlig Bevægelse og de fine Blodkar udvides. En Fornemmelse
af Lethed, Varme og Kraft opstaar. Og medens Hjernens
Blodfattigdom under Sorgen viser sig i aandelig Træthed og
Sløvhed, saa viser Hjernens rigelige Blodforsyning under
Glæden sig i Tankens og Indbildningkraftens Flugt. --- \emph{Skrækken}
staar fysiologisk Sorgen nær, kun at Symptomerne ere stærkere,
og at ikke blot Karmusklerne, men ogsaa alle andre
uvilkaarlige Muskler synes at trække sig sammen. --- \emph{Vreden} staar
omtrent i samme Forhold til Glæden som Skrækken til Sorgen.

Af disse forskellige Symptomer anser Forf. nu \emph{Indsnæv}\emph{%
% --- p. 54 ---
\opage{54}ringen af Blodkarrene} for det vigtigste, det primære, hvis
fysiologiske Konsekvenser alle de andre ere. Vel kan denne
Antagelse endnu ikke fuldstændigt bevises, og Forf. henstiller
den som en Hypotese, der stiller fremtidig Forskning en
betydningsfuld Opgave. Men der er dog vigtige Kendsgerninger,
som foreløbigt tale for den. Navnligt véd man, at Hjerne,
Rygmarv og det hele Nervesystem i høj Grad paavirkes i
Henseende til Virkeevne af Forandringen i deres Blodrigdom.
Forstyrrelsen i Karsystemet vil ad den Vej kunne faa
Indflydelse paa alle de vigtigere Funktioner i Legemet. Individer,
hvis vasomotoriske System let sættes i Bevægelse, vise sig
ogsaa særlig tilbøjelige til Sindsbevægelser.

Naar det nu skal godtgøres, at vi i de her beskrevne
Fænomener have hele Sindsbevægelsens Natur for os, saa
lægger Forf. især Vægt paa, at \emph{de samme Tilstande, som man
i Regelen betegner som rent sjælelige Følelser, kunne fremkaldes
ved rent legemlige Paavirkninger}, naar disse blot paa en vis
bestemt Maade faa Indflydelse paa det vasomotoriske Apparat.
Vin gør jo glad, lige saa vel som godt Nyt; Vrede kan
bevirkes ved Nydelsen af Fluesvamp; Ipekakuanha frembringer
en deprimeret Stemning, der ligner Sorg eller Frygt; Angest
kan opstaa under Nervesygdomme o. s. v. Her hen hører
ogsaa saadanne Tilfælde, hvor Følelsesudbrud, f. Eks. Latter,
fremkaldes af Aarsager, som ikke synes at passe til denne
Virkning\footnote{Forf. har ikke set dette Fænomen omtalt af andre Forfattere, men
nævner to Tilfælde fra sin egen Praksis. Jeg har i min Psykologi
(2den Udg. p. 337, 4. Udg. p. 324), samlet en Del Eksempler herpaa.
--- Det maa maaske være mig tilladt med det Samme at bemærke,
at jeg har lagt det af Forf. her fremhævede Slægtskab mellem de
af legemlige og de af „sjælelige“ Aarsager fremkaldte Følelser til
Grund for min Fremstilling af Følelsernes Psykologi, idet jeg gaar
fra de simpleste Former til de mere sammensatte. Jeg følger f. Eks.
Frygten (Angesten) gennem dens Former, som blot Udtryk for en
organisk Tilstand (ved hindret Aandedræt eller Kardialgi), som
motiveret ved Sanseindtryk (instinktmæssig Sammenfaren ved pludselige
Indtryk) og som motiveret ved Forestilinger % PRINTED AS IS (for Forestillinger)
 (altsaa ved centrale
Processer i Hjærnen).}. Paa den anden Side kunne »sjælelige«
Sindsbevægelser lindres eller hæves ved legemlige Midler.

% --- p. 55 ---
\opage{55}Her har nu Forf. sikkert Ret i, at det er uberettiget,
at kalde de ad legemlig Vej fremkaldte Sindsbevægelser for
tilsyneladende, de ad sjælelig Vej fremkaldte for sande
Sindsbevægelser. Hvorfor skulde hine ikke være lige saa sande
som disse? Sandt er jo kun hvad der virkeligt \emph{er}, det vil her
sige \emph{føles}. En Glæde eller en Sorg kan i og for sig hverken
være sand eller falsk. En Følelse er tilstede eller er ikke
tilstede; de Forestillinger, Individet knytter dem til, kunne
være sande og falske, men paa selve Følelsen har denne
Forskel ingen Anvendelse. Følelsen kan have en mere eller
mindre solid Grundvold. Den Glæde og den sympatiske
Stemning, der vækkes ved Vin eller Opium, har ikke den
sunde og faste Grund som den, der skyldes et lykkeligt
Temperament eller selverhvervet Klarhed og Dygtighed eller
lykkelige, af Tilfældets og Øjeblikkets Gunst uafhængige
Livsvilkaar. Den ved Opium frembragte Lykkefølelse afløses
(naar den i længere Tid er bleven dyrket) siden af den
forfærdeligste Kval, medens den ad naturlig Vej frembragte
Lykkefølelse kan holde sig Livet igennem. Kun i dette
Forhold kan man hente Maalestokken til en Værdsættelse af
Sindsbevægelserne.

Men jeg tror alligevel, at Forf. overvurderer den
fysiologiske Betragtningsmaade. Først og fremmest kende vi dog
Sindsbevægelserne af \emph{direkte, subjektiv Iagttagelse}, og uden
denne vilde Forf. ikke have kunnet give sin Fremstilling.
Hans Mening maa egentligt være denne: det, vi hos os selv
føle, naar vi sige, at vi sørge, vilde for en ydre Iagttager,
der kunde se Alt, hvad der foregik i vore Nerver og Muskler,
tage sig ud som Spænding af nogle Muskler og Slappelse af
andre; Andet kan den ydre Iagttager i Følge Sagens Natur
ikke faa Øje paa. Men en saadan Iagttager slutter sig stedse
til Noget han ikke \emph{ser}, men \emph{forestiller} sig i \emph{Analogi} med visse
Erfaringer, han har gjort hos sig selv. Kun hvis han af \emph{egen,
subjektiv Erfaring} kender Sorg, kan han \emph{som klinisk
Iagttager} se Sorgen i de fysiologiske Fænomener, han har for sig.
Ellers vilde de for ham kun være ligegyldige Trækninger i
et ligegyldigt Stof. Hin Analogislutning gøres i Regelen let
og hurtigt, saa vi oftest ikke blive os den klart bevidst; men
% --- p. 56 ---
\opage{56}det er i teoretisk Henseende af afgørende Betydning at gøre
opmærksom paa den.

Man kan let tage fejl i denne Fortolkning af de
fysiologiske Fænomener. Skulde Forf. saaledes ikke tage fejl,
naar han af Overensstemmelsen mellem nogle vigtige
Symptomer slutter, at Vrede staar nærmere ved Glæde end ved
Sorg? De Fleste ville sikkert kalde Vreden en Ulystfølelse
og af den Grund sætte den nærmere ved Sorg end ved Glæde.
Det er kun for saa vidt Vrede er forbunden med indledende
Skridt til Hævn, at en Lystfølelse kommer til at spille en
væsentlig Rolle i Vreden. Vrede er Sorg, men Hævn er
Lyst, siger allerede \emph{Aristoteles}; Hævnen faar Vreden til at
høre op, idet den fremkalder Lyst i Stedet for Sorg. Det er
jo ogsaa klart, at Ingen vredes, uden der er tilføjet ham en
Ulyst eller en Sorg. Vreden er forskellig fra Sorgen ved, at
den er af aktiv, ikke passiv Karakter, og derved nærmer den
sig ganske vist til Glæden. --- Saa usikker den subjektive
Iagttagelse ogsaa kan være, saa er den dog en nødvendig
Forudsætning.

Hvad nu \emph{Forholdet} angaar \emph{mellem det, der træder frem
for den subjektive Iagttagelse, og det, den ydre Iagttager kan
se}, er jeg fuldstændigt beredt til at gaa ind paa Forfatterens
Opfattelse, der ikke tilsteder at henlægge Sindsbevægelsen
til Sjælen og dens Virkning til Legemet. Jeg betragter
enhver Appel til Sjælens Indgriben i den fysiologiske
Sammenhæng som uvidenskabelig. Det er almindelige teoretiske
Betragtninger, som have ført mig til at hævde \emph{den fysiologiske
Kontinuitet}. Dersom nemlig den Proces, som foregaar i
Nervesystem og Hjerne, paa nogle Punkter afbrødes for at
forplante sig over i Sjælen og senere atter genoptages ved
Sjælens Reaktion, saa vilde vi her have et Brud paa Loven
om den fysiske Energis Bestaaen, idet der først vilde gaa
fysisk Energi til Grunde og senere skabes fysisk Energi.
Derfor har jeg sluttet mig til den Hypotese, i Følge hvilken
Bevidsthedsfænomenerne, saaledes som subjektiv Iagttagelse
lærer os dem at kende, ere indre sjælelige Former for den
samme Virksomhed, som for den ydre Iagttagelse fremtræder
som materielt Fænomen. Paa denne Maade kommer baade
% --- p. 57 ---
\opage{57}Psykologien og Fysiologien til deres fulde Ret. Det er
egentligt en meget gammel Antagelse, som vi her hylde.
»Naturforskeren og Filosofen,« siger allerede Aristoteles, »ville give
forskellige Definitioner af ethvert [sjæleligt] Fænomen, f. Eks.
af Vreden. Den Ene vil nemlig sige, at den er en Attraa
efter at gengælde en ubehagelig Følelse eller noget Lignende;
den Anden vil sige: det er en Kogning i Blodet eller i det
Varme omkring Hjertet.« Aristoteles udtrykker her ved
Hjælp af sin Tids Fysiologi en Teori, der søger at lade baade
den fysiologiske og den psykologiske Side ved Fænomenet
komme til sin fulde Ret. I den nyere Tid har fremfor Alle
Spinoza hævdet denne Synsmaade, som i vore Dage har saa
mange Tilhængere. ---

Hævdelsen af \emph{den fysiologiske Kontinuitet} er en Lære,
der kan bestaa uafhængigt af den specielle Antagelse, som
Forf. hævder: at \emph{det Væsentlige ved Sindsbevægelsen netop er
de organiske Virkninger,} i hvilke den sædvanlige Opfattelse
kun ser Bifænomener, --- at \emph{Sindsbevægelse altsaa egentligt kun
er Fornemmelse af de organiske Virkninger}. En amerikansk
Psykolog, \emph{William James}, har i Tidsskriftet »Mind« (April
1884) fremsat en lignende Teori. For at den skulde godtgøre
sin Sandhed, maatte det kunne vises, at Sindsbevægelsen ikke
opstaar, \emph{før} de organiske Virkninger fornemmes. Dette vilde
stemme godt med, at Lyst- og Ulystfølelse i det Hele synes
at behøve længere Tid til sin Optræden end Fornemmelser
og Forestillinger. Men Betingelsen for, at hine organiske
Virkninger kunne indtræde, er, at der foregaar en Udstraaling
i Hjernen fra de Centra (for Sanseiagttagelse eller Forestilling),
som vække Sindsbevægelsen. Kun naar denne Udstraaling
naar til visse andre Centra (i Følge Forf. til de vasomotoriske),
opstaar Sindsbevægelse. Mon der nu ikke i Bevidstheden
skulde gøre sig Noget gældende, som svarer til denne
Udstraaling? Og skulde det ikke være berettiget at skelne
mellem to Stadier inden for Sindsbevægelsens Forløb, saaledes
at det ene svarer til, hvad der foregaar i Hjernen under hin
Udstraalingsproces, det andet til den Maade, hvorpaa denne
Hjerneproces forstærkes og modificeres ved de organiske
Virkningers Indtryk paa Hjernen? Her vilde da foregaa
% --- p. 58 ---
\opage{58}et Kredsløb: fra Hjernen ud i Legemet og tilbage til Hjernen.
I sin fulde Ejendommelighed have vi først Sindsbevægelsen,
naar denne Tilbagevirkning er fuldbyrdet. Men det kunde
dog være, at en Lyst- eller Ulystfølelse opstod allerede før
Virkningen forplanter sig ud over Hjernens Omraade, selv
om Tidsforskellen kun er meget ubetydelig. Omvendt kunne
undertiden efter en stærk Sindsbevægelse de organiske
Virkninger endnu vare ved, skønt Sindet ellers er i Ro. Her
vedvarer altsaa det andet Stadium, medens det første er
afsluttet. Saaledes f. Eks. naar vi »ikke kunne komme os«
efter en Forskrækkelse eller efter et Latterudbrud.

Til Ideen om den \emph{fysiologiske Kontinuitet} og til Teorien
om \emph{de organiske Reflekser som det fysiologiske Hovedfænomen
ved Sindsbevægelsen} kommer endnu hos Forf. hans specielle
Teori om \emph{den vasomotoriske Refleks som den primære}, af
hvilken alle andre fysiologiske Virkninger hidføres og
betinges. Den samme Teori er nyligt kommen frem fra anden
Side, nemlig i et Skrift om Frygtens Fysiologi af Italieneren
\emph{Mosso} (La peur. Paris 1886). En saadan Antagelse ligger
vor Tids Fysiologi nær, dels fordi saa mange Kendsgerninger
godtgøre Karinnervationens Indflydelse paa de øvrige Nervers
Funktioner, medens en omvendt Indflydelse ikke kan paavises,
dels fordi der i de senere Aar er paavist Centralorganer i
Hjernen for Karmusklernes Nerver, ikke blot i den forlængede
Marv, men efter Nogles Mening i selve den store Hjernes Bark,
altsaa mellem de højeste Nervecentra. Dersom denne
Antagelse vinder fortsat Bekræftelse, vil Sindsbevægelsernes
Fysiologi afrunde sig paa en smuk og overskuelig Maade.
Om Udsigterne hertil kunne naturligvis kun Fysiologer af
Faget have nogen begrundet Mening, og jeg skal derfor ikke
gaa nærmere ind herpaa. Kun vil jeg bemærke, at selv om
det maatte have nok saa stor Betydning for Sindsbevægelsens
fysiologiske Symptomer, at Bevægelsen inden for Hjernen naar
hen til det vasomotoriske Centrum, saa er derved ikke
udelukket, at der sker en ligesaa umiddelbar Paavirkning af
de i den forlængede Marv beliggende Centra for Nervus
vagus, hvorved den stærke Bevægelse i Hjernen straks faar
Indflydelse paa Hjertets Funktioner og derved igen paa Blodets
% --- p. 59 ---
\opage{59}Cirkulation, først og fremmest paa Blodtilførselen til selve
Hjernen. \emph{Claude Bernard} lagde Hovedvægten paa denne
Omstændighed i sin Afhandling om Hjertets Fysiologi (paany
optrykt i den under Titel »La Science Expérimentale«
udgivne Samling af hans Afhandlinger). Ogsaa ad denne Vej
vilde det være Forandring i Blodtilførselen til de forskellige
Organer, som blev det vigtigste fysiologiske Fænomen ved
Sindsbevægelsen. Jeg har i min Psykologi lagt Hovedvægten
paa det af Bernard fremhævede Punkt, men har dog ogsaa
fremhævet de vasomotoriske Centras Betydning, og jeg vil,
uden at behøve at forandre Noget i min psykologiske Teori,
kunne gøre Brug af Forf. smukke Fremstilling, dersom den
vinder fortsat Bekræftelse. ---

\emph{Sindsbevægelsernes Fremtid} ser det efter Forf.s Anskuelse
trist ud med. Baade Slægtens og Individernes Opdragelse
og Dannelse gaar jo ud paa at trænge Sindsbevægelserne
tilbage, at hæmme de organiske Reflekser, som de ere knyttede
til. Vil dette da ikke ved Nedarving føre til, at
Karinnervationen bliver saa træg, at der tilsidst ikke mere kan være
Tale om Sindsbevægelser? Dertil kommer, at Forstandslivet
spiller en stedse større Rolle under Kulturens Udvikling.
Forstandslivet og Følelseslivet staa i Modsætning til hinanden.
De tære begge paa Hjernens Blodforsyning, og jo mere
den ene Art af Hjernevirksomhed forbruger, des mindre
bliver der tilovers til den anden. Vil da ikke tilsidst
Forstandslivet lægge Beslag paa det Hele, saa at Følelsernes
Tid er forbi? --- Med dette pessimistiske Resultat slutter
Forf. sin Bog.

Jeg kan ikke tro Andet, end at han kun mener dette
som en Advarsel, --- en Advarsel, vi ogsaa højligt kunne
behøve. Det er ganske vist en stor Ensidighed ved hele vor
moderne Kultur, at den overvejende har udviklet
Forstandsevnerne. Det er dens Styrke, men ogsaa dens Mangel.
Følgen er dog, tror jeg, ikke saa meget bleven en Svækkelse
af Følelseslivet i det Hele, som en Disharmoni mellem vor
Erkendelse og vor Følelse. Vore Tanker omspænde langt Mere
end vore Følelser endnu kunne naa at fylde og varme. Vor
Viden gaar altfor ofte i én Retning, vor Tro i en anden.
% --- p. 60 ---
\opage{60}Deraf de vaklende og delte Karakterer. Hele Livet er ikke
længere, som fordum, med al sin Energi rettet mod ét Maal,
der formaaede at tilfredsstille alle sjælelige Fornødenheder.
For at Følelseslivet skal være energisk nok, maa Sympatien
være ganske anderledes levende og fyldig end den endnu
har vist sig. Det Hang til Individualisme og Egoisme, som
den moderne Forstandsudvikling ikke kan siges fri for at
have begunstiget, er det, der indsnævrer Følelsen. Den rette
Opdragelse vil gaa ud paa at fortrænge de egoistiske
Sindsbevægelser, men derimod søge at begunstige de sympatiske
og lede dem ad de gavnligste Veje. Det kan kun være en
stor Misforstaaelse, dersom Slægten vilde mene, at den nu var
kommen saa vidt, at den kunde undvære Sindsbevægelserne.
Disse ville stedse vise sig som de levende Kræfter, der drive
Værket og overvinde Modstanden.

Den indbyrdes Vekselvirkning mellem Forstandsudvikling
og Følelsesudvikling vil sikkert mere og mere forandre
Sindsbevægelsernes Karakter. Deres elementære Heftighed og
Voldsomhed vil afløses af mildere Former. Det er til Dels
sket, og mon det altid er et Tab? Er det et Tab, at
Menneskene nu kende Elskov i Steden for Dyrets brutale Instinkt,
eller at Hunger og Tørst efter Retfærdighed kan afløse det
blinde, hensynsløse Hævninstinkt? Det er ganske vist ikke
altid, at Vindingen er saa let at paavise, og en Advarsel
kommer derfor altid med en vis Berettigelse. Vi maa i hvert
Tilfælde erindre, at jo flere og jo mere omfattende Genstande
der lægge Beslag paa vor Følelse, des større Energi maa
der kunne anvendes, dersom ikke Livet skal svækkes. Jo
flere forskellige Udgifter vi have, des større samlet Kapital
maa vi raade over, hvis vi ikke skulle komme til kort. ---

Den psykologiske Forskens Fremgang sker ved Bidrag
fra saare mange forskellige Sider. Forf. har fra sit
Standpunkt givet et betydningsfuldt Indlæg, betydningsfuldt ikke
blot ved det Resultat, han deri er kommen til, men ogsaa
ved de Tanker, det i forskellig Retning sætter i Bevægelse.
Og det maa endnu føjes til, at hans Bog ogsaa ved sin
tiltalende Form og sit smukke Sprog fortjener at vinde sig
mange Læsere.

\streg{}

% --- p. 61 ---
\opage{61}
\essay[FREDERIK CHRISTIAN SIBBERN]{FREDERIK CHRISTIAN SIBBERN.}

\begin{center}1785---1872.\end{center}

\streg{}

Der gives fremragende Mænd, hvem det synes at være
lykkedes fuldstændigt at give sig selv i deres Værker. Sibbern
hører ikke til disses Tal. Hvad han selv ofte klagede over,
sande ogsaa Læserne af hans videnskabelige Skrifter, at han
havde Vanskelighed ved at forme og udtrykke hvad der stod
for hans indre Øje. Dette hang vel til Dels sammen med
en Ufuldkommenhed ved hans videnskabelige Standpunkt,
som han aldrig, trods sit stadigt fornyede Arbejde, naaede
ud over; men det hang ogsaa sammen med, at hans
Personlighed var saa rig, varm og ejendommelig, ogsaa
ejendommeligt kantet, at den ikke kunde faa Luft gennem en enkelt
Form. Tænkeren, Digteren, den religiøst Grebne, den
hjertelige Menneskeven og den begejstrede sociale Profet vare
forskellige Jeger i ham og krævede hver sin Art Tilfredsstillelse.
Han arbejdede sit Liv igennem med disse forskellige Kræfter,
af hvilke snart én, snart en anden skød sig frem i
Forgrunden; Enheden, Sammenhængen mellem disse forskellige
Bestræbelser har han overladt til Eftertiden at gribe, om
den formaar det. Og der maa gøres et Forsøg derpaa, medens
Mindet om ham endnu lever iblandt os.

Han hører ikke til dem, der stærkt og gennemgribende
have paavirket det aandelige og det offentlige Livs Gang i
vort Fædreland. Men han er et af de sundeste og tillige
dybeste Mennesker, som har levet her hjemme. Og han
gennemlevede sin Tid fuldt og selvstændigt, optog ikke blot
i sig hvad den frembød, Stort og Smaat, men tog det
Altsammen paa sin ejendommelige Maade, gav sit Besyv med,
% --- p. 62 ---
\opage{62}privat eller offentligt, og ved sin varme Sympati og
indtrængende Forstaaelse har han været overordentligt Meget
for saa mange af de Bedste her hjemme, som han da selv
maa regnes til vort Lands bedste Mænd. I ham spejler hele
det aandelige og offentlige Liv i Danmark fra 1820 til 1870
sig paa en lærerig Maade. Lærerig --- ikke blot fordi han
levede det med under varm Deltagelse, men ogsaa fordi han
selv stedse var en Stræbende, der lige til sin høje Alderdom
gik frem i Søgen efter Sandhed. Hans Ungdom og første
Udvikling faldt i Reaktionsperioden mod det attende
Aarhundrede. Han følte ret den kvægende Virkning af den
Følelsens og Fantasiens Fylde, der afløste den
forstandsmæssige »Oplysning« og den tørre Fornuftmorals Herredømme.
Men han tog kun hvad han kunde bruge, og han kunde kun
bruge det, som havde sand menneskelig Kærne i sig; Skallerne
lod han ligge. Det blev derfor ikke noget Tilbageslag, men
Resultatet af en naturlig og sund Udviklingsgang, naar han
senere udtalte religiøse og politiske Anskuelser, som ved
deres Frihed og Radikalisme forundrede og forargede mange
af hans tidligere Venner, og vilde have forundret endnu
Flere, hvis de vare blevne mere bekendte i større Kredse.
Det var hans levende Sans for det Levende, der fik ham til
at gennembryde enhver Form, Tidsaanden stillede op som
ene saliggørende, hvad enten det var paa Religionens, paa
Politikens eller paa Filosofiens Omraade. Der er en indre
Forbindelse mellem hans Polemik mod den hegelske Filosofi,
hans Løsrivelse fra al dogmatisk Tro og hans Bekæmpelse
af Overtroen paa Papirskonstitutioner og juridiske
Formaliteter.

Ved sine afsluttende religiøse og sociale Ideer mødes han
med friere Bestræbelser paa disse Omraader i den nyeste
Tid. Ved sin levende Sans for det virkelige Liv, dets
Vilkaar og dets Krav, mødes han med det Bedste i den moderne
realistiske Aandsretning. Men Alt er hos ham vokset ud af
det dybe Sinds Trang til at opsøge og fordybe sig i Livet
under alle dets Former, paa alle dets Trin. Man kan neppe
paavise ydre Paavirkninger, som skulde have fremkaldt og
bestemt hans Udvikling i den senere Del af hans Liv. Han
% --- p. 63 ---
\opage{63}blev forberedt til Revolutionen netop ved fuldt og helt at
gennemleve Reaktionen. Og han \emph{voksede} sig fra Reaktionen
ind i Revolutionen uden mærkeligt Brud og uden pludselig
Katastrofe. Han, som selv stedse advarede mod at plukke
den umodne Frugt af Livets Træ, lod ogsaa sin egen Frugt
have den fornødne Tid. Des værre blev det i Mellemtiden
Aften for ham selv; det, der var hans Livs og Stræbens
sundt udviklede Resultat, formaaede han nu ikke at vække
Opmærksomhed og Anerkendelse for. Det kom, saa længe
han selv levede, mest til at staa som en gammel Mands sære
Ideer, og hans gamle Venner og Beundrere lod helst, som
om de ikke vidste af dem at sige. Maaske ønskede han det
heller ikke selv anderledes. Han var den stille Udviklings
Mand og Fjende af al Forhastelse. Des mere opmuntrende
og lærerigt er hans Billede for den yngre Slægt, som har
den Opgave at føre Traaden videre, hvor hans gamle Haand
maatte slippe den. Det opfordrer dem baade til at gaa
fremad og til at lægge en dyb og fast Grundvold for
Fremskridtsarbejdet. Naar vi i vore Dage se saa mange
uhyggelige Eksempler paa, at gamle Frihedsmænd --- ofte uden selv
at ville være ved det --- fornægte deres Ungdoms Idealer,
saa maa det komme af, at Opstillingen af disse Idealer har
været Hastværksarbejde uden solidt Grundlag. Det er
Arbejdet, det stille, ihærdige og taalmodige Arbejde, først og
fremmest Arbejdet med sig selv, der har manglet. Men der
har ogsaa manglet den levende Sans for, og dermed den
levende Tro paa Livet som det, der aldrig staar stille. Den
Form, hvorunder man først har tænkt sig Livet, forlanger
man, at det skal iføre sig alle Tider. Den ældre Slægt
kommer derved til at tynge som en Mare paa den yngre
Slægt, det ene sociale Lag paa det andet, i Stedet for at
bane Vejen for sin Efterfølger og hjælpe ham frem. Men
Livet har stedse Ret; det var netop en af Sibberns
Hovedsætninger.

Naar jeg nu her vil afbetale en Taknemmelighedsgæld
ved at forsøge paa en Karakteristik af F. C. Sibbern, saa
maa Begyndelsen ske med at skildre hans Personlighed, saa
vidt det kan lykkes at samle de spredte Træk om visse
% --- p. 64 ---
\opage{64}Hovedpunkter. Vi ville derved se, hvorledes hans Natur
forudbestemte ham til hans Gerning. Hvad han har virket
som Psykolog, som Filosof, som religiøs og social Tænker,
vil kun ret forstaas, naar vi først --- for at bruge et af hans
egne Yndlingsudtryk --- opsøge »den indre Kilde«, af hvilken
det Altsammen er udvældet.

\section*{I.\\  Personlig Karakteristik.}

Vi lære Sibbern at kende gennem hans Breve, baade de
virkelige og dem, han har udgivet under Navn af Gabrielis’s
Breve. Det Billede af ham, som herved kan vindes, bekræftes
ved hans andre Skrifter, men indeholder tillige adskillige
Træk, som tjene til bedre Forstaaelse af disse.

Et Vendepunkt i hans Ungdom var det langvarige
Ophold i Tyskland fra 1811 til 1814, fra hvilket han vendte
hjem for at overtage en Professorpost i Filosofi ved
Universitetet. Vi følge ham paa Rejsen gennem hans Breve, og vi
lære hans Stemning og Tilstand i den nærmeste Tid efter
Hjemkomsten at kende gennem første Del af Gabrielis’s
Breve, som er skreven 1813---1814, skønt den først blev
udgiven 1826.

Han var rejst ud for at studere Filosofi og uddanne sig
til akademisk Lærer. Han gjorde Fichtes, Schleiermachers,
Steffens’ og Schellings Bekendtskab og paavirkedes især stærkt
og varigt af de to sidst Nævnte. Han opsøgte Goethe, der
for ham stod som en ophøjet Genius og var Genstand for
hans Kultus, ja næsten for hans Forelskelse. Han foretog
Fodvandringer gennem Tysklands skønneste Egne, levede
med Bjerge, Skove og Søer og kvægede sig ved Naturens
Skønhed. Men under alt det Nye og Herlige, han saa og
lærte at kende, hang hans Hjerte ved Hjemmet. Han hang
ved den danske Natur, ved sin Fædrenestad, hvis store
Ulykker han havde oplevet, og hvor han haabede, at en ny
og betydningsfuld Periode for Videnskab og Kunst var
indledet. Men der var dog Et, der bandt ham fastere end
noget Andet til Hjemmet, nemlig hans Kærlighed til Sofie
% --- p. 65 ---
\opage{65}Ørsted, til hvem de fleste og interessanteste af hans
Rejsebreve ere rettede. »Ingen skulde nogen Sinde rejse«, skriver
han til hende, »for end Fædrelandet havde bundet ham med
de fasteste, stærkeste Baand«. --- Paa Galleriet i Berlin traf
han en ung Pige, med hvem han ofte talte, og hvis Billede
han bevarede og ofte dvælede ved i sine Tanker. »Men at
jeg kan gøre dette \emph{med Ro}, derfor takker jeg mit Fædreland,
takker det inderligt derfor; det lod mig tidligt skue og fatte
høj Skønhed, og det skænkede mig derved den lykkelige
Gave: overalt at kunne frydes ved hver Skønhed, jeg møder
paa min Vej, fryde mig uforstyrret af Attraa og Begær,«

Han kom fra ung af meget i A. S. Ørsteds Hus. For
den store Retslærde nærede han en overordentlig Beundring
og en varm Hengivenhed, som han ofte har givet et smukt
og hjerteligt Udtryk, smukkest maaske i sine »politiske
Intelligensblade«, i Anledning af Ørsteds Tale ved Aabningen
af de raadgivende Stænders Forsamling i Roskilde. Han
var selv Ørsteds Discipel, idet hans første Studier vare
juridiske. Ørsteds Hustru, Oehlenschlägers smukke og begavede
Søster, samlede ikke mindre end sin Mand en stor Kreds
om sig. Sibbern elskede hende med hele sin Naturs Varme
og Entusiasme. Endnu paa sine gamle Dage bevægedes han
i en Samtale med Sofie Ørsteds Biograf, Nikolaj Bøgh, ved
Mindet om sin Ungdoms Kærlighed. Han spadserede med
hende og læste for hende i Odysseen, Sofokles, Plato og
Goethe. Han var idel Glæde og Beundring over hendes
Skønhed og over hendes livlige Interesse og klare Blik.
Han var rejst ud i jublende Glæde over hvad han her havde
fundet. Det tilhørte ham ikke, men Glæden derover
tilhørte ham.

Allerede af Rejsebrevene ser man dog, at det ikke faldt
ham saa ganske let blot at tage Forholdet fra denne Side.
Sofie Ørsted elskede ham ikke; hun behandlede ham som en yngre
Broder; og desuden var hun knyttet til en Mand, hun agtede
og elskede saa højt som nogen Anden. Fortvivlelsen over
det Haabløse i hans Kærlighed maatte bryde frem, og det
var først, da Sofie Ørsted var død, at han kunde udtale, at
»det var sødt og glædeligt blot at vide, at hun var til, og
% --- p. 66 ---
\opage{66}at være kendt af hende«. Der var Tider, hvor hans Hjerte
forlangte Mere end det. I et af sine sidste Rejsebreve til
hende omtaler han, hvorledes han i Breslau, hvor han
besøgte Steffens, havde følt sig overvældet af sin Smerte og
havde fundet Lindring i at græde af ganske Hjerte. »Og
virkeligt er det«, tilføjer han, »som om jeg dér, hos min
kære, dyrebare Landsmand, først har fundet den Fred og
Hvile, som Orest fordum fandt i Gudindens Lund, i den rene
Søsters Nærhed«. --- Der synes her at være hentydet til en
Kamp mellem den opblussende Attraa og den glade, af alt
Begær uhildede Følelse ved hvad han havde i den, han
elskede. Han har kæmpet sig frem til den Beslutning, som
han lader sin Gabrielis udtale: »Er der Noget, jeg med
Bestemthed vil, saa er det det, ikke at se hende igen, førend
jeg véd, at jeg er rolig og uden Begær vil kunne glæde mig
ved Beskuelsen«.

»Den Time, jeg vender hjem,« havde han skrevet, »den
vil bringe mig alle Glæders Blomst«. Men hvorledes det
blev for ham, lader han sit andet Jeg, Gabrielis, skildre den
Ven, til hvem han skriver. Man vil deri genfinde nogle
Vendinger, som allerede ere anførte fra et af de virkelige
Breve til Sofie Ørsted.

»Du véd, hvor længe jeg var borte, men Du véd ikke,
hvor tro jeg hængte ved hende. Alle mine Længsler
tilhørte hende og mit Hjem, og tusinde Gange sagde jeg til
mig selv: Lykkelige! ja du kan rejse; ti dig holder
Fædrelandet fast ved de skønneste Baand, og hvad der saa møder
dig, det Herligste, Jorden har at give, gav Hjemmet dig. ---
Ak! og ogsaa det Bitreste. Jeg kom tilbage; mine Følelser
luede endnu i mig med ufortæret Inderlighed og Styrke.
Jeg fandt ogsaa hos hende den vante Godhed for mig.
Uudsigeligt takkede jeg hende for, at hun modtog mig saa
god, saa glad. Jeg levede atter nogle glædefulde Dage. Men
det blev snart mattere. Nej, jeg tabte, hvad jeg nu først
haabede at tilvinde mig. --- Med hver Dag blev jeg fattigere.
I en endeløs Lidenskabs, jeg kunde sige Skinsyges aldrig
hvilende Ængstelser gik mine Følelser op, og hvad jeg
endnu besad, derved skilte mig snart min forvildede Adfærd.«

% --- p. 67 ---
\opage{67}Det var den første store Kamp, han maatte tage op med
Livet. Han gennemlevede dens hele Bitterhed, skønt det
ofte forekom ham, som hans Liv var forspildt. Det var en
tung Vinter for ham, den første efter hans Hjemkomst, den
i hvilken »Gabrielis« blev skreven. Han skriver selv til sin
Søster den 5te Marts 1814: »At jeg skal tilstaa dig det,
endnu har jeg ikke tilbragt nogen Vinter med saa nedstemte,
ofte bitre Følelser som denne«. Der kom foruden den
Aarsag, hvorved vi her have dvælet, og som han ikke omtaler
over for sin Søster, endnu to andre Aarsager til. --- Han havde
overtaget den filosofiske Lærerpost, men følte, at det ikke
gik for ham i hans nye Stilling, som han ønskede det. Han
kæmpede, som vi allerede se af Rejsebrevene, med Tvivl om
sig selv. Han frygtede, at der ikke ret vilde blive Noget af
ham, følte sig mat, aandsforladt; Alt inden i ham syntes at
gaa i Staa. I den Forladtheds- og Mismodsperiode, som var
kommen over ham, maatte denne Tvivl og Afmagtsfølelse
vokse, og maatte tillige føles saa meget des stærkere, jo mere
han trængte til at fordybe sig i Virksomhed for at glemme
sin Smerte, og jo større Fordringer hans nye Stilling rettede
til ham. --- Endeligt var det jo ogsaa denne Vinter, hvor
den ulykkelige Syvaarskrig afsluttedes med den lige saa
ulykkelige Kielerfred. Han havde paa nært Hold oplevet
Frihedskrigens Udbrud i Tyskland og Glæden over Sejrene;
Usselheden og Ulykken her hjemme traadte i Modsætning
dertil i et saa meget des grellere Lys. »Medens hele det
øvrige Europa foruden Frankrig kan udbryde i høj Jubel,
skulle vi alene lide paa en saa skammelig Maade!«

Saadanne vare de Erfaringer, gennem hvilke Sibberns
Karakter først træder frem for os. Naar der nu i hans
Breve, i hans Skrifter og i saa mange mundtlige Ytringer
af ham, saa vel som i hans hele Færd, fremtræder en saa
frejdig Glæde ved Livet, en urokkelig Optimisme, der suger
Lyksalighed af Alt, saa er det allerede af det, vi nu have
omtalt, klart, at denne Livsglæde og Optimisme ikke var
bygget paa Sand, men vandtes gennem Kamp med megen
og dybt følt Modstand. Han lader sin Gabrielis trøste sig
med, at der paa det Hav af Taarer, der trængte til Gennem%
% --- p. 68 ---
\opage{68}brud hos ham, en Gang vilde følge et Hav af Lys. Saaledes
er der altid for Sibbern en dyb og stærk Skygge, der giver
Lyset dets fulde Virkning. Den gode Tid jublede han over,
fordi han kendte den onde, og især, naar han havde holdt
sig sin Lærer Goethes Formaning efterrettelig: ikke at hidse
sig i den onde Tid.

I den Krise, som er omtalt, se vi ham særligt kæmpe
for at bevare sin Glæde og Beundring selv dèr, hvor
Besiddelsen er umulig. Beundring og Sympati vare i det Hele
de Følelser, som havde den dybeste Grund i Sibberns Natur.
Livet var for ham, naar han blot kunde neddæmpe den
Bitterhed, der kunde stige op i hans eget Indre, væsentligt en
Kilde til Glæde. I sit 24. Aar skriver han til sin Søster:
»Dan dig af det, der omgiver dig, en Himmel fuld af Glæde!«
og endnu i sin høje Alderdom følte han (som han skriver i
et mig velvilligt overladt Brev fra Aaret før hans Død)
»Sommerlighed« hos sig. Ydre og indre Tryk og Lidelser
kunde formørke og forbitre Sindet; og Tryk var efter hans
Erfaring værre end Lidelser; --- men selv i slige Tider kan
man, paastod han, skaffe sig Glimt af noget Lysere eller en
»Forudfølelse af Musens Duft.« Den værste Fjende for
Livspoesien, for Bevarelsen af Ungdommen, fandt han i Begæret,
som er knyttet til noget bestemt Enkelt, og som afføder
Bitterhed og Smerte, naar dette ikke kan opnaas. »Nej, tag
til det, som overalt er at have. Der er underfuld Livsmusik
nok overalt; man skal kun vorde stille i sig selv og lytte
efter.« Saaledes lader han Gabrielis (i Brevenes anden Del)
besvare det Spørgsmaal, hvorledes Ungdommen og Livspoesien
bevares. Og nu følger hele den vidunderlige smukke
Udvikling, som Enhver bør kende fra første Haand, om hvorledes
man kan finde Glæde ved Beskuelsen af Solglimtet paa
Taget, af Spurvene, som hoppe omkring, af Scenerne paa
Gaden, af Menneskelivet i dets forskellige Former. Man
skal leve med »Naturens hele store Livsbevægelse«, »slutte et
fortroligt Bekendtskab med Vinteren som med Sommeren, saa
at man kan føle i sig de Melodier, som de spille op for os,
naar de komme til os i Besøg«. Vi forstaa her Sibberns store
Glæde ved Fodture. Som han havde travet rundt i Tyskland
% --- p. 69 ---
\opage{69}paa sin store Rejse, travede han i senere Aar omkring i
Sjælland og paa Fyen, og i Gabrielis’ Breve findes adskillige
Træk af hvad man kunde kalde Fodturenes Epikuræisme,
den Vuggen i Stemninger og Tanker, som Ensomheden i
Mark og Skov begunstiger, og den hyggelige Hvile ved
Aftenstid. Det var ikke blot Naturen, han svælgede i, og
hvis smaa og store Scener han belurede og ofte skildrer med
stor Friskhed. Det var, som allerede berørt, ogsaa
Menneskelivet. Han fandt det smukt, naar Lavater havde udgivet sit
store Værk over Menneskefysiognomierne »til Befordring af
Menneskekærligheden«. Man skal elske Menneskene, som
Botanikeren elsker Planter: for ham er der intet Ukrudt.
Man skal sætte sig hos Haandværkeren, se paa hans Dont
og tale med ham derom. »Mon der ikke i Haandværkerens
Arbejde, i Kræmmerens Travlhed og i al menneskelig Færd
paa Torve og i Gader ere store Livsmagter til Stede?« I
sine senere Aar nævnte han især Husmandsfolk og
Landboere med Forkærlighed. Han, som levede i vort Folks
finest og ædlest dannede Kredse og selv havde umiddelbar
Adgang til Videnskabens og Poesiens Skatte, ender sin (efter
hans Død udkomne) Moralfilosofi med de Ord: »Jeg kunde
anføre mere end ét Eksempel af Husmandsstanden paa
Landet, ligeledes af Tjenestefolk, med hvilke Samlivet kan virke
højst velgørende paa den Agtsomme«. Ogsaa i andre Skrifter
kommer han ofte tilbage til Husmandsstanden paa en Maade,
som viser, at han taler af personligt Bekendtskab.

Man ser: den Livspoesi, han prædiker og vil have
Øjnene oplukkede for, bæres af en hjertelig og varm Sympati.
Det er intet Fantasilefleri som i saa mange Former af Romantik,
og (maa jeg føje til) i saa mange Former af »Realisme«, der
i denne Henseende kun ere forklædt Romantik. Han søgte
Næring, ikke Underholdning. Livet i alle dets Brydninger
har Værdi for ham, og det for dets egen Skyld. Det var
hans Tro, at det samme Liv, hvis mindste Rørelser uden for
ham kunde have saa stort Værd i hans Øjne, ogsaa rørte
sig i ham selv, --- og at det var derfor han kunde faa Øje
for det. Livspoesiens Evangelium hænger her, som det senere
vil vise sig, nøje sammen med hans filosofiske Anskuelse og
% --- p. 70 ---
\opage{70}hans religiøse Tro, som han da i det Hele ikke gør den
store Forskel mellem disse tre Ting, som vi stakkels moderne
Mennesker ere nødte til.

Det var ikke blot Natur og Mennesker, han fandt Næring
og Vederkvægelse ved. Han levede ogsaa i Digterne, i
Bibelen, i gamle, kærnefulde Opbyggelsesskrifter, som han
endnu satte Pris paa og anbefalede, efter at Dogmetroen
forlængst var bleven ham aldeles fremmed. Yndlingssteder fra
saadanne Forfattere havde han ofte hos sig i udtagne Ark
eller opskrevne paa Sedler til Brug paa Fodture for ret at
kunne gøre sig dem levende og nærværende. Endnu Aaret
før sin Død opskrev han en Del saadanne Yndlingssteder,
især af Schiller og Goethe, tillige med nogle af sine egne
Yndlingstanker, til en ung Pige, der skulde rejse ud som
Lærerinde.

Som flere Gange sagt, den Livsglæde, som var Sibbern
saa ejendommelig, var ikke overfladisk og letsindig; den var
ikke blot ofte dyrt købt, men den var ham en alvorlig Sag,
en Hjertes- og Viljessag. Han havde jo, som vi have set,
Tider, hvor selve Smerten var det Eneste, der holdt Sindet
sammen og hindrede, at det sprængtes. »Intet glædeligt«,
siger Gabrielis, »havde jeg, der kunde give min Eksistens
Kontinuitet; saa blev da Smerte og Fortvivlelse det, der
holdt min Tilværelse sammen og bevarede mig for at give
min Sjæl hen i Søndersplittelse«. Der ligger en sund og
mandig Tanke i disse Ord. Den aandelige Smerte har vel, som
al Smerte, en Tendens til at opløse vor Sjæl. Men den
er selv, hvor den melder sig, et naturnødvendigt Led i
Rækken af vore indre Erfaringer, og at skyde den ud, at
ville fortrænge den for enhver Pris, det er at hengive sig
til en aandelig Morfinisme og er den visse Vej til
Undergang. Smerten maa gennemleves, dersom man ikke kan
forandre de Forhold, som fremkalde den. Kun derved vokser
og hærdes Sindet. Det var det, Sibbern atter og atter, lige
til det Sidste, indskærpede, at »hvad der hæmmer, det
fremmer, naar det tages retteligen«. Han kendte af egen Erfaring,
hvad Goethe kalder Smerternes »hemmeligt dannende Magt«.
Man maa kæmpe sig frem med Udholdenhed og Taalmodig%
% --- p. 71 ---
\opage{71}hed fra den onde Tid til den gode; først saa virker
Modsætningen mellem Lys og Mørke paa rette Maade. Sibberns
Livsanskuelse var væsentligt en Tro paa Livet, og den sidste
Tanke hos ham var vel den, at man maatte lade Livet have
dets Gang og ogsaa hos sig selv vente paa, at de gode Tider
kom. De onde Tider, især Aandsforladthedens Tider, som
han saa ofte klagede over, betragtede han som Tider, i hvilke
Sindet laa Brak og kunde samle Kraft til ny Virksomhed.
Men han hyldede ingenlunde en Kvietisme, der førte til at
lade Hænderne ligge i Skødet. Hvor energisk hans
Opfattelse tvært imod var, kan ses af, at han ofte opponerede
mod det Skriftsted, som siger, at man ikke kan lægge en
Alen til sin Vækst. »Men en Tomme«, skriver han endnu
Aaret før sin Død, »kan man dog nok føje til den, naar man
holder sig rask og rank.« Han vilde ikke have Grænserne
skarpt afstukne for, hvor langt man kan komme. Og skønt det
Højeste for ham var Hengivelse til Livet i os og uden for
os, saa opfattede han ikke denne Hengivelse som noget blot
Passivt, men som Betingelsen for, at Livets Fylde kan blive
en Magt i os. Den maa gaa i Et med vor egen Villen.

Dertil kom, at der maaske nok var noget Sentimentalt i
hans Natur, men at han ikke hørte til de Blide og Sagtmodige.
Han var et afgjort kolerisk Temperament. En psykologisk
Karakteristik af ham vilde være meget mangelfuld, naar den
ikke fremhævede dette, baade fordi det præger Mandens hele
Fysiognomi, og fordi vi her have en anden hæmmende og
mørk Magt at kæmpe med. Det er bekendt af flere
Anekdoter, hvor han kunde fare op i Heftighed. Han skal en
Gang, medens Martensen var Professor, have overfuset ham
paa Universitetet i Studenternes Paahør, da de vare i Strid
om et Auditorium; og dog var Martensen hans gode Ven.
Vi have forskellige Selvbekendelser af ham i den Retning.
--- Et af de smukkeste Breve i den udgivne Samling er det
til Lektor C. F. Molbech i Sorø, som var bleven greben af
en Sindssygdom, der havde forandret hans forhen kærlige og
sagtmodige Karakter og gjort ham bitter og ustyrlig. Da
den værste Tid syntes at være ovre, men den Lidende endnu
havde at kæmpe med Efterdønningerne, skrev Sibbern til
% --- p. 72 ---
\opage{72}ham for at hjælpe og styrke ham. »Just jeg vil skrive til
Dem, fordi ogsaa jeg véd og af mit eget Indre kender, hvad
alt der bor i vor Sjæls Afgrund og kan stige op af den og
formørke vort Sind. Et kender jeg især, som især er ondt,
fordi det er mest ødelæggende for dem, som ere os de
Nærmeste .... Dette ene er den bitre Vredagtighed, det mørke
Misnøje og Mismod. Spørg mig ikke, hvad der dog kan
ægge den frem og nære den hos mig. Det Mindste er
undertiden nok dertil, naar Stemningen er derefter. Men Grunden
er at søge tidligt i mit Liv. Allerede fra min Moder har
jeg arvet den; og gik ikke vor hele Barndoms Dannelse og
i vor Ungdom Alt ud paa at oplue os i en Ærgerrighed og
en Selvfølelse, som den Gang maaske drev os fremad, men
nu standser og hilder?« Der var altsaa en stadig Kilde til
Bitterhed og Vrede, selv med forsvindende ydre Aarsager.
Men hvad Sibbern saa omkring sig i det offentlige Liv, fik
ofte den indre Ild til at flamme voldsomt op i ham, i senere
Aar især de nationalliberale Politikeres og Journalisters
Færd. Han svælgede da i Fantasier om, hvad han vilde
gøre ved dem, hvis han havde Aladdins Lampe eller kunde
gøre sig usynlig. Man maatte ikke ønske sig at være i
Lehmanns, Plougs og Billes Sted, dersom disse Fantasier vare
blevne til Virkelighed. --- I »Meddelelser af Indholdet af et
Skrift fra Aaret 2135« fortælles om en Bog, som var
udkommen i den sidste Fjerdedel af det nittende Aarhundrede,
og som indeholdt en gammel Mands Levnedsbeskrivelse af
ham selv. Heri havde Forfatteren fortalt, hvorledes han,
uden at Andre havde mærket til det, havde lidt af en
forfærdelig Sjælesygdom, som han kaldte Hævnhidsighedssygen.
Den opstod kun ved hvad han læste og udmalte for sig i
sine Forestillinger. Ofte kunde han af den Grund ikke holde
ud at læse Dagbladene. Sygdommen næredes ved Alt hvad
han læste eller selv oplevede af teologisk og politisk
Fanatisme baade i Fortid og Nutid; han havde ofte derved faaet
Hidsighedsanfald, som han endog havde følt legemlige
Virkninger af. »Bogen«, hedder det, »udkom først efter hans
Død, og Alle erklærede, at de aldrig havde tænkt sig, at han
% --- p. 73 ---
\opage{73}havde kunnet bære en saadan Hævnlystshidsighed i sig.«
Denne Mand er tydeligt nok Sibbern selv.

Dette Hang vidste han vel var begrundet i hans Natur;
men han fremhæver skarpt og bestemt, at det næres og plejes
mere eller mindre hos os Alle ved den Maade, hvorpaa de
offentlige Forhold behandles i Stat og i Kirke. Paa hans
gamle Dage, som Tilskuer ved det politiske Liv, især efter
1864, var det ofte vanskeligt for ham at holde fast ved hvad
han havde ladet sin Gabrielis udtale, at der egentligt intet
Ukrudt findes, heller ikke i Menneskeverdenen. Han gjorde
den Erfaring, at »for at bevare sin Menneskevelvilje,
Menneskekærlighed, Menneskeagtelse, maa man lade den stundom
ty ind under Menneskeforagten og her søge sig et
Tilflugtssted til Ly«. Der er altsaa, saaledes synes det at maatte
forstaas, visse Udvækster, man maa forkaste, for at kunne
bevare sin Tro og Kærlighed til Menneskeheden i det Hele.
Men denne Art Foragt, som kan tjene til
Konservationsmiddel for Menneskekærligheden, passer kun for de Gamle,
for dem, som have gennemgaaet Livets Skole. Hos de Yngre
maa der findes »en anderledes stærk og til Had hengrænsende
Foragt end den, hvorunder Menneskekærligheden skal søge
Ly«. (Samfundsbetragtninger. Andet Hefte 1870.) Det var
især Harmen over ubillige og af Partifanatisme og
Uvidenhed udspringende Domme, som fremkaldte saadanne Udtalelser
af ham.

Denne Mand har ikke taget sig Livet let. Det var
ogsaa en af hans Grundsætninger, »at man skal tage sig Livet
\emph{tilbørligen} nær«. Han havde nok at kæmpe med, og at han
kunde vidne om Livets Værdi, som han har, er et
betydningsfuldt Faktum. At hans Tro paa Livet vandtes gennem en
Kamp, gør hans Resultat saa meget mere vederhæftigt. Og
endeligt byggede han paa en Indsigt i det aandelige Livs
Love som sit teoretiske Grundlag. Han førtes tidligt til den
Overbevisning, at det aandelige Liv, ligesom alt Liv, udvikler
sig sporadisk, det vil sige ud fra forskellige, ofte
tilsyneladende modsatte Udgangspunkter. Meget forskellige Anlæg
og Drifter kunne paa én Gang ytre sig, og idet de drage
Sindet hver i sin Retning, opstaar den Gæring, Uro og
% --- p. 74 ---
\opage{74}Smerte, som især saa let kommer i de betydningsfulde
Ungdomsaar hos dybere Naturer. Ingen har beskrevet dette med
saa klar Forstaaelse som Sibbern. Det var en af hans
Specialiteter baade i den teoretiske og i den praktiske
Psykologi. At han kunde være saa Meget for saa mange Yngre
i deres Gæringstid, skyldtes ikke mindst denne Indsigt.
Han forstod hvad der brødes i dem og kunde indgyde dem
Taalmodighed til at lade de forskellige Kræfter udfolde sig i
Ro og Tillid til, at der ud af Striden og Uroen vilde udvikle
sig en saa meget des rigere Harmoni.

Selv følte han sig ikke lykkeligst i sin Ungdom. Han
saa hen til den Tid, da Aftenens Ro skulde afløse Dagens
Uro og Pinagtighed, og han ønskede sig inderligt at blive
Olding for at kunne se paa Livet uden »Brynde og Begær«.
Hans Ønske blev opfyldt. Det beroede vel, som tilstrækkeligt
bekendt, paa en Illusion, der er saa sædvanlig hos gamle
Folk, naar han i et Brev fra Aaret 1863 skrev, at han holdt
sine Forelæsninger med lige saa megen Aandskraft som før.
Det holdt haardt at faa ham til at trække sig tilbage fra
Universitetet; han følte sig saa rask, at han ikke kunde
tænke sig Andet end, at det var hans politiske Modstandere,
der vilde fjærne ham. Dog burde man for hans egen Skyld
have grebet mere energisk ind. Det var en Skandale, at
man udsatte den udmærkede gamle Mand, hvis Sanser vare
svækkede, og som ikke kunde gøre sin Myndighed gældende,
for de unge Tilhøreres kaade Drengestreger. Men havde han
saaledes Illusioner med Hensyn til sin Evne til udad gaaende
Virksomhed, saa var der en anden Henseende, i hvilken
ingen Illusion var mulig, nemlig Følelsernes Liv og Friskhed,
den »Sommerlighed«, som han selv udtrykte sig, der holdt
sig hos ham til det Sidste.

Denne Følelsernes Friskhed hang sammen med hans hele
barnlige Natur. En stor, kostelig Naivitet prægede Alt hvad
han tænkte, følte og sagde. Ejendommeligt for Naiviteten er
det jo at gaa lige løs paa Sagens Kærne uden alle
Formaliteter, med aabent Blik for Livets virkelige Sammenhæng.
Denne Naivitet maatte naturligvis ofte virke komisk. Da
han f. Eks. ved en Doktordisputats, som i sin Tid gjorde
% --- p. 75 ---
\opage{75}megen Opsigt, endte med meget hjerteligt at ønske Doktoranden
til Lykke med, at han nu dog endeligt en Gang havde opnaaet
det, som han saa mange Gange forgæves havde prøvet paa,
kendte Tilhørernes Jubel ingen Grænse. Det var ikke ironisk
ment. Ti Ironi laa aldeles ikke for Sibbern. Han skrev i
Tyverne et »ironisk« Forsvar for Englændernes Tog til
København 1807; men de halve Læsere toge det for Alvor,
og han maatte offentligt forklare, at det var ironisk ment. ---
Var denne Evne nægtet ham, saa havde han fuld Erstatning
derfor i den umiddelbare Varme og Begejstring, hvormed han
tog Livet og stræbte at faa Andre til at tage det.

Des værre lykkedes det ikke Sibbern i noget enkelt Værk
at give sig selv helt, heller ikke at give noget enkelt Værk
den fulde Afslutning hvad Tankeindholdet angaar. Der er
ingen Bog af ham, man kan pege paa som den, i hvilken
man har ham helt og holden. Maaske nogle udvalgte Stykker
at % PRINTED AS IS (for af)
»Gabrielis«, af »Psykologisk Patologi« (»Læren om de
menneskelige Følelser og Lidenskaber«), og ikke mindst af
»Meddelelser af Indholdet af et Skrift fra Aaret 2135« vilde
egne sig bedst til at karakterisere ham fra den Side, fra
hvilken han sikkert har sin største Betydning. Ti hvilken
Mening man end har om ham som videnskabelig Filosof, saa
var han i Ordets smukkeste Betydning Livsfilosof, en Mand,
der baade har levet rigt og dybt og formaaet at give os
værdifulde Bidrag til at forstaa vort Liv og til at øve den
vanskelige Kunst: at leve.

\section*{II.\\  Sibbern som Psykolog.}

De fleste og største af Sibberns videnskabelige Skrifter
behandle Psykologien, og han medbragte til dette Fag i flere
Henseender fortrinlige Forudsætninger. En frisk og levende
Natursans og Iagttagelsesevne, især for Fænomener, der høre
hjemme paa Følelses- og Driftslivets Omraade, og ikke ringe
Kendskab til Naturvidenskab og Fysiologi. Men hans
psykologiske Arbejder (af hvilke det første udkom 1819) bære
stærkt Præget af, at man den Gang i langt højere Grad end
% --- p. 76 ---
\opage{76}nu var uklar over Psykologiens Stilling og Forhold til andre
Videnskaber. Sibberns Bestræbelse gaar vel ud paa at
opfatte Psykologien som en Erfaringsvidenskab; men han
udsondrer den ikke tilstrækkeligt hverken fra den spekulative
Filosofi eller fra Moralfilosofien. Nogle af de vigtigste
psykologiske Begreber staa hos ham endnu belemrede med og
bestemte ved Elementer, der skrive sig fra spekulative eller
metafysiske Antagelser, som Psykologien i og for sig ikke
har noget at gøre med. Den blot iagttagende og beskrivende
Metode var ikke nok for Sibbern; han fandt i de vigtigste
psykologiske Iagttagelser saa betydningsfulde Data, at de
efter hans Mening kun kunde finde deres Forklaring i Lyset
af en højere Idé. Han skelner vel mellem »en Psykologi
paa Iagttagelsens Standpunkt« og »en spekulativ Psykologi«,
men efter at have anerkendt denne Distinktion føjer han til, at
»medens man paa Naturvidenskabens Omraade bestandig bør
holde sig nær til Erfaringen, saa kan man selv i den
iagttagende Psykologi ikke hurtigt nok hæve sig til den højere
Fornuftidé for at betragte alle Menneskets indre sjælelige
eller aandelige Ytringer i dens Lys«. Dels skulle vi kun
ved at gaa ud fra en saadan Idé forstaa, hvorledes de
forskellige Kræfter og Drifter kunne udvikle sig i Overensstemmelse
med de Vilkaar, under hvilke Individerne leve, dels skal
Menneskets Natur, dets Anlæg og Evner hænge saa nøje
sammen med dets moralske og religiøse Bestemmelse, at det
Ene ikke kan forstaas uden det Andet. Mange Partier af
hans Psykologi kunde lige saa godt findes i en Etik. Da
Etiken stedse bør bygges paa Psykologien, er der i og for
sig ingen Fejl heri; men Fejlen kommer, naar man uden
videre gør etiske Distinktioner til psykologiske Distinktioner.
Hvad der i etisk Henseende staar milevidt fra hinanden, kan
derfor dog psykologisk høre meget nært sammen. Derfor er
det af stor Vigtighed, at Psykologien bevarer sin
Selvstændighed som Erfaringsvidenskab, uafhængig af Etik og
Spekulation.

Dersom der gives en spekulativ Filosofi, kan det kun
være i den Betydning, at der gives visse Grundantagelser og
Grundforudsætninger, der danne de sidste Konsekvenser af
% --- p. 77 ---
\opage{77}al vor Viden. Det kan have sin Betydning at forsøge at
fremdrage og forme dem; derved faa vi først klart at se,
hvad der ligger i vor Viden, og komme altsaa først derved
til fuldstændig Selvbevidsthed. Men hvilken Betydning man
end vil tillægge saadanne Spekulationer, der aldrig kunne
blive andet end Hypoteser, saa maa Erfaringsvidenskaben ---
og det Erfaringsvidenskaben om de aandelige Fænomener
lige saa vel som Erfaringsvidenskaben om de materielle
Fænomener --- saa vidt muligt holdes uafhængig af dem. De maa
bygges paa Konsekvenserne af Erfaringsvidenskaberne, ikke
omvendt. Saaledes opfattede man nu ikke Sagen paa den
Tid, i hvilken Sibberns filosofiske Udvikling fandt Sted. Man
vilde først udlede og paavise den højeste Idé, der udtrykte
Tilværelsens inderste Væsen, og saa af den igen udlede, eller
i hvert Tilfælde i dens Lys betragte de specielle Fænomener.
Og i denne højeste Idé fandt man da tillige det sidste
Formaal, henimod hvilket Livet maa stræbe, naar det skal naa
sin rette Fuldkommenhed.

Som god Discipel af Schelling og Steffens stod Sibbern
ogsaa i det Væsentlige paa dette Standpunkt, og det har
hæmmet ham i den klare, naturlige Gennemarbejdelse af
Psykologien. Et enkelt, men meget vigtigt Eksempel vil
vise dette.

Der er neppe noget psykologisk Begreb, ved hvis
Behandling de forskellige psykologiske Anskuelser bedre
karakteriseres end ved Driftens Begreb. I Ordets videre
Betydning er Drift enhver dunkel Tilskyndelse, især naar denne
gaar i en bestemt Retning. Hvorledes forklares nu saadanne
Tilskyndelser? De opstaa jo ofte forud for al Erfaring om
og Nydelse af det, de gaa ud paa! --- Sibberns Forklaring er
den: »Livet vil frem i os og vil i os udtale sig i den Fylde
af Virkninger, i hvilken det hører Livet til at træde frem,
og det kan da ikke heller Andet end at medføre og fremkalde
Hentydninger til alt det, som er Betingelse for alle dens
Virkninger«. Her er nu en spekulativ Antydning sat i en
virkelig Forklarings Sted. At Livet vil frem i os, vide vi
kun af selve de uvilkaarlige Livsytringer, som vi kalde
Naturtilskyndelser. Men hvoraf kommer det, at Livet vil frem i
% --- p. 78 ---
\opage{78}en vis bestemt Retning? Det er intet Svar at sige, at denne
bestemte Retning hører med til Livet. Det gør den ganske
vist, da enhver Naturtilskyndelse i en vis Retning er en
Livsytring. Man slutter da: alle Livsytringer, altsaa ogsaa alle
Naturtilskyndelser, maa have deres Grund; nu er der
Naturtilskyndelser, som gaa ud paa Ernæring, andre, som gaa ud
paa Tilfredsstillelse af højere Fornødenheder, som Samliv med
andre Væsener, eller føre til moralske Handlinger o. s. v.;
altsaa maa den almindelige Drift, i Følge hvilken Livet
arbejder sig frem, fremtræde dels som en Ernæringsdrift, dels
som en Meddelelsesdrift, dels som en moralsk Drift, dels som
en religiøs Drift o. s. v. Man udstyrer altsaa den menneskelige
Natur med lige saa mange oprindelige Drifter, som der er
forskellige Retninger, i hvilke Livet udfolder sig. Saaledes
siger Sibbern: »Der er Intet, som henhører med til Livets
Indhold eller til det Indbegreb af Virkninger og Former,
hvori det fuldendte Liv udtaler sig, uden at det omfattes
med af Naturtilskyndelsen, og der gives derfor Intet, hvorpaa
Menneskets Liv og Stræben gaar ud, lige til det Allerhøjeste,
uden at en dertil svarende Drift bliver at nævne«.

Ved denne Maade at fremstille Sagen paa drejer man
sig i en Kreds, idet man skaber lige saa mange oprindelige
Drifter, som der er specielle Livsretninger. Man staar her i
Psykologien paa samme Standpunkt, som man i Fysiologien
stod paa, saa længe man antog en særegen Livskraft til
Forklaring af det, der udgør Forskellen mellem den organiske
og den uorganiske Naturs Fænomener. Kun da vilde man
have naat videre end til en saadan Kredsgang, hvis man
af Livets almindelige Idé havde kunnet udlede de specielle
Kræfter og Drifter i deres indbyrdes Forhold. Men dette
har stedse vist sig umuligt, og Sibbern prøver heller ikke
derpaa; han kunde, som vi senere skulle se, efter sin hele
Opfattelse af Filosofien heller ikke prøve derpaa. Naar en
saadan Udledning nu ikke er mulig, saa staar kun tilbage at
give en katalogiserende Opregning af de forskellige Drifter
med stadig Paamindelse om, at de alle have deres Plads og
Betydning i Livet. Sibbern slaar især i sine senere Skrifter
og Udgaver stærkt ind paa en saadan Katalogiseren og Ru%
% --- p. 79 ---
\opage{79}briceren, og det er for en stor Del dette, der giver hans Skrifter
noget Afskrækkende for mange Læsere.

Den nyere Psykologis Fremgangsmaade er i Modsætning
hertil den, for det første at holde Beskrivelsen af de enkelte
psykologiske Fænomener uafhængig af alle spekulative og
etiske Synspunkter, og dernæst at søge tilbage til de simpleste,
mest usammensatte Fænomener for at prøve, om de højere,
det vil sige: mere sammensatte Fænomener kunne forklares
som opstaaede af dem ved en sammenhængende
Udviklingsgang. Den nyere Psykologis Standpunkt kan betegnes som
analytisk-genetisk. Ved Analyse søger den at finde de fælles
Elementer i de forskellige Bevidsthedstilstande, og søger
dernæst at vise, at Forskellighederne bero paa de
forskelligartede Forbindelser af disse Elementer. Derved bliver det
ogsaa først muligt at paavise et naturligt Slægtskab mellem
de forskellige psykologiske Livsytringer. Dersom Talen f. Eks.
er om Forholdet mellem Egoisme og Sympati, saa er det ikke
nok at paavise, at der virkeligt gives Drifter, der føre i begge
Retninger; men det hjælper os heller ikke til dybere
psykologisk Forstaaelse, naar Sibbern begrunder Muligheden af
en ren sympatisk Drift ved, »at Individet jo overhovedet ikke
er til for sin egen Skyld, men for dets Skyld, hvis Organ
det skal være«. Derimod vil Betragtningen af de Instinkter,
til hvilke Slægtens Forplantning og Bevarelse er knyttet,
vise, at der i det enkelte Individ, selv om det ikke kommer
til dets fulde Bevidsthed, rører sig Kræfter og Tendenser,
som pege ud over dets egen Selvopholdelse, og det vil da
være muligt at vise, hvorledes der paa dette naturlige Grundlag
og efter det Stød, Naturen her selv giver, i Kraft af
almindelige psykologiske Love, kan finde en Udvikling Sted, der
gør en uselvisk Sympati forstaaelig. Sibbern har
fortræffelige Bemærkninger om de her omtalte Instinkter og Drifter;
men disse Bemærkninger gøres paa Grund af hans
almindelige psykologiske Standpunkt ikke frugtbringende for
Forstaaelsen af Følelsernes virkelige Tilblivelse.

Alligevel har Sibberns Psykologi en naturalistisk Karakter.
Det ligger allerede i den Maade, hvorpaa den lader Livet
fra dets højeste til dets laveste Trin arbejde sig frem under
% --- p. 80 ---
\opage{80}Driftens Form. Det ses ogsaa af den Iver, hvormed han
benyttede hvad hin Tids Fysiologi frembød til Belysning
af Bevidsthedslivet. Men en af Grundene til
Ufuldkommenheden ved Datidens Psykologi var netop det uklare
Standpunkt, hvorpaa Fysiologien da endnu stod. Den havde endnu
den Gang ikke konstitueret sig som mekanisk Naturvidenskab,
der følger samme Metode og anerkender samme almindelige
Principer som Fysiken og Kemien. Læren om »Livskraften«
afgav endnu det almindelige Synspunkt for de organiske
Fænomener. Ved dette Ubestemte, halvt Mystiske i de ledende
fysiologiske Principer kunde Teorien om Forholdet mellem
Sjæl og Legeme heller ikke forme sig klart og bestemt.
Problemet kom til at mangle sin egentlige Brod. Først naar
Fysiologien opstiller den Fordring, at \emph{alle} Fænomener i
Organismen, i Hjerne- og Nervesystemet saa vel som i Blodomløbs- og
Ernæringssystemet, skulle forklares efter de almindelige
mekaniske Principer, springer Vanskeligheden ved at forstaa
Sammenhængen mellem Bevidsthedsliv og organisk Liv ret i Øjnene.
Gælde de mekaniske Principer for denne Sammenhæng? Hvis
ikke, have vi da her ikke et Brud paa den materielle Naturs
almindelige Love? Det er netop Læren om den fysiske Energis
Bestaaen, som i vor Tid paany har givet det gamle
Spørgsmaal om Forholdet mellem Sjæl og Legeme en stor
Interesse. Sibberns psykologiske Arbejder vare endnu ikke
paavirkede af dette Synspunkt. Selv hans Bog »Om Forholdet
mellem Sjæl og Legeme« (1849) udarbejdedes endnu paa en
Tid, da Mayers, Joules og Helmholz’s Arbejder endnu ikke
havde faaet Indflydelse paa Fysiologien. Alligevel er det
interessant at se, hvor meget hans Opfattelse nærmer sig til
Teorier, som i den nyeste Tid gøres gældende af adskillige
Forskere paa dette Punkt.

Han erklærer sig imod enhver dualistisk Opfattelse og
lægger den Opfattelse til Grund, at Livet er ét som
Bevidsthedsliv og som materielt Liv. Kun denne Livets Enhed gør
det muligt at forstaa, hvorledes der til visse
Bevidsthedsfænomener kunne svare visse Hjernefænomener og omvendt:
Livet i Bevidsthedens Form er i Følge sin Natur i
Overensstemmelse med Livet i Materiens Form, dersom det er ét
% --- p. 81 ---
\opage{81}og samme Liv, der fremtræder under de to Former.
»Legemsvirksomhed og Sindsvirksomhed ere væsentligen
sammenhængende \emph{sideordnede} Virkninger af én fælles, samtidigt i
begge virkende Aarsag«. Han lægger især Vægt paa at det
legemlige Liv kun bestaar under en uophørlig
Reproduktionsvirksomhed, idet der stedse udskilles Dele af det organiske
Væv og optages og dannes andre i deres Sted. Det er
egentligt kun et Skin, naar Legemet synes at have en solidere
Basis for sin Bestaaen end Sjælen. Begge Steder have vi
kun en Virksomhed for os: ogsaa Sjælen lægger sig kun for
Dagen i en bestandig Frembringen og Bearbejden af
Forestillinger, Følelser og Bestræbelser.

Ikke alle Sibberns Udtalelser om dette Punkt ere lige
tydelige og konsekvente, men den almindelige Opfattelse, han
peger hen imod, synes at være klar nok. Vel er han ved
denne Teori mere behersket af den almindelige Interesse for
at hævde Livets Enhed under alle Former end for at udfinde,
hvilken Hypotese der bedst tilfredsstiller Erfaringsvidenskabens
vigtigste Resultater og ledende Tanker; men han har for sin
Tid forstaaet med stor Sikkerhed at gribe nogle af de
væsentligste Synspunkter. Han er bl. a. en bestemt Modstander af
den spiritualistiske Lære om Sjælen som en Substans eller
et Væsen for sig: »Sjæl er kun sjælelig Leven«. Den
Dogmatiseren, hvormed man slaar \emph{det} fast som værende og hvilende,
der stedse kun er givet i Vorden og Udvikling, har bidraget
Meget til at gøre de psykologiske Problemer vanskeligere, end
de i Virkeligheden ere. En saadan Fastslaaen og
Substantialisering stred desuden mod Sibberns hele Verdensanskuelse.
»Vor hele Tilværelse bestaar i at være en Funktion i
Alorganisationen«. »Den individuelle Virken og Leven er kun
en Funktion af det Alvirkende«. Individet fremgaar af den
hele store Naturvirksomhed og kan ikke rive sig saaledes
løs fra denne, at det skulde blive en selvstændig Helhed for
sig. Det staar i samme Forhold til hele Naturvirksomheden
som den enkelte Forestilling til hele Bevidsthedsvirksomheden.
Derved bliver det, efter Sibberns Opfattelse, ogsaa
forstaaeligt, hvorledes der i det enkelte Individ kan røre sig Noget,
som er af universel Natur: Instinkter og Tendenser, der pege
% --- p. 82 ---
\opage{82}langt ud over Individets Omraade. Men før jeg gaar nærmere
ind paa Sibberns almindelige filosofiske Verdensanskuelse,
vil jeg fremdrage nogle Punkter af hans Psykologi, som sikkert
have blivende Interesse og Betydning.

Hans Psykologi betegner tydeligt Reaktionen mod den
rationalistiske Periode ved den Vægt, der lægges paa
»Hjertet« i Modsætning til »Hovedet«. Han skildrer under
forskellige Former og Vendinger Forskellen mellem den rent
intellektuelle, forstandsmæssige Tilslutning til hvad man
erkender og anerkender, og den levende Hengivelse, ved hvilken
vi blive Ét med selve Sagen. «Det, som gaar ud over
Fornuften«, nemlig den kraftige og varme Følelse for det,
Fornuften har grebet, hører til det, han skildrer bedst. Mange
Sider ved Følelsernes Psykologi, især den store Rolle,
Kontrast- og Blandingsforhold her spille, har han i de store Træk
klart fremhævet\footnote{Betydningen af Sibberns Opfattelse paa dette Punkt anerkendes
nu ogsaa i udenlandske Værker (som \emph{Ribot:} La Psychologie des
Sentiments. Paris 1896. p. 265), efter at der i min Psykologi
(1882) og i A. Lehmanns „Hovedlovene for det menneskelige
Følelsesliv“ (1892) var henvist til den.}.

Hermed hænger ogsaa den Betydning sammen, som han
tillægger det ubevidste Sjæleliv, eller, som han udtrykker sig:
»det, som foregaar i Sjælens dybe Grund«. Kun en ringe Del
af vor Natur træder klart frem for vor Bevidsthed. Ligesom
vi fra først af blive til som Individer formedelst en
Naturvirksomhed, vi ikke kunne blive os bevidst, saaledes maa
denne for os ubevidste Naturvirksomhed stedse antages at
fortsætte sig, saa længe vort aandelige Liv varer. Der er
jo heller ikke paa det fysiske Omraade nogen absolut Grænse
mellem Avling og Vækst. I Instinktet og i Geniet have vi
slaaende Eksempler paa Noget, som virker i os, men først i
sine Virkninger, ofte fjerne Virkninger, kommer op for
Bevidstheden. Vor aandelige Vækst foregaar overvejende
ubevidst, saa det gaar med os som med Manden i Parablen, der
»kaster Sæd i Jorden og sover og staar op Nat og Dag, og
Sæden vokser og bliver høj, han véd ej selv hvorledes.« Der
er bestandigt Noget, som danner Bevidsthedens Grundlag uden
% --- p. 83 ---
\opage{83}at komme frem for Bevidstheden; og der er et stadigt
Vekselforhold mellem det Ubevidste og det klart Bevidste. Det er
sikkert en fuldstændigt rigtig Bemærkning, Sibbern gør om
Angeren (hvilken han i det Hele skildrer godt og uden den
indeterministiske Dogmatisme, som den ældre Psykologi var
tilbøjelig til), --- at dens psykologiske og etiske Betydning
overvejende beror paa dens Indflydelse paa hele det ubevidste
Sjæleliv; Menneskets Karakter kan gennem Angerfølelsen
og ved dens lutrende Kraft blive en anden. Man mindes
her Gabrielis’s Ytring om, at Smerten var det Eneste, der
holdt hans Bevidsthed sammen. Saaledes kan Angeren være
et uundgaaeligt Mellemled, uden hvilket Karakterudviklingen
ikke kan gaa for sig.

Hvad der nu fra dette ubevidste Grundlag paavirker og
bestemmer Bevidsthedslivet, melder sig ofte spredt ud fra
flere isolerede Punkter og i meget forskellig Retning paa
én Gang. Og ligeledes faar Bevidstheden fra Omverdenen
en Mængde forskellige, indbyrdes usammenhængende
Iagttagelser og Forestillinger. Der vækkes endeligt ved Erfaringerne
en Mængde forskellige Drifter, Ønsker og Forsætter. Alt
dette kan nu i broget Forvirring kæmpe i Bevidstheden, og
det netop, jo rigere Individets Natur er. Det vil da komme
til »et vist Livets Røre, en Livets Gæring, en Livets Debat,
en Livets Kamp«, inden der kan danne sig et bestemt Centrum
og en harmonisk Orden i Sindet. »Skønt stærkest i Livets
ungdommelige Periode (Ynglingsalderen), der for Resten hos
forskellige Individer kan falde i forskellige Tidspunkter,
vedvarer den dog, mer eller mindre, hele Livet igennem. Der
gives ogsaa Individer, som i denne Henseende blive som
bestandigen unge, ja nogle gives der, som endog hele Livet
igennem blive saa hildede i den med denne Livets Kamp
naturligen i Ungdommen forbundne Gæring, at det er, som
om hele Livet for dem var en aandelig Fødselskamp.«

Dette er Læren om den sporadiske Udvikling, som i Sibberns
Verdensanskuelse spiller en ikke mindre Rolle end hans Tro
paa Livets Enhed og Helhed, og som han ikke mindst finder
klart bekræftet paa det psykologiske Omraade. Han gør
rigtigt opmærksom paa, at dersom Udviklingen ikke foregik
% --- p. 84 ---
\opage{84}sporadisk, vilde den ikke foregaa gennem Kamp. Hans
Iagttagelse og hans Tankegang fører ham her lige til den
moderne realistiske Lære om Kampen for Tilværelsen. Dette
Synspunkt ligger ham saa nær, at han bestemt indskærper
(allerede i en Afhandling fra Aaret 1829), at den Sætning:
Livet er en Kamp, ikke maa misforstaas , som om det kun
var de ydre Forhold, der nødte Livet til at kæmpe; nej,
Livet er kun til i Kampen, kan ikke bestaa i og for sig
uden den. Ved den Energi, hvormed han, ikke blot paa det
psykologiske Omraade, men, som vi snart skulle se, i sin
Verdensanskuelse i det Hele, fastholder dette Synspunkt, har
han givet interessante Bidrag til den almindelige Lære om
Udviklingen.

\section*{III.\\  Sibberns filosofiske Verdensanskuelse.}

Sibbern var stærkt paavirket af den tyske Filosofi, og
hans Lære er en af de sundeste Former, i hvilken denne
Filosofi har vundet Indflydelse paa Tankegangen i den danske
Aandsverden; men hans Sans for det Konkrete, hans
Iagttagerinteresse og hans varme Sympati for det individuelle
Liv gjorde, at han ikke kunde blive en Tilhænger af den
abstrakte og aprioriske Deduceren og Konstrueren, hvori den
ældre Skole endte. Han traadte bestemt i Marken mod Hegel,
og allerede Fichte og Schelling tog han paa sin egen Maade.
Han tager bestemt Afstand fra den Schellingske Naturfilosofi.
I et Brev fra 1824 hedder det: »Det er raadeligt for enhver
Naturbetragter at holde sig Erfaringen saa nær som muligt
og at være varsom med ethvert Forsøg paa at ville, sættende
sig over paa det højere Standpunkt, oplyse og udtyde
Naturtingene ved en Betragtning \emph{fra oven}: ti dertil ere
Naturbetragtningerne endnu langt fra at være modnede«. Som man erindrer,
gennemførte han ikke dette Princip for Psykologiens
Vedkommende. Og dog er det karakteristisk for hele hans
Opfattelse af Filosofien.

Han saa’ ganske vist det højeste Ideal i en Erkendelse,
ved hvilken vi ud fra en højeste Idé begrunde og forklare
hele vor Viden om Tilværelsen. En Konstruktion a priori
% --- p. 85 ---
\opage{85}var ogsaa for ham den filosofiske Stræbens højeste Formaal.
Men Filosofien maa efter hans Overbevisning stedse gaa ud
fra et givet Indhold, og først gennem Bearbejdelsen af dette
kan efterhaanden den sidste, ledende Idé udvikle sig.
Filosofien maa, som han udtrykker sig, være eksplikativ, før den
kan blive spekulativ, det vil sige, den maa gaa ud fra
Betragtninger over det Givne for i dette at finde Hentydninger
til en Idé, i hvis Lys det maatte kunne finde sin Forklaring.
Og den spekulative Erkendelse maa selv igen stedse søge
sin Bekræftelse gennem stadig Sammenholden (»Konferens«)
med det virkeligt Givne.

Spørge vi, hvorledes det er muligt, at en saadan Idé
kan findes, og en saadan Forklaring til en vis Grad naas,
da vi jo dog kun kende og overskue saa ringe en Del af
Tilværelsen, saa svarer Sibbern, at vi selv jo ere et Led af
Tilværelsen, og at det Samme, som udgør Tilværelsens inderste
Væsen i det Hele, ogsaa gør sig gældende i os, naar Idéerne
under Betragtning og Bearbejdelse af de givne Erfaringer
arbejde sig frem i os. Men at Maalet ikke fuldstændigt vil
kunne naas i en overskuelig Tid, ligger i det Sporadiske i
Tilværelsen. Verden træder frem for os fra sin Fyldes mange
Punkter paa én Gang; det Centrum, omkring hvilket alle
disse Punkter kunne ses i deres Sammenhæng, vil kun kunne
findes gennem »en Debat af Alt mod Alt«, som vanskeligt kan
tænkes afsluttet. Det er da intet Under, at der er saa mange
Stridigheder mellem Filosoferne; vi ere jo endnu ved at
prøve os frem, og derfor maa snart et, snart et andet
Hovedpunkt komme til at træde frem med ensidig Styrke.

De, som kende Filosofiens nyere Historie, ville finde, at
Sibberns hele Opfattelse af Filosofiens Væsen og Metode ikke
lidet minder om Lotze. Ogsaa ellers fremtræder der et vist
Slægtskab mellem de to Tænkeres Anskuelser.

Sibbern baner sig Vej fra sin Erkendelseslære til sin
metafysiske Teori ved at vise, at Erkendelsens Mulighed beror
paa, at Tilværelsen frembyder en Sammenhæng af det Ene
med det Andet, saa at Alt tilsidst danner en altomfattende
Helhed. Dersom der ikke i Tilværelsen gives en altomfattende
Tingenes Orden, i Følge hvilken hver enkelt Ting bestemmes
% --- p. 86 ---
\opage{86}ved de andre, saa kan ingen enkelt Ting finde sin fulde
Forklaring. Der maa gives en Totalkonsekvens, som bærer de
enkelte Ting og Fænomener. Det er den filosofiske
Grundtanke, endnu i ubestemt og indholdsløs Form. Den er allerede
Idéen om et Alværende og Alvirkende, som gennem den
altomfattende Tingenes Orden bærer og bestemmer de enkelte
Fænomener.

Men denne Forudsætning er ikke tilstrækkelig. Ti det
Alværende og Alvirkende, den almindelige Tingenes Orden
maa paa hvert enkelt Punkt virke forskelligt efter det Givne,
som bestemmes ved det. Al Opstaaen og Tilbliven fremkommer
ved to Faktorer: den ene er den i Alt virkende Tingenes
Orden, den anden er noget vist Givet, som medfører, at hver
Gang noget vist Bestemt og ikke Andet bliver til. Der er
bestandigt i Naturen forskellige Udgangspunkter, ud fra hvilke
Processerne tage deres Løb. Og disse Udgangspunkter ere
hvert for sig endelige og relative; de staa i indbyrdes
Vekselvirkning, saa at de ere ikke blot Aktionspunkter, men ogsaa
Reaktionspunkter. Derfor fremtræder Naturen for os som
en stor Vekselvirkningsproces, en stor Debat og Kamp af
Alt imod Alt, og først gennem denne store Verdensdebat
udfolder Tilværelsesprincipet sin hele Fylde.

Denne anden Forudsætning kunde man kalde den
realistiske Forudsætning, ligesom man kunde kalde den første
den idealistiske, skønt Sibbern selv bruger andre Udtryk.
Ved denne Forudsætning træder han fra en ny Side i en
bestemt Modsætning til Hegel og den tyske spekulative Skole.
Han bebrejder Hegel, at han paa Grund af sin ensidigt
idealistiske Synsmaade ikke kunde faa nogen virkeligt \emph{historisk}
Opfattelse af Tilværelsen. Historisk Udvikling bliver først
mulig, naar der er sondrede Udgangspunkter i indbyrdes
Vekselvirkning. Udviklingen gaar da ud paa at bringe
Harmoni mellem de sporadisk fremtrængende Processer, at
organisere det i spredt Form givne Materiale. Erfaringen
viser os, som vi allerede for Psykologiens Vedkommende have
fremhævet, at Alt i Naturen fra først af danner sig sporadisk.
Ved Krystallisationsprocessen, som f. Eks. den, der finder Sted,
naar Vandet fryser, opstaa Naalene ud fra forskellige Punkter,
% --- p. 87 ---
\opage{87}men gaa efterhaanden sammen og danne en kompakt Enhed.
Fosterdannelsen sker i flere Henseender ud fra forskellige
Organisationspunkter, hvorfra Processerne mødes for tilsidst
at ende i et Resultat, der har et slaaende Enhedspræg. Den
sociale Verden viser os noget Tilsvarende: inden for den enkelte
Stat virker Bestræbelser af forskellige Individer og fra
forskellig Side til at danne en fælles Samfundsorganisme, og
de forskellige Stater ere igen Udgangspunkter for Virkninger,
som gennem mange Sammenstød føre os fremad mod et
almindeligt Menneskesamfund. Gaa vi tilbage til Jordklodens
og hele Solsystemets Oprindelse, saa maa enhver Partikel i
Urtaagen have spillet sin Rolle, have været Udgangspunkt
for Tiltrækken og Frastøden, og gennem det Mylr af
Vekselvirkninger, som derved bragtes til Veje, fik vor Klode og
vort Solsystem sin nuværende Form. Paa ethvert Omraade
kan der spores en stor, fremskridende Organisationsproces i
ubrudt Sammenhæng. »Det allerførste Røre i
Jordklodeprocessen var Begyndelsen til den hele nu paa Jorden
subsisterende Planteverdens, Dyreverdens og Aandeligheds
Konstitution«. Tilværelsen er en stor, kontinuerlig Udviklingsproces,
indenfor hvilken Menneskets Tilblivelse betegner et enkelt
Stadium. »Hvorledes mon det første Menneske blev til?
Uden Tvivl derved, at en vis Dyreformation, fra en første
ringe Begyndelse gennem en Række af Generationer eller
Affødninger, idet i Afkommet bestandigen Uddannelsesprocessen
gik videre og videre, naaede til det Punkt, da nu dets sidste
Affødning blev et Væsen, som saa at sige slog
Aandelighedens Øje op, idet det sagde Jeg til sig selv og følte sig
i sin Jeghed«\footnote{En lignende Anskuelse af Menneskets Tilblivelse havde allerede
Kant udtalt (i sin Antropologi), og hos os Treschow (i sin „Historiens
Filosofi“).}.

Der er altsaa trods de sporadiske Udgangspunkter en
Fremgang til stedse større Harmoni og til højere Livsformer.
Men paa de mellemliggende Stadier kan der ofte komme Former
og Fænomener frem, der synes at pege i ganske anden Retning
end imod en Fremgang til højere Fuldkommenhed, ligesom
% --- p. 88 ---
\opage{88}Fosteret paa visse Stadier af sin Udvikling forekommer den
ukyndige Betragter at være en Vanskabning. Naturen synes
ofte at gaa Absurditetens Veje. Alt Sligt er dog kun noget
Foreløbigt og Forbigaaende, en lille Krumning paa en Vej,
som Verdensudviklingen maa tilbagelægge. Endnu er
Tilværelsen paa Vejen, ja endog paa Begyndelsesstadiet af Vejen
til Maalet. Der viser sig endnu i vor Erfaring et
Modsætningsforhold mellem det centrale Princip, Enhedsprincipet i Verden,
og de mangfoldige periferiske Udgangspunkter. Men det vil
mere og mere blive afløst af et Harmoniforhold. Kun for
en begrænset Del af Verden kan der dog være Tale om en
fuldstændig Virkeliggørelse af Verdensformaalene; da der i
Verden i det Hele skal virkeliggøres en uudtømmelig
Indholdsfylde, vil der dertil behøves en uendelig Tid.

Disse vare de filosofiske Ideer. Sibbern udtalte i sit
mærkelige lille Skrift »Spekulativ Kosmologi« , som udkom 1846,
det eneste, i hvilket han i Sammenhæng har fremstillet sin
almindelige Verdensanskuelse. Det lønner sig endnu for dem,
der interessere sig for filosofiske Spekulationer, at studere
dette Skrift, der, som man vil have set, ikke mangler
Tilknytningspunkter for den nyeste Tids Teorier og Synsmaader.
Til nærmere Karakteristik vil jeg her kun tilføje nogle faa
Bemærkninger.

En af de vigtigste Sider, fra hvilken en filosofisk
Verdensanskuelse kan karakteriseres, er det Forhold, den statuerer
mellem Enhed og Mangfoldighed i Tilværelsen. Sibbern søger
at give begge hvad der tilkommer dem, hver for sig, og i
hans Fremstilling ses meget klart de Motiver, som give dem
Interesse og Betydning i filosofisk Henseende. Til at
forudsætte Enhed føres vi ved det Grundfaktum, at Tilværelsen,
hvor vi formaa at trænge ind i den, fremtræder for os som
en lovordnet Sammenhæng. Der maa da, synes det, være et
eneste Princip, som virker i Alt. Men paa den anden Side
viser Erfaringen os lige saa vel en Mangfoldighed af virkende
Elementer, af sondrede Udgangspunkter. Uagtet nu dette
sidste Synspunkt betones meget stærkt af Sibbern, især i den
for ham saa ejendommelige Lære om den sporadiske Udvikling,
saa underordner han dog afgjort Mangfoldigheden under
% --- p. 89 ---
\opage{89}Enheden. De reale Udgangspunkter betegnes vel som noget
»Givet«, paa hvilket »det Alvirkende« øver sin Indflydelse; men
dette Givne skal selv kun være det Alvirkendes egne
primitive Virksomheder, og det Alvirkendes Indflydelse skal ganske
gaa i Ét med de sondrede Elementers Virksomhed. En
Opfattelse som Leibniz’s Monadelære eller som den absolute
Atomlære, der ender i et Mylr af selvstændige Enkeltvæsener som den
sidste Tanke, afviser Sibbern bestemt. Ethvert Enkeltvæsen
er kun et Moment i det uendelige Væsen, enhver Molekule
kun Sæde for en af det Alvirkendes Funktioner. Det skal
kun være vor Fantasi, der ikke kan lade være at tillægge
Elementerne eller Molekulerne en anden Realitet end den at
være blotte Virksomhedsudgangspunkter for det Alvirksomme,
ligesom vore enkelte Forestillinger ere lige saa mange Former
for Aandsvirksomheden og ikke have nogen Eksistens afset
fra denne.

Men hermed er det ikke at forlige, at Indbegrebet af
Elementerne kaldes et \emph{Modstaaende}, et \emph{Andet}, et \emph{Objektivt},
som skal bearbejdes og formes af det Alvirksomme. Her
fremtræder en Dobbelthed, en Modsætning, som ikke er
overvunden. Naar en Idé skal virkeliggøres, siger Sibbern i
denne Sammenhæng, kan Alt ikke fremgaa af dens indre
Væld, men Stof og Middel maa tages ude fra. Her er altsaa
en Dobbelthed til Stede fra først af; de to Forudsætninger,
hvorpaa der bygges, kunne ikke føres tilbage til en Enhed,
og der ligger her en truende Fare for Sibberns optimistiske
Verdensanskuelse.

Hvorledes kan han være saa sikker paa, at der vil komme
stedse større Harmoni mellem Enhedsprincipet og de sporadiske
Udviklingsprocesser? Kun saa kort og ringe en Del af
Udviklingen kunne vi overskue, og vi kunne ikke herefter stille
Horoskopet for hele Udviklingen. Vi kunne trøste os med,
at Absurditeterne og Disharmonierne kun skyldes vor Mangel
paa Kendskab til Udviklingens egentlige Love og Betingelser:
men det er da kun en Tro eller et Haab, vi her tage vor
Tilflugt til, ingen videnskabelig Indsigt. Det er, om man vil,
en filosofisk Tro, vi kunne have paa, at Spliden og
Disharmonien gennem Kampen, gennem den store Debat af Alt mod
% --- p. 90 ---
\opage{90}Alt, fører over til en saa meget des rigere Organisation;
men Mere kan det aldrig blive. Det er en energisk
Overbevisning, der rører sig hos den begejstrede Tænker, der saa
stærkt har betonet det Sporadiske og Brudstykkeagtige, naar
han stedse fastholder Ideen om en sig gennem Striden
fremarbejdende Harmoni; men den kan ikke videnskabeligt
begrundes og gennemføres.

Maaske hænger det sammen hermed, at Sibbern ikke
helt ud fastholder sin Udviklingslære. Hans Tro paa det
Alvirksomme fører ham ikke blot til at tilskrive det en
uendelig Overlegenhed over for de enkelte Elementer, der
skulle organiseres; men han undtager det endog fra at være
Udviklingsprocessen underlagt. Den egentlige Udvikling
bestaar for ham i, at Elementerne organiseres; den organiserende
Magt, det Alvirksomme er ikke selv vordende men evigt.
Eller som han udtrykker det med en teologisk Vending: Gud
vorder ikke, men kun Guds Rige vorder. En saadan
Distinktion kan ikke fastholdes, og Sibbern føler det selv.
Hvorledes kan Guden være evig og uden Udvikling, naar hans
Rige ikke er det? Guden og hans Rige høre saa nøje
sammen, at de kun blive til med hinanden; der kan ikke
være et Ligegyldighedsforhold mellem dem. Vorder Guds
Rige, saa er Gud ogsaa stedt i Vorden. Sibbern har her
optaget Mere af Teologien end de blotte Ord. Vel peger
den kristne Lære om en lidende Gud med storartet
Dybsindighed hen til hvad den gennemførte Udviklingsfilosofi
fordrer; men Teologiens sædvanlige Tilbøjelighed til at halvere
Tankerne har ladet den gyse tilbage for denne dristige
Konsekvens.

Uagtet Sibbern ynder at bruge teologiske Udtryk, er
hans Lære aldeles kættersk. Gud er for ham ikke noget fra
»Verden« sondret og forskelligt Væsen. Han skelner mellem
Tilværelsens »Verdensside« og »Guddomsside«: hin have vi i
Elementernes Mangfoldighed, denne i Enhedsprincipets Virken.
Guddommen virker gennem de fra Elementerne udgaaende
Virksomheder. Dog tilfredsstiller Panteismen (hvilken han iøvrigt
omtaler i varme Ord som en alvorlig og berettiget Livsanskuelse)
ham ikke, fordi den lærer, »at der ikke gives anden Gud end Guds
% --- p. 91 ---
\opage{91}Rige, eller at der bag det Rige, jeg her kalder Guds Rige,
ikke er nogen personlig Gud at søge«. Det er alligevel et
Spørgsmaal, hvor Mange af dem, der tro paa en personlig
Gud, der ville finde sig tilfredsstillede ved den Sætning, i
hvilken Sibbern slutteligt udtrykker sin religionsfilosofiske
Anskuelse: »Fører dog ikke Alt til, at man tilsidst maa
erkende, at Verdenstilværelsens inderste Grund maa være
centraliseret i sig, at Verden saa at sige er en stor Individualitet?«

Vi føres ved det nys Omtalte naturligt over til Sibberns
religiøse Anskuelser.

\section*{IV.\\  Sibberns religiøse Udvikling.}

Sibberns Ungdom faldt i Tiden for den store religiøse
Reaktion mod det attende Aarhundrede. Hans filosofiske
Lærere, Fichte, Schelling og Steffens, vare, da han kom under
deres Indflydelse, allerede paa Veje til at teologisere og
befandt sig i hvert Tilfælde i en vis Reaktion mod deres egne
Ungdomsanskuelser. Her hjemme stod han i nært personligt
Forhold til Grundtvig og Mynster. Han stod endog en lang
Tid som den eneste Lærer ved Universitetet, hos hvem de,
der følte sig utilfredsstillede ved Rationalismen, kunde finde
Tilhold. Hans Stilling i denne Henseende var bestemt ved
hans hele Personlighed og ved hans psykologiske
Grundanskuelse. Han følte dybt Trangen til et rigere Følelses- og
Fantasiindhold end Rationalismen havde kunnet byde, og han
saa klart, at deri den menneskelige Aand rørte sig langt
flere Kræfter og Drifter, end den flade, forstandsmæssige og
uhistoriske »Fornuftreligion« havde anerkendt.

Dog tog han stedse sine Forbehold. Uagtet han
protesterede mod Hejbergs Sætning, at Frelsen var at søge i
Filosofien, og uagtet den filosofiske Forsken for ham kun var
et Middel, en forberedende Prøvelse, medens den egentlige
Livsanskuelse hviler paa personlig Grebethed og hjertelig
Tilslutning, saa gjorde han alligevel enhver dogmatisk
Antagelses Gyldighed afhængig af den filosofiske Drøftelse. Et
er, at en Lære er grundet i kirkelig Overlevering, og at den
sætter vor Følelse i stærk Bevægelse, et Andet er det, om
% --- p. 92 ---
\opage{92}den indeholder Noget, der har Gyldighed og Realitet.
Filosofien skal undersøge dette uathængigt % PRINTED AS IS (for uafhængigt)
 af al Tradition og alle
Følelsesindflydelser. »Overhovedet staar,« hedder det i en
Afhandling om Filosofiens Forhold til Teologien, som findes
i »Filosofisk Arkiv« (1829), »Filosofien i samme Forhold til
det i religiøs Henseende positivt eller de facto Givne, som til
det paa anden Maade i Virkeligheden Givne; og det hvad
enten hint er givet som en udvortes, saakaldet positiv
Aabenbaring, eller som den af Nogle saakaldte indvortes
Aabenbaring.« Begge Steder er det den reale Betydning og
Grundethed, som skal godtgøres. Og det er efter Sibberns Opfattelse 
ikke en Anmasselse af en udenfor Staaende, naar Filosofien
vil prøve og undersøge de dogmatiske Lærdomme. Tvert
imod: i disse Lærdomme er der indeholdt en gammel Filosofi.
»Jeg tror at kunne paastaa,« siger han, »at dersom man af
Dogmatiken, ej blot saaledes som den nu er, men og som
den fra de første Tider af har dannet sig, vilde uddrage Alt
hvad der er udsprunget af Filosofering, forfølgende det i dets
fineste Forgreninger ind i Præmisserne for de dogmatiske
Lærdomme, saa vilde den falde sammen.« Den historiske
Kritik giver Sibbern Ret i denne Paastand, idet den
paaviser Platonismens store Indflydelse paa de kristne
Centraldogmers Udvikling og Formulering. Men han var sig, da
han nedskrev de anførte Sætninger, ikke Rækkevidden af
deres Indhold bevidst. Han mente sig endnu væsentlig i
Harmoni med den gældende Teologi og med Kirkens Lære.

Kun paa enkelte Punkter havde han faaet
Betænkeligheder, og disse Punkter indeholdt Spiren til den afgjorte
Ændring i hans religiøse Standpunkt. Det var især
Lindbergs Optræden i den grundtvigske Strid, som henledte hans
Opmærksomhed derpaa. Der var to Ting, der stødte ham.
For det Første den teologiske Fanatisme, hvorom hin Strid
syntes ham at vidne; dernæst den Vægt, Lindberg lagde paa
Djævletroen og Troen paa den evige Fordømmelse. Om denne
Tro maatte Sibbern indrømme, at den fandtes i Skriften, men
han forkastede den alligevel paa det Allerbestemteste. Dette
maatte fjerne ham afgjort fra Ortodoksiens Forkæmpere.
Den heftige Strid paa det teologiske Omraade her hjemme i
% --- p. 93 ---
\opage{93}Tyverne, Trediverne og Fyrrerne drejede sig i dogmatisk
Henseende væsentligt om, hvor vidt en Mand kunde være
Præst eller Lærer i den danske Kirke, naar han ikke troede
paa en personlig Djævel og paa en evig Fordømmelse. Den
ortodokse Retning har, som vitterligt for Alle, væsentligt sejret;
de nævnte Dogmer ere de mest karakteristiske for den i vort
Fædreland nu herskende teologiske Anskuelse; Sibberns
Tilslutning til den religiøse Bevægelse hørte op paa det Punkt,
da han mærkede, hvilken Vægt man lagde paa dem.

Efter at han i længere Tid ikke havde ytret sig om
religiøse Spørgsmaal, fremsatte han i en Række
Universitetsprogrammer (1846, 1847 og 1849) det nye Standpunkt,
hvortil han nu var naaet. Om den Forandring, der var
foregaaet med ham, har han selv udtalt sig saaledes i et Brev
fra Aaret 1864: »Alt Aanden Tiltalende, Sjælefødende, alt sin
Kraft fra Aandens indre Samstemning Hentende i
Kristendommen er for mig ganske det, det var. Kun det, som
skulde hente sin Gyldighed fra Positiviteten, \emph{respekterede} jeg
den Gang som et Givet ...... Denne Athængighed % PRINTED AS IS (for Afhængighed)
 vil jeg
nu slet ikke mere indrømme nogen Gyldighed.« --- Hvad han
i de anførte Programmer havde udtalt i sin Almindelighed,
gennemførte han videre dels i »Meddelelser af et Skrift fra
Aaret 2135«, dels i den efter hans Død udgivne Moralfilosofi.

Man kan spore flere forskellige Motiver, der hos Sibbern
havde virket sammen til at fremkalde og forme hans nye
Standpunkt. --- For det Første rører der sig hos ham en
filosofisk-historisk Interesse for at bringe Kristendommens
Opstaaen og Kristi Personlighed i naturlig Sammenhæng med
Verdensudviklingen. Allerede i sin »Kosmologi« havde han
udtalt, at ligesom Menneskets Tilblivelse fra først af maatte
forklares som Resultat af en videre og højere Udvikling fra
lavere Former, saaledes maatte Kristendommens Opstaaen igen
inden for den menneskelige Udvikling opfattes som et særeget,
men i Kontinuitet med den foregaaende Udvikling fremkommet
Vendepunkt. Vil man tale om en Inkarnation, maa derfor
hele Verdensudviklingen, særlig Jordudviklingen, opfattes
som en fremskridende Inkarnation. Paa det Punkt under
Urtaagens Udvikling, hvor Jordkloden begyndte at udsondre
% --- p. 94 ---
\opage{94}sig fra den øvrige Verdensmasse, maa vi søge den første
Spire til det aandelige Liv paa Jorden i det Hele og derfor
ogsaa til de højeste Former, hvorunder det aandelige Liv
historisk er fremtraadt. Den historiske Kristus er fremgaaet
af Menneskehedens Midte som en Mand, i hvem Livsvældet
rørte sig paa en særegen Maade, saa han kunde blive en
aandelig Fører for Aartusinder. Kun naar han opfattes paa
rent menneskelig Maade, har hans Kamp og Lidelse, hans
Stræben og Virken sand og naturlig Betydning. Man skulde
ikke have forvansket denne ved at grunde den paa noget
Andet end hans Virksomheds egen indre Magt og Kraft. ---
Dertil slutter sig den Overbevisning, at det religiøse Liv
væsentligt bestaar i et indre Forhold til Tilværelsens
Grundkilde, en Følelse af inderlig Enhed med og Tilslutning til
denne »indre Giver«. Man mener ofte, at det er en usikker
Grund at bygge paa, og at man maa have et objektivt
Grundlag i en dogmatisk Bekendelse og en i det Ydre fremtrædende
og organiseret Kirke, naar det religiøse Liv ikke skal tabe
sig i blot subjektive og vilkaarlige Indskydelser og Illusioner.
Men hertil maa svares, at alt Aandeligt, som skal komme til
Herredømme i Verden, netop gaar Subjektivitetens Vej, rører
sig med indefra kommende, gribende og besjælende Magt
hos de enkelte Personligheder. Det egentlige Objektive er
her netop den indre besjælende Magt, som rører sig i de
enkelte Individer. Den historiske Udvikling og de historiske
Former tabe derved ikke deres Betydning, ti der maa ydre
Paavirkning og Hjælp til for at det Indre kan blive sat i
Bevægelse. Derimod er det en Illusion, som Kirken og
Teologien have hildet sig i, at Noget skulde blive objektivt,
det vil sige i og for sig gyldigt, fordi det anerkendes af store
Kredse og lægges til Grund for Institutioner i den ydre
Verden. Den indre Magt, som fører det aandelige Liv fremad,
rørte sig nok saa kraftigt i Rationalismen og Pietismen som
i den ortodokse Kirke. »Brød ikke Rationalismen vældigen
frem, og tør Nogen sige, at ikke Humanitet, Menneskeomsorg,
Omsorgen for at bespise de Hungrige og læske de Tørstige
lige saa vel og lige saa varmt har fundet Sted paa denne
Side som paa hin?« --- Man ser her, hvorledes Sibbern er
% --- p. 95 ---
\opage{95}kommen ud over Reaktionen mod Rationalismen og formaar
at stille denne paa dens berettigede Plads blandt de med
hverandre kæmpende Retninger. Han finder ogsaa i denne
Sammenhæng Anvendelse for sit Yndlingssynspunkt: Debatten
af Alt mod Alt. Gennem den indbyrdes Brydning mellem alle
de forskellige Retninger og Personligheder gaar det religiøse
Liv fremad. Man skal ikke stille sig \emph{ligegyldig} over for disse
forskellige Anskuelser, men man skal tillægge dem \emph{lige
Gyldighed}, for saa vidt de hver for sig vidne om det indre
Livsvæld og have Ret til at gøre sig gældende og yde deres
Bidrag til Menneskelivets Udviklingsproces. I anden Del af
Gabrielis’ Breve, der ellers har et saa kirkeligt Præg, og
hvor Helten ender med at blive Præst, er det dog den stadige
Understrøm, at enhver stærk og levende Tro i og for sig er
berettiget som kraftigt Udslag af Livets indre Magt. Selv
om dem, der efter fuld Overbevisning forkaste Troen paa
Gud og Udødeligheden, maa man --- hedder det --- sige, at
de have Ret dertil, naar det er »en levende Sans for
Idealet eller dog for det, som skulde og burde være«, der rører
sig hos dem; de tænke ogsaa paa Sjælenes Vel og Frelse.
Den Form for Religiøsitet, som Sibbern selv skildrer med varm
Tilslutning, er en Tro paa en personlig Guddom, »et Liv i
Gud som i den indre Grundvirker og Grundgiver, ved hvis
indre Virkelser og Givelser enhver god oplivende og
frugtbringende Tanke og Følelse opstaar,« saaledes, at Individets
egen Stræben og den indre Givers Leven i det bestandigt
gaa i Et. Det er en Art Mystik, som reducerer de
dogmatiske Antagelser til noget forholdsvis Lidet. I Bibelen
finder han vel stadigt mange dybe og tiltalende Tanker, men
langt fra fuld Klarhed og Sandhed. Med al sin Dybde, hedder
det, giver den hellige Skrift ofte kun halve Tanker. Den
giver os f. Eks. ikke nogen klar og fyldestgørende Oplysning
om den hinsidige Tilværelse. Det kan da ikke hjælpe at
støtte den saa kaldte Aabenbarings Betydning paa den
Oplysning, den i denne Henseende skulde give os, især da det,
som Skriften tydeligt nok lærer derom, nemlig den evige
Fordømmelse af saa Mange, er nok til, at vor moralske
Følelse maa forkaste dens Autoritet. I Stedet for at holde
% --- p. 96 ---
\opage{96}sig til en saadan enkelt, overleveret Bog, holde man sig til
Naturens Bog, »den store jordiske, baade i det Legemlige
og i det Aandelige fremtrædende Livsfyldes store Bog«. Og
spørges der nu, om vi da ved vore naturlige Erfaringer kunne
finde Oplysning om Alt hvad vi kunne ønske at vide, særligt
om et andet Liv, saa svarer Sibbern, at hvad enten der er
en Udødelighed eller ikke, ville vi leve Livet som vi leve
det; vor moralske Overbevisning er ikke bygget paa nogen
Lære om det Hinsidige.

Den Tendens, som gaar igennem Kirkens Historie, til at
dogmatisere, det vil sige slaa den Sandhed, man mener at
have vundet, fast i bestemte, foreskrevne Formler, er i flere
Henseender fordærvelig. Man rykker derved Sandheden løs
fra dens levende Rod, dens Livsgrund, og gør en Mumie af
den, for at den saa i denne Mumieskikkelse kan blive
antagen af Alle. Denne Tendens ytrer sig maaske endnu stærkere
i den ortodokse Protestantisme end i Katolicismen. Man har
i Protestantismen sat en Dogmahellighed i Stedet for
Katolicismens Værkhellighed. I Stedet for paa denne Maade at
omhegne de Veje, ad hvilke Aanden virker, lade man den
frit Spil! Man trænge sig ikke ind i den Enkeltes
Helligdom med Polemik eller med Spot! Man lade Enhver beholde
sin Fetisch, det, hvori han har fundet sit aandelige Tilhold!
Men saaledes har Kristendommen ikke opfattet Sagen. Uagtet
den skulde være Kærlighedens Religion, har den dog i al
sin Tid langt mere opfostret og næret Anmasselsens,
Herskesygens og Hadets Aand end Menneskekærlighedens Aand.
Den har afskyet og bekæmpet Hedenskabets Guder og
Gudevæsen, uden at tænke paa, om der ikke ogsaa blandt
Hedningerne kunde findes saadanne Smaa, som man fremfor alt
ikke maatte forarge, og uden at undersøge, om ikke ogsaa
Flerguderi kunde have sin gode Berettigelse. Hvad det
gælder om, er, at Menneskekærligheden udfolder sig, og man
maa derfor raade Enhver til at blive ved sin Tro, naar den
er ham til Fortrøstning og til Fremme for det, som Pligt og
Menneskeomsorg kræve.

Inden for Kristendommen er man efter Sibberns Mening
i denne Henseende gaaet ad saa fordærvelige og forkerte
% --- p. 97 ---
\opage{97}Veje, at han lader de Fremtidseuropæere, han skildrer i
Bogen fra Aaret 2135, vægre sig ved at kalde sig Kristne.
»Skulde vi kunne ville bære Navn sammen med dem, som
i Kristendommens Navn og under dens Banner have fyldt
Aarhundreder med Anmasselser og Grueligheder?« --- Han
skelner mellem Kristelighed og Kristenskab. Det Kristelige
er den sande og sunde Kærne i Kristendommen, men til
Kristenskabet hører alt det dogmatiske og kirkelige Væsen.
Han slutter sig væsentligt til Nyrationalismen, hvilken han i
sin Moralfilosofi karakteriserer paa følgende Maade: »Den
vil have det kirkelige Kristendomsherredømme stødt bort
for i dets Sted at sætte den frie Gudhengivenheds Sindelag,
som jeg i Modsætning til Kristenskabet har kaldet
Kristeligheden, ladende Ordet Kristendom staa hen som ubestemt«. ---

Hvad der forekommer mig at være det mest
Betydningsfulde ved Sibberns religiøse Standpunkt, er den storartede
Tro paa Personlighedens Ret. Han har i langt fyldigere
Maal og paa langt naturligere Maade end S. Kierkegaard
gennemført den Sætning, at Subjektiviteten er Sandheden.
S. Kierkegaard gav denne Sætning en altfor snever
Anvendelse og Udførelse. Vi møde atter her Sibberns faste
Tro paa Livet i dets sporadiske Virken i Forbindelse med
hans varme Sympati for enhver personlig Livsytring. Det
er ikke et paa Forlegenhed og Skepticisme grundet Forslag
om Tolerance, vi her have for os; men en paa Grundlag af
en stort % PRINTED AS IS (for stor)
og dybt anlagt Verdensanskuelse opstillet Fordring.
Her er heller ingen ængstelig Søgen efter, om der dog ikke
skulde findes et Skjul, hvor vi kunne gemme vore Dogmer
saa godt, at Kritiken ikke skal kunne finde dem. Der er
tvert imod en Bestræbelse for at bringe saa stor
Overensstemmelse som muligt mellem Tro og Viden, mellem
Livsanskuelse og Verdensanskuelse. Og vi have set, at
væsentlige Træk af hans Verdensanskuelse finde Bekræftelse ved
de nyere Teorier og Synsmaader; derfor vil ogsaa hans
religiøse Standpunkt være af Interesse for Nutiden, afset fra den
store Interesse det har som Resultat af et sundt tænkende
og varmt følende Menneskes Forsken og Søgen. De religiøse
Spørgsmaal ere her hjemme i de senere Aar fra deres Side,
% --- p. 98 ---
\opage{98}der ikke kunne slutte sig til Ortodoksien, sikkert altfor meget
blevne behandlede kritisk og polemisk. Hos Sibbern er
Kritiken og Polemiken stedse hidført og baaren af den
positive Fylde, han virker for. Han kan heri være et Forbillede
ogsaa for dem, der komme til Resultater, som afvige fra hans.

\section*{V.\\  Sibberns politiske og sociale Anskuelser,}

Sibberns politiske og sociale Anskuelser bære i høj Grad
Præget af hans hele Personlighed. Dette berøver dem
ingenlunde deres Betydning og blivende Interesse, ti som en af
hans Hovedejendommeligheder have vi lært at kende den
levende Sans for og Kærlighed til Livet i dets Virkelighed
og Individualitet. Han søger overalt Indholdet og Kærnen,
og har han nogen Fejl, saa er det den at overse, at der i
denne Verden nu en Gang behøves visse bestemte Former
og Skaller, ofte haarde og usmagelige Skaller, selv for dets
værdifuldeste Indhold.

Hvad der saa ofte gjorde ham upraktisk, naar han
deltog i den politiske Strid, var ikke mindst det, at han
optraadte som psykologisk Sjælesørger. Han havde et skarpt
Blik for det Opirrende og Aandsfortærende i den politiske
Strid, og ikke mindst for, hvorledes der i Stridens Løb kan
opdynge sig Meget, som skjuler det, som det egentligt kommer
an paa. Den Form for politisk Virksomhed, som mest
tiltalte ham, var den, der svarede til Herakles’s Maade at
rense Augias’s Stald paa. Ligesom Herakles udførte sit
Værk ved at lede Strømmen ind og saa overlod til den at
skylle ud i alle Kroge, saaledes ansaa Sibbern det for sit
Kald, naar i Stridens Løb dels Polemik, dels Fortabelse i
Enkelthederne havde formørket Blikket hos de Fleste, da at
arbejde saa direkte som muligt paa at skaffe selve
Hovedsagen Indgang. Derved mente han ikke blot at arbejde
bedst for Sagen, men ogsaa at skaffe Sindene den Forfriskelse
og Hvile, hvortil de trængte. Godt var det og, mente han,
om man i Politiken kunde lade f. Eks. hver
halvtredsindstyvende Uge være en Jubeluge, hvor man lod Stridighederne
% --- p. 99 ---
\opage{99}hvile under fornøjeligt personligt Samkvem. Det havde, efter
Sibberns Mening, aldeles ikke nogen Hast i Politiken. Han
glædede sig i sine »Politiske Intelligensblade« (1835) over,
at de raadgivende Stænder først traadte sammen fire Aar
efter, at Forordningen om deres Indførelse var kommen. I
dette store Tidsmellemrum var der Plads for alsidig og
grundig Overvejelse. Derfor betragtede han ogsaa de
særskilte Stænderforsamlinger for Øerne og Jylland som en
heldig Indretning, da enhver Sag derved blev overvejet to
Gange. Senere glædede han sig over, at den frie
konstitutionelle Forfatning ikke kom før 1848, da Forholdene egnede
sig for den. Og han optraadte baade under Striden om
Helstatsændringen (1854) og under Striden om Grundlovens
Revision (1865) med det Raad, at stille Sagen i Bero eller
i Hvilestand i nogle Aar, for at den kunde overvejes og
bearbejdes i Stilhed, og for at de vigtigere Statsformaal, navnlig
Omsorgen for »Lavmandsfolkets« Velfærd og borgerlige Frihed,
kunde fremmes.

At man i Stridens Hede ikke var tilbøjelig til at lytte
til slige Raad, kan ikke undre. Og ligeledes følger det af
sig selv, at de store statsretlige Afgørelser, der ere
bestemmende for hele den Maade, hvorpaa, og de Betingelser,
hvorunder alle de vigtigste Statsanliggender i det Enkelte
behandles og afgøres, ikke kunne opsættes for længe, og at en
saadan Opsættelse i hvert Tilfælde neppe vil fri Sindene fra
Spænding og Uro.

Det er i det Hele karakteristisk for Sibberns politiske
Anskuelser, at han er saa omhyggelig for, at Overvejelsen
kan blive grundig og alsidig. Hans Politik hænger her
sammen med hans Filosofi. I hele Tilværelsen gaar jo efter
hans Lære Udviklingen frem gennem en stor Debat af Alt
mod Alt. Om den politiske Udvikling gælder det Samme.
En Statsforfatnings Godhed beror derfor paa, om den aabner
Adgang for alle Kræfter og Retninger i Folket til at komme
til Orde og virke medbestemmende i Udviklingen. Deri ser
han det Karakteristiske for den \emph{frie} Forfatning. Derimod er
det efter hans Mening af mere underordnet Betydning, hvor
den endelige Afgørelse, Viljesbestemmelsen efter Overvejelsen,
% --- p. 100 ---
\opage{100}ligger. Han anser det for heldigst, at den ligger hos Kongen.
Den frieste mest alsidige Overvejelse vilde da kunne forenes
med den mest energiske Afgørelse, den gennemførte
Decentralisation med den gennemførte Centralisation. Sibberns
Ideal af en Forfatning vedblev stedse at være »en
Eneraadighedskonge med en blot raadgivende Folkeformandsforsamling,
der skulde søge al sin Vægt og Styrke i sine Grundes Godhed
og Velbetænkthed«. Netop naar Folkerepræsentanterne ikke
faa Del i Afgørelsen, kan der herske den største Frihed i
den offentlige Drøftelse, og Folkets Skøn vil kunne faa en
langt større Indflydelse paa den endelige Bestemmelse, end
naar denne beror paa Afstemning i Repræsentanternes
Forsamling. Regeringsmagten skal staa under Stændernes Kontrol.
En Kontrol af Alle over for Alle er i enhver Forfatning det
sidste Middel, og de saakaldte statsretlige Garantier ere til
Syvende og Sidst alle af moralsk Natur, i en
»Fristatsforfatning«, hvor Repræsentationen har Stemmeret ved Afgørelsen,
saa vel som i en rent monarkisk Forfatning. Det er langt
fra, at en saadan fuldstændig Frihed i Meningstilkendegivelse
og offentlig Drøftelse skulde skade Monarkiet, at netop Intet
støtter det saakaldte absolute Monarki mere end en
republikansk Aand saa vel i Folket som i Regeringen. Folket føler
sig derved staaende sin Konge nær.

Ved en Konstitution forstaar man i Reglen et paa
Papiret nedskrevet Grundrids af en Statsorden. Men naturligere
er det at tage Ordet i dets medicinske Betydning, altsaa her
om Folkets og Statens virkelige Tilstand og Livsorden. Kun
en ringe Del af denne kan til enhver Tid være formuleret
og nedskreven. Og en saadan skreven Konstitution vil stedse
være et Baand paa Livet. »Jo flere Bestemmelser en
Grundlov indeholder, desto større Formynderskab har man tiltaget
sig over Fremtiden og dens Slægter, desto mindre er selve
Staten i Sandhed fri, desto mere er derimod Livet i Staten
og hele Nationaludviklingen, Nationalbestræbelsen og Nationens
friske Rørighed belemret med Baand, som let kunne blive til
besværlige Ulemper og Hildelser« (Dikaiosyne, 1843, p. 82).
En saadan skreven Konstitution var efter Sibberns Mening
lige saa fordærvelig for det politiske Liv, som en skreven og
% --- p. 101 ---
\opage{101}foreskreven Trosbekendelse for det religiøse Liv. Af alle
Former for Absolutisme ansaa han den, som overdroges til
et Stykke Papir, for den værste. ---

Det er en forunderlig Modsætning, Sibbern statuerer
mellem Overvejelse og Beslutning, og den er sikkert lige saa
lidt holdbar i Politiken som i Psykologien. Ethvert Motiv,
der kommer frem under Overvejelsen, faar ogsaa sin
bestemmende Indflydelse ved Beslutningen, og kunde vi helt
gennemskue, hvad der gaar for sig i vort Indre under Viljesakten,
vilde vi ofte finde, at Beslutningen staar som Afslutning af
en Proces, hvor det Tidligere helt igennem bestemmer det
Senere. Det er altid et Tegn paa indre Splid og Disharmoni,
naar Beslutningen bestemmes ved Noget, som ikke har kunnet
sejre under Overvejelsen, og det vil i Reglen være til det
Onde, naar det sker. Sibberns politiske Ideal vil paa
tilsvarende Maade føre til en Disharmoni, eller ogsaa til at
Kongemagten drager Slutningen af, hvad der i den
offentlige Debat er kommet frem. --- I første Tilfælde --- altsaa
naar Kongemagten og den offentlige Mening ere uenige ---
vil det Spørgsmaal opstaa, om i Reglen Kongen eller Folket
har rigtigst Blik paa Tingene. Antager man, at Kongen i
Reglen har Ret, hvor bliver saa egentligt »den republikanske
Aand« af? Saa kan Kongen i det Højeste være forpligtet
til at tage Forsigtighedshensyn, eller et pædagogisk Hensyn
til den offentlige Mening, men Debatten og Kontrollen af
Alt med Alt bliver illusorisk, og det bliver da vel dog
hensigtsmæssigst at stanse den ørkesløse Drøftelse. Frederik
den Sjettes berømte Sætning: »Vi alene vide, hvad der tjener
til Undersaatternes sande Vel«, i hans afslaaende Svar paa
et Andragende om udvidet Trykkefrihed, er den logiske
Konsekvens af denne Opfattelse. Antager man derimod, at
den offentlige Mening i Reglen har Ret, vil man omvendt i
Kongemagten faa en stadig Hindring for Fremskridt og en
Fare for Staten. --- Man vil da drives over til den anden
Mulighed, at Kongemagten skal udøve og fuldbyrde
Resultatet af den offentlige Debat. Men i saa Fald vil det være
hensigtsmæssigst, at dette Resultat udtrykkeligt formuleres og
slaas fast, og det kan kun ske ved en Afstemning, altsaa
% --- p. 102 ---
\opage{102}ved at Folkerepræsentanterne (eller Folket gennem sine
Repræsentanter) have Del i Statsmagten.

\emph{En} Sætning maa man give Sibbern Ret i, og det er den,
at alle statsretlige Garantier ere moralske. Det giver ingen
Sikkerhed, at Folkets valgte Repræsentanter faa Andel i
Statsmagten, naar det kun beror paa Regeringens egen Vilje,
om den skal tage noget Hensyn til deres Afstemning.
Dersom Sibbern havde oplevet Hundredaaret efter sin Fødsel,
vilde han have faaet en Erfaringsbekræftelse paa sin Sætning.

Men for ret at vurdere hans politiske Opfattelse, maa
man lægge Hovedvægten paa den republikanske Aand, han
fordrer skal herske i Monarkiet. Hvad han vil udelukke
og hindre, er Standsforskellenes Skærpelse og Uddannelse af
et Aristokrati. Han vilde hverken have et Adels- eller et
Embedsaristokrati, og han frygtede for, at i en
»Fristatsforfatning« vilde det blive »Velstandsfolket«, der kom til at
regere over »Lavmandsfolket«. Det havde efter hans Opfattelse
været de danske Enevoldskongers Bestræbelse at arbejde
paa Folkets Emancipation. Selve Regeringsforandringen 1660
betragtede han som en Frigørelse af Borgerstanden. Fra
dette Aar af regnede han derfor Friheden i Danmark. Især
virkede Frederik den Tredje og Kristian den Femte, og senere
igen Kronprins Frederiks Regering fra 1784 overordentligt
meget til Bedste for Lavmandsfolket og for almindelig Lighed
mellem Landets Borgere. Intet harmede Sibbern mere, end
naar han hørte de liberale Politikere tale om Tiden før 1848
som en Ufrihedstid. Han havde selv foretaget et grundigt
Studium af det 17. og 18. Aarhundredes danske Historie.
Endnu i sine »Samfundsbetragtninger« (Første Hefte 1865)
siger han: »Dersom jeg i min høje Alder turde tænke paa
Sligt, vilde jeg nok forsøge en historisk-filosofisk
Gennemgaaelse af Danmarks Historie fra 1536 af ..... Det er især for
dem, som deltage i Politiseringerne, at jeg ønskede Skildringer
af Danmarks indre \emph{Stats- og Folke-Udviklingshistorie}, skreven
saaledes, at de vigtigste Kilder, de kongelige Forordninger,
og hvad der blandt Reskripterne maa stilles sammen med
dem, retteligen blive gennemgranskede, og det Betydelige i
dem fremdraget.« Han fremhæver med Rette, at heri ligger
% --- p. 103 ---
\opage{103}Kontinuiteten i det danske Folks Historie. Han vilde højligt
have glædet sig ved Edv. Holms forskellige Arbejder over
Danmarks indre Udvikling under Enevælden og vilde have
fundet en Bekræftelse for sine Anskuelser i dem. Han har
her et grundigere historisk Blik end de nationalliberale
Politikere, som vare saa optagne af deres egen Statsvisdom, at
de saa lutter Mørke forud for deres egen Tid, ligesom de
da ogsaa have ment, at der vilde følge lutter Mørke efter
dem. Derimod saa han ikke, at Enevælden uddannede et
Embedsaristokrati, som paa sin Vis maatte blive en ny Hæmsko
for Udviklingen. Men deri saa de nationalliberale Politikere
ikke klarere end han. ---

Nu vilde Skæbnen, at Danmark alligevel skulde have
en Fristatsforfatning, og skønt Sibbern ikke ansaa denne
for den ideale Forfatning, saa han dog, at det ikke kunde
være anderledes. Han protesterede kun mod, at det var
Friheden, der 1848 blev skænket Folket. Friheden kan ikke
skænkes. Det var Magten, der blev delt mellem Konge og
Folk. At det kom dertil, skyldtes efter Sibberns Opfattelse
flere forskellige Omstændigheder. Den herskende politiske
Bevægelse i Europa forplantede sig ogsaa til os, og de ivrige
liberale Agitatorer (»Røropsmændene«, som Sibbern i sit Sprog
kaldte dem) vare ikke fornøjede med den gamle Forfatning.
Derhos var Frederik den Syvende ikke skikket til at hævde
Enevoldsmagten. »Man maa sige, at for ham var den nye
Statsindretning en stor Lykke; ti dels befriede den ham for
Regeringspligter, som han aldeles ikke var skabt til --- ej
heller havde søgt at uddanne sig til --- at kunne opfylde
blot nogenlunde, som Sligt af Mennesker kan fordres, dels
satte den ham i et saare godt og tillige menneskeligen frit
Forhold til Folket.« Kun Skade, at denne »lidet
kongeforstandige« Mand ikke forstod at bruge den Del af Magten,
han beholdt tilbage! (Meddelelser af et Skrift fra Aaret 2135.
Første Del, p. 645 og flere St.).

Folket var \emph{egnet} til den nye Forfatning. Under
Enevælden var Ligheden bleven fremmet, og Emancipationen
gennemført. Det var et stort Held, at den nye
Foratning % PRINTED AS IS (Wikisource corrects to "Forfatning")
 ikke kom, før den kunde grundlægges paa en
% --- p. 104 ---
\opage{104}meget bred Basis, med saa udstrakt en Valgret. Man
slap derved i Danmark for den Modsætning, som sætter
andre Lande i saa stor og stadig Uro: mellem de
Stemmeberettigede og Ikke-Stemmeberettigede. »Det saakaldte
Demokrati staar overalt bagved Stændersalene og vil ind;
lykkeligt det Land, som har aabnet Dørene helt for det, og
med det Samme aabnet en Tumle- og Kampplads ej blot
for det, men tillige for dets Modstandere og Modkæmpere.«
\emph{Egnet} var Folket, ikke modent; ti modent bliver et Folk kun
ved Hjælp af selve Forfatningen. Alt kunde være gaaet
godt, dersom kun ikke de nationalliberale Ledere havde
manglet den Modenhed, som maa fordres af dem, der ville
styre Statsskibet; Modenhed til at affatte Søkortet havde de
maaske nok« (Samfundsbetragtninger. Første Hefte, p. 18, 42 f.).

Hvor oprigtigt Sibbern havde sluttet sig til Aanden i
den nye Forfatning, da den historiske Udvikling havde ført
den med sig, fik han to Gange Lejlighed til at vise, under det
Ørstedske Ministerium og under Grundlovskampen 1865---1866.

Det Skrift, han udgav under det Ørstedske Ministeriums
Kamp med Rigsdagen (»Om og i Anledning af Kampen
mellem de tvende højeste Statsmyndigheder i Danmark«) er
af ikke liden Interesse netop i vore Dage. Man finder der
en Række af Klagepunkter, som paany ere komne frem her
hjemme under den nuværende politiske Strid. --- Sibbern
klager over, at Ministeriet har draget Kongen ind i Striden.
Den frie Adgang til Kongen staar ikke længer aaben,
saaledes som den i lang Tid havde været i Danmark. »Det
var,« siger Sibbern, »næppe hverken statsrigtigt eller
statsklogt, at den frie Adgang lukkede sig, da man vilde komme
med Adresser fra mange baade Kongens Person og Kongens
Kongeanseelse meget hengivne Mænd. Regeringens Formaal
fremmes neppe saaledes; langt bedre uden Tvivl paa modsat
Maade .... Man har derved rørt Folket paa et meget ømt
Sted, maaske især Folket i Jylland.« Det var jo netop hvad
Sibbern stedse havde fremhævet som Kongedømmets Opgave:
at staa i nær og levende Berøring med det hele Folk. Han
citerer i et andet Skrift med Begejstring den bekendte Sang
af »Kong Salomon og Jørgen Hattemager«:

% --- p. 105 ---
\opage{105}
\begin{verse}
Bestandig var der Vej i Herthas Lund \\
Fra Folkets Hytter op til Fyrstens Sale. \\
Den Mindste selv sig nærmed Folkets Drot \\
Og bragte did sit Hjertes Bønner \\
Og følte selv i Kongens Slot, \\
At han var en af Landets Sønner.
\end{verse}

Det var nu forbi, og under en Fristatsforfatning! Under
Enevælden vilde det, efter Sibberns Mening, ikke kunne gaa
saaledes!

Ved de mange Pressesager var Frygten for, at Regeringen
vilde Trykkefriheden til Livs, vaagnet. Embedsstandens Frihed
og Anseelse rokkedes ved adskillige Afskedigelser. Man
rokkede ved Rigsdagens Anseelse, og dog hører den lige saa
vel som Regeringen til de Myndigheder, danske Borgere ere
undergivne, og som de Kristne i Følge Romerbrevets 13.
Kapitel skulle adlyde! --- Sibbern gør endeligt opmærksom
paa, at alle lovlige Midler til at faa Ende paa Striden ikke
ere forsøgte, saa længe det ikke er prøvet, hvad Indflydelse
et Ministerskifte vilde faa\footnote{Hele denne Strid pinte særlig Sibbern, fordi han saa A. S. Ørsted,
hvem han saa højt beundrede, paa den anden Side. „Jeg kan ikke
sige Dem,“ skriver han til Wegener, „hvor stærkt dette Reskript
[som indeholdt en ministeriel Tilrettevisning til Wegener for hans
Optræden i Arvefølgesagen, uagtet Ministeriet forgæves havde
anklaget ham for Højesteret] har afficeret mig, allermest for Ørsteds
Skyld. Jeg har nu i 48 Aar staaet i det bedste Forhold til ham,
nydt meget Godt i hans Hus, holdt paa ham, troet paa ham, indtil
forrige Aar; og nu igen dette!“ --- Efter Ministeriets Afgang skriver
Sibbern (til Ørsteds Søster): „Han [A. S. Ørsted] er ved Ministeriets
Entledigelse kommen ud af en ikke god Stilling. Det vides
almindeligen, at meget af det, som skete, og som han maatte give sit Navn
til, ja stundom sætte sit Navn under, ingenlunde havde hans Bifald,
men var i Strid med hans Anskuelse. Men han var i Ministeriet
i Minoriteten og blev overstemt.“}.

En halv Snes Aar senere tog Sibbern Ordet mod et
»Grundlovsattentat« af en anden Art. Der udkom i Aaret
1865 tre forskellige politiske Smaaskrifter af ham, alle om
Grundlovskampen. Det oprørte ham at se, at man straks
efter den sørgelige Krig benyttede sig af juridiske
Spidsfindigheder til »ad en Omvej, der ser ud som en Krogvej, at an%
% --- p. 106 ---
\opage{106}taste Danmarks gode frie, folkelige Grundlov«. Det var for
ham kun »Finessemageri«, naar man ikke vilde indrømme, at
den Del af sin Magt, Rigsdagen i sin Tid havde afstaaet
for at gøre det muligt at gennemføre en Fællesforfatning,
der ogsaa omfattede Hertugdømmerne, --- at denne Magt efter
Hertugdømmernes Afstaaelse straks maatte falde tilbage til
Rigsdagen. Man raabte paa Lovlighed, men gik Ulovlighedens
Vej, fordi man havde sine bestemte Formaal, man vilde
opnaa: »Hvorfor vilde man, da Betingelserne for
Indskrænkningen af Grundloven af 1849 vare rent borte, … lade et
Spilleværkeri indtræde, som maatte være stødende for hele
det danske Folks sunde Sans? Uden Tvivl var det, fordi
man deri saa et Middel til at faa Landstingsreformen frem,
og saa maatte det øvrige af Komedien opføres med«
(Samfundsbetragtninger. Første Hefte, p. 70).

Der var her rokket ved den eneste Betingelse, under
hvilken Sibbern kunde sympatisere med en »Fristatsforfatning«,
nemlig at den byggedes paa et demokratisk Grundlag.
Landstingsreformen var en Reaktion mod den Udvikling af
Lighedsfølelsen, som især i det sidste Aarhundrede var gaaet
for sig. Man deler alle Landets Sønner i to Hovedklasser,
Velstaaende og Ikke-Velstaaende, og man ser ikke, at man
derved skærper Modsætningen mellem Folketing og
Landsting, saa at den nye Forfatning indeholder Spiren i sig til
lange Stridigheder. Der er i den opdynget et betydeligt
Stof til fortsat Agitation, og de tage fejl, der mene, at
Forfatningsforandringens Vedtagelse har bilagt Striden. --- Sibbern
stod i denne Sag ganske paa samme Standpunkt som
Tscherning og Grundtvig og sympatiserede med de Skridt, de gjorde
for at hævde deres Anskuelse. I det Hele er der ikke ringe
Lighed mellem disse tre Mænds politiske Ideer. Et
Kongedømme paa meget bred demokratisk Basis, med Udjevnen
af Standsforskellene, var alle tres Ideal. Under de skiftende
Forhold forfølger Sibbern dette Ideal med stor Sikkerhed og
sund Bedømmelse af Personligheder og Situationer. Hans
politiske Pjecer ere af ikke ringe psykologisk og historisk Interesse.
At de ikke virkede paa deres Samtid, laa dels i Forfatterens
idealistiske Opfattelse, dels i de Snurrigheder, hvoraf de vrimle,
% --- p. 107 ---
\opage{107}og som det var let at slaa sig til Ridder paa. Men for den, der ser
tilbage paa Begivenhederne, har det sin store Interesse at se,
hvorledes de opfattes og bedømmes af en sund, ærlig og med
varm Følelse begavet Aand, der tillige medbringer en ikke
almindelig historisk og filosofisk Sans. Jeg vil ikke gøre noget
Forsøg paa at angive, hvorledes Sibbern vilde have stillet sig,
hvis han havde oplevet den nuværende politiske Strid. Han
var en erklæret Modstander af Parlamentarismen; men han
havde ogsaa i sin Tid været en erklæret Modstander af
»Fristatsforfatningen« og dog lært at betragte denne som en
naturlig og nødvendig Konsekvens. Hans ægte demokratiske
Sindelag vilde have stillet ham langt anderledes i denne Strid
end de Mange, som ledes af Standshovmod og bureaukratiske
Fordomme. Han vilde smerteligt have følt, hvorledes de to
Elementer i hans Statsideal, Kongemagten og Demokratiet,
traadte i skarp Modsætning til hinanden; men han vilde ikke
have fundet nogen Udvej i et Godsejer- og Embedsaristokratis
Herredømme\footnote{For dem, der mene, at det paa Grund af Partistridighederne gaar
tilbage med alt Godt og Ædelt i Danmark, vil det være af Interesse
at vide, at det i hvert Tilfælde er en gammel Skade. Den 5. Januar
1825 skriver Sibbern til sin gamle Rektor, Brorson, paa Herlufsholm:
„I Sandhed, naar jeg tænker paa Dem, saa træder Billedet af en
Tid mig i Møde, da der var i mange Henseender et skønnere Liv i
Danmark end nu. Belevenhed og et tækkeligt Væsen var endnu
regnet for Dyd; Landets store Mænd respekterede Enhver; kom noget
ret fortrinligt og godt ud, saa fejrede ethvert Hjerte en Nationalfest . . .
Nu derimod er Alt saa uhyre kritiserende, at for denne kritiske
Syge næsten Intet kan komme op; saaledes at sætte sin Egoisme
mod hele Verden, forstod man ikke, da De var ung.“ --- I Aaret
1828 taler han (i sin „Psykologiske Patologi“) om „det Hang til at
nære Partiaand og at bevæges af Partiaand, hvilket er saa
fremherskende i vor Tid“. --- I et Brev fra December 1835 udtaler
Paludan-Müller (den senere Biskop) sin Harme over „Dagbladenes
Tone“ saa vel som over, at man bruger den nye Stænderforsamling
til at gøre Opposition, „som man overhovedet gør sig Flid for at
bringe Folket til at betragte alle bestaaende Forhold i en opponerende
Aand“. --- Om de Nationalliberale var det Sibberns Mening, at de
havde Meget tilfælles med Schleswigholsteinerne. „De vare oplærte
i samme Skole. Agtelse for Folkets Sjælelighed, Omsorg for den
jevne gode Folkeaands Bevaring havde de ikke. De gik saa vidt, at de reve ned paa Landets mest kundskabsrige, indsigtsfulde,
højtagtede Mænd, og undergrove derved selv Agtelsen for Lærdom
saaledes, at, da siden Lærdmændsfolkene, de saakaldte doktrinære,
kom for Had, kunde man sige til dem: I have jo selv arbejdet
derpaa af al Magt“ (Meddelelser af Indholdet af et Skrift fra Aaret
2135. Første Del, p. 695 f.).}. ---

% --- p. 108 ---
\opage{108}Alle de forskellige politiske og konstitutionelle
Spørgsmaal angik nu egentligt saare lidet det, som laa Sibbern
inderst paa Hjerte. Allerede i et Universitetsprogram fra
1849 hævder han, at det sociale Spørgsmaal var langt
vigtigere end det konstitutionelle, og at det var det, Statsmændene
først og fremmest skulde sætte sig ind i. Det konstitutionelle
Væsen har, mente han, Blændværk eller Halvhed for og bag,
medens Kommunismens Ide vel har Mørke for og Mørke bag,
men dog har Realitet til Grundlag og Realitet til Formaal.
Det Spørgsmaal, hvis Afgørelse Alt her kom an paa, var for
ham, om det hører til Statssamfundets Væsen at virke ved
ydre Magt, eller om dette ikke maa betragtes som noget
Foreløbigt, Noget, der gælder for Menneskeslægtens nuværende
Livstrin, men vil blive afløst af et idealt Statsliv, hvor der
vel findes et fast organiseret Samfund til Opnaaelse af fælles
menneskelige Formaal, men hvor den ydre, tvingende Magt
er falden bort, fordi den ikke behøves. Den paa Magt grundede
og med Magt virkende Stat fører os ikke ud over
Naturstanden. »Vore Stater ere saa langt fra at ophæve den Krig
af Alle mod Alle, som man har ment at herske i Naturstanden,
at Statsforfatningerne og de dermed forbundne Folkeforfatninger
jo tvert imod paa mange Maader have gjort saadan Krig af
Alle mod Alle stærkere og skarpere, hvassere, mere omgribende
og indgribende, ej alene udad til, men ogsaa i hver Stats og
Nations Indre midt under den saakaldte Fredstilstand. Man
er % PRINTED AS IS (Wikisource corrects to "ser")
 i denne vistnok ikke Blodsudgydelse, men desto mere
ser man, naar man har hele Europa for Øje, Kampen for
Erhvervet, ja Kampen endog for det simple Livsophold, at
føre til utallige Tærelser og Fordærvelser af mangfoldige
Strømme af Nervekraft, ja af hele moralske Tilværelser.«
Skulde denne Kamp om Magt eller Rigdom eller endog om
den blotte Eksistens være evig? Vil ikke en saa stor Fælles%
% --- p. 109 ---
\opage{109}følelse udvikle sig, at den ydre Magt kan falde bort og den
blinde, haardhudede Konkurrence kan afløses af en fornuftig
Ordning? Er det dog ikke Statens Opgave at varetage
Omsorgen for Alle? Kunde man ikke bestemme, at hvad den
Enkelte erhvervede ud over et vist Maal skulde afgives til
Statens fælles Driftskapital, saa at Staten for Alvor kunde
optræde som den fordelende Retfærdigheds Organ, som en
sand Menneskeorganisation, der ikke lader Tingene gaa deres
Gang under Kræfters og Drifters vilde Kamp, men ordner
og hjælper paa alle Omraader? ---

Det er karakteristisk for Sibbern, at han tillægger det
sociale Spørgsmaal større Betydning end det politiske. Vi
have ogsaa set, hvorledes han gentagne Gange opfordrer til
at stanse den politiske Strid for at sørge for »Lavmandsfolket«.
Dog ser han klart Sammenhængen mellem begge Spørgsmaal,
og det var ikke hans Mening, at man kunde springe lige
ind i en kommunistisk Ordning. Tvert imod, saa længe
Tvangsstaten bestaar, kan Kommunismen ikke virkeliggøres.
Sibbern saa hen til en Tid, hvor den europæiske
Menneskehed var kommen ud over sine nuværende Livsvilkaar og
kunde indrette sig paa en helt ny Maade. Denne
Fremtidsfantasi er det, han udfører i »Meddelelser af Indholdet af et
Skrift fra Aaret 2135«. Sibbern antyder selv, at denne Bog
maa betragtes som et Sidestykke til Platons »Stat«. Han
har altsaa villet give sin ideale Fantasi Luft ved at skildre,
hvorledes et Menneskesamfund vilde se ud under saa gunstige
indre og ydre Forhold som muligt. Bogen indeholder tillige
adskillige »Tilbageblik« paa det nittende Aarhundrede. I det
Foregaaende have vi benyttet den ved Fremstillingen af Sibberns
religiøse og politiske Anskuelser. Nu ville vi betragte det
sociale Utopia, han skildrer. ---

Henimod Slutningen af det nittende Aarhundrede steg
den europæiske Civilisations Onder og Skyggesider til en
Højde, der gjorde Livet til en Byrde. Al Livsglæde var
forbi; end ikke Musiken formaaede længere at oplive Nogen.
Pessimismen herskede. Den eneste Vederkvægelse, man fandt,
var at udtrykke sin Følelse af Elendighed og sin Fortvivlelse
med saa stærke Ord som muligt. En Fyrste forlod i Aaret
% --- p. 110 ---
\opage{110}1896 det ulykkelige Europa med de Ord: »Jeg kan ikke
udholde det længere, det Spil med ormstukne Sjæle om
ormstukne Goder; jeg har længe nok trællet i denne Trædemølle
og tærsket dette kærneløse Straa.« Og man foreslog, at der
over alle Rigsdags- og Regeringsbygninger skulde staa disse
Ord: »Kommer hid alle I, som ere livsfriske og glade, her
skulle I vorde mødige og besværede«. Men som Nøden,
Fortvivlelsen og Livstrætheden var paa sit Højeste, kom Frelsen.
En Dvaletilstand udbredte sig over den europæiske
Menneskehed, efter at først Døden paa hemmelighedsfuld Maade
havde raset gennem Landene. Der kom en lang Hviletid, i
hvilken alle Civilisationens Former saa godt som opløstes
og de faa Efterlevende vendte tilbage til Naturlivets simple
Former. Tillige mildnedes Følelserne; der kom Ro og
Menneskekærlighed i Sindene. Træthed var dog ikke den eneste
Aarsag til, at »Søvnens store Frelse« kom. Store Forandringer
i Naturen, underjordiske Dampes Frembrud o. s. v. virkede
med. Tiden syntes at være kommen for et Gennembrud baade
i det fysiske og i det aandelige Liv paa Jordkloden. Da
Genopvaagningens og Genoplivelsens Tid kom, kunde den
ungdomsfriske, genfødte Menneskehed begynde sin Tilværelse
under nye Vilkaar. Statsforfatning, Lovgivning,
Kirkeforfatning, Skolevæsen, Pengevæsen, Handelsvæsen, Ejendom,
Straffevæsen, Menneskeforskelle og Menneskeuligheder --- Alt
var man bleven befriet for, og man ønskede det ikke tilbage,
men følte Gru og Rædsel ved at læse derom eller se det
hos andre Folk. Man levede ud af den indre
Fællesskabsfølelse. Alle toge Del i legemligt og aandeligt Arbejde. Først
og fremmest gjaldt det om ikke at kue eller hæmme det
sjælelige Liv hos Nogen. --- Hver bragte til de offentlige Forraad
hvad han producerede mere end han selv behøvede; og deraf
tog de Alle, hvad de trængte til. Naar et Hus var bygget,
flyttede den ind, som mest havde Lyst. Regering behøvede
man ikke. Drengenes Optøjer var det Eneste, man havde
at værge sig imod. Lige saa lidt som man vilde have nogen
hellig Bog eller nogen Kirke mellem sin Følelse og den Livets
indre Kilde, man troede paa, lige saa lidt vilde man have
nogen Lovforfatning som kunstigt Gærde for Livet. Man
% --- p. 111 ---
\opage{111}mente, at al den Skarpsindighed og Eftertænksomhed, som
var bleven anvendt paa Retsvidenskab, Lovgivning og
Lovanvendelse, havde været til saare liden virkelig Fremme for
Aandslivet. Forbrydelsen havde man allerede længe betragtet
som en Sygdom, der saa snart den var godtgjort ved Dom,
skulde helbredes ved menneskekærlig Omsorg og Forbedring.
Straffeanstalterne bleve Lægeanstalter eller Sygehjem. Men
med Ophøret af Særejendom og af alle kuende og æggende
Institutioner hørte ogsaa Forbrydelserne op. ---

Bag det Fantastiske i disse Forestillinger og deres nærmere
Udførelse mærker man et varmt og menneskeligt Sind og en
stærk Harme over al Undertrykkelse og Forkuelse. Det var
Sorg og Lede ved det Virkelige, der drev den for menneskelig
Lykke saa bekymrede Tænker til at søge i en indbildt Verden,
hvad han ikke fandt i den virkelige. Ved Fortællingen om
Dødsperioden, om den store Søvn og Genopvaagningen har
han indrømmet, at han er ude af Stand til at vise, hvorledes
den kommunistiske Lyksalighedstilstand skulde kunne udvikle
sig ad naturlig Vej ud af de nærværende Samfundstilstande.
Der maa foregaa en indre Omdannelse af Menneskenaturen,
saaledes at Følelsen af Fællesskab fuldstændigt faar
Overherredømmet over Egoismen. Og selv da maatte der være
en lykkelig Harmoni mellem de enkelte Menneskers
Tilbøjeligheder, siden f. Eks. et nybygget Hus skulde bebos af den,
som havde mest Lyst dertil, og Enhver skulde kunne tage
af det offentlige Forraad hvad han trængte til. Hvorledes
mon man skal afgøre, hvem der har mest Lyst til Noget,
eller hvorledes vil man kunne faa de forskellige
Fornødenheder til at passe sammen? 

Det har dog neppe været Sibberns Mening at ville
henvise Alt hvad han udtaler i denne Bog til Fantasiens, evigt fra
Virkeligheden fjernede Land. Det ses allerede af, at han i
sit Universitetsprogram fra 1849 foreslaar, at Staten skal
samle sig en Driftskapital og optræde som fordelende
Myndighed. Han er her inde paa hvad man i vore Dage kalder
Statssocialisme. Den tyske Nationaløkonom Wagner har
videre udført, hvad Sibbern antyder i sit Program. Wagner
betragter det nemlig som et væsentlig Synspunkt og Øjemed
% --- p. 112 ---
\opage{112}for Skattevæsenet, at Staten gennem dette virker regulerende
paa Formuernes Fordeling i Landet. I Statssocialismen
fremtræder den fordelende Retfærdighed, Sibbern fordrer af Staten.
Men Tvangsstaten vilde han dog ikke betro dens Udøvelse,
og hvorledes faa vi en Stat, som ikke er en Tvangsstat? Og
hvorledes faa vi hos Statens Ledere den Indsigt, den
Menneskekundskab og den Menneskekærlighed, som maa være
til Stede, for at vi kunne tilstede dem at gribe fordelende
ind? Selv om Statssocialismen skulde have Fremtiden for
sig, vil Udviklingen i denne Retning kun kunne gaa meget
langsomt og med fornyet Prøvelse paa hvert enkelt Punkt.
Sibbern har jo selv lært os, at al Udvikling foregaar sporadisk,
ud fra forskellige individuelle Udgangspunkter. Saaledes vil det
vist vedblive at gaa, og det vilde allerede være en stor
Lykke og et stort Fremskridt, om vi kunde komme saa vidt,
at de forskellige Individer ikke støde for haardt sammen
under deres forskellige og ofte modsatte Bestræbelser. Mere
end at virke fremhjælpende og understøttende i denne Retning
formaar Staten neppe, saa længe den ikke er bleven noget
helt Andet end den nu er, --- og naar den er bleven det, er
det jo ikke sagt, at der er Noget at raade Bod paa. Sibberns
faste og urokkelige Tro paa Udviklingen kommer her igen.
Han var jo, som vi have set, sikker paa, at det Sporadiske
og Disharmoniske vilde forsvinde, og at den Absurditet og
Pinagtighed, Erfaringen frembyder, kun kommer af, at
Verdensudviklingen endnu er paa et af sine første Stadier. Denne
Tanke holdt ham oppe og gav hans menneskekærlige Sind
Ro. Den har vel ogsaa gjort, at han selv ikke strengt har
skelnet mellem Utopier og virkelige Idealer i sin Skildring
af Fremtidens Europa.

\streg{}

% --- p. 113 ---
\opage{113}
\essay[ALEXANDER HAMILTON OG DEN NORDAMERIKANSKE UNIONSFORFATNING]{ALEXANDER HAMILTON OG DEN NORDAMERIKANSKE UNIONSFORFATNING.}

\streg{}

\section*{I.}

Det er ofte blevet bemærket, at den historiske Opfattelse,
som i vort Aarhundrede er den herskende, staar i en skarp
Modsætning til den, der raadede i forrige Aarhundrede. Man
var i det 18. Aarhundrede tilbøjelig til at forklare historiske
Begivenheder, især Grundlæggelsen af Institutioner, der havde
vist sig at være betydningsfulde, af bevidst Overlæg og
Hensigt hos enkelte Personligheder. Og i naturlig Forbindelse
hermed nærede man en stor, næsten ubegrænset Tillid til,
hvad den ret »oplyste« menneskelige Fornuft kunde gøre for
at ordne Samfundsforholdene paa bedste Maade. Det Ædleste
og det Bedste, som det 18. Aarhundrede har frembragt, skyldes
denne Tillid. Den nu herskende historiske Betragtningsmaade
gaar derimod ud fra den Tanke, at alt Stort sker, eller i
hvert Tilfælde indledes, i Tavshed, i Ubevidsthed. Grundigere
Kendskab til Historien har vist, hvorledes sociale og politiske
Institutioner udvikle sig langsomt, og uden at der fra først
af behøver at være nogen Bevidsthed om deres Betydning
og deres fjernere Virkninger. Historien staar for os som
en stor ubevidst Udviklingsgang, ved hvilken de bevægende
Kræfter, som fra først af raade, kunne være helt andre end
de, der senere føre til at fastholde og paaskønne, hvad der
har udviklet sig. Den Betydning, en Institution faar paa
Højdepunktet af sin Udvikling, kan ikke antages at have
gjort sig gældende i deres Bevidsthed, som lagde den første
Grund til den. Og vi pleje at tilføje, at det er sundt, naar
Udviklingen saaledes gaar overvejende ubevidst for sig, lige%
% --- p. 114 ---
\opage{114}som Barnets Udvikling foregaar heldigst, naar man lægger saa
lidt Mærke til det som muligt, og først og fremmest naar det
lægger saa lidt Mærke til sig selv som muligt. Naar Kraften
og Formerne ere udviklede, saa er det Tid for Selvfølelsen og
Selvbevidstheden at vaagne, for at det Udviklede kan fattes og
fortsættes med fuld Forstaaelse af dets Værd. Men denne
grundigere og sundere historiske Opfattelse, ved hvilken vort
Aarhundrede staar saa højt over det forrige, har paa den anden Side
berøvet os dettes store Tillid til Tankens Magt paa det sociale
og politiske Omraade. Vi have set saa mange Eksempler paa,
at de snildest udtænkte Ordninger ikke bestaa, hvor den
naturlige Jordbund i Folkets Erfaringer og dets gennem
lange Tider grundfæstede Karakter mangler. Det er intet
Under, at man er bleven skeptisk over for den menneskelige
Aands Evne til med Bevidsthed at gribe ind i og paa gavnlig
Maade at ordne de sociale og politiske Forhold. --- En saadan
Tvivl, en Tvivl, der synes saa vel begrundet ved hele den
moderne Opfattelse af Historien, har noget Svækkende og
Trøstesløst ved sig. Er vi da henviste til at lade os glide
med Strømmen? Kunne vi lige saa godt lade være at tænke
og at ville?

Allerede en simpel Betragtning (som allerede er antydet
i det Foregaaende) viser, at vor Opfattelse let kan blive lige
saa ensidig som det 18. Aarhundredes paa sin Maade var.
Den bevidste Tanke og Vilje hjælper ganske vist ikke, naar
ingen naturlige Betingelser og Evner ere til Stede. De
kunne Intet bygge eller skabe ud af sig selv. Men naar
Tanken netop først og fremmest benyttes til at forstaa, hvilke
Betingelser der ere til Stede og hvortil de kunne benyttes,
--- hvilke Farer der true, og hvorledes de kunne undgaas
og forminskes, saa er den selv et vigtigt Led i Udviklingen,
et Led, som ikke kan undværes, naar Livsgangen skal være
stor og sund.

En interessant Bekræftelse heraf frembyder den
amerikanske Unionsforfatnings Grundlæggelse. Her blev med klar
Bevidsthed Konsekvensen dragen af en lang, i mindre Kredse
forløbende Udvikling. Slutstenen blev sat paa en Bygning,
som var vokset frem i det Skjulte. Med fremskuende Blik
% --- p. 115 ---
\opage{115}blev der lagt Midler til Rette for at forebygge Farer, om
hvis Alvor det følgende Aarhundredes Historie noksom har
vidnet. Og der var en bestemt Anelse om, at det
grundlagte Værk vilde komme til at afgive Ramme for en
storartet Udvikling, som vilde blive af Betydning ikke blot for
det enkelte Folk, men for hele Menneskeheden\footnote{Paa Konventet i Filadelfia, hvor Unionsforfatningen blev udarbejdet,
var især \emph{James Madison} virksom for dens Vedtagelse. Han stod
tilbage for Hamilton som Tænker, men havde mere Blik for de
Kompromis’er, som de politiske Forhold gjorde nødvendige.
Forhandlingerne paa Konventet holdtes hemmelige, men Madison
skrev om Natten en udførlig Beretning om dem, for --- som han
sagde --- „at bevare Materiale til en Konstitutions Historie, paa
hvilken \emph{et --- endog i sin Barndom stort --- Folks Lykke} beror, ja
paa hvilken maaske \emph{Frihedens Sag i hele Verden} beror“. (Curtis:
History of the Constitution of the United States. I. p. 427). Det
første Nummer af „The Federalist“ (27. Oktbr. 1787) aabnes med
følgende Udtalelse af \emph{Hamilton:} „Efter tilstrækkelig Erfaring om
den bestaaende Forbundsforfatnings Utilstrækkelighed opfordres I
nu, Medborgere, til at forhandle om en ny Forfatning for de
forenede Fristater i Amerika. Emnets Betydning er klar. Det
indeslutter i sig ikke mindre end \emph{Unionens Eksistens}, dens enkelte
Deles Sikkerhed og Velfærd, \emph{det Riges Skæbne som i mange
Henseender er det mærkeligste i Verden}. Det er ofte blevet bemærket,
at det synes at være blevet forbeholdt for dette Lands Befolkning
ved deres Adfærd og Eksempel at afgøre \emph{det vigtige Spørgsmaal,
om menneskelige Samfund virkeligt ere i Stand til at oprette en god
Regering ved Overvejelse og Valg, eller om de for evigt ere bestemte
til at afhænge af Tilfældet og Magten}, hvad deres politiske
Forfatninger angaar .... Denne Forestilling maa, ved at forene
Menneskekærlighedens Tilskyndelser med Fædrelandskærlighedens, forøge den
Bekymring, som alle betænksomme og gode Mennesker maa føle
for, hvad Udfaldet bliver.“}. Ogsaa
denne forudskuende Bevidsthed har Historien bekræftet. Den
har vist os det i nyere Tid enestaaende Syn af en demokratisk
Stormagt. Der er nu gaaet hundrede Aar, siden
Unionsforfatningen traadte i Kraft; og det vil sikkert have sin
Interesse at dvæle ved de Tanker, og de ydre Forhold, ved
hvis Vekselvirkning den blev til.

Æren for dette store Værk tilkommer først og fremmest
en hel Kreds af udmærkede Mænd, som vare besjælede af
% --- p. 116 ---
\opage{116}de nye Tanker, det 18. Aarhundredes engelske og franske
Literatur havde sat i Bevægelse, og som ikke mindre vare
fortrolige med de nordamerikanske Staters Forhold og med
den Folkeaand, som der havde udviklet sig, --- Mænd, der
i Krig og Fred havde kæmpet for Frihedens og Selvstyrelsens
Sag. Men blandt dem er der igen én Mand, hvis Navn paa
en særlig Maade er knyttet til Unionsforfatningen, fordi han
var den første, der undfangede den Grundtanke, paa hvilken
den beror, og med størst Klarhed og Energi begrundede og
udviklede denne Tanke, ligesom han ogsaa med stor Kraft
virkede for dens Udførelse. Denne Mand var \emph{Alexander
Hamilton}, en af de største Statsmænd og tillige en af de
største politiske Filosofer, som nogen Sinde har levet. Det
er paa Tide, at hans Navn drages frem for den almindelige
Bevidsthed. Han fortjener en Plads ved Siden af Washington
og Franklin i Nordamerikas Historie, og i de politiske Ideers
Historie fortjener han en Plads i første Række.

Ogsaa ved sin Personlighed er han værdig til en mere
almindelig Opmærksomhed, end der hidtil er bleven ham til
Del. De betydningsfulde Tanker, han kæmpede for, fandt i
ham ikke blot en af de mest geniale og energiske, men
tillige en af de ædleste Forkæmpere, som nogen Sinde en stor
og god Sag har haft. Han satte sit Liv ind for sin Sags
Renhed. Under al Døgnets Strid, under den Forplumring af alle
Forhold, al den Løgn og Smaalighed, som Nutidens politiske
Liv saa mange Steder frembyder, er det styrkende at vende
Blikket mod et Syn, som Historien ingenlunde frembyder
paa mange af sine Blade: en stor og praktisk Idé, fattet med
klar Bevidsthed og gennemført af en ren og energisk Vilje,
hvem det lykkes i det afgørende Øjeblik at faa Folket til at
lytte til sin Røst.

\section*{II.}

De nordamerikanske Fristater skylde
Uafhængighedsdriften og Frihedsfølelsen deres Oprindelse. De grundedes af
Mennesker, som vilde undfly det religiøse og politiske Tyranni
i Europa. Ny-England (de nordligste af Fristaterne) blev
f. Eks. til ved de saakaldte »Pilgrimsfædres« Udvandring
% --- p. 117 ---
\opage{117}(1620). En Flok Puritanere var 1612 udvandret fra England
til Holland, men besluttede senere at drage over til den nye
Verden for at kunne leve efter deres egen Overbevisning.
Under store Møjsommeligheder og Farer fæstede man Bo der.

Den nye Kolonis Forfatning var i den første Tid aldeles
uafhængig af Englands. Alle Menighedens Medlemmer samledes
fra Tid til anden for at holde Raad, og man valgte sig selv
en Guvernør, der med nogle faa Raadgivere ledede de løbende
Forretninger. Da senere Folkemængden tiltog, indrettede man
(1639) en Repræsentation, som skulde give Love. Hvad der
mest udpræget fremtraadte i Ny-England, udviklede sig i
beslægtede Former i de andre Kolonier. Den engelske Regering
søgte forgæves at kue denne saa naturligt opstaaede
Selvstyrelse. Gennem Deltagelsen i den indre Styrelse af de
offentlige Anliggender paa de forskellige Trin udviklede der sig
en Almensans og en Borgeraand, som paa den Tid var
enestaaende i Verden. Selv de mest fremskredne Aander i
Gammel-England forstod ikke ret hvad der gik for sig der
ovre. Man ser dette f. Eks. af den Forfatning, som John
Locke udkastede for Karolina. Vel blev flere af Udkastets
mere frisindede Bestemmelser ændrede af de Lorder, hvem
Kolonien var overdragen af Karl II; men selv afset herfra
stred Udkastet mod den hele Retning, Udviklingen havde
taget der ovre, og det kom derfor aldrig til virkelig Anvendelse.

Ligesom Koloniernes Borgere saaledes havde Lejlighed til
at øve sig i indre Selvstyrelse, saaledes frembød de hyppige
Kampe med Indianere og Franskmænd Anledning til at erfare
Betydningen af Sammenslutning og Enhed mellem de enkelte
Stater. Allerede 1643 sluttede saaledes fem Kolonier sig
sammen som »the united Colonies of New-England« til et
offensivt og defensivt Forbund. Medens enhver Koloni
beholdt sin Selvstyrelse, skulde alle Spørgsmaal om Krig og
Fred og i det Hele alle Sager af fælles Interesse afgøres af
en Forsamling, der bestod af to Udsendinge fra hver Koloni.
Et Aarhundrede senere (1754) udkastede Benjamin Franklin
en Plan til et Forbund mellem alle Kolonierne under Ledelse
af en af den engelske Regering udnævnt Generalguvernør.
Men denne Plan blev forkastet i Kolonierne, fordi man mente,
% --- p. 118 ---
\opage{118}at den indeholdt for store Forrettigheder for Kronen, og af
det engelske Ministerium, fordi den var for demokratisk.
Franklin mente (som det ses af hans Selvbiografi), at hvis
hans Plan var bleven fulgt, vilde Kolonierne have været
stærke nok til at forsvare sig selv; den engelske Regering
havde da manglet ethvert Paaskud til at beskatte dem, og
den blodige Kamp mellem Kolonierne og England vilde være
undgaaet. »Men,« tilføjer den gamle Vismand, »saadanne
Misgreb ere intet Nyt; Historien er fuld af Staters og Fyrsters
Vildfarelser. Se dig omkring i Verden! Hvor Faa forstaa
deres eget Vel eller følge det, naar de forstaa det! De, der
staa ved Statens Ror, have for det Meste Meget at gøre og
gøre sig ikke den Ulejlighed at overveje og udføre nye
Planer. De bedste offentlige Forholdsregler antages derfor
sjeldent paa Grund af forudgaaende Visdom, men paatvinges
af Omstændighederne.«

Saaledes havde det amerikanske Folk baade Lejlighed
til at gøre Erfaringer og øve sine Kræfter ved selvstændig
Deltagelse i indre og ydre Anliggenders Ordning. Det havde
faaet den bedste Opdragelse, et Folk kan faa, den, som
bestaar i at handle selv, Ansigt til Ansigt med de virkelige
Forhold, uden at en arvelig Regerings Formynderskab eller
et selvklogt Bureaukrati i enhver Sag træder imellem.

Den amerikanske Revolution fik derfor ogsaa en anden
Karakter end den franske. Man havde i Amerika udøvet
Friheden, før man officielt fik den eller tog den. Gennem
lang historisk Udvikling var der bragt en solid Underbygning
til Veje, som kunde holde, selv om Spiret toges bort. Med
den engelske Styrelse i Amerika faldt ikke enhver Styrelse;
man havde lært at styre sig selv. Kontinuiteten blev ikke
afbrudt. Folket var ikke vant til at faa sin Parole fra London,
og det gjorde da Intet, at Forbindelsen afbrødes.
Vanskeligheden ved Unionsforfatningens Tilblivelse beroede tvert
imod paa, at man var saa vant til næsten ubegrænset
Selvstyrelse i de mindre Kredse, at man følte en Uvilje ved at
oprette en Centralmyndighed. Men Livet opløstes ikke, fordi
den øverste Autoritet stødtes bort. Anderledes i Frankrig.
Her var Centralisationen trængt i den Grad igennem, at Alt
% --- p. 119 ---
\opage{119}tilsidst hang i én Traad, saa at det maatte falde til Jorden,
naar denne ene Traad overskares. Kontinuiteten i Folkets
Liv blev paa én Gang afbrudt og er aldrig bleven fuldt
genoprettet siden. Udviklingen af det enevældige Kongedømme
havde ført til at trykke alle Klasser ned og berøve alle Kredse
ethvert Initiativ. Intet Fællesskab føltes. Kongen og hans
Embedsmænd maatte gribe ind overalt, hvis ikke Alt skulde
gaa i Staa. I en Indberetning til Ludvig XVI skrev Turgot:
»Deres Majestæt er nødt til at afgøre Alt selv eller ved Deres
Embedsmænd. Man venter Deres udtrykkelige Ordrer, før
man bidrager til det offentlige Vel, før man agter Andres Ret,
undertiden endog før man handler i sine egne Anliggender«\footnote{Smlg. \emph{Tocqueville:} L’ancien régime et la revolution. 7. ed. p. 158.}.
Ved Revolutionen beholdt man den centraliserede
Styrelsesmaade, og de jakobinske Magthavere saa vel som senere den
militære Usurpator fandt derved deres Vej banet.

Hvad Amerikanerne kæmpede for i deres Frihedskrig,
var den gammel-engelske Frihed, det Princip, at et Folk ikke
kan beskattes, naar det ikke selv gennem sine Repræsentanter
har bevilget Skatten. Dette erkendtes ogsaa offentligt af
mange Englændere. Den ældre Pitt erklærede saaledes i
Overhuset (20. Januar 1775): »Den Aand, som nu gør
Modstand mod Beskatningen i Amerika, er den samme, som forhen
modsatte sig Laan, Privilegier og Skibsafgift i England,
den samme Aand, som fik hele England til at rejse sig og
ved »Bill of Right« hævdede den engelske Konstitution, det
samme Princip, som fastslog den store, fundamentale og
væsentlige Grundsætning, paa hvilken vore Friheder hvile: at ingen
engelsk Undersaat kan beskattes uden efter eget Samtykke.
Whiggismens hæderfulde Aand besjæler tre Millioner
Mennesker i Amerika, som foretrække Fattigdom med Frihed og
ville dø under Forsvaret for deres Rettigheder som Mænd,
som frie Mænd.« --- Selv i deres Revolution havde
Amerikanerne historisk Grund under Fødderne. Vel beraabte de
sig i deres berømte Uafhængighedserklæring paa almindelige
og oprindelige Menneskerettigheder, og man beskyldte
Erklæringens Forfatter, Jefferson, for at have skrevet den ud
% --- p. 120 ---
\opage{120}af Locke’s Skrift »On Government«. Men de amerikanske
Ledere havde en anden Forestilling om deres
»Menneskerettigheder« end de franske Revolutionsmænd. De mente ikke, at
Alt først maatte raseres, for at man saa paa bar Grund
kunde opføre en ny Bygning efter visse abstrakte Principer.
De »oprindelige« Principer betragtede man som de inderste
Motiver og Tendenser i selve den historiske Udvikling, og
hvad man krævede, var blot, at denne Udvikling ikke skulde
stanses ved noget Magtbud.

Frihedskrigens Historie ville vi her ikke dvæle ved.
Gennem skiftende Sejre og Uheld førte den til et Punkt, hvor
Englands Kræfter vare udtømte. Trods sine uhyre Udgifter,
trods de Grusomheder og de til Dels forræderske Midler, som
anvendtes, maatte England give tabt. --- Amerikanerne stode
ved Maalet. Men for at blive et Folk, er det ikke nok at
afkaste det ydre Aag. Store indre Vanskeligheder vare endnu
at overvinde, og med Rette har en amerikansk Forfatter
nyligt kaldt Tiden fra Fredsslutningen til Unionsforfatningens
Vedtagelse (1783---1789) den mest kritiske Periode i Amerikas
Historie.

\section*{III.}

Medens det engelske Herredømme varede, havde
Kolonierne ikke haft nogen stadig Forbindelse med hverandre.
Enhver af dem havde havt at gøre med den engelske
Regering og kæmpet for at hævde sin Selvstyrelse overfor den;
men afset fra enkelte Forsøg paa Sammenslutning var
Unionstanken ikke meget levende. Lige saa lidt som der var
nogen stadig politisk Forbindelse, var der nogen synderlig
Handelsforbindelse mellem Kolonierne indbyrdes. Nordamerika
var væsentligt et agerdyrkende Land. Rejserne fra den ene
Koloni til den anden vare desuden besværlige og sjældne;
medens nu Tusender dagligt færdes paa Rejse mellem Boston
og New-York, vare to Postkareter den Gang tilstrækkelige til
at rumme alle de Rejsende. Efter Pacifikbanens Anlæg kan
man rejse hurtigere fra Boston til de nordvestligste Stater
end man for hundrede Aar siden kunde rejse fra Boston til
Filadelfia.

% --- p. 121 ---
\opage{121}Under Frihedskrigen havde der (ligesom tidligere under
Kampene med Indianere og Franskmænd) været følt Savn af
en fast Centralregering. En Kongres, bestaaende af
Delegerede med Fuldmagt fra de forskellige Stater, var
sammentraadt. Denne Kongres havde udstedt
Uafhængighedserklæringen og ledet Sagernes Gang under Krigen og ved
Fredsslutningen. Men Enkeltstaterne havde stadigt følt en vis
Skinsyge overfor Kongressen, og man var bange for at give den
for stor Myndighed. Kongressen kunde nok fatte Beslutninger,
men ikke udføre dem. Den kunde kun \emph{henstille} til de enkelte
Staters Regeringer, om de vilde udføre dem! Egentlig
Regeringsmagt kendtes kun for de enkelte Staters, ikke for hele
Forbundets Vedkommende. I hver enkelt Stat var der en
fri Regeringsform, som skrev sig fra Koloniernes
Grundlæggelse og som havde udviklet sig under stadig Kamp med
de fra England sendte kongelige Guvernører. De i hver
enkelt Stat valgte lovgivende Forsamlinger vare i Folkets
Øjne de eneste Myndigheder, som havde Ret til at paalægge
Skatter. Det kom derfor an paa dem, om de af Kongressen
ønskede Bidrag til Krigens Førelse bleve udredede. Hver
Gang en Bevilling skulde gives, maatte de alle spørges, og
ofte gaves slet intet Svar. Da Kongressen i Aaret 1781
til Gældens Betaling foreslog en Indførselstold paa 5 p. c.,
vakte det stor Harme i de enkelte Stater. Man spurgte,
hvorfor man havde modsat sig Stempelakten og Teskatten,
naar der nu alligevel blev paalagt Skatter af en Magt
udenfor Staten (Enkeltstaten)! Kun formedelst sin store
Taalmodighed og Uegennyttighed formaaede Washington at
udholde de pinagtige Forhold, i hvilke han ofte bragtes af Mangel
paa Penge. Han udtalte senere, at Krigen vilde være endt
tidligere og have kostet mindre Blod og Penge, hvis Unionens
Regering havde haft Ret til at paalægge Skatter. Den 1. Januar
1777, ganske kort efter Washingtons berømte Overgang over
Delaware og Overfaldet paa Trenton (Juledag 1776), maatte
Robert Morris, som under Krigen havde den ubehagelige
Post at sørge for de fælles Finanser, gaa fra Dør til Dør i
New-York for at indsamle Penge til Felttogets Fortsættelse.

Værre blev det endnu efter Krigen, da den Begejstring,
% --- p. 122 ---
\opage{122}som stedse paany flammede op og gjorde det muligt at ende
Krigen paa en saa heldig Maade, var forbi. Man vendte sig
nu hver til sin Gerning, og Sansen for de store og fælles
Anliggender tabte sig. Man var jo paa den engelske Tid
vant til ikke at have med Noget at gøre, som laa udenfor
Koloniens Grænser. Der gjorde sig saa ofte en kortsynet
Egenkærlighed gældende, som vægrede sig ved at opgive Noget
af Enkeltstatens Suverænitet til Bedste for Unionen. Dog
rørte der sig tillige hos Mange, til Dels blandt de mest
fremragende Mænd, en principiel Overbevisning, som frygtede
for, at Oprettelsen af en fast Centralmagt skulde kue Livet
i de mindre Kredse og tilsidst endog kue det frie,
selvstændige Liv hos de enkelte Personligheder. Saaledes betragtede
General George Clinton, hvem New-York ni Aar i Træk
valgte til sin Guvernør, de enkelte Staters Suverænitet som
nødvendig Betingelse for Folkets Frihed, og ansaa dem, der
virkede for en stærk Unionsregering, for »Despotiets
Talsmænd«. En anden anset amerikansk Statsmand, Thomas
Jefferson, som dog billigede den senere Unionsforfatnings
Vedtagelse, skriver i sin Selvbiografi: »Det er ikke ved
Konsolidering eller Koncentration af Magten, men ved dens
Fordeling, at god Regering bevirkes. Var dette store Land
ikke allerede delt i Stater, maatte en saadan Deling foretages,
for at hver Stat for sig; kunde besørge hvad der direkte
vedkommer den, og hvad den kan gøre langt bedre end en fjern
Myndighed. Enhver Stat er igen delt i Grevskaber, af hvilke
ethvert skal sørge for hvad der ligger indenfor dets Grænser;
ethvert Grevskab atter i Kredse eller Distrikter til
Varetagelsen af mindre Enkeltheder, og ethvert Distrikt i Gaarde,
der hver styres af sin individuelle Besidder. Naar vi fra
Washington [Byen W.] skulde have Ordrer om, naar vi skulle
saa og høste, vilde vi snart mangle Brød. Kun ved denne
fra det Almindelige til det Specielle gradevis nedad stigende
Deling af Arbejdet kunne de menneskelige Anliggender
besørges til Alles Held og Lykke«\footnote{\emph{Th. Jeffersons} Works. I. p. 82.}. Jefferson var saa ivrig
Demokrat, at han efter eget Sigende ikke var Ven af nogen
% --- p. 123 ---
\opage{123}energisk Regering nogensteds i Verden. Efter de indre
Uroligheder, der i Begyndelsen af Firserne fandt Sted i
Massachusetts, skrev han til en Ven, at saadanne Uordener vare
nødvendige, naar der skulde herske offentlig Frihed. Selv
om Friheden skulde udarte til Tøjlesløshed, maatte der drages
snevre Grænser for Regeringens Magt. Enhver enkelt Stat
burde have udelukkende Myndighed i sine egne Anliggender,
særligt Beskatningsret\footnote{Smlgn. \emph{Curtis}: History of the Constitution of the United States.
II. p. 507.}. --- Naar saadanne Anskuelser
hyldedes af fremragende Mænd og havde Rod i de historiske
Forhold, var det intet Under, om de Erfaringer, man havde
gjort om Sammenslutningens og Unionens Nødvendighed ikke
fik tilstrækkelig Paaagtelse. Der syntes at foreligge et
\emph{Enten}---\emph{Eller}: enten Enkeltstaternes Suverænitet eller Unionens
Suverænitet, og af disse to tilsyneladende uforenelige
Modsætninger havde foreløbigt den første Folkestemningen for sig.

Overfor Udlandet var man ved Fredsslutningen optraadt
som én Stat. Man havde paadraget sig Forpligtelser overfor
England; men de enkelte Stater tog Forholdsregler, som
aldeles stred mod disse Forpligtelser. Saaledes bleve de, der
havde sluttet sig til England under Krigen, og som i Reglen
kaldtes Toryer, forfulgte og mistede deres Ejendomme, uagtet
disse vare garanterede dem ved Fredstraktaten. Kongressen
manglede Midler til at gennemføre sine egne Forpligtelser.
--- Det gjaldt om at bringe den udenlandske Handel paa Fode.
Men de fremmede Magter vægrede sig ved at afslutte
Handelstraktater med et Folk, hvor der ikke var nogen Myndighed,
som kunde skaffe sig Anseelse. Udlandet spekulerede endog
i Enkeltstaternes indbyrdes stridende Handelsinteresser. Der
opstod en formelig Told- og Handelskrig mellem Nabostater
indbyrdes. --- Endvidere skulde Hæren hjemsendes, og Officerer
og Soldater skulde have deres Tilgodehavende. Men Kongressen
kunde ikke skaffe de fornødne Midler til Veje. Oprørske
Soldater samlede sig da udenfor Kongressens Lokale i
Filadelfia og udstødte truende og haanende Udraab mod
Kongresmedlemmerne. Guvernøren i Pennsylvanien vilde ikke
% --- p. 124 ---
\opage{124}samle Militsen, fordi det kun var Kongressen, som
fornærmedes! --- Ogsaa mod Enkeltstaternes Regeringer skete flere
Steder Oprør, og samtidigt truede det med en Indianerkrig.

Det var da intet Under, at man baade i Udlandet og
Indlandet begyndte at nære Tvivl om Nordamerikas Fremtid.
De udenlandske Magter, ikke blot England, men ogsaa Frankrig
og Spanien, der i Krigen have kæmpet paa Nordamerikas
Side, vilde helst se Landet delt. Eller man mente, at den
eneste Redning laa i Indførelsen af den monarkiske
Forfatning, hvad enten det nu blev en fransk eller en engelsk Prins,
som skulde over for at holde sammen paa Landet.

Og det var heller intet Under, at der blandt de
fremragende Mænd i Nordamerika herskede Mismod og
Betænkelighed i de nærmeste Aar efter Fredslutningen. Man følte
hvad der manglede: en stærk Centralmagt; men Opgaven var,
hvorledes den skulde oprettes, og hvorledes den kunde
forenes med Folkets demokratiske Tankegang og med
Selvstyrelsen. Det gjaldt, om det af Kongressen vedtagne
Valgsprog: E pluribus unum! skulde blive til Virkelighed. Og
det gjaldt, om der i Nordamerika skulde danne sig en ny
Stormagt, hvis folkelige Institutioner kunde danne et Hjemsted
for Frihed og Selvstyrelse, eller om dets frie Stater skulde
gaa til Grunde i indbyrdes Kiv, hver optagen af sine begrænsede
Interesser, et let Bytte for fremmede Magters Intriger.

\section*{IV.}

En af de Mænd, som bidrog mest til dette Spørgsmaals
lykkelige Afgørelse, var \emph{Alexander Hamilton}.

Han var født paa den vestindiske Ø Nevis den 11. Januar
1757. Hans fædrene Familje var skotsk; han % PRINTED AS IS (for hans)
mødrene Slægt
var en fransk Huguenotfamilje. Hans Moder, som skal have
været en begavet og højsindet Kvinde, døde tidligt, og da
hans Faders Anliggender gik tilbage, toge nogle af hans
mødrene Slægtninge sig af ham og lod ham opdrage paa
St. Croix. Her gik han i Skole til han var tolv Aar gammel,
og har vel kun faaet de aller nødvendigste Skolekundskaber.
Derefter blev han anbragt paa et Handelskontor paa samme
Ø, og viste sig her saa dygtig, at Forretningens Chef over%
% --- p. 125 ---
\opage{125}lod den trettenaarige Knøs Ledelsen af alle Sagerne, medens
han selv af Helbredshensyn gik paa en længere Rejse. Man
har bevaret Forretningsbreve fra ham fra denne Tid, som
bære Præget af tidligt udviklet Klogskab og Energi. Sin
Fritid anvendte han til Læsning, især af nationaløkonomiske,
historiske og filosofiske Skrifter, og lagde derved Grunden
til de flersidige Kundskaber, han fik Brug for i sit senere
Liv. Hans Tragten gik allerede nu ud over de snevre
Forhold, under hvilke han levede. I et Brev fra Aar 1769 skriver
han til en Ven: »Jeg maa bekende, at Ærgerrighed er det
Fremherskende hos mig, saa at jeg foragter en Kontorists
ringe Stilling, hvortil min Skæbne fordømmer mig, og gerne
vilde sætte mit Liv, skønt ikke min Karakter, paa Spil for
at vinde en højere Stilling. Jeg véd vel, at min Ungdom
udelukker mig fra ethvert Haab om øjeblikkelig Forfremmelse,
ikke heller ønsker jeg en saadan; men jeg vil berede Vejen
for Fremtiden. Jeg er, som du ser, ikke Filosof, og man
vil sige, at jeg bygger Luftkasteller; min Taabelighed faar
mig til at skamme mig, og jeg beder dig, ikke at omtale
den. Dog har man set saadanne Planer føre til Maalet, naar
Ophavsmanden var udholdende. Jeg vil slutte med at sige,
at jeg ønsker, der vilde komme Krig.« Sjeldent udtales vel
den Unges Ærgerrighed paa en saa klar og bestemt Maade
og tillige med en saa tydelig Bevidsthed om, hvad det koster
at naa Maalet. Intet bestemt Maal staar for ham, men han
føler sin Kraft. Han slaar med Vingerne, skønt han ikke
véd, hvor Flugten skal gaa hen.

Guvernøren paa St. Croix blev opmærksom paa ham i
Anledning af en Beskrivelse af et Jordskælv, han havde givet
i en Avis. Der bleve nu sørget for, at han kunde fortsætte
sin saa energisk begyndte Uddannelse. Han kom til
New-York og studerede ved Kings College. Da Frihedskrigen
udbrød, bleve de unge Studenter naturligvis stærkt grebne af
Bevægelsen. Paa et Møde, som holdtes udenfor Byen den
6. Juli 1774 for at øve et Tryk paa de ledende Mænd og
faa dem til at opgive deres Tøven, optraadte Hamilton første
Gang som Taler. Den syttenaarige Mands fyrige og klare
Veltalenhed fængslede Forsamlingen. Man vedtog, at New-York
% --- p. 126 ---
\opage{126}skulde følge Bostons Eksempel og afbryde al Handelsforbindelse
med England. Hermed havde Hamilton betraadt den politiske
Bane. I den nærmeste Fremtid forfattede han flere
Smaaskrifter imod en Gejstlig, der havde forsvaret det engelske
Ministeriums Fremgangsmaade. Han hævdede her borgerlig
Frihed som stemmende med Menneskets Natur og som
nødvendig for Samfundets Velfærd. Det engelske Parlaments
Ret skriver sig fra dets Vælgere og er derfor begrænset til
Storbritanien. En engelsk Vælger har Intet at sige i Amerika.
Et vilkaarligt Regimente føres der overalt, hvor et Folk
regeres ved Love, i hvis Tilblivelse det ikke selv har haft
nogen Del. Storbritanien er kun et frit Land, fordi dets
Indbyggere have Del i Lovgivningen. Vel maa der i ethvert
borgerligt Samfund være en højeste Magt; men derfor
behøver ikke en Del af Riget at herske over de andre Dele\footnote{Disse Smaaskrifter findes i 2. Del af Hamiltons samlede Værker.}.
--- Denne Pennefejde fik en ejendommelig Afslutning. Nogle
ivrige Frihedsmænd bortførte Hamiltons Modstander, Præsten
Seabury, som Fange og overfaldt det Trykkeri, fra hvilket
hans Skrifter vare udgaaede. Dette vakte stor Harme, ikke
blot hos det engelske Parti, men ogsaa hos Oppositionen.
Hamilton sluttede sig selv til Militsen, da den forfulgte dem,
som vilde ende Debatten paa saadan Vis. Men iøvrigt var
nu den fredelige Forhandlings Tid forbi. Det første Blod var
allerede flydt ved Lexington og Bunkers Hill (April og Juni 1775),
og Kampen var snart i fuld Gang. Hamilton og flere andre
Studerende fra Kings College vare med at erobre nogle
Kanoner fra et Fort udenfor Byen. De deltog i de Frivilliges
Øvelser, og ivrige militære Studier afløste for Hamilton de
tidligere fredelige Studier. Efter at have underkastet sig en
Eksamen blev han udnævnt til Kaptajn ved et Batteri, som
Staten New-York udrustede. Den unge Kaptajn (ikke tyve
Aar gammel!) tildrog sig snart ved sin Iver og Dygtighed
de kommanderende Generalers, særligt Washingtons,
Opmærksomhed. Han blev Oberstlieutenant, og Washington gjorde
ham til sin Adjutant og Sekretær. I Hæren gik han under
Navn af »den lille Løve«; Washington kaldte ham gærne
% --- p. 127 ---
\opage{127}»my boy«. Det nøje Forhold, som nu var knyttet mellem
de to udmærkede Mænd, varede hele Livet igennem, skønt
et lille Sammenstød gjorde, at Hamilton i Aaret 1781 forlod
Washingtons Stab for igen at overtage en aktiv Kommando.

Hamiltons Stilling i Washingtons Nærhed gav ham
Lejlighed til at følge de militære og de politiske Forholds
Udvikling paa nært Hold og fra et ophøjet Stade. Han fik et
dybt Indtryk af Nødvendigheden af en fast Ordning af Staternes
fælles Anliggender. Ofte maatte --- under de ydre og de
indre Vanskeligheder --- den Mulighed komme frem for ham,
at Udfaldet ikke vilde blive godt. I et Brev til sin Forlovede
skrev han paa denne Tid, at hvis det skulde gaa galt, vilde
de bosætte sig i Sveitz; kun i et frit Land kunde de tænke
sig at bo.

Midt under Krigens Sysler og Anstrengelser drejede
hans Tanker sig om, \emph{hvorledes} Unionen mellem Staterne kunde
ordnes og befæstes. I et Brev til Duane, Medlem af
Kongressen for New-York, af 3. September 1780, udvikler han,
at Kongressen mangler Magt paa Grund af Enkeltstaternes
af misforstaaet Frihedsaand udsprungne Skinsyge. Kongressen
er bleven afhængig af Enkeltstaterne i Stedet for at staa over
dem. Det er en Selvmodsigelse, at den har udøvet nogle af
de vigtigste Suverænitetsrettigheder: Udstedelsen af
Uafhængighedserklæringen, Oprettelsen af Hær og Flaade,
Udsendelsen af Mønt, Alliancen med Frankrig o. s. v. --- og saa
dog ikke har besiddet Myndighed til at skaffe det Fornødne
for at føre Krigen til Ende! Hvis Kongressen ikke selv vil
tiltage sig den Myndighed, som aabenbart i Grunden allerede
er overdragen den i Følge Forholdenes Natur, maa der
sammenkaldes en Forsamling af Udsendinge fra alle Staterne, som
med besluttende Myndighed kan vedtage en ny Ordning,
særligt Oprettelsen af en kraftig udøvende Magt, og Folket
maa beredes til at gaa ind paa dens Afgørelse ved vækkende
populære Skrifter. Unionsregeringen skulde da have
Suverænitet med Hensyn til Krigsvæsen, Handelspolitik, Toldvæsen
og Møntvæsen.

Her blev for første Gang den Tanke udtalt, paa hvilken
Unionsforfatningen kom til at hvile. Det Geniale i dette
% --- p. 128 ---
\opage{128}Brev, som Hamilton efter eget Sigende skrev i sit Telt midt
i en Hær, der snarere kunde kaldes en Pøbelflok, en Hær, som
manglede Klæder, Proviant, Penge, Moral og Disciplin, ---
det Geniale ligger i Tanken om, at Unionen grundes paa
\emph{hele Folkets}, ikke blot paa \emph{Staternes}, Tilslutning. Elementerne
i den nordamerikanske Unionsstat skulde være de samme som
Elementerne i Enkeltstaterne. »Den nationale Magts Strøm
skal«, som Hamilton et andet Sted udtrykker sig, »flyde
umiddelbart fra al lovlig Autoritets Kilde«. I hver Borgers Sind
skal Unionen slaa Rod. Det var denne Tanke, der førte til
at grunde Unionens, lige som Enkeltstaternes Suverænitet
paa den almindelige Valgret, og som førte til Dannelsen af
en \emph{Forbundsstat} i Stedet for det tidligere \emph{Statsforbund}.
Enkeltstaten skulde ikke være Mellemled mellem Individet og
Unionen, men Individet skulde som Borger leve umiddelbart i
begge, ligesom det lever umiddelbart baade i Familje,
Kommune og Stat. --- Denne Tanke betegner en stor politisk
Opdagelse. Hverken Oldtid, Middelalder eller den senere
Tid kendte noget ganske tilsvarende. Siden er den sveiziske
Forbundsforfatning (af 1848) og det tyske Riges Forfatning
(af 1867) grundede paa samme Princip\footnote{\emph{Tocqueville} erklærede i sin Bog om „Demokratiet i Amerika“, at
Sproget manglede et Ord til at betegne den ejendommelige Blanding
af national og føderativ Forfatning, som Unionsregeringen hviler
paa. --- \emph{Bluntschli} har (i sin „Politik als Wissenschaft“ 1876)
sammenstillet den med den sveiziske og tyske Forfatning, efterat allerede
tidligere en sveizisk Forfatter \emph{(Rüttimann)} havde anstillet en
Sammenligning mellem nordamerikanske og sveiziske Statsforfatninger.}.

Men der var endnu et strengt Arbejde at gøre, inden
Tanken traadte ud i Virkeligheden. Aaret efter det omtalte
Brevs Affattelse deltog Hamilton i Stormen paa Yorktown,
en af de smukkeste Bedrifter i Krigen. Han anførte højre
Fløj og var første Mand i den fjendtlige Skanse. Kort efter
traadte han ud af Hæren og nedsatte sig som Advokat. En
kort Tid var han Medlem af Kongressen, men han traadte
ud af den, fordi han ikke formaaede at hævde sin Værdighed.
Nu kom (efter Freden) de sørgelige Aar, hvor Mathed,
Egennytte og Kortsynethed afløste Krigstidens Energi og Begejstring.
% --- p. 129 ---
\opage{129}»Jo mere jeg erfarer«, sagde Hamilton i denne Tid, »des
mere Grund finder jeg for Alle, som elske dette Land, til at
græde over dets Blindhed«. Forgæves havde han allerede
før Freden (1781---82) i en Række Smaaskrifter, som udkom
under Titel »The Continentalist«, søgt at paavirke den
offentlige Mening. Og forgæves udtalte Washington sig, efter
Opfordring af Hamilton, i samme Retning i en Rundskrivelse
til Staternes Guvernører, da han nedlagde Kommandoen. Som
Hamilton skrev i et Brev til Washington: de centrifugale
Kræfter vare stærkere end de centripetale.

Dog var det ikke uden Betydning, at Tanken var kommen
frem. Den maatte vente paa, at Betingelserne for dens
Udførelse kom til Stede; men naar Opmærksomheden for en Tanke
er vakt, opdages mulige Udgangspunkter for dens
Virkeliggørelse lettere. Et Udgangspunkt fandtes nu i de materielle
Interesser, og i en Sammenhæng, hvor man ikke skulde have
anet det. Efter Fredsslutningen beskæftigede Washington
sig med Planer om en Forbindelse mellem Ohiodalen og
Østkystens Floddale, idet han klart saa Betydningen deraf for
Forbindelsen mellem Østen og Vesten, især da Missisippis
Udløb endnu var i Spaniens Besiddelse. Udsendinge fra de i
en saadan Kanalforbindelse interesserede Stater samledes paa
Mount Vernon, Washingtons Herresæde, i Begyndelsen af
Aaret 1785. Han benyttede Lejligheden til at indlede en
Raadslagning angaaende Told- og Møntspørgsmaal. Disse
Forhandlinger førte til, at alle Stater bleve enige om at sende
Delegerede til Annapolis 1786 for at drøfte en ny, fælles
Ordning af Handelslovgivningen. Mødet blev ikke
tilstrækkeligt fuldtalligt; men de Tilstedeværende besluttede at rette
en af Hamilton forfattet Opfordring til alle Staterne om at
sende Delegerede til en Forsamling i Filadelfia i Maj 1787,
som skulde udarbejde Forslag til en saadan Forandring i
Unionsforfatningen, ved hvilken denne kunde blive i Stand til
at fyldestgøre de Krav, Forholdene stillede til den. Saaledes
kom Sagen da i Gang.

Konventet i Filadelfia sammentraadte den 14. Maj 1787.
Nordamerikas berømteste Mænd mødte som Delegerede fra
deres Stater. Washington blev valgt til Præsident. Foruden
% --- p. 130 ---
\opage{130}ham vare bl. a. Benjamin Franklin, James Madison og Hamilton
til Stede. Modsætningerne traadte her skarpt overfor
hinanden. Medens Hamilton stillede et Forslag, der trykkede
Enkeltstaterne stærkt ned og indrettede en kraftig udøvende
Magt, fordredes fra den modsatte Side Unionsregeringens
Myndighed indskrænket til et Minimum. Af New-Yorks
Delegerede var Hamilton den Eneste, der var stemt for en
virkelig Unionsforfatning; hans to Kolleger modsatte sig den
og forlode Konventet, da de saa, at det blev til Alvor med
den. James Madison viste sig her som den praktiske
Statsmand, der lemper sig efter de virkelige Forhold. Kun gennem
forskellige Kompromiser kunde det lykkes at faa et Udkast
vedtaget. Af disse Kompromiser kan nævnes det, ved hvilket
det bestemtes, at medens Kongressens Repræsentanthus skulde
vælges af Folket, skulde dens Senat bestaa af to Delegerede
fra hver Stat. Senatet udtrykker altsaa alle Enkeltstaternes
Ligeberettigelse, medens Repræsentanthuset fremgaar af Folket
i det Hele, uden Hensyn til Enkeltstaterne. Desuden kom
man overens om, at Slavehandel kunde forbydes af Unionen
fra 1808 af, medens Slaveriet vedblev at bestaa og endog
skaffede Slavestaterne større politisk Magt, idet en Slave
regnedes for \textsuperscript{3}/\textsubscript{5} Menneske ved Fastsættelsen af hvor mange
Repræsentanter enhver Enkeltstat skulde sende. Ved denne
sidste Bestemmelse kom man foreløbigt ud over det
vanskelige Slavespørgsmaal, som allerede nu delte Nord- og
Sydstater i to Partier, som havde et helt forskelligt Syn.
Stridsspørgsmaalet blev udsat\footnote{En Hovedgrund til Udsættelsen var vistnok ogsaa den Forventning,
man nærede om, at Slaveriet snart vilde falde bort af sig selv.
Man anede ikke, at der forestod et uhyre Opsving af
Bomuldsdyrkningen, foranlediget ved Opfindelsen af Maskiner i den engelske
Bomuldsindustri. Der sendtes paa den Tid endnu slet ikke Bomuld
fra Sydstaterne til England, og der dyrkedes i det Hele ikke megen
Bomuld. Med Bomuldsdyrkningens Opsving bleve Sydstaternes
Politikere snart fanatiske Modstandere af Negeremancipationen.
Smlgn. \emph{John Fiske:} The critical Period of Amercian % PRINTED AS IS (for American)
 History.
London 1888. p. 266 f. --- Et interessant Eksempel paa, hvorledes
økonomiske Forhold faa Indflydelse paa etiske Anskuelser!}, indtil det kunde finde sin Afgørelse
gennem en klar og energisk offentlig Mening. Og selv da
% --- p. 131 ---
\opage{131}kom jo Afgørelsen kun efter en haard Kamp, der truede med
at sprænge Unionen --- og maaske havde sprængt den, hvis
Grunden ikke var bleven lagt paa en saa klog og solid Maade.

Mange havde været tilbøjelige til at gaa endnu mere paa
Akkord, fordi de ikke troede, at Folkestemningen var for saa
megen Centralisation af Magten, som det endeligt vedtagne
Udkast fastslog. Da der faldt Ytringer i denne Retning,
rejste Washington sig fra sin Formandsstol og sagde: »Det
er kun altfor sandsynligt, at ingen Plan, vi foreslaa, vil blive
antagen. Maaske maa endnu en haard Kamp udholdes. Men
hvorledes kunne vi senere forsvare vort Værk, dersom vi,
for at behage Folket, tilbyde, hvad vi selv misbillige?«

Efter lange Debatter og skarpe Sammenstød vedtoges
Forfatningsudkastet. Konventet havde da siddet i fire Maaneder.
Det var et højtideligt Øjeblik, da Udkastet skulde underskrives.
Da Washington greb Pennen, sagde han: »Dersom denne
udmærkede Konstitution forkastes af Staterne, vil der neppe
være Lejlighed til at sætte en anden i dens Sted; --- den
næste vil blive skreven med Blod!« Men den gamle Benjamin
Franklin havde lysere Forhaabninger. Idet han pegede paa
en Sol, som var afbildet over Formandsstolen, sagde han:
»Medens vi sad og forhandlede, var jeg ofte i Tvivl, om det
var en opgaaende eller nedgaaende Sol dér; nu er jeg vis
paa, at det er en opgaaende Sol!« ---

Endnu var der dog et stort Arbejde at gøre. Udkastet
skulde vedtages af folkevalgte Forsamlinger i de enkelte
Stater. Her var flere Steder en meget stærk Modstand
at overvinde. Hamilton traadte nu igen i Skranken for
sin Yndlingstanke. Sammen med Madison og Jay (en
amerikansk Retslærd, der havde udmærket sig under
Fredsforhandlingerne i Versailles) udgav han fra Oktober 1787 til
August 1788 en Række Flyveskrifter under fælles Titel af
»The Federalist«, i hvilke alle Forfatningens Bestemmelser
orklaredes og begrundedes, dels ved almindelige, fra den % PRINTED AS IS: „orklaredes“
menneskelige Natur udgaaende Betragtninger, dels ved historiske
Eksempler, dels ved Henblik paa de faktiske Forhold. Hamilton
har selv forfattet den aller største Del af dette Skrift, som
man har kaldet Republikanismens Bibel, og som fortjener
% --- p. 132 ---
\opage{132}en Plads ved Siden af Aristoteles’, Lockes og Montesquieus
Skrifter. Det udmærker sig endog fremfor disse derved, at
det staar den praktiske Virkelighed nærmere, --- udgaar fra
den og igen vil føre over i den. Den forholder sig til
Unionsforfatningen som Motiverne forholde sig til et Lovudkast.
Udkastet dækkede vel ikke Hamiltons oprindelige Tanke;
men hans Patriotisme og hans Virkelighedssans fik ham til
at sætte sig saa levende ind i det vedtagne Udkast, at han
nu kunde optræde som dets vigtigste Talsmand. I hans egen
Stat New-York var der en heftig Kamp. Paa det af
Befolkningen valgte Konvent, som sammentraadte i Juni 1788, vare
i Begyndelsen to Tredjedele af de Delegerede imod Udkastet.
Ved sin Begejstring og sin Veltalenhed fik Hamilton dog et
Flertal vundet for Sagen, og den 26. Juli blev Udkastet
vedtaget. Ved New-Yorks Tilslutning var Forfatningen sikret,
skønt Nordkarolina først tiltraadte 1789 og Rhode Island
først 1790. I Foraaret 1789 traadte Forfatningen i Kraft,
og Washington blev valgt til Unionens første Præsident.

Det store Tog, som efter Udkastets Vedtagelse drog
gennem New-Yorks Gader, bragte særligt Hamilton sin Hylding.
Hans Billede fandtes paa Fanerne, dels alene, dels sammen
med Washingtons. Det var en velfortjent Hæder for den
Mand, hvis Tanke og hvis Ord havde udrettet saa store Ting.

\section*{V.}

Inden vi følge Hamilton videre paa hans Bane, ville vi
dvæle lidt ved det mærkelige Skrift, hvis Hovedforfatter han
var, og fremhæve dets Hovedtanke.

Den politiske Opdagelse, Hamilton allerede 1780 udtalte,
var Unionens Grundlæggelse paa folkeligt Grundlag, i Stedet
for et statsligt Grundlag. I noget mere abstrakt Form kan
den udtrykkes saaledes: der fandtes en grundigere Form for
\emph{Forbindelsen af Enhed og Frihed}, end der hidtil var fundet.
I Oldtidens Republiker var der borgerlig og politisk Frihed;
men det syntes, som om Historien havde godtgjort, at denne
Frihed kun kunde bestaa med et begrænset Landomraade.
Saa snart et stort Rige skulde dannes, maatte Friheden
forsvinde for Enheden. Ja, selv i de antike Republiker var den
% --- p. 133 ---
\opage{133}individuelle Frihed paa mange Maader stærkt beskaaret, for
at Statens Enhed ikke skulde lide. \emph{Montesquieu} lærte
udtrykkeligt, at den republikanske Forfatning kun var mulig ved et
begrænset Landomraade. Kun dér kunde den Enhedsfølelse,
den umiddelbare Opgaaen i de fælles Interesser opstaa, som
den republikanske Regeringsform forudsætter. Alle Borgere
kunde tage Del i Sagernes Afgørelse\footnote{Esprit des Lois. III, 3, kfr. IX, 1.}. \emph{Frederik den Store},
ytrede i Aaret 1782 til den engelske Gesandt i Berlin, at han
var overbevist om, at den amerikanske Union ikke kunde
bestaa længe i sin republikanske Form; allerede Landets store
Udstrækning vilde være en Hindring, ti en republikansk
Regering har aldrig eksisteret længe, hvor Territoriet ikke
var begrænset og koncentreret\footnote{\emph{Bancroft:} History of the Constitution of the United States. I p. 71.}. Endnu i vort Aarhundrede
er det ofte af europæiske Forfattere blevet udtalt, at kun
Monarkiet kunde være Form for en omfattende Statsorganisme\footnote{Se f. Eks. \emph{Hegel:} Philosophie der Geschichte. 3. Aufl. p. 379. ---
\emph{J. N. Madvig:} Blik paa Oldtidens Statsforfatninger.
(Universitetsprogram). København 1840. p. 12 og 21.}.

Det var en Indvending, som Unionens Modstandere ikke
undlod at benytte sig af. Frihedens Sag var for disse knyttet
til Enkeltstaternes absolute Suverænitet. Hamilton svarer,
at Montesquieu kun kendte Oldtidens smaa Republiker, og
at Historien endnu ikke havde vist nogen Anvendelse af
Repræsentationssystemet i et stort Land med republikansk
Forfatning. Oldtidens Republiker forudsatte alle Borgernes
\emph{personlige} Deltagelse i Sagernes Afgørelse og kendte kun
Repræsentation i meget indskrænket Omfang. Hvis man altsaa
vil beraabe sig paa Montesquieu, maa man være saa
konsekvent at forlange flere af de største nordamerikanske
Enkeltstater delte; ti de er for store til at være Republiker i den
Betydning, i hvilken Montesquieu tog dette Ord. Tillige overse
de, at denne udtrykkeligt anbefaler et Forbund mellem frie
Stater som eneste Middel til at udstrække den republikanske
Forfatning over et større Omraade. Et saadant Forbund er
nu prøvet i Amerika og har vist sig utilstrækkeligt. Det har
vist sig, at hvis det er Alvor med at give Unionen Styrke
% --- p. 134 ---
\opage{134}og Varighed, maa Forbundsmagten ikke baseres paa den gode
Vilje hos Enkeltstaternes Regeringer, men den maa kunne
udøve sin Suverænitet umiddelbart over alle Borgere, lige
saa vel som Enkeltstaterne gøre det. Det gælder blot om at
fordele Rettighederne mellem Unionen og Enkeltstaterne, saa
at intet Sammenstød bliver muligt. Men Unionens
Myndighed maa tilsidst udspringe af samme Kilde som
Enkeltstatsregeringernes Myndighed: fra selve Folkets Vilje. Enheden
maa udspringe af selve Friheden\footnote{I det tyske Rige er den almindelige Valgret indført som nødvendigt
Middel til at styrke Rigets Magt overfor Enkeltstaterne. Smlgn.
\emph{Bismarck:} Reden. III. p. 194. Men medens man i Amerika har
ladet Enheden bestaa, for saa vidt den var forenelig med Friheden,
er man i Tyskland gaaet den omvendte Vej, idet Friheden kun
faar Lov til at bestaa, for saa vidt den er forenelig med den
stramme Enhed, som det nye Riges Bestaaen og Magt synes at fordre.}; saa vil den ikke vise sig
at være fjendtlig mod denne, og saa vil den være stærk nok
til at kunne lade de mindre Kredse bestaa uden at absorbere
dem. Et stort Samfund vil meget lettere lade Selvstyrelsen
blomstre end et lille. Et enkelt Parti kan ikke saa let
tilrive sig en udelukkende og stadig Magt i et stort som i et
lille Land, naar Institutionerne blot ere tilstrækkeligt frie.
Kun Retfærdigheden og det almene Vel ville kunne samle
meget store Majoriteter i et Land af stor Udstrækning.

En anden Hovedtanke, som fremtræder i »The Federalist«,
er, at en konservativ Tendens i Forfatningen er meget vel
forenelig med frie og republikanske Institutioner. Ligesom
den amerikanske Unionsforfatning er karakteriseret ved
Forbindelsen af \emph{Enhed og Frihed}, saaledes er den ejendommelig
ved Forbindelsen af \emph{Fasthed og Frihed}, og har netop i den
nyeste Tid vakt Opmærksomhed hos konservative Forfattere.

Unionsforfatningen hviler paa Folkets Suverænitet.
Hamilton lagde stor Vægt paa, at det i Indledningen til
Forfatningen kom til at hedde: »Vi, de forenede Fristaters Folk,
anordne og oprette denne Forfatning«. Forfatningen var for
saa vidt grundet ikke blot paa Folkets Magt, men ogsaa paa
Tillid til den Maade, hvorpaa Folket vilde bruge denne Magt.
»Ligesom der«, siger Hamilton, »gives en Fordærvelse i
% --- p. 135 ---
\opage{135}Menneskenaturen, som gør Forsigtighed og Mistro nødvendig,
saaledes er der andre Egenskaber i den menneskelige Natur,
som retfærdiggøre en vis Grad af Agtelse og Tillid. Den
republikanske Forfatning forudsætter Tilstedeværelsen af disse
sidste Egenskaber i en højere Grad end nogen anden
Forfatning. Dersom de Skildringer, som nogle Forfattere have
givet, vare sande Gengivelser af den menneskelige Karakter,
maatte man drage den Slutning, at der ikke er tilstrækkelig
Dyd hos Menneskene til Selvstyrelse, og at kun Despotiets
Lænker kunne hindre dem i at ødelægge og fortære
hverandre«\footnote{At Hamilton selv ikke just havde høstet gode Erfaringer om
Menneskene, ses af en karakteristisk Udtalelse af ham i et Brev til hans
Ven Meade (i Aaret 1782): „Intet skal afbryde vort Venskab. Mit
Venskab for dig er bygget paa det solide Grundlag af en fuld
Overbevisning om, at du fortjener det, og at det er gensidigt; og det
er saa meget des fastere grundet, som du har faa Medbejlere.
Erfaring er en stadig Kommentar over den menneskelige Slægts
Uværdighed, og de faa Undtagelser, vi finde, have saa meget des større Ret
til at skattes, jo sjeldnere de ere. Jeg kender Faa, som ere
agtværdige, Færre endnu, som fortjene at elskes; og naar jeg møder
Nogen af sidstnævnte Art, staar det ikke i min Magt at holde mine
Følelser tilbage.“ (Works I. p. 298).}. Det republikanske Princip hviler altsaa paa en
Tillid til Folket. »Men det fordrer ikke ubetinget Eftergiven
for enhver opblussende Lidenskab, eller enhver forbigaaende
Tilskyndelse, som maatte røre sig hos Folket ved
Kunstgreb af saadanne Mennesker, der smigre dets Fordomme
for at forraade dets Interesser. Det er en rigtig Iagttagelse,
at Folket i Almindelighed \emph{tilsigter} det offentlige Vel. I selve
dets Vildfarelser viser dette sig ofte. Men dets sunde Sans
vilde foragte den Smigrer, der vilde paastaa, at det stedse
\emph{tænker rigtigt om Midlerne} til at fremme det offentlige Vel.
Det véd af Erfaring, at det undertiden farer vild; og et Under
er det, at det farer saa sjeldent vild, som det gør, skønt det
stadigt er udsat for Snyltegæsters og falske Angiveres Kunster,
for de Ærgerriges, de Havesyges og de Hensynsløses Snarer,
for deres List, som besidde dets Tillid uden at fortjene den.
Naar der kommer Tilfælde, i hvilke Folkets Interesser stride
mod dets Tilbøjeligheder, er det en Pligt for dem, som Folket
% --- p. 136 ---
\opage{136}har udvalgt, at værne om dets Interesser, at modsætte sig
den øjeblikkelige Vildfarelse for at give det Tid og
Lejlighed til rolig og besindig Overvejelse«\footnote{The Federalist. Philadelphia 1882. p. 427, 533.}.

Af denne Betragtning udledes Nødvendigheden af
konservative Elementer i en republikansk Forfatning. De vigtigste
af disse Elementer ere følgende.

1) En udøvende Myndighed, som kan optræde med Kraft
og Fasthed baade udadtil og indadtil, er den vigtigste
Betingelse. Kraft og Fasthed hos den udøvende Magt beror
igen paa, at Magten lægges i én Mands Haand, at den har
en vis Varighed, og at selve den Magt, den raader over, er
tilstrækkelig stor. Naar Republikens Præsident vælges af
Folket (indirekte, gennem Valgmænd, der kunne raadslaa om
den Værdigste) paa fire Aar og kan drages til Regnskab
for sin Embedsførelse, saa ville de to Betingelser være
forenede, som synes Mange aldeles uforenelige, nemlig udstrakt
Myndighed og Athængighed % PRINTED AS IS (for Afhængighed)
 af Folket. (Dog havde Hamilton
oprindeligt helst villet have en paa Livstid valgt Præsident,
der kun kunde afsættes ved Dom af en Rigsret). Naar de
frie Forfatninger i Europa indskrænke den udøvende Magt
og nære en stadig Mistanke overfor den, kommer det af, at
de ere Overenskomster mellem en arvelig Herskerslægt og
Folket. De have en vis negativ Karakter og bestaa
væsentligt i Opregning af det, som Regenten \emph{ikke} maa gøre, af de
Rettigheder, han \emph{ikke} har. Saaledes Magna Charta, Petition
of Rights, Bill of Rights. Den amerikanske Forfatning er
derimod udtrykkeligt grundet paa Folkets Magt og skal hævdes
af dets umiddelbare Repræsentanter og Tjenere. Derfor
behøves ikke saa stor Ængstelighed. Hamilton vilde endog
finde det betænkeligt, om Unionsforfatningen opregnede
særlige »Menneskerettigheder«, som Regeringen ikke maatte
krænke; ti man kunde da let drage den Slutning, at hvad
der ikke fandtes i denne Opregning, var tilladt. Dette er
Hamiltons Svar til dem, der i Udkastet savnede en udtrykkelig
»bill of rights«\footnote{The Federalist, p. 360. --- Denne Bemærkning er træffende, især
naar man tænker paa den i mange europæiske Lande herskende Antagelse, at hvad der ikke med rene Ord er forbudt en Regering,
er tilladt den.} Han har her haft et klart Blik for de
% --- p. 137 ---
\opage{137}ejendommelige Forhold, under hvilke den amerikanske
Forfatning blev til: som positiv Fastsættelse af Folkets samlede
Magt, ikke som en Akkord eller en Kontrakt mellem stridende
Magter i Folket. Den Grund til Skinsyge overfor den
udøvende Magt, som fandtes, og stadigt findes i Europa, og som
har gjort saa mange frie Forfatninger ufrugtbare, fandtes ikke i
Amerika. Lockes og Montesquieus Teori om Magtens Deling
kunde derfor gennemføres dér i renere og konsekventere Form.
Den amerikanske Republik er ingen »parlamentarisk« Republik.

2) Det andet konservative Element ligger i
Tokammersystemet. Ved den Bestemmelse, at Senatets Medlemmer
skulde vælges af Enkeltstaterne, to fra hver Stat, lykkedes
det at lade Staternes Selvstændighed træde frem i
Forfatningen. Men tillige giver Tokammersystemet Lejlighed til
grundigere Overvejelse end Etkammersystemet. En enkelt,
talrig Forsamling giver let efter for en pludselig Tilskyndelse.
Tillige kan i Senatet Sagkundskaben blive særligt repræsenteret.

3) For det Tredje ere Forfatningsforandringer
vanskelige at gennemføre. De skulle enten foreslaas af Kongressen,
naar to Tredjedele af begge Huse ere enige derom, eller af
et Konvent, der sammenkaldes af Kongressen efter Opfordring
af de lovgivende Forsamlinger i to Tredjedele af Enkeltstaterne.
Til Forandringernes Vedtagelse udkræves, at tre Fjerdedele
af Enkeltstaternes lovgivende Forsamlinger, eller tre
Fjerdedele af Konventer, der sammenkaldes i Enkeltstaterne, stemme
for dem. Den Bestemmelse, at Senatet skal bestaa af et lige
stort Antal Medlemmer fra de forskellige Stater, kan ikke
ændres. --- Medens det engelske Underhus ved simpel Lov
kan ændre den engelske Forfatning, saa at denne befinder
sig i en uafbrudt Vorden, er der her sat større Hindringer
for Forandringer.

4) Domstolene ere aldeles uafhængige af den udøvende og
den lovgivende Magt. Den øverste Unionsdomstol kan erklære
Love, som stride mod Konstitutionen, for ugyldige. Den har
altsaa Ret til at fortolke Konstitutionen. Hensigten med
% --- p. 138 ---
\opage{138}denne Bestemmelse er at hindre, at den i enkelte, specielle
Love udtalte Folkevilje skal komme til at stride mod den i
Konstitutionen udtalte permanente og almindelige Vilje. Den
Autoritet, som den lovgivende Magt har, maa være
underordnet den Autoritet, som har sat selve Forfatningen i
Kraft. --- Ligesom Sondringen mellem den udøvende og den
lovgivende Magt saaledes er ogsaa Sondringen mellem den
dømmende Magt og de to andre Myndigheder skarpere
gennemført i Nordamerika end i europæiske Lande. Det er et
Element i Unionsforfatningen, som vi her have regnet blandt
dens konservative Elementer, men som ogsaa kunde regnes
til Frihedselementerne. Montesquieu (Esprit des lois XI, 6)
siger med Rette, at der ingen Frihed er, hvor den dømmende
Magt ikke sondres fra den lovgivende og udøvende. ---

De vigtigste af de her nævnte Grundtanker har
Hamilton selv fra først af undfanget. De andre (særligt den om
Senatets Sammensætning) har han antaget og søgt at begrunde
nærmere. Kun nogle faa Bestanddele af det rige Indhold i »The
Federalist« ere her fremdragne. Om en hel Række mere
underordnede Punkter finder man udførlig og klar Begrundelse i
det mærkelige Skrift. Man maa ønske de europæiske
Folkeslag, at den Tid maa komme, da deres konservative Partier
ville støtte sig til en saa klar, sund og folkelig Tankegang
som den, der gaar igennem »Federalisten«. Bogen er en i
sin Art enestaaende Forbindelse af et statsfilosofisk Arbejde
og et Agitationsskrift. Men sjældent er et Agitationsskrift
fremtraadt i en saa ren og ædel Form. For at give en
Forestilling om det ophøjede Standpunkt, fra hvilket dets
Hovedforfatter saa paa den saa stærkt spændte Situation, under
hvilken det ikke manglede paa grove og hadefulde Ord fra
begge Sider, skal jeg anføre en Udtalelse fra det første
Nummer.

»Blandt de frygteligste Hindringer, som den nye
Forfatning vil have at overvinde, vil man let bemærke den
aabenbare Interesse, som en vis Klasse af Mænd i enhver
Stat har af at modsætte sig alle Forandringer, som kunde
medføre en Forminskelse af den Magt, Indtægt og Indflydelse,
de besidde ved deres Embeder under Statsinstitutionerne, saa
% --- p. 139 ---
\opage{139}vel som den fordærvede Ærgerrighed hos en anden Klasse
af Mænd, der enten haabe at stige højere under de forvirrede
Forhold i deres Land eller smigre sig med at opnaa større
Magt ved Rigets Deling i forskellige mindre Forbund.

»Det er dog ikke min Hensigt at dvæle ved
Betragtninger af denne Natur. Jeg indsér, at det vilde være uædelt
uden Forskel at forklare Oppositionen fra nogen som helst
Klasse af Mænd ved egenkærlige eller ærgerrige Motiver,
blot fordi deres Stillinger kunde gøre dem til Genstand for
saadan Mistanke. Oprigtighed nøder os til at indrømme,
at ogsaa saadanne Mænd kunne være besjælede af ærlige
Hensigter, og det kan ikke betvivles, at Meget af den
Opposition, som allerede er fremtraadt, eller herefter vil fremtræde,
udspringer fra Kilder, som ikke kunne lastes, om de ikke
kunne agtes: hæderlige Vildfarelser, som ere fremkaldte ved
forudfattet Mistanke og Frygt. Saa talrige og mægtige ere
jo de Aarsager, der kunne føre Dømmekraften vild, at vi ved
mange Lejligheder se vise og gode Mænd saa vel paa den
vrange som paa den rette Side ved Spørgsmaal af aller største
Betydning for Samfundet. Denne Omstændighed vil, naar
den ret overvejes, stedse afgive en gavnlig Lærdom for dem,
der ere indviklede i en eller anden Strid, hvor overbeviste
de saa ere om at have Ret. Og en Grund til Forsigtighed
ligger ogsaa deri, at vi ikke ere sikre paa, at de, der
forsvare Sandheden, besjæles af renere Principer end deres
Modstandere«.

Bedre vil man ikke kunne tydeliggøre, hvad der taler
for at holde en Debat i rent saglig Form. Og bedre er denne
Form næppe bleven overholdt i nogen politisk Debat.

\section*{VI.}

I Washingtons første Ministerium traadte Hamilton ind
som Finansminister. Det faldt derved i hans Lod at befæste
det Værk, hvis Grundtanke han først havde fattet, og hvis
Natur han grundigst havde udviklet. Han konsoliderede
Statsgælden, saaledes at Unionen overtog Enkeltstaternes
Gæld. Den fælles Statsgæld skulde tjene som et nyt Middel
til at holde Unionen sammen ved at skabe fælles Interesser.
% --- p. 140 ---
\opage{140}Desuden søgte han at skaffe Unionen Indtægter ved en Lov
om Accise, og han fik oprettet en Nationalbank. Det er
anerkendt fra alle Sider, at det skyldtes hans kloge
Foranstaltninger, at Landets Anseelse og Kredit hævedes, og at
Handel og Industri kom i Flor.

Da Talleyrand i 1794---1795 opholdt sig i Amerika (fordi
han hverken maatte opholde sig i England eller Frankrig),
omgikkes han meget fortroligt med Hamilton, for hvem han
nærede den største Beundring. Han betragtede i sin Alderdom
Napoleon, Fox og Hamilton som de største Mænd, han havde
mødt paa sin mærkelige og bugtede Vej, og skulde han sætte
en af dem højst, var han tilbøjelig til at give Hamilton denne
Plads. Han sagde om ham, at »han havde gættet Europa«,
hvormed han hentydede til Hamiltons skarpe Blik for den
politiske Situation og de Muligheder, den rummede.

Men Hamiltons Personlighed egnede sig dog ikke for
den specielle Politik med dens Partikampe og dens
personlige Sammenstød. Han var uheldig som Partifører, og han
forstod bedre at kæmpe paa egen Haand end at faa en sluttet
Skare til at følge sig gennem den politiske Kamps mange
Svingninger. Men lige saa uheldig, som han i denne Henseende
var, lige saa rolig og bestemt gik han den Vej, ad hvilken
hans Overbevisning førte ham, hvad enten han var ensom
eller blev fulgt af Mange.

Partikampen begyndte snart igen efter Forfatningens
Vedtagelse. I Washingtons Kabinet sad overfor Hamilton
Jefferson som hans ivrige Modstander, der hævdede den
diametralt modsatte politiske Opfattelse. Han saa med
Betænkelighed paa Hamiltons Bestræbelser for at befæste Unionen.
Disse Bestræbelser maatte, frygtede Jefferson, føre i Retning
af et Monarki. Han hævdede Enkeltstaternes Ret til at
nægte Lydighed mod Unionsregeringen, naar de mente, at
denne overskred sin Kompetence. Det var det saakaldte
Nullifikationsprincip, der i de følgende Tiders Kamp mellem Norden
og Syden kom til at spille saa stor en Rolle.

I Aaret 1795 traadte Hamilton ud af Ministeriet og
levede nu Resten af sit Liv som Sagfører. Kun én Gang kom
han endnu i offentlig Tjeneste, da han under den Troppe%
% --- p. 141 ---
\opage{141}samling, der med Washington som Overgeneral foretoges i
Aaret 1798 i Anledning af det spændte Forhold til Frankrig,
efter Washingtons Ønske udnævntes til Næstkommanderende.
Da Washington døde, blev han Overgeneral i hans Sted.

Da Hamiltons Parti, de saakaldte Føderalister, ikke
havde Held med sig overfor de stedse mægtigere
»Republikanere« (saaledes kaldte Jeffersons Parti, der svarer til de
senere »Demokrater«, sig som en Protest mod Føderalisternes
foregivne monarkiske Sindelag), viste mange af dets Tilhængere
en upatriotisk Fanatisme, som Hamilton var langt fra at dele.
Han støttede saaledes Jeffersons Valg til Præsident, da man
kun havde Valget mellem ham og den forbryderiske
Eventyrer Oberst Burr, hvilken Sidste Føderalisterne stemte paa
af Had til Jefferson. Og da det lykkedes Jefferson at
erhverve Louisiana fra Frankrig, billigede Hamilton dette
Skridt, medens hans fanatiske Partifæller vare imod det,
fordi de frygtede, at Syden, hvor »Republikanerne« stedse
vandt flere Tilhængere, derved skulde blive mægtigere. Nogle
Føderalister, bl. A. Burr, fattede endog den Tanke at
opløse Unionen og danne en særlig Union for Nordstaterne.
Intet kunde stride mere mod Hamiltons Overbevisning, og
han bekæmpede derfor af al Magt og med stærke Ord Burrs
Forsøg paa at blive valgt til Guvernør i New-York, da han
frygtede for, at Burr vilde benytte denne Stilling til at virke
for sine landsforræderiske Planer. Dette gav Burr Anledning
til at udfordre Hamilton, og skønt denne ansaa Dueller for
utilbørlige, mente han dog ikke at kunne afslaa Udfordringen,
Kun 47 Aar gammel blev han haardt saaret af Burrs Skud
den 11. Juli 1804 og døde Dagen efter. Han ofrede sit Liv
for sin Sags Renhed.

\streg{}

% --- p. 142 ---
\opage{142}
\essay[SOCIAL PESSIMISME]{SOCIAL PESSIMISME.}

\begin{center}\emph{Ferdinand Tönnies:} Gemeinschaft und Gesellschaft. Leipzig 1887.\end{center}

\streg{}

\section*{I.}

Naar en optimistisk og en pessimistisk Anskuelse staa
overfor hinanden, er det vel i Regelen Stemninger, som ere
det Bestemmende, Stemninger, der enten have deres
Oprindelse fra Individernes Temperament eller fra de begrænsede
og tilfældige Erfaringer, de have gjort eller mene at have
gjort. Dog vil der undertiden kunne være Tale om mere
end blot Stemning. Det er saaledes klart, at Processerne i
et organisk Legeme kunne tage et Forløb, der gør det
muligt for den Kyndige at forudsige en Opløsning af Livet.
Og Fysikere have søgt at vise, at den Tid vil komme, da
alt Liv og al Bevægelse i vort Solsystem vil ophøre, idet
der vil indtræde en ligelig Fordeling af Varmen, og ingen
Omsætning af Varme til Bevægelse længere vil finde Sted.
De udlede dette af, at medens al Bevægelse kan omsættes
til Varme, kan ikke al Varme igen omsættes til Bevægelse,
men en Del af den vil udbrede sig til Omgivelserne. Tænker
man sig denne Proces fuldført, faa vi Stilstanden. Skulde
det nu ikke paa lignende Maade være muligt at afgøre, om
der paa det menneskelige Samfundslivs Omraade finder en
Opløsning og en Tilbagegang eller en Fremgang Sted?

Man kunde mene, at Sagen maatte være klar i deres
Øjne, som hylde Udviklingshypotesen. Men man vilde da
forveksle Udvikling og Fremskridt. Udvikling er en
Tillempelsesproces, ved hvilken Individer og Arter antage
saadanne Former og blive i Stand til at udøve saadanne
Virksomheder, som Livsforholdene, de ydre Omstændigheder kræve.
% --- p. 143 ---
\opage{143}Men en saadan Tillæmpelse kan under visse Betingelser føre
til, at Former og Evner forminskes og forringes i Stedet
for at udfolde sig rigere og mangfoldigere. Hos Snyltedyr
findes saaledes Spor af en afgjort Tilbagegang, idet Organer,
der hos beslægtede Dyrearter fremtræde med en rig Bygning,
ere indskrumpede til svage Antydninger, og de tilsvarende
Evner ligeledes ere forsvundne. Ogsaa Mennesker kunne
tillempe sig efter saa ringe og simple Forhold, at de miste
de højere Evner, som vare anlagte hos dem. Det er dog
ikke blot ved Overgang fra mangfoldigere og større Forhold
til snevrere og ringere, at en saadan Tilbagegang kan ske.
Ved Overgang fra simple og jevne, men let overskuelige
Forhold til mere sammensatte og omfattende, men tillige
ogsaa mere kunstige Forhold kunne Evner af stor Værdi
ogsaa gaa tabte. Der kan gaas frem i Mangfoldighed, og
dog finde et væsentligt Tilbageskridt Sted i Koncentration,
i indre Sammenhæng og Kraft. Dette sker saaledes let, naar
det begrænsede, men sikre og trygge Instinkt afløses af den
mere omfattende, men ofte tvivlende og fejlende Fornuft. Hos
Mennesket blive derfor Udskejelser og Unaturligheder mulige,
som det af Instinktet ledede Dyr ikke kender. Efter Darwin’s
Opfattelse er der sket et Tilbageskridt kort efter »det meget
tidligere Tidspunkt, da Mennesket lige var blevet Menneske«,
og da Fornuften paa saa mange Punkter skulde til at afløse
Instinktet\footnote{Se „Menneskets Afstamning og Parringsvalget“. Dansk Overs. II
p. 364, 379, 400. --- Der kan i denne Sammenhæng mindes om Kants
og Schillers Fortolkning af Myten om Syndefaldet som Udtryk for
Overgangen fra Instinkt til Fornuft.}. Den menneskelige Udvikling, baade hos det
enkelte Individ og hos hele Folkeslag, viser ofte den
Tilbagegang, der i hvert Tilfælde paa visse Punkter --- kan
være forbunden med en Overgang fra simple og begrænsede
til sammensatte og mangfoldige Livsbetingelser. De
Indsigelser mod Kulturens Værdi, som under Tiden fremkomme
med stort Eftertryk, have deres relative Berettigelse ved
denne Omstændighed.

De senere Aars Undersøgelser af Samfundsforholdenes
og Retsforestillingernes Udvikling hos de indoeuropæiske
% --- p. 144 ---
\opage{144}Folkeslag have nu endog ført til Opstillingen af en hel social
Patologi. Man har søgt at begrunde den Paastand, at hele
den moderne Kulturudvikling i Virkeligheden er en stor
Opløsningsproces, der skrider stille og umærkeligt frem, kun
aabenbarende sig for den sammenlignende og konsekvente
Forsker. Det ovenanførte Skrift af \emph{Ferdinand Tönnies} (fra
Husum, Docent i Kiel) er et interessant og, som det
forekommer mig, meget lærerigt Indlæg i denne Retning.
Forfatteren, der tidligere har gjort sig bekendt ved sine Studier
over Hobbes og Spinoza, Arbejder, der høre til det
Fortrinligste, som er ydet paa Filosofiens Histories Omraade, har
den samme objektive Ro og Uforstyrrethed i Betragtningen af
det menneskelige Livs Fænomener, som udmærker hine ældre
Tænkere, han har beskæftiget sig saa meget med. Man mærker
ikke hos ham en Rousseaus eller en Schopenhauers
Lidenskabelighed. Han undersøger og konstaterer paa rent
teoretisk Vis, og fordrer at imødegaas paa samme Vis. Han
forbinder paa en ejendommelig Maade Sociologi og Psykologi,
idet han viser, hvorledes den sociale Udvikling nødvendigvis
hænger sammen med og har sit Sidestykke i en tilsvarende
Udvikling af de menneskelige Aandsevner. Det rige Stof,
han bygger paa i sin Undersøgelse, har han saaledes gjort
til sin Ejendom, at han har omsat det til en Kreds af Tanker,
der udvikles ofte i meget abstrakt Form og med Benyttelse
af nydannede Begreber og Udtryk, men stedse saaledes, at
man mærker fast Grund under Fødderne. Skønt Bogen kræver
indtrængende Studium for at forstaas, lønner den det Arbejde,
som anvendes paa den. Meget som hidtil var kendt, kommer
til at staa i ny Belysning, og adskillige nye Synspunkter
fremføres. Den, der for Alvor vil forstaa sin Tid og dens
Tendenser, bør ikke sky det Arbejde at prøve en Dyst med
denne Forfatter. Jeg vil her indlade mig paa en saadan Dyst,
idet jeg vil gøre Rede for hans Standpunkt og dernæst
undersøge, hvad der kan gøres gældende imod det.

\section*{II.}

De europæiske Folkeslags nuværende moralske og sociale
Forhold have i Følge de nyere Undersøgelser udviklet sig af
% --- p. 145 ---
\opage{145}en tidligere Tingenes Orden, i hvilken den skarpt udprægede
Retsorden med bestemt Fastsættelse af enhver Persons Krav
og Pligt ikke eksisterede, lige saa lidt som den store
Interessekamp fandt Sted, under hvilken hver Enkelts Vinding
eller Tab bestemmes ved Forholdet mellem Tilbud og
Efterspørgsel. Der eksisterede hverken den juridiske eller den
økonomiske Sondring mellem Individerne, som har udviklet
sig hos de indoeuropæiske Folkeslag, der fra deres Hjem i
det fjerne Østen efterhaanden bosatte sig i de vestlige Egne.
Levninger af den gamle Samfundsorden findes nu til Dags
paa Landet mere end i Byerne, hos Almuen mere end hos
de saakaldte Dannede, i de østlige Lande mere end i de
vestlige. I Indien bestaar saaledes endnu den gamle, paa
Familje- og Naboforholdet grundede Landsbyordning. Det var
Blodets og den rumlige Nærheds Baand, som forenede
Menneskene. Umiddelbar Sympati og Forstaaelse, personligt
Kendskab og en ved Samliv styrket Samvirken vare de
Kræfter, der bar Samfundet. Det var et naturligt,
selvfølgeligt, uvilkaarligt Samliv, ikke indgaaet med klar Bevidsthed
og i Haab om bestemte individuelle Fordele. De enkelte
Individer følte sig umiddelbart som Dele af et Hele og
tillagde sig ingen Eksistens udenfor dette. Reglerne for deres
Færd vare givne i urgamle Skikke, Anordninger og
Overleveringer, der ikke drøftedes, men fulgtes som en
Naturorden og betragtedes med Ærefrygt. Ogsaa det større
Verdensløb, af hvilket Familiens og Stammens Skæbne var afhængigt,
opfattedes og vurderedes ikke med prøvende Forsken og
Beskuen, men med levende Indbildningskraft, Følelse og Tro.
Ogsaa her saa man en Tingenes Orden, der traadte de
enkelte Individer i Møde, paa én Gang som noget Nært og
Hjemligt, og dog som noget Ærværdigt og Raadende.

Kulturens Gang har nu --- som Tönnies viser --- mere
og mere ført ud over denne Tilstand. Der er foregaaet en
fremskridende Emancipation. Individet har faaet en
selvstændig Ret, afset fra det snævre, naturlige Samfund, det
hører til. Der er sat bestemte Grænser mellem de enkelte
Individers Raaderum. De ere revne ud af deres umiddelbare
Sammenhæng, ere blevne Centra hver for sig i Stedet for at have
% --- p. 146 ---
\opage{146}et fælles Centrum. Der finder vel Samkvem og Forbindelse
Sted imellem dem, men Samkvem og Forbindelse bestemmes
nu mere og mere ved hvad den Enkelte mener at kunne
opnaa af de Andre. \emph{Selskabet} (Gesellschaft) træder mere og mere
i Stedet for \emph{Fællesskabet} (Gemeinschaft). I Stedet for
Instinktet og den ved Tradition, Autoritet og Broderskab
umiddelbart bestemte Følelse træder nu det forstandige,
beregnende Overlæg. Ydre Interesser forbinde Menneskene,
uden at de derfor staa i nogen væsentlig indre Forbindelse.
Den Ene vil med det mindst mulige Offer fra sin Side opnaa
det mest mulige af det for ham Attraaværdige, som den Anden
besidder. \emph{Handelsforholdet} er det mest fremtrædende Eksempel
paa Selskabet i Modsætning til Fællesskabet, og Handelen
er tillige det vigtigste Middel til, at Selskabet afløser
Fællesskabet\footnote{Ogsaa i Følge \emph{Leist} (Alt-arisches Jus Gentium. Leipzig 1889 p.
523---525) er den gammelariske Husordning gaaet til Grunde ved
Handelens og Pengevæsenets Opsving, der havde den skarpe
Udvikling af det rent privatlige Ejendomsbegreb til Følge. Det er „den
dæmoniske Kapitalmagt“, der omdanner Forholdene. I Stedet for
Husfællesskabet og den til det knyttede Erhvervsvirksomhed trænger
nu Egeninteressen sig frem, som er fjendtlig mod ethvert varigt
Fællesskab. Leist gør træffende opmærksom paa, hvor forstaaelig
og historisk betydningsfuld i denne Sammenhæng Aristoteles’s
Polemik (i hans Statslære) mod enhver Stræben efter Erhverv, der gaar
ud over det for Husstanden Nødvendige, kommer til at staa for
os. Det var dog en reaktionær Tankegang, der ikke formaaede at
stanse Udviklingen. --- Paa lignende Maade har Handelen virket
ved Overgangen til den nyere Tid. Det var de forandrede
Afsætningsforhold, Dannelsen af et Verdensmarked i Stedet for de tidligere
begrænsede Afsætningsforhold, der førte til Bruddet med den gamle
Lavsorden, til Haandværkets Afløsning ved Storindustrien, til
Opfindelsen af Maskiner og til hele den moderne Konkurrence. Smlgn.
\emph{L. Brentano:} Über die Ursachen der heutigen socialen Noth. Leipzig
1889 p. 16. Brentano bebrejder Marx (hvem Tönnies paa dette
Punkt følger vel meget), at han vender Forholdet om, idet han
lader den store Industri udvikle sig af sig selv og derpaa
frembringe Verdensmarkedet, ligesom Marx i det Hele i sine Teorier
overser Efterspørgselens Betydning.}. Ved »Selskabet« ses der bort fra alle oprindelige
eller naturlige Forhold mellem Menneskene indbyrdes. Et
Selskab forudsætter kun en Flerhed af nøgne Personer, der
% --- p. 147 ---
\opage{147}ere i Stand til at yde, følgeligt ogsaa til at love. Mod selv
at faa Midler til \emph{sine} Formaal yder den Ene den Anden
Midler til \emph{hans} Formaal. Ved Overgangen fra Fællesskab
til Selskab nærme vi os en Tilstand, hvor Enhver er
Købmand.

I \emph{Fællesskabet}, i det gamle Samfund, var Forbindelsen
mellem Individerne given, før de gensidigt kunde yde
hverandre Tjenester, og før de kunde gøre Regning paa saadanne
gensidige Ydelser. I \emph{Selskabet} derimod faa Individerne for
det Meste kun med hverandre at gøre, naar deres Interesser
mødes paa harmonisk eller disharmonisk Maade. En kunstig,
mekanisk Sammenhæng træder i Stedet for den indre.
Aktieselskabet er den karakteristiske Type for Selskabet, ligesom
Familien for Fællesskabet. I Stedet for Agerbrug,
Haandværk og Kunst træder Industri og Handel. I Stedet for
det personlige Forhold mellem Haandværkeren og hans Kunde
træder det upersonlige Forhold mellem Masseproducenten og
hans tilfældige, ubekendte Købere. Storstaden med sit
brogede, internationale Liv træder i Stedet for Huset,
Familien og Landsbyen. Selskabet tilspidser sig tilsidst i
Forholdet mellem Kapitalist og Arbejder som et rent Købs- og
Kontraktforhold. For at bruge Sir Henry Maine’s Udtryk:
den naturlige, overleverede Tilstand \emph{(status)} afløses mere og
mere af den vilkaarligt, med bevidst Overlæg tilvejebragte
\emph{(contractus)}.

Der foregaar her, i Følge Tönnies, en patologisk
Udviklingsproces, som vel kan ændres, men ikke hindres af
Filantroper og Lovgivere. Om hin »dualistiske
Konstruktion af Selskabet«, der viser sig i Modsætningen mellem
Kapitalister og Arbejdere, er den eneste mulige, --- hvad
der vil ske, naar Linen er løben ud, og alle Forhold ere
blevne Kontraktforhold, --- det er Spørgsmaal, som
Forfatteren erklærer for at være ham uvedkommende.

At det er en \emph{patologisk} Udvikling, der her er beskreven,
ses især, naar vi betragte den psykologiske Side af
Udviklingen, ikke blot den sociale. Som Selskabet forholder sig
til Fællesskabet, saaledes forholder \emph{den vilkaarlige Villen} sig
til \emph{Væsensviljen}. Ved Væsensvilje forstaar Tönnies Instinkt,
% --- p. 148 ---
\opage{148}umiddelbar Tilskyndelse, der med indre Nødvendighed gaar
imod sit Maal. Den virker overvejende ubevidst, efter
Naturanlæg og Vane eller bestemt ved Efterligning og
umiddelbar Sympati. Ved vilkaarlig Villen skelnes derimod
bestemt mellem Formaal og Middel, og det bliver Genstand
for Overvejelse, hvor meget det er værdt at ofre paa Midlerne
for at opnaa Formaalet. Enhver vilkaarlig Villen er et Køb,
ofrer Noget for at vinde Noget. Det Ejendommelige ved
den vilkaarlige Villen træder tydeligt frem, naar Midlet ene
og alene villes for Formaalets Skyld. Der er derfor ved
den vilkaarlige Villen noget Unaturligt, en vis Selvtvang, en
Falskhed. Naar man gør sig klart Rede for hvad man ofrer,
og hvad man derved vil opnaa, vil tillige Egoismen ligge
nær og vil fortrænge den umiddelbare Hengivelse. Det er
ved Vilkaarligheden ikke Godheden, men kun Klogskaben og
Konsekvensen, som er at beundre. Livet staar ikke længere
som et Kald, i hvilket der skal arbejdes, men som en
Forretning, der skal gøres.

Selskabets Udvikling begunstiger Tendensen til
vilkaarlig og egoistisk Villen paa Bekostning af Væsensviljen og
Hengivelsen. Under den stadige Kappestrid bliver
Vilkaarlighedens hensynsløse Brug en Betingelse for Individets
Bestaaen. Den klare Forstand træder i hele det aandelige
Liv i Stedet for den umiddelbare Følelse, Fantasi og Tro.
Der finder en fremskridende Intellektualisme Sted; den
videnskabelige Refleksion træder mere og mere i Stedet for den
barnlige Tro. Det religiøse og det sociale Problem ere
saaledes Grene af et og samme store Problem, som Tönnies
behandler dels fra den sociologiske, dels fra den
psykologiske Side.

\section*{III}

Ingen vil kunne negte, at der foregaar en saadan
sociologisk og psykologisk Udvikling, som Tönnies skildrer.
Kulturen trænger ind, hvor der før kun var Natur;
Livsformerne blive stedse mere sammensatte og i højere Grad
Frugten af vilkaarlig Ordning; der stilles stedse større For%
% --- p. 149 ---
\opage{149}dringer til Kløgt og Beregning. Vi synes for saa vidt stedse
mere at fjerne os fra det Primitive, Naturlige og Uvilkaarlige.
Vil Linen da ikke løbe ud, saa det menneskelige Liv
forstenes eller opløses, ligesom det efter en fysisk Teori vil
gaa med Jordklodens og vort Solsystems Liv og Bevægelse i
det Hele? Eller skulde den af Tönnies skildrede Udvikling
blot være én Side af Sagen, idet der for det opmærksomme
Blik vil fremtræde andre Udviklingstendenser, som kunne
holde dem Stangen, saa at den ikke kommer til at løbe
Linen ud?

Lad os i Korthed se, hvilke Erfaringer og Betragtninger
der kunne føre til at afvinde Sagen en lysere Side. ---

Pessimismen er her som saa ofte beslægtet med
\emph{Romantiken}. Den ser Fortiden i et straalende Lys, medens
Nutiden synes kold, tom og mørk. De gamle Samfundsforhold
forud for den store Emancipationsproces, der udfylder saa
stor en Del af Historien, vare dog ingenlunde ideale. De
fremtræde for os med Enhedens og Sluttethedens Karakter
og tillige med et vist idyllisk Præg. Menneskelivet og de
menneskelige Følelser fremtræde i simple, intensive og
elementære Former, der ere velgørende for det af Nutidens brogede,
delte og stridende Liv trættede Øje. Hos Tönnies
forekommer det mig særligt, at man mærker den pietetsfulde
Vedhængen ved Livet som det former sig paa Landet og i de
smaa Flækker; han skildrer det og dets Former med en
Varme, som denne klare Tænker ikke formaar at
vedligeholde, hvor det gælder Forstaaelsen af Storstadens og det
sammensatte Kulturlivs Forhold, og som ganske sikkert
ogsaa er vanskeligere at bevare dér. Men Idyllen er ofte kun
et Skin --- i den menneskelige Verden som i den øvrige
Natur. Der kunde f. Eks. ved det gamle patriarkalske System
gaa megen Grusomhed og Undertrykkelse i Svang.
Husfaderen herskede uden Kontrol, kun underkastet Gudernes
Straf; men de lod ikke altid høre fra sig. Om det Værste
tier endda maaske Historien, ti de aldeles Retsløse faa ingen
Historie. Det var i Virkeligheden intet Fald, da det
gammelariske, til Dels paa præstelig Tradition og
Præsteherredømme byggede Familiesystem afløstes af det klart udviklede
% --- p. 150 ---
\opage{150}græsk-romerske Borgersystem, indenfor hvilket Privatretten
undergik sin grundlæggende Udviklingsproces. Det var ikke
blot Egoisme og selvisk Erhvervsinteresse, som sprængte den
gammelariske Husorden og i en senere Tid den
middelalderlige Lavsorden. Det var ikke mindst en --- ved de forandrede
økonomiske Forhold vakt og næret --- udtrykkeligere
Anerkendelse af den enkelte Personligheds Værd og Evne.

Tönnies gør vistnok med fuld Ret opmærksom paa den
store Forskel mellem Trællene indenfor de gammelariske
Familier og senere Tiders Slaver. Trællens (der Knecht)
Stilling nærmede sig ofte til Sønnens eller Vennens; han var
et Medlem af Husstanden saa vel som Hustru og Børn. »Idet
han deler Familiens Glæde og Sorg, yder sin Herre en
voksen Søns Ærefrygt og nyder en Hjælpers eller endog en
Raadgivers betroede Stilling, er han efter hele sin
Beskaffenhed et frit Menneske«. Kun undtagelsesvis trykkes han ved
Mishandling ned til en virkelig slavisk Stilling. Senere har
Leist fremhævet, at Trællevæsenet i Oldtiden, endnu hos
Grækerne, hverken maa betragtes og bedømmes fra et rent
privatretligt Synspunkt eller fra den offentlige Rets
Synspunkt. »I Husordningen have vi et arisk Urretsinstitut for
os, som er ældre end Sondringen mellem privat og offentlig
Ret. Den er den juridiske Mikrokosme, af hvilken
Privatretten og den offentlige Ret først efterhaanden have udviklet
sig«. Ligesom ved Erhvervsvirksomheden saaledes forstaa vi
ogsaa her Aristoteles’s Standpunkt (særligt hans bekendte
Forsvar for Trælleinstitutionen) som et Forsøg paa at
fastholde den gamle Husordning overfor den fremskridende
Individualisme.

Dette er en vigtig Betragtning, hvorved Meget i
Historien bliver bedømt paa en retfærdigere Maade end en abstrakt
filantropisk og naturretlig Anskuelse er i Stand til. Men
netop naar Trællens Stilling ofte og maaske i lange Tider
har været forholdsvis god og betrygget ved patriarkalsk
Omsorg, vil det være aldeles uforstaaeligt, hvorledes der kunde
opstaa Ønsker om at forlade denne Stilling, hvis der ikke
fandt Misbrug og Undertrykkelse Sted. Det er formodentligt
gaaet her som ved Udviklingen af den frie Arbejderstand
% --- p. 151 ---
\opage{151}ved Middelalderens Slutning, der blev fremkaldt ved den
Egoisme, hvormed Lavsmestrene søgte at holde deres
Fordele og Rettigheder indskrænkede til saa lille en Kreds som
muligt, samtidigt med, at (som allerede berørt i en anden
Sammenhæng) Markedet udvidedes og stedse flere
Arbejdskræfter maatte tages i Beslag\footnote{\emph{L. Brentano} skildrer (i sit Værk: Die Arbeitergilden I p. 59) denne
Sag saaledes: „Industriens Opblomstring og den større Anvendelse
af Kapital i den, hvilke gensidigt fremkaldte og betingede hinanden,
trak nødvendigvis de Livegne i Mængde til Byen. Derved sørgedes
der vel paa den ene Side for det større Antal Arbejdskræfter, som
Industriens Fremskridt fordrede, men paa den anden Side opstod
der i enhver ny Arbejder en mulig fremtidig Konkurrent. Udbyttet
af Kapitalanlæggelse i Industrien truedes altsaa; Bekymringen
herfor vakte Haandværkernes Monopolaand, og der opstod en Mængde
Bestemmelser, der indskrænkede hine ny fremtrædende Familiers
Konkurrence. Betalingen for Indtrædelse i Lavet forhøjedes; paa
Fastlandet opstod den Skik, at der af Enhver, der vilde drive
selvstændigt Haandværk, forlangtes et for det Meste kostbart
Mesterstykke, hvorimod Sønner og Svigersønner saa vel som de, der
ægtede en Enke, som hørte til Lavet, nød Lettelser.“}.

Det gamle System var ikke blot et \emph{patriarkalsk
Hussystem;} det var tillige et \emph{Tvangssystem}, eller det havde i
hvert Tilfælde hvert Øjeblik Muligheden for at gaa over til
et Tvangssystem. Men hvad enten det virkede ved
omsorgsfuld Mildhed eller ved Tvang, én Ting var det ikke egnet
til: at udvikle \emph{Selvstændighed} i Evne og Drift og den dertil
knyttede \emph{Selvfølelse} hos de enkelte Individer. Disse stode ---
med Undtagelse af dem, der var Autoriteter (Fædre, Præster,
Profeter osv.) --- som mere eller mindre selvløse Dele af det
Hele. Men de vare ikke hver for sig smaa Helheder.
Overgangen til den mere individualistiske Ordning (fra Fællesskabet
til Selskabet), hvad enten den nu mest er fremkaldt ved de
autoritære Magthaveres Egoisme eller ved de formedelst ydre
Forhold ændrede økonomiske Betingelser, har da i hvert
Tilfælde bevirket det Gode, at der blev Mulighed for og vaktes
Drift til Udfoldelsen af Evner og Kræfter, som hidtil laa
slumrende hos de Enkelte. Derved muliggjordes ganske vist
ogsaa Isolation og Udvorteshed i Forholdet mellem Indivi%
% --- p. 152 ---
\opage{152}derne. Men der stilledes dog den store Opgave for hver
Enkelt at blive til en selvstændig og ejendommelig
Personlighed; der blev anlagt en Mængde selvstændige Udgangspunkter
for Slægtens Liv. En større Inderlighed i det menneskelige
Liv, helt nye Erfaringsomraader af stor Værdi banedes der
derved Vej til. Selv om Individernes Virken ikke gaar saa
let og instinktmæssigt for sig som i det gamle Samfund, saa
er der dog opnaaet et stort Gode. Under Udviklingens Gang
tabes der uundgaaeligt Noget, hver Gang der vindes Noget.
Det er nu én Gang saa. Men det er uberettiget blot at se
paa hvad der tabes. Det Højeste vilde ganske vist være at
forene det Værdifulde fra de tidligere Udviklingstrin med
det Værdifulde fra de følgende Trin. Men det er et Ideal,
som ikke fuldt kan naas, skønt det har sin store Betydning
at holde sig det klart for Øje. Den pessimistiske Anskuelses
Mission er at hindre os i at glemme det.

Ikke blot en større Selvstændighed og Inderlighed bliver
mulig, men ogsaa et \emph{større Samfund}. Staten og
Menneskeheden, det er Begreber, der først kunne udvikles, naar den
patriarkalske Ordning løslader de enkelte Individer og
tillader dem at føle sig i menneskelige Forhold, selv om de
skilles fra det oprindelige Hjemsted. Ogsaa denne Side af
Udviklingen frembyder sine Farer; de snevrere Kredse, med
deres umiddelbare og derfor hjerteligere Berøringer, skydes
let til Side for de større, men mere abstrakte Forhold, og
Livet kan derved lide i Kraft og Varme. Men i sig selv
var det dog en stor og betydningsfuld Udvidelse af
Menneskelighedsfølelsen, der fandt Sted, da Stat og
Menneskehed bleve ledende Ideer i det menneskelige Liv. Skranker
nedbrødes, som paa mange Maader medførte Hæmmelse og
Strid\footnote{I „\emph{Dansk Bondeliv}, saaledes som det i Mands Minde førtes, navnlig
i Vestjylland“ (København 1889) har Pastor \emph{H. F. Feilberg} givet
en interessant Skildring af Livet, ogsaa Samfundslivet i
Landsbyen, saaledes som det førtes, før de moderne Reformer, de
udvidede Færselsmidler, Oplysningen og den politiske Bevægelse bragte
nye Forhold og en ny Aand til Veje. Jo større den Kærlighed og
Forstaaelse er, med hvilke Forfatteren dvæler ved den gamle Tingenes Orden, des interessantere er det at høre hans Dom om
den Forandring, der er i Færd med at foregaa og til Dels allerede
er foregaaet. Han fremhæver, hvilket snevert Syn og hvilke
Standsfordomme der maatte opstaa under den Isolation og Begrænsning,
under hvilken Livet førtes. Man saa kun „til sin egen Marks
Grænse, eller naar det kommer højt til Sognets Skel; underligen
vilde det have været og hartad ubegribeligt, om man havde truffet
Folk med Syn for det store Samfund, Land og Folk“. Der er nu
som Følge af de nævnte Aarsager, ikke mindst af den politiske
Bevægelse selv i de usleste Hjem kommet „\emph{noget af en
Fællesskabsfølelse, som vist aldrig har været der før}“. Arbejderen og
Husmanden have lært, at „de høre et Folk til, at de har fælles
Interesser med deres Folk, at de har en Ret og et Ansvar“ (p. 130,
137). --- Hvis Forfatteren har Ret heri, se vi, hvorledes Opløsningen
af de gamle Bykommuner med deres Fællesskab og deres paa
Naboforholdet hvilende humane og sociale Forbindelser kan være
en Overgang til en større og mere omfattende Fællesfølelse og til
en Bevidsthed om større Opgaver og Synspunkter, end den gamle
Ordnings begrænsede Horizonter gjorde mulige.}. ---

% --- p. 153 ---
\opage{153}Vanskeligheden ved at forene de tidligere Udviklingstrins
Goder med de seneres lettes nu for en stor Del ved, at
\emph{Slægten i en vis Forstand stedse begynder forfra}. Det gaar
her som i Naturen, hvor vi mere og mere have lært, at de
Kræfter, der have virket i de tidligere Jordperioder, ikke
ere væsensforskellige fra dem, der den Dag i Dag ere i
Virksomhed rundt om os. Ogsaa paa det menneskelige og sociale
Omraade gælder Aktualitetsprincipet. Al mulig Emancipation
og Individualisme kan ikke ophæve de simple Grundforhold,
som stedse paany sætte Fællesskabet i Selskabets Sted.
Familien maa stedse vedblive at være et Fællesskab.
Forholdet mellem Forældre og Børn og mellem Mand og Hustru
lægger Grunden til de stærkeste og varigste Følelser, som
Menneskelivet kender. Indenfor det Samfund, som bestemmes
ved disse Følelser, begynder ethvert Menneske sit Liv. Her
finder et Samliv Sted, som er Medfølelsens og
Menneskelighedsfølelsens stadige Hjemsted. Fra det udgaar stedse paany
de Kræfter, der holde Menneskelivet sammen. Her bøjes og
lempes Egoismen; her undergaar Brutaliteten en stille
Ændring. Ingen Retsmekanik, selv den kløgtigst udtænkte, og
ingen Oplysning, selv den mest grundige, kan erstatte det
% --- p. 154 ---
\opage{154}Væld, som her har sit Udspring. Og denne inderliggørende,
mildnende og sammenholdende Magt er ikke svækket, fordi
Livets ydre Orden er bleven en anden end den var, da man
ikke kendte andre Samfundsformer end dem, der grundedes
paa Blod og Naboskab. Familien kommer netop først til sin
Ret, naar den ikke længere er det eneste, eller saa godt
som det eneste Samfund. Magtforholdet træder da tilbage;
de Elementer udskilles, som i den gamle Familie altfor meget
mindede om et Forhold mellem Hersker og Beherskede\footnote{Den russiske Dramatiker \emph{Ostrovsky} har skildret Familielivet,
saaledes som det var for 40 Aar siden i Byerne i det indre Rusland.
Man ser her, i hvilken Grad selv Familien under den gamle
Tingenes Orden var et Tvangssystem. Hans „Oeuvres dramatiques“
have derfor ikke blot en overordentlig stor poetisk, men ogsaa en
meget stor kulturhistorisk og etisk Interesse.}.
Den almindelige Frigørelsesproces er ogsaa kommen Familien
til Gode; dens egentlige Væsen har først kunnet udfolde sig,
efter at det \emph{gamle}, patriarkalsk ledede Fællesskab var
ophørt. Ikke blot virke altsaa endnu de samme Kræfter i det
menneskelige Liv, som for Aartusinder siden; men der er
endog Grund til at tro, at de virke under heldigere
Vilkaar! ---

Og selv om Individerne udenfor Familien ere komne til
at staa mere i Kontraktforhold end i naturligt
Sammenhøringsforhold til hverandre, --- selv om contractus breder sig paa
Bekostning af status, --- saa føre dog atter disse forandrede
Forhold til nye Former for Fællesskab gennem \emph{fri
Association}. Ikke alle Associationer have nemlig Aktieselskabets
løse og selviske Karakter. Hvor der ikke blot er fælles
Formaal, men hvor der arbejdes og stræbes, lides og taales i
Fællesskab, vil der efterhaanden udvikle sig en
Fællesskabsfølelse, et Kammeratskab, som kan faa et større Sving og
besidde en større Styrke, end den rent egoistiske
Forventning om bestemte ydre Formaals Opnaaelse kan forklare.
Selv om Associationen fra først af indgaas fra de Enkeltes
Side af overvejende egoistiske Grunde, saa vil der under
Samlivet og Samvirksomheden foregaa en Forskydning af
Motiverne, saaledes at det oprindelige Formaal træder til%
% --- p. 155 ---
\opage{155}bage for Følelsen af Fællesskab. Ad denne Vej bliver den
Enkelte delagtig i Slægtens fælles Liv; naar han lægger
sit lille Lod i Vægtskaalen, er det ikke i Tanken om, hvad
han selv vil opnaa derved, men det er den fælles Skæbne,
han derved vil gøre Sit til at bøje i den Retning, han
anser for gavnligst for den Helhed, til hvilken han hører.
Det bedste og mest kendte Eksempel herpaa frembyder
Arbejderfagforeningernes Historie. Hvad vi kunne lære af
denne, har jeg et andet Sted (Etik p. 244)\footnote{2. Udgave p. 336.} udtrykt saaledes:
»Den kameratlige Følelse, Sympatien med Andre vokser. Den
Enkelte ser, at hans Adfærd under og udenfor Arbejdet, hans
Dygtighed, Flid, Selvbeherskelse og Sparsommelighed blive
af Betydning for langt Flere end ham selv. Der dannes saa at
sige en etisk Sfære omkring ham, en stor Familie, af hvilken
han føler sig som Led. Han lærer Broderskabet at kende,
som før næsten glemtes for Frihed og Lighed. Han lærer at
underordne sine egne Interesser under de fælles. Og han faar
gennem de fælles Erfaringer og det Overblik over Handels-
og Fabrikforholdene, der bestemmer Fagforeningernes Politik,
en klarere Opfattelse af Arbejdernes Stilling til de andre
Samfundsklasser, lærer baade sine Rettigheder og Pligter som
Led af Slægten bedre at kende. Det er kort sagt en hel
Opdragelse fra Egoisme til Sympati, fra blind Raahed til klart
seende Kraft, fra Kamp til fredelig Forhandling, som her
foregaar. Og alt dette sker ad Frihedens Vej. Der gives
intet bedre Svar til dem, der betragte vor Tid som en blot
Opløsningens Tid, end at henvise dem til den Samfundsdannelse
og den etiske Udvikling, som her er gaaet for sig.« ---

Kun med nogle faa Ord kan jeg her omtale den
\emph{psykologiske} Udvikling, som i Følge Tönnies gaar parallelt med og
tillige staar i real Sammenhæng med den sociale. Ogsaa den
var jo af patologisk Beskaffenhed: Forstand og Refleksion
afløse Indbildningskraft og Tro; i Stedet for Følelsens
umiddelbare Varme træder Tænkningens kølige Klarhed; det
aandelige Liv bevæger sig ikke længere efter naturlige Tilskyndelser,
men efter kloge Beregninger. Men ogsaa her gælder det, at
% --- p. 156 ---
\opage{156}nye naturlige og umiddelbare Tilstande og Motiver kunne
danne sig, efter at man har løsrevet sig fra det Instinktmæssige
og Overleverede. Hvad vi først kun omfatte med Erkendelsen,
kan tilsidst igen sætte Følelsen og Fantasien i Bevægelse. Jo
mere vi leve os ind i den vundne Erkendelse, des mere vil
ogsaa den kunne blive en Ramme eller et Grundlag for et nyt
umiddelbart aandeligt Liv. Den moderne Verdensopfattelse
har sin Storhed og Ophøjethed, som --- naar den forberedende
Forsken og Kritik, der foreløbigt lægger Beslag paa de bedste
Kræfter, har gjort sit Værk --- vil kalde Fantasien og Følelsen
til Virksomhed. Naar vi betænke, hvor kort den
videnskabelige Erkendelse har virket, og hvor snevre Kredse den endnu
er indskrænket til, er det med fuld Ret, vi antage, at vi leve
i en Overgangstid. Strid mellem Viden og Tro er karakteristisk
for en saadan Overgangstid. Men der er ingen Grund til at
antage, at denne Strid --- i sin nuværende Form --- skal vare
bestandigt. Den Harmoni, som nu den Enkelte møjsommeligt
tilstræber i sin Livsanskuelse, naar Bruddet med Traditionen
er sket, vil til sin Tid træde ud i fælles Symboler, i Tanker,
om hvilke store Kredse kunne samle sig. Vi leve i en »kritisk«
Periode; men den vil --- det er ikke en ubegrundet Tro ---
blive afløst af en »organisk« Periode. ---

Ville nu alle disse Betragtninger gøre noget Indtryk paa
den sociale Pessimist?

Som i Begyndelsen omtalt, ligger der bag Striden mellem
Optimisme og Pessimisme to forskellige Grundstemninger. Det
gælder selv her, hvor vi staa overfor en klar og energisk
Tænker, der raader over et rigt Materiale af Kendsgerninger
og over Analysen med stor Indsigt. Bag de Grunde, der
lade sig formulere, skjule sig dog stedse Stemningsgrunde,
der ikke lade sig formulere. Hvad et enkelt Punkt angaar
(en vis romantisk Forkærlighed for de primitive, naturlige,
begrænsede Livsformer), er der henpeget derpaa i det Foregaaende.

Desuden ere de Forhold, det her gælder om, saa
indviklede, at Ingen kan overskue eller beregne de forskellige
samvirkende og modvirkende Tendensers Magt. Hverken
Sociologien eller Psykologien vil her kunne opstille nogen
aldeles sikker Deduktion, der kunde afgøre Striden.

% --- p. 157 ---
\opage{157}Vor Modstander vil svare os: »Selv om Alt, hvad
her er fremført, er rigtigt, selv om der f. Eks. stadigt sker
Nydannelser i Retning af umiddelbart socialt og aandeligt
Liv, --- den modsatte Tendens, der gaar i Retning af
Mekanisering, Isolation og Intellektualisering, virker i stedse større
Omfang; maa den da ikke tilsidst løbe Linen ud, saa at Livet
opløses?«

Nu ja, det gælder en Prøve. Ingen alvorligt Tænkende
kan være blind for Faren. Livet er en stadig Kamp, og
dennes Resultat kunne vi ikke paa Forhaand kende.
Pessimismen har den store Fortjeneste at indskærpe, at vi ikke
over en Fremgang paa Overfladen glemme og miste det, der
dog tilsidst er Livets Kærne og Kraft. Deri har Tönnies
Ret, at hvis »Væsensvillie« og alt »Fællesskab« opløstes,
vilde Livet med det Samme blive et Skinliv, en tom Form.
Den udtrykkelige Bevidsthed om, hvad det gælder, vil kunne
skærpe Sansen for de virkelige Værdier og Viljen til at
bevare dem. Tilbage kunne vi ikke komme, og staa stille
kunne vi heller ikke; vi maa fremad og ville fremad, men en
af de vigtigste Betingelser for den virkelige Fremgang er,
at vi ikke miste hvad der én Gang er naat, og hvad der er
en stadig Fornødenhed for et sundt Liv.

\streg{}

% --- p. 158 ---
\opage{158}
\essay[„HEDENSKE SANDHEDSSØGERE“]{„HEDENSKE SANDHEDSSØGERE“.}

\begin{center}\emph{F. W. Farrar:} Hedenske Sandhedssøgere. (Seneca --- Epiktet --- Marcus
Aurelius). Oversat af Louise Hoffmeyer. Med et Forord af Rektor
Dahlenborg. København. 1891.\end{center}

\streg{}

\section*{I}

Gennem denne Bog ville danske Læsere kunne stifte
Bekendtskab med nogle af de mærkeligste Aander fra Oldtidens
Slutning. Den skildrer Tanke- og Følelseslivet hos Mænd,
der levede paa den Tid, da Kristendommen var i Færd med
at udbrede sig i Romerriget, men hvis hele Livsanskuelse
dog var aldeles uberørt af den, --- Mænd, som netop
repræsentere en af de Hovedformer, i hvilke den antike Livsanskuelse
fremtræder ved Oldtidens Slutning. Og da der i deres Maade
at tænke og føle paa viser sig en Inderlighed, en Renhed og
en Alvor, som Mange ikke vente at finde hos Forfattere, der
have udviklet sig uden for al Berøring med Kristendommen,
have deres Skrifter Betydning af et Eksperiment, ved hvis
Hjælp man kan prøve, hvor langt det aandelige Liv kunde
naa i Udvikling paa Oldtidens Vilkaar. Hvad enten man
betragter Kristendommen som overnaturlig Aabenbaring eller
som uvilkaarligt Udtryk for hvad der har rørt sig i
Menneskeaanden, har det dog sin store Interesse, at det er muligt at
anstille en saadan Sammenligning.

De Forfattere, med hvilke Bogen beskæftiger sig, tilhøre
Filosofiens Historie. Men de ere ikke rent teoretiske Tænkere.
Skøndt de ere udgaaede af en bestemt filosofisk Skole, den
stoiske, og hævde dennes Grundtanker, er det Karakteristiske
for dem deres Tænknings personlige Karakter. De frembyde
netop et interessant Eksempel paa, hvorledes en og samme
almindelige Tankeretning former sig forskelligt hos de for%
% --- p. 159 ---
\opage{159}skellige Personligheder. Deres Interesse er Livsfilosofi. De
ere, hvad Søren Kierkegaard kaldte subjektive Tænkere.

Deres Forskel fra den klassiske Tids græske Filosofer
beror især paa, at den umiddelbare Harmoni mellem
Elementerne i Livet og særligt i Menneskenaturen ikke længere
raader. De have ikke Sokrates’ Forvisning om, at klar
Indsigt er Et med Villen af hvad man har erkendt. Der har
vist sig at være større Afstand mellem Ideal og Virkelighed
end den klassiske Græcitet antog. Disse Tænkere ere i en
vis Forstand Udtryk for en Opløsningsproces, for den
Kendsgerning, at den klassiske Harmoni, der havde sat sig saa
store Mindesmærker i Kunst, Tænkning og Statsliv, maatte
opgives. Men fra den anden Side set betegne de et
Fremskridt, idet deres hele Opfattelse har en Inderlighed og en
Dybde, hvortil den græske Tænkning ikke naaede i sin
klassiske Tid. Mellem Tanken og Viljen skyder Følelseslivet sig
frem, en Verden af indre Bevægelser, der for den klassiske
Opfattelse med Lethed føjede sig ind i de Rammer, Tanken
kunde danne og Viljen omfatte i sin Handlen, men som nu
trængte til at vinde sit ejendommelige Udtryk og krævede
sin særskilte Tilfredsstillelse. Fra Betragtningen af den
ydre Tilværelse og fra Deltagelsen i Statslivet trak \emph{den
Enkelte} sig da nu tilbage for at fordybe sig i en indre Verden,
der var ikke mindre rig paa Begivenheder og Opgaver end
hin ydre. Alt kom nu for ham an paa, hvorledes denne indre
Verden formede sig, eller rettere, hvorledes han kunde forme
den. Hvad vedkom egentligt ham den ydre Verdens
Herlighed eller Ulykke, Uendelighed eller Sneverhed, naar den
ikke greb ind i denne \emph{hans} indre Verden? --- Stemningen,
Sindelaget blev nu det Afgørende, ikke den teoretiske Tænken
eller den udad gaaende Handlen. Deraf følger en ganske
særegen Kolorit over disse Forfatteres Skrifter, en
Stemningsnuance, som er ukendt i den klassiske Oldtids Literatur.

Men samtidigt med, at Synskredsen saaledes indsnevres,
for saa vidt den Enkelte trækker sig tilbage i sig selv,
udvides den, idet det gaar op for ham, at et saadant indre Liv,
som han selv fører, er der Mulighed for hos \emph{ethvert} andet
Menneske. Med Isolationen af den Enkelte gaar ved Oldtidens
% --- p. 160 ---
\opage{160}Slutning Udviklingen af en almindelig menneskelig Enheds-
og Fællesskabsfølelse Haand i Haand. De ydre Forskelle
mellem Trællen og den Frie, mellem Barbaren og Helleneren
blive forsvindende og betydningsløse i Forhold til Slægtskabet
i det indre Liv. Der gennembrydes Skranker, som i den
klassiske Oldtid antoges satte af Naturen selv. Det er en
i historisk og etisk Henseende overordentligt vigtig
Kendsgerning, at vi allerede før Kristendommens Fremkomst havde
\emph{almindelig Menneskekærlighed} (caritas generis humani)\footnote{Om det mærkelige Sted hos \emph{Cicero} (de finibus V, 23, 65), hvor dette
Udtryk første Gang forekommer, skal jeg paa et senere Punkt gøre
nogle historiske Bemærkninger. Ordet Menneskekærlighed eller
Menneskevenlighed (philantropia) forekommer vel tidligt hos Grækerne,
men faar først hos Stoikerne en dybere Betydning, og det er mere
ubestemt end det hos Cicero forekommende Udtryk: Kærlighed til
Menneskeslægten.} som
en klart bevidst Følelse.

Ogsaa i en anden Henseende udvides Sindet. Det
omfatter ikke blot alt Menneskeligt med sin Medfølelse, men
ogsaa hele Tilværelsen, hele Verden (Kosmos) med en inderlig
Følelse af Beundring og Slægtskab; det føler sig
gennemstrømmet af den samme Kraft, der rører sig i det store Hele,
opfatter Samvittighedens Lov som en enkelt, individuel Form
for den store Verdenslov og bøjer sig med villig
Resignation for de Skranker, Verdensordenen nu en Gang sætter.
Den stoiske Etik har fra sin første Tilblivelse (omtrent 300
Aar f. Kr.) denne \emph{religiøse} Karakter, der allerede hos en af
Stoicismens første Repræsentanter, Kleanthes, fandt Udtryk i
en begejstret Hymne.

Elementer og Spirer til den Udvikling, den stoiske
Filosofi saaledes er Udtryk for, fandtes ganske vist allerede
hos Sokrates i hans grundlæggende Opfattelse af det Etiske.
Men den Selverkendelse og Selvomsorg, som Sokrates havde
taget Ordet for, fremtræder dog hos Stoikerne, især de
senere Stoikere, om hvem den foreliggende Bog handler,
betydeligt uddybet. Det viser sig ikke mindst deri, at medens
Sokrates vel havde indskærpet Nødvendigheden af at føle sin
Uvidenhed og af at drage en skarp Grænse mellem, hvad
% --- p. 161 ---
\opage{161}man forstod, og hvad man ikke forstod, hvormed tillige fulgte
Grænsen mellem hvad man kunde gøre, og hvad man ikke
kunde gøre, --- saa indskærpes nu paa langt mere indtrængende
Maade Følelsen af egen Ufuldkommenhed som Betingelse for
Fremgang. For at tage et enkelt Eksempel: blandt Tegn
paa, at man gaar frem, anfører Epiktet (i sin »Haandbog«
48, 2), at man i sig selv ser en Fjende og Forræder. Dette
udsiges med en Betoning, der er langt skarpere og tyder
paa langt større Erfaring end den, der fandt Udtryk i den
sokratiske Ironi.

Der er i alt dette saa Meget, der kunde minde om
de kristelige Forestillinger, at det ikke er underligt, at man
har ment at maatte antage kristelig Paavirkning paa de
senere Stoikere. Man har f. Eks. antaget, at Seneca skulde
have staaet i Forbindelse med Apostelen Paulus. Den
Brevveksling, som gaar under disse to Mænds Navn, er dog et
Makværk, og den hele Antagelse usandsynlig og unødvendig.
I den Titel, Oversætterinden har givet det foreliggende Skrift,
og som er heldigere end Originalens Titel »seekers after God«,
da det Religiøse dog ikke er det eneste eller en Gang
Hovedsagen hos disse Mænd, er der ligesom en Genklang af
Forundring over, at disse Mænd, skøndt de stod udenfor
Kristendommen, dog ere naaede saa vidt som de ere, eller rettere
over, at de i det Hele vare inde paa denne Vej: Hedninger
--- og \emph{dog} søgende efter Sandhed! Denne Modsætning er i
Virkeligheden ikke tilstede. Iveren for at finde Sandheden
har været nok saa stor udenfor Kristendommen som indenfor
den. Ogsaa den ældre Oldtid vidner derom. De her
omhandlede Mænd fortsætte kun paa deres Vis de tidligere
Tænkeres Værk, drage deres ejendommelige praktiske
Konsekvenser af dette. --- Den valgte Titel har dog den Fordel, at
den vil vække Opmærksomhed for et Problem, der vil vinde
stedse stigende Betydning: Sandhedsbestræbelsens
Uafhængighed af religiøse Bekendelsers Skranker.

Den engelske Bog, som Oversætterinden med stor
Kærlighed og Omhu har gengivet paa Dansk, yder god Hjælp
og let Adgang til at stifte Bekendtskab med nogle af
Oldtidens ædleste Aander. Skøndt Forfatteren er Teolog, stiller
% --- p. 162 ---
\opage{162}han sig med stor Fordomsfrihed og Sympati overfor de gamle
Tænkere. Han er mild og skaansom i sin Dom over en
Mand som Seneca, hvis Karakter har været Genstand for
Angreb allerede i Oldtiden, og hvis Liv ganske vist ikke
altid stod i Overensstemmelse med hans Lære. Han er fuld
af Beundring for Epiktets Selvbeherskelse, Nøjsomhed og
frejdige Resignation og for Marcus Aurelius’ inderlige og
aandfulde Selvbetragtninger. Lidt naiv og ofte
forstyrrende er hans Sammenstilling af (ofte uheldigt valgte)
Skriftsteder med de gamle Filosofers Sætninger, og hans kritiske
og sammenlignende Slutningsafsnit er ikke meget dybtgaaende.
Men Bogens Styrke ligger i de interessante og udførlige
Biografier af de tre Mænd og i Uddragene af deres Skrifter,
der ere saa rigelige, at man derigennem lærer deres
Hovedtanker at kende. Bogen maa derfor anbefales til Enhver, der
interesserer sig for det menneskelige Aandslivs Udvikling.

\section*{II.}

Et enkelt Punkt i den foreliggende Fremstilling skal jeg
dvæle et Øjeblik ved, fordi det i sig selv har stor Interesse,
og fordi Forfatteren her baade befinder sig i en historisk
Vildfarelse og tillige ledes af en misforstaaet kristelig
Interesse. --- Han omtaler Seneca, Epiktet og Marcus Aurelius
som aldeles enestaaende Fænomener paa deres Tid, som
Draaber i »den almene Fordærvelses skummende Hav«. Men
han har aldeles intet Materiale, der kan berettige ham til at
antage en særligt stor sædelig Fordærvelse ved Oldtidens
Slutning. Berettigelsen kan ikke søges i Senecas og Andres
Deklamationer, lige saa lidt som i Apostlen Paulus’ Udtalelser
i Begyndelsen af Romerbrevet. Tiden var i en vis Forstand
en Opløsningsperiode. Men skønt den ikke frembød mange
store Skikkelser, frembød den dog mange rene og ædle
Karakterer, der under ugunstige og ofte pinagtige Forhold
bleve deres Livsideal, som ikke var lavt, tro. Og fra selve
den Anerkendelse, der blev de ogsaa af Forfatteren beundrede
Mænd til Del, kan der hentes et Modbevis imod hans
Antagelse. Saaledes siger en Forsker, der med Hensyn til et
Spørgsmaal som dette hører til de mest sagkyndige, \emph{Ludwig}
% --- p. 163 ---
\opage{163}\emph{Friedländer} (Darstellungen aus der Sittengeschichte Roms in
der Zeit von August bis zum Ausgang der Antonini. 5. Aufl.
III. p. 676): »En Tid, der ved egen Kraft hævede sig til
højere og renere sædelige Anskuelser end hele den tidligere
Oldtid, --- der ikke blot frembragte en Musonius [en med
Seneca samtidig Stoiker], Epiktet og Marcus Aurelius, men
i hvilken disse Forkyndere af en mild, ægte menneskelig
Sædelære ogsaa fandt den mest almindelige Beundring og
deres Lærdomme almindelig Udbredelse, kan ikke have været
en Tid af det dybeste sædelige Forfald, saaledes som den
ofte er bleven kaldet. Naar der overhovedet ikke gives nogen
Maalestok for Sædeligheden i en Periode, selv om denne er
nok saa godt kendt, saa gælder det allermest for disse
Aarhundreder [de to første efter Kristus], fra hvilke der kun
foreligger enkelte, dels til bestemte Omraader begrænsede,
dels farvede eller ensidige Beretninger.« --- Og den samme
Forskers Undersøgelser godtgør, at som det forholder sig
med Fordommen om et særligt stort sædeligt Forfald ved
Oldtidens Slutning, saaledes forholder det sig med Antagelsen
af et almindeligt Forfald af det religiøse Liv paa denne Tid.
Mange Vidnesbyrd godtgøre en inderlig Fastholden af den
gamle Tro og de gamle religiøse Skikke hos Folket i det
Hele. Man kan ved denne Periode lige saa lidt slutte fra
visse literært dannede Kredse, som man ved det 18.
Aarhundrede fra de af de franske Encyklopædister paavirkede
Kredse kan slutte til hele Folkets religiøse Tilstand. Der
viste sig endog netop paa den Tid, der her er Tale om, en
stærk religiøs Bevægelse, der dels førte til Genoplivelse af
den gamle Kultus, saa at f. Eks. Oraklerne igen begyndte
at tale efter længe at have været stumme, dels til Udbredelse
af orientalske Religionsdyrkelser i Romerrigets vestlige Egne.
Selve Kristendommen og dens Udbredelse er, historisk set,
kun en enkelt, skønt den vigtigste Side af hele dette
Fænomen. Den haarde Kamp, Kristendommen havde at bestaa,
viser ogsaa, at den gamle Tro ikke var borte. Selv efter at
Kristendommen var bleven Statsreligion, var flere
Aarhundreders grusomme og fanatiske Forfølgelse endnu nødvendig
for at fortrænge den gamle Religion --- fra Overfladen. I
% --- p. 164 ---
\opage{164}mangehaande Forklædninger vedbleve de gamle Forestillinger
at virke og modificerede paa mange Maader de kristelige
Lærdomme, i mere filosofisk Form hos de Lærde, i mere
mytologisk Form hos de Ulærde.

En Misforstaaelse, mener jeg, gør sig her gældende ved
Antagelsen af et stort sædeligt og religiøst Forfald, for saa
vidt man mener derved at forherlige Kristendommen. Da
denne efter den ortodokse Opfattelse ikke blot skal træde i
Forhold til en bestemt Tid, da den heller ikke er Frugt af
bestemte historiske Betingelser, --- og da efter samme
ortodokse Opfattelse den menneskelige Syndighed til alle Tider
er lige bundløs og hjælpeløs, saa vilde det netop for dem,
der hylde denne Opfattelse, være betænkeligt at antage et
særligt Forfald i den Periode, da Kristendommen fremtraadte.

\section*{III.}

De tre Mænd, hvis Liv og Tanker fremstilles i den
foreliggende Bog, tilhørte, som allerede bemærket, den stoiske
Skole. Det var især den praktiske Livsanskuelse, som denne
Filosofi begrundede, der skaffede den saa mange Tilhængere
i den romerske Verden. Den hævdede den Enkeltes
Personlighed overfor den ydre Verden og indskærpede til samme
Tid mandig Resignation overfor Skæbnens Gang. Begge
Dele hang nøje sammen; ti naar Livets højeste Værd ligger
i det Indre, bliver alt Ydre tilsidst ligegyldigt. Der var
Noget her, som tiltalte den romerske Nationalkarakter, og vi
finde da ogsaa en hel Række af fremragende Statsmænd fra
Republikens Tid blandt den stoiske Skoles Tilhængere. Den
stoiske Filosof Panaitios var Scipio Afrikanus den Yngres
og Lælius’ Ven, og Cæsars ædleste Modstander, Cato den
Yngre, tilhørte afgjort denne Retning. I Kejserdømmets
første Tider dannede en hel Skare stoisk dannede Mænd en
Opposition mod det afsindige Tyranni paa Herskersædet, en
Opposition, der ofte førte til Døden. Den Mandighed,
hvormed saa Mange af disse Mænd gik Døden i Møde, og som
Tacitus har skildret i en Række berømte Scener, var et
Udslag af den stoiske Filosofis Aand. Men Stoicismen fik
ogsaa anden Indflydelse end den at lære, hvorledes man kan
% --- p. 165 ---
\opage{165}trodse en ugunstig Skæbne og dø paa ædel Vis. Den virkede
snart i ikke ringe Grad paa Regeringens og Lovgivningens
Aand. Et stort Antal humane Foranstaltninger fra den ældre
Kejsertid udledes vistnok med Rette af denne Skoles
Indflydelse. Og af de tre Mænd, den foreliggende Bog handler
om, øvede de to en bestemt Indflydelse paa Statsmagtens
Anvendelse.

\emph{Lucius Annæus Seneca} (død 65 e. Kr.) var Kejser Neros
Lærer og havde --- sammen med Prætorianpræfekten Burrhus ---
i Begyndelsen af sin Elevs Regeringstid en stor og god
Indflydelse paa ham. Nogle Aar efter blev han styrtet og
maatte paa Kejserens Befaling dræbe sig selv. Senere kom
der en Reaktion mod Filosoferne; de bleve gentagne Gange
(under Vespasian og under Domitian) forjagne fra Rom. Men
der kom atter bedre Tider, og det næste Aarhundrede saa
endog en Filosof paa Kejsertronen: \emph{Marcus Aurelius Antoninus}
(født 121 e. Kr., død 180). --- At det dog ikke blot var i de
højere Kredse, at Stoicismen vandt Indgang, ses af, at den
tredie af de store stoiske Forfattere fra Kejsertiden, \emph{Epiktetos}
(der levede i sidste Halvdel af det første Aarhundrede), var
en Slave, der vandt aandelig Frihed ved at leve sig ind i
denne Lære, hvis Forkynder han siden blev.

Ikke mindre stor end Forskellen mellem de tre Mænds
ydre Kaar var Forskellen mellem deres Aand og deres Sprog,
indenfor den fælles stoiske Ramme. Seneca er den
frugtbareste Forfatter af de tre. Han har ofte Mere af Retoren end
af Filosofen, og en vis Sentimentalitet kan han ikke frikendes
for. Men en rig Erfaring paa det aandelige Omraade, en virkelig
Sjælesørgeromhu for sine Venner og en alvorlig Fordybelse i
etisk Tænkning gøre mange af hans Skrifter interessante og
lærerige, ligesom de ere af stor Betydning for Kendskabet
til den stoiske Lære, da vi ingen ældre stoiske Skrifter have.
Selv om han maaske har vist Karaktersvaghed i sin
offentlige Optræden, og selv om hans yppige Omgivelser efter
Nogles Mening ikke ganske stemmede med den strenge Lære,
han forkyndte, saa gaar der dog gennem hans Liv en
alvorlig Stræben efter det Sande, og den Værdighed, der maaske
undertiden svigtede ham under Kampen for Tilværelsen paa
% --- p. 166 ---
\opage{166}Hoffets glatte Gulve, fandt han i ethvert Tilfælde, da det
gjaldt om at kunne dø.

I stor Modsætning til den rige Hofmand Seneca, hvis
yppige Omgivelser havde en Analogi i hans blomstrende og
retoriske Fremstilling, staar den sygelige og fattige Epiktet,
der indskrænkede baade sine literære og fysiske
Fornødenheder til det mindst Mulige. Han har selv Intet skrevet,
men en af hans Disciple optegnede hans Foredrag og gjorde
igen af dette Værk (hvis ene Halvdel er bevaret) et Uddrag,
»Epiktets Haandbog«, den korteste og simpleste Fremstilling,
vi have af den stoiske Lære. Det er den Enkeltes Etik, der
her læres. Gennem Udholdenhed og Afholdenhed vinder
Mennesket sig selv. Den sikreste Vej til aandelig Frihed
er den at indskrænke Fornødenhederne. Vær ikke bitter, om
Retterne ved Livets Bord gaa dig forbi, --- men vis dig endog
saa overlegen, at du lader en Ret gaa forbi, selv om den
bydes dig! Og dog er det ingen mørk Askese, her prædikes.
Livets og Tilværelsens Herlighed priser Epiktet med
inderlig Begejstring; ti netop for den, der gør sig aandeligt fri i
Selvbeherskelsens store Skole, bliver det klart, at der intet
virkelig Ondt gives i Verden. Og ligeledes læres den
inderligste Kærlighed til Andre, selv til Fjender og til dem, der
mishandle os.

Den tredie Forfatter, Marcus Aurelius, forbinder Epiktets
Klarhed og Simpelhed med Senecas glimrende Forstand og
retoriske Kraft, ligesom han formaaede at bevare
Selvbeherskelsen og det milde Sind paa Kejsertronen og at
fastholde Begejstringen over Livet trods den Modstand og de
Skuffelser, han erfarede, og trods alt det Smaa og det Smaalige,
hans Ærlighed og skarpe Blik ikke oversaa. Hans
Selvbetragtninger ere i en vis Henseende den mærkeligste Bog, vi
have fra Oldtiden. Der er noget Moderne i, at en Mand
saaledes stille for sig selv gør Rede for sine Tanker og
Stemninger ved Livet. Og der er en Dybde i Stemningen,
en ejendommelig Tone over hans Udtalelser, som vi ikke finde
i noget andet Oldtidsskrift. Stoicismen fremtræder tillige
med større Mildhed hos ham end endnu hos Seneca og Epiktet,
der begge have noget af Prædikantens Strenghed. Marcus
% --- p. 167 ---
\opage{167}Aurelius havde kun med sig selv at gøre, og naar han i
Paladset eller i Lejren trak sig tilbage til sit Lønkammer,
var det --- trods den Skarphed, hvormed han saa paa Tingene
i deres materielle Nøgenhed --- ikke for at give sit Tungsind
eller sin Harme Luft, men for at bestyrke sig i Følelsen af
det menneskelige Fællesskab og dets Opgaver. Hans stoiske
Kosmopolitisme hindrer ham ikke i energisk at føle sig som
Romer: »som Antoninus har jeg Rom, som Menneske hele
Verden til Fædrenestad; kun det, der gavner disse to Steder,
kan være godt for mig.«

\section*{IV.}

Det vil maaske have sin Interesse at undersøge,
hvorledes den Tankegang og Livsanskuelse, vi møde hos disse
Mænd, og som i mange Træk afvige i saa høj Grad fra den
klassisk græske, har udviklet sig. Kun korte Antydninger
kunne her gives, da, som allerede bemærket, hele den ældre
stoiske Filosofis Literatur er gaaet til Grunde, saa at vi ere
indskrænkede til at benytte andenhaands Kilder.

Den stoiske Filosofi grundlagdes af \emph{Zenon} fra Cypern,
der lidt efter Aar 300 f. Kr. lærte i Athen i en Buegang
(stoa), efter hvilken Skolen fik Navn. Det var en Tid, der
var ikke saa lidt forskellig fra den, da Sokrates og Platon
virkede i Athen. Grækenland havde mistet sin Frihed og
var kommen under makedonisk Herredømme, efter at de
enkelte græske Stater vare blevne svækkede og Sansen for
de offentlige og nationale Interesser havde tabt sig. Samtidigt
havde Alexanders Erobringer aabnet en større Horizont og
banet Vejen for et nyt og mere omfattende Syn paa de
menneskelige Forhold. Den unge makedoniske Hærfører havde
nedbrudt Skrankerne mellem Hellener og Barbarer, ikke blot
ved at forene dem under ét Scepter, men ogsaa ved at
anlægge græske Byer rundt om i Orienten og foranledige
Vekselvirkning mellem græske og orientalske Skikke og
Forestillinger. Man lærte nu fra græsk Side i Praksis, at
der ikke var den principielle Forskel mellem Hellener og
Barbarer, som endnu Alexanders store Lærer Aristoteles havde
hævdet. Karakteristisk er den Fortælling (der findes i et
% --- p. 168 ---
\opage{168}Skrift, som, maaske med Urette, tillægges Plutark), at
Aristoteles skal have raadet Alexander til at behandle Hellenerne
som Hellenere, Barbarerne som Barbarer, men at den unge
Erobrer, hvis ærgerrige politiske Planer faldt sammen med
hvad Humaniteten fordrede, ikke adlød sin Lærer heri. Gennem
Samlivet og den fælles Skæbne uddannedes der nu en fælles
Følelse; Forestillingen om en almindelig Menneskehed, der
omfattede alle Nationer trods deres Forskelligheder, blev nu
til. De overvundne Nationer fandt Trøst i den Tanke, at de
ikke vare bukkede under for at gaa op i et enkelt andet
Folk, men for at staa som Led af den hele Menneskehed\footnote{Paa et enkelt Punkt, der historisk gaar langt tilbage i Tiden, førte
ogsaa den klassiske græske Etik udover den lokale og nationale
Begrænsning af Menneskelighedsfølelsen, nemlig i
Gæstevenskabsforholdet og i det Hele i Forholdet til de Fremmede, der fra gammel
Tid betragtedes som staaende under Gudernes særegne Beskyttelse.
Pligterne mod den hjælpeløse Fremmede stode som nogle af de
højeste Pligter. At vise den vildfarende Vandrer Vej, at give Andre
Ild og Vand, naar de mangle, og ikke at raade dem Andet end
man selv holder for gavnligt, var Pligter, der indskærpedes ved
særlig, ærværdig Overlevering. Her gjaldt altsaa ikke den Skranke,
der ellers sondrede Menneskene i forskellige etiske Verdener. Den
Fremmede og den Bedende stod som Menneske, hverken som
Hellener eller Barbar. Der var her en Spire, som kunde udvikles videre,
og ogsaa blev udviklet.}.
--- Den samme Kreds af Erfaringer blev der Lejlighed til
at gøre i endnu større Kredse og under varigere og fastere
Forhold, da hele den civiliserede Verden gik op i Romerriget.

Den stoiske Filosofi, som betoner det Indres, Fornuftens,
Villiens, Uafhængighed af det Ydre, var særligt kaldet til at
udvikle den ad praktisk Vej forberedte Forestilling om en
almindelig Menneskehed. De nationale Forskelligheder maatte
staa som noget Ydre og Tilfældigt, naar man fandt Menneskets
egentlige Væsen ikke i ydre Livsformer og Livsvilkaar, men
i det, der bevæger sig i Sind og Tanke. Der viser sig her
et interessant Vekselforhold mellem historisk Erfaring og
filosofisk Tænkning. Opdagelsen af den fælles Menneskenatur
betingedes og lettedes ved den faktiske Nedbrydelse af de
nationale Skranker, og den filosofiske Tænkning udtalte for
% --- p. 169 ---
\opage{169}saa vidt kun den historiske Erfarings Resultat. Men paa
den anden Side var denne filosofiske Tænkning allerede
tidligere grundlagt. Zenon, Stoicismens Stifter, var paavirket
af den store, fra Sokrates et Aarhundrede forud udgaaende
Bevægelse. Efter en Fortælling skal Zenon have faaet
Interesse for Filosofi ved at læse Xenofons Skrift om Sokrates.
Det var den sokratiske Tanke om Selverkendelsen som Grundlag
for al Erkendelse og Livsførelse, der hos Stoikerne udvikledes
videre i Konsekvenser, der ikke forud vare dragne. Og den
historiske Erfaring, der gjordes efter Alexanders Erobringer,
fik nu først sin højere Betydning ved, at Ideen om en
almindelig Menneskehed udvikledes. Stoikerne uddybede denne
Tanke ved den hele Verdensanskuelse, ifølge hvilken
Tilværelsen udgør en Enhed, idet Verdenssjælen lever og rører
sig i de enkelte Væsener. Hele Verden, lærte de, er en stor
Stat, fælles for Guder og Mennesker, besjælet af en og samme
uendelige Kraft og underlagt en og samme store Lov, der
fremtræder for hvert enkelt Væsen som dets egen indre Lov.
De søgte herved at forbinde den enkelte Personligheds
Selvstændighed med Hengivelsen til det Universelle: idet den
Enkelte følger sin Natur, følger han den almindelige Natur,
i hvilken han kun er et enkelt Led.

Hos \emph{de første} Stoikere betones den personlige Selvhævdelse
og Uafhængighed mere end Hengivelsen til det Universelle,
og, hvad der hænger sammen hermed, indenfor den enkelte
Personlighed fremhæve de Fornuften og Villien som det
Vigtigste. Deres Etik faar derved en streng og negativ Karakter.
Den Vise, deres Ideal, skal være hævet ikke blot over alle
ydre Omskiftelser, men ogsaa over Bevægelserne i hans eget
Indre, over Lyst og Smerte, Kærlighed og Vrede, over alle
Affekter, der betragtedes som noget Ondt, da de ophævede
Selvstændigheden, den absolute Hvilen i sig selv, og førte
til Afhængighed af Noget, der laa udenfor Villien.
Hoveddyden var Retfærdighed, der opfattedes som Modsætning
til Kærlighed. --- Senere Stoikere (\emph{den mellemste stoiske Skole}
i 2. og 1. Aarhundrede f. Kr.) ændrede denne Opfattelse og
uddybede Humanitetsfølelsen, ja gjorde egentlig først
Menneskelighedsbevidstheden til en Menneskelighedsfølelse. Des værre
% --- p. 170 ---
\opage{170}vide vi saare lidet om denne Udvikling, og kun paa anden
eller tredie Haand fra senere Forfattere. Den stoiske Tænker,
hvis Indflydelse her især har været af Betydning, er den
allerede som Afrikanus den Yngres Ven nævnte \emph{Panaitios}.
Paa Grundlag af den naturlige Trang til Samliv, som allerede
Forgængerne havde antaget, og som Aristoteles havde udtrykt
i den bekendte Sætning om Mennesket som et til Samfund
anlagt Væsen (zoon politikon), udviklede Panaitios den etiske
Betydning af Menneskekærligheden som sammenholdende
Kraft. Mellem Retfærdigheden og Kærligheden stillede han
Billigheden, der ikke holder sig til den formelle Ret, men
tager alle individuelle Hensyn med i Betragtning. Og medens
de ældre Stoikere egentligt kun antog et sandt
Venskabsforhold mellem de ideale »Vise«, hos hvem Fornuften var
naat til fuld Udvikling, fandt Panaitios i
Menneskekærligheden et levende Baand mellem alle Mennesker, selv om de
ikke have naat det absolute Ideal. Af særlig Interesse er
det, at han ikke blot har holdt sig til psykologiske Beskrivelser
og etiske Postulater, men har givet Grundtrækkene til en
Sympatiens Udviklingshistorie, der minder om den moderne
Udviklingslære. I sit Skrift »om det højeste Gode og Onde« (de finibus)
siger Cicero (V, 23, 65): »Af alt Godt er intet saa
beundringsværdigt og omfattende som Menneskenes indbyrdes
Forbindelse, som deres Fællesskab i det, der er dem nyttigt, og
som Kærligheden til Menneskeslægten, en Kærlighed, som
har sin Kilde i Forplantningen, idet Afkommet elskes af
Forældrene og hele Huset holdes sammen ved Forholdet
mellem Ægtefæller og mellem Forældre og Børn, men som
derpaa efterhaanden udbreder sig udad til, først ved nærmere
og fjernere Slægtskabsforhold, saa ved Venskabsforbindelser,
ved Naboforhold og ved Sammenslutning i borgerlige og
offentlige Forhold, indtil den tilsidst omfatter hele Menneskeslægten.«
Der er i disse Ord antydet en successiv Udvikling af
Menneskelighedsfølelsen fra snevrere til videre Kredse, betinget
ved den Maade, hvorpaa Menneskene af indre Drift og ved
Forholdenes Magt komme til at leve, virke og lide sammen,
en Udvikling, der ganske stemmer med den Maade, paa
hvilken den moderne Sociologi opfatter Sagen. Den første
% --- p. 171 ---
\opage{171}Del af denne Udviklingsproces har allerede Aristoteles skildret
paa klar og interessant Maade, og han har angivet den Grundlov,
som gør sig gældende ved hele denne Udviklingsgang: at
den uvilkaarlige Indøvelse i en Virksomhed eller et Livsforhold
frembringer en tilsvarende Følelse af Tilfredsstillelse derved.
Jeg har derfor (i min »Etik«) kaldet denne vigtige Lov »det
aristoteliske Princip.« Den sidste Del af den i Stedet hos
Cicero antydede Udviklingsproces, Overgangen fra den
nationale Sammenhøren til hele Menneskeslægtens Enhed, skyldes
Betingelser og Ideer, der først fremkom efter Aristoteles’ og
Alexanders Tid, først i den østlige Del af den da bekendte
Verden, senere, ved Romerrigets Udbredelse, i endnu større
Kredse. Og det er sandsynligt, at Fremstillingen hos Cicero
bygger paa Tanker, der første Gang saa Lyset i den stoiske
Skole, vistnok hos Panaitios\footnote{Den græske Filosof, som Cicero følger i det Afsnit, hvorfra det
anførte Sted er taget, er \emph{Antiochos} fra Askalon, hvem han selv havde
hørt i Athen. Denne Mand tilhørte den akademiske Skole, som var
grundlagt af Platon, men han søgte, ofte paa meget udvortes Maade,
at forbinde forskellige Skolers Læresætninger og har optaget meget
baade fra den aristoteliske og den stoiske Skole. Der kunde da
være Spørgsmaal om, hvilken af disse to Kilder vi nærmest maa
antage for den anførte Tankegangs Udspring. Madvig synes (i sin
Udgave af de finibus, Noten til III, 19, 62) nærmest at antage en
stoisk Oprindelse. Senere har Bernays søgt at vise, at det er fra
den aristoteliske Skole, Ideen om den almindelige
Menneskekærlighed skriver sig. Men de nyeste Undersøgelser (O. Apelt: Beiträge
zur Geschichte der griechischen Philosophie 1891 p. 354, --- A. Schmekel:
Die Philosophie der mittleren Stoa. 1892 p. 398 kfr. p. 220) pege
hen paa, at Ciceros Hjemmelsmand Antiochos her nærmest er
afhængig af Panaitios og den Udformning, han gav den ældste stoiske
Lære.}. Hos denne Filosof, og i endnu
højere Grad hos hans Efterfølger \emph{Poseidonios}, fremhævedes i
det Hele Følelseslivet og dets Betydning langt mere end hos
de første Stoikere. Man søgte at mildne den skarpe
Modsætning, disse havde fastslaaet mellem Fornuften og den aktive
Villen paa den ene Side, Følelse og Affekt paa den anden
Side. Dette hang naturligt sammen med, at de søgte at
ophæve den absolute Modsætning, Skolens Stiftere havde fast%
% --- p. 172 ---
\opage{172}slaaet mellem de ideale »Vise« og alle andre Mennesker.
Der har i det Hele været Noget i denne Bevægelse indenfor
den stoiske Skole i de sidste Aarhundreder f. Kr., som er
beslægtet med den Indflydelse, Rousseaus Hævden af Følelsens
Magt og Ret overfor de andre Sider af Sjælelivet øvede i
forrige Aarhundrede baade paa Psykologiens og Etikens
Udvikling og paa Literaturen i det Hele.

Den Ændring og Udvikling, som har fundet Sted i den
mellemste stoiske Skole, og ved hvilken der lagdes en stigende
Vægt paa det Universelle, Almenmenneskelige i Modsætning
til det blot Individuelle og Særskilte, paa Følelsen i
Modsætning til Erkendelsen og paa Kærligheden i Modsætning til
Retfærdigheden, hører til det Allervigtigste, der foregaar i
Historien i de sidste Aarhundreder f. Kr. Baade paa Grund
af den successive Maade, hvorpaa denne Udvikling er gaaet
for sig, og paa Grund af, at den foregaar i Menneskenes
Sind og Tanke, trænger den sig ikke frem i Historiens
Forgrund, hvor de larmende politiske Begivenheder foregaa.
Men derfor kan dens Betydning være nok saa stor og varig.

\emph{De sidste Stoikere}, som repræsenteres ved de tre Romere
Seneca, Epiktet og Marcus Aurelius, staa nu ikke for os
som et uforklarligt historisk Fænomen. Det er let at
forstaa, hvorledes de, saaledes som Tidsforholdene udviklede
sig efter den romerske Republiks Undergang, kom til at
lægge en endnu større Vægt paa det, som allerede den
mellemste stoiske Skole (Panaitios og Poseidonios) havde
betonet. De vilde ikke med Rette kunne kaldes Stoikere,
dersom de ikke med Energi hævdede den indre Personligheds
Uafhængighed af Skæbnens Gang og de ydre Goder og
Onder. I den Henseende mangler det hos dem ikke paa
skarpe og krasse Udtalelser. Men ejendommelig for dem er
den store Inderlighed, den dybe Stemning, hvormed de udtale
de stoiske Tanker. Et enkelt Eksempel vil vise Karakteren
af denne Aandsretning, som især kommer for Dagen paa en
smuk og ofte gribende Maade hos den af disse Mænd, der
baade som Personlighed og Forfatter staar højest, Marcus
Aurelius. Han udtaler et Sted i sin Dagbog, at Menneskene ere
anlagte til fælles Virksomhed, til Samarbejde, ligesom Legemets
% --- p. 173 ---
\opage{173}Lemmer, og fortsætter: »Denne Tanke vil træde klarere frem for
dig, naar du hyppigt siger til dig selv; jeg er \emph{Led} (melos) i de
aandelige Væseners Samfund, end naar du (ved en lille
Forandring i Bogstaverne) blot kalder dig en \emph{Del} (meros). Hvis
du blot er Del, ikke Led, \emph{saa elsker du endnu ikke
Menneskene af Hjertet;} det at gøre Andre vel glæder dig endnu
ikke saaledes, at det helt betager dig, men du gør det blot
af Pligt, endnu ikke med den Følelse, at du derved gør dig
selv vel.« (Marcus Aurelius’ Optegnelser VII, 13). Den
samme Inderlighed fremtræder, som allerede omtalt, i
Fordringen om at elske sine Fjender og i den religiøse Opfattelse
af Verdenslivet. Det er den indre Rigdom og Fylde, der
hæver disse Mænd over Menneskers Overgreb og lader dem
overse Tilværelsens mørke Sider. Den Kraft, de følte hos
sig til at overvinde det Ydre, Skiftende og Smertefulde, stod
for dem ikke i Strid med den Kraft, der rørte sig i
Verdenslivet, men var en Forgrening af denne; derfor blev Harmonien
deres sidste Tanke, og de vedblev at være Grækere, trods
den skarpe Modsætning mellem det Indre og det Ydre, der
var en saa væsentlig Del af deres Lære. Ikke med Trods,
men med frejdig Beredvillighed venter den kejserlige Filosof
hvad Skæbnen vil bringe ham, og hans Frejdighed slaar
undertiden ud i en Hymne, som paa det berømte Sted, hvor
det hedder: »Alt er harmonisk for mig, som er harmonisk
for dig, o Verden. Intet kommer for tidligt eller for sent
for mig, naar det kommer betimeligt for dig. Alt, hvad dine
Aarstider bringe, er Frugt for mig, o Natur. Fra dig kommer
Alt, i dig er Alt, til dig vender Alt tilbage. Naar Digteren
siger: kære Kekrops’ Stad\footnote{Kekrops var efter Sagnet Athens Grundlægger.}, skulde jeg da ikke sige: kære
Zeus’ Stad!« (Optegnelser IV, 23). --- Den inderlige
Optimisme hvilede ikke paa nogen Betragtning af den ydre Natur.
De ældre Stoikere havde endnu haft Interesse for Naturforsken
og for historisk Videnskab, og de havde kun ved skrøbelige
og sofistiske Grunde (som til Dels gaa igen i nyere Tiders
teologiske Ræsonnementer) formaaet at hævde deres
optimistiske Opfattelse overfor Erfaringen. Fra al saadan udad%
% --- p. 174 ---
\opage{174}gaaende Interesse have de seneste Stoikere trukket sig
tilbage; derfor slippe de for store og æggende Problemer. De
løse ikke den Opgave, hvorledes det indre Livs Inderlighed,
Frejdighed og Kraft kan hævdes sammen med Erfaringen om
Lidelsen og Kampen i Verden. Men denne Ufuldkommenhed
opvejes ved den faste Overbevisning, hvormed de have
henpeget til det indre Liv som vor nærmeste og væsentligste
Verden og hævdet det som et Omraade for sig selv. »Se
indad«, siger Marcus Aurelius (VII, 59); »der inde er det
Godes Kilde, som altid kan sprudle frem, naar der blot graves!«

Mænd som Epiktet og Marcus Aurelius staa som smukke
Vidnesbyrd om den Evne, kraftige og klare Aander have
til at blive sig selv tro og bevare deres Frihed under
Tidsstrømninger, der førte Mange til at kaste sig i Armene paa
Mystik og Overtro. De fra Orienten kommende gamle og
nye Religioner lokkede Mange til sig. Mange følte sig trætte
--- Nogle af Udsvævelser, Andre af Kunst og Literatur, Andre
igen af Politik. Og naar man er træt, bøjer man villigt det
Hoved, som man maaske før bar noget mere knejsende, end
den indre Kraft virkeligt berettigede til I Epiktet og i Marcus % PRINTED AS IS: „til“
Aurelius se vi, at der var Mænd baade i de laveste og i de
højeste Kredse, som ikke kunde og ikke vilde opgive sig
selv, og som heller ikke behøvede det. Deres mandige
Livsanskuelse har ikke blot i sit Indhold Meget, der endnu fra
vor Tids Synspunkt set er af stor Interesse og af blivende
Betydning; men den staar særligt ved sit Forhold til Datidens
Strømninger som et mærkeligt og lærerigt Fænomen.

\section*{V.}

Det vilde være en vanskelig Opgave at paavise, hvori
egentligt Følelseslivet hos Mænd som de her omtalte var
forskelligt fra det, som Kristendommen førte med sig. Naturligvis
bliver Opgaven let, naar man blot laver en ny Slags Følelse
efter enhver Rubrik i Dogmatiken. Paa den Maade vil det
være let at vise, hvilken Rigdom der med Kristendommen
kom ind i Følelseslivet; dette blev paa denne Maade ganske
vist beriget med adskillige nye Arter. Langt vanskeligere
er Opgaven for den, der er klar over, at saa vigtigt og af%
% --- p. 175 ---
\opage{175}gørende det kan være, at der kommer nye Forestillinger frem
om Livets Formaal og Betydning, saa har Følelseslivet dog
sine egne Love, sine naturlige Grundforhold, som ikke kunne
ændres eller udvides i det Ubestemte. Det er derfor langt
lettere at paavise, hvilke nye Forestillinger Kristendommen
indførte, end hvilken Ændring i Følelseslivet den medførte.
Og Opgaven bliver saa meget des vanskeligere, naar man
lægger Mærke til, hvilken Inderlighed og hvilket Omfang
Følelseslivet kunde udvikle sig til hos det »naturlige«
Menneske. Dog er jeg overbevist om, at det var en vigtig, i sin
Art enestaaende Ændring, der i sin Bestaaen viser sig at
være uafhængig af de dogmatiske Forestillinger, som fra først
af fremkaldte den. Jeg skal forsøge at fremstille, hvorledes
jeg tror, at denne Ændring kan beskrives, uden at jeg tør
paastaa, at jeg dermed angiver den eneste Maade, hvorpaa
dette Spørgsmaal kan besvares. En Sammenligning mellem
Stoicismen i dens sidste Udviklingsform og Kristendommen
vil føre til hvad jeg anser for det Afgørende.

Stoicismens Hovedfejl var den megen Reguleren, den
stærke Betoning af Nødvendigheden af vilkaarlig Indgriben
i Sjælelivets Gang. Til Grund herfor laa en Mistillid til
de uvilkaarlige Bevægelser i Sindet. Stoikeren erklærede
alle Sindsbevægelser for onde. Den Vise er hævet over
Stemningernes Bølgegang og handler ud af klar og stræng
Bevidsthed. At følge Stemningen, at lade sig lede af
Sindsbevægelsen førte efter Stoikernes Mening uundgaaeligt til
Fald. Og hvad det gjaldt om, var jo dog det, at staa som
sin egen Herre, selv sidde til Rors i sin Fart over Livets
urolige Hav. Som saa ofte hænger ogsaa her Fejlen sammen
med selve det Store og Udmærkede, med det, der gør de
stoiske Karakterer enestaaende og giver dem Ophøjethedens
Præg. Men Sagen er netop den, at Stoikerne kun troede at
kunne naa denne ideale Stilling ved stadigt at hæmme den
uvilkaarlige Trang til Hengivelse. En Hovedsætning i den
stoiske Etik var og blev det, at man skal hæmme
Forestillingernes Følelsesvirkninger. Tingene vække Forestillinger
hos os; det kan ikke undgaas, men vi kunne undgaa at lade
disse Forestillinger faa Magt over os ved stadigt at sige: det
% --- p. 176 ---
\opage{176}er kun Forestillinger (eller som Stoikerne ofte sagde: Dogmer)!
Hindre vi Forestillingerne i at rive os med, saa have Tingene
ingen Magt over os. Dette Forhold til Forestillingerne er i
Virkeligheden det Eneste, som staar i vor Magt, og derfor
det eneste Afgørende.

Blandt de sidste Stoikere fremtræder dette
Karaktermærke skarpest hos Epiktet. Hos Seneca og især hos Marcus
Aurelius er det i høj Grad mildnet, uden dog at være
opgivet. Marcus Aurelius’ Optegnelser vise, hvorledes de
uvilkaarlige Stemninger beherske Sindet; der er en Inderlighed
i Hengivelsen, som mangler hos de andre Forfattere. Den
Opgave at skildre Forskellen mellem den sidste Stoicisme og
Kristendommens Følelsesliv vilde derfor blive betydelig
vanskeligere, hvis vi blot vilde holde os til ham. For at faa
Forskellen tydeligt frem holder jeg mig da til Epiktet\footnote{Dog maa det vel mærkes, at Epiktets Ytringer i „Haandbogen“
fremtræde med en Skarphed og Strenghed, der for en stor Del skyldes
den kortfattede Form. Den udførligere Fremstilling, af hvilken
„Haandbogen“ er et Uddrag, viser os Epiktets Lære i en mildere,
menneskeligere Form. Denne Iagttagelse er gjort af \emph{Martha:} Les
moralistes sous l’empire romain. 2. ed. Paris 1866. p. 162 f.}.

Da Alt er ligegyldigt, hvad der ikke staar i vor egen
Magt, og da det Vigtigste af Alt, vore Opgaver og Pligter,
høre til det, der staar i vor egen Magt, saa er Stoikeren
ikke nogen stærkt interesseret Tilskuer ved eller Deltager
i hvad der foregaar i Verden uden om ham. Han ved jo,
at det altid er ham muligt at tage, hvad der kommer, paa
rette Maade; Fremtiden vækker derfor ikke hans Spænding
eller Nysgerrighed; \emph{hvad} der vil ske, ved han ganske vist
ikke, men \emph{af hvad Art} det vil være, ved han, nemlig, at
hvis det er Noget, der ikke staar i hans Magt, er det hverken
godt eller ondt, og hvis det er Noget, der er hans Magt
underlagt, saa beror det Væsentlige paa hans egen Villie, og
hvad den bør gøre, det ved han paa Forhaand. (Epiktets
Haandbog, Kap. 32). Hvad Epiktet udtaler om, hvorledes
man skal forholde sig ved et Skuespil, kan anvendes paa
Stoikerens Forhold til Livets store Skuespil. »Det er,« siger
han, »for det Meste ikke nødvendigt at gaa i Teatret. Men
% --- p. 177 ---
\opage{177}hvis det en Gang er betimeligt at gaa derhen, saa maa du
ikke være optaget af noget Andet end dig selv, det vil sige,
du skal blot ønske, at det sker, som virkeligt sker, og kun
at den maa sejre, som virkeligt sejrer; ti du vil da ingen
Hindring møde. Men hold dig ganske fra Raab, Hoveren og
stærk Bevægelse, og naar du er gaaet derfra, skal du ikke
tale meget om hvad der er gaaet for sig, uden for saa vidt
det kan tjene til din egen Bestyrkelse; ti af saadan Adfærd
vil det vise sig, at du havde følt Undren over, hvad du saa.«
(Haandbog, Kap. 33, § 10). Ogsaa hvad det virkelige Liv
angaar, vil Stoikeren kun, at det skal ske, som sker, og
saaledes som det sker, og derpaa beror hans indre Velfærd
(Haandbog, Kap. 8). Den grunder sig paa hans egen
Opfattelse. Og den samme Betragtning finder Anvendelse paa
Andre: ogsaa om dem ved han, at et virkeligt Onde ikke
kan vederfares dem. »Naar du ser, at En græder og er i
Sorg, saa vogt dig, at Forestillingen ikke river dig hen, saa
du mener, han er stedt i Ulykke; men gør straks en
Distinktion i dit stille Sind, og hav den Tanke paa rede Haand, at
ikke det, der er hændt ham, piner ham (en Anden vilde det
maaske ikke pine), men hans Opfattelse af det er det, der
piner ham. I Ord skal du ikke tøve med at vise ham
Deltagelse, og om saa skal være, sukke med ham; \emph{men vogt
over, at du ikke ogsaa kommer til at sukke i dit Indre!}«
(Haandbog, Kap. 16).

Hvad der her advares imod, er den fulde Hengivelse til
og Betagethed ved det, som virkeligt sker. For at Sindet
ikke skal blive sat i saa stærk Bevægelse, at det ikke kan
finde sin Ligevægt igen, holder man det borte fra det levende
Forhold til Virkeligheden og gør denne til et ligegyldigt
Spil. Den Vej, Stoikeren gaar for at bevare sin Frihed, er
at indskrænke sine aandelige Fornødenheder, sin
Følelsestrang; des sikrere er han paa at faa den tilfredsstillet. Han
faar Brøken i Livets Regnestykke til at vokse ved at
forminske Nævneren, ikke ved at forstørre Tælleren. Her
fremtræder det Negative og Passive i den stoiske Etik
tydeligt. En vis Frygt for Livet gør sig gældende. Jo mere
% --- p. 178 ---
\opage{178}dette Element kom til at raade, des mere maatte Stoicismens
Begrænsning vise sig.

Psykologisk set er Kristendommen i Modsætning hertil
karakteriseret ved, at den dristigt sætter Følelseslivet i den
stærkeste Bevægelse, lader det svinge ud i de mest ekstreme
Former. Den højeste Lyst og den dybeste Smerte har
Kristendommen Plads for. Dens Dogmer udtrykke den
Opfattelse, at kun den kender Livet, som har gennemlevet
Ekstremerne. Den Trang, som den menneskelige Følelse
har til at faa fuldstændig Udladning, til at skaffe sig Luft
i hensynsløs Hengivelse, fyldestgøres i den kristelige Kultus.
Jesus selv var bedrøvet indtil Døden, og han forudsiger sine
Disciple, at de skulle sørge og bedrøves. Og paa den anden
Side bebudes en Herlighed og en Glæde, som ingen Sans
kan fatte. Her er ingen Ængstelighed over for Følelsen, ingen
Frygt for at sukke i sit Inderste, saa lidt som for at lade
sig henrive af Undren. Derfor førte Kristendommen ikke
til at skille sig ud fra Tilværelsen, men opfordrede til at
leve og lide under dens Gang. Alle Tanker rettedes mod
ét stort Maal, der snart skulde naa sin Virkeliggørelse, og
paa Forholdet til hvilket den Enkeltes højeste Ve og Vel
beroede. Den Enkelte var altsaa virkeligt draget ind med
i de største Verdensbegivenheder, i den store historiske
Sammenhæng, hvis Grundtræk Aabenbaringen havde afsløret.
Som ligegyldig Tilskuer, der havde nok i sig selv, kunde
den Enkelte her ikke staa. I Forventningen om den store
og nære Afslutning af Verdenslivet var han revet ind i saa
voldsom en Bølgegang af Haab og Frygt, som den stoiske
Vise, hvem Tilværelsen Intet kunde bringe, og fra hvem
den Intet kunde tage, ikke kendte og i hvert Tilfælde gyste
tilbage for. Derved vandt den en Magt over Sindene, som
Stoicismen var udelukket fra, afset fra en lille Kreds af
energiske og inderlige Karakterer. Kristendommen lærte at
vinde sig selv ved at glemme sig selv, medens det var
Stoikerens Ulykke, at han altid havde saa travlt med at
regulere Sjælelivets uvilkaarlige Gang, at han ikke kunde
naa den sande Selvforglemmelse.

Men saa mangesidige ere Forholdene paa det aandelige
% --- p. 179 ---
\opage{179}Livs Omraade, at neppe have vi fremdraget og betonet en
afgørende Modsætning mellem Stoicismen og Kristendommen,
før vi nødes til at blive opmærksomme paa den Lighed, der
er imellem dem. \emph{Til Stoicismens skarpe dualistiske Forskel
mellem det Indre og det Ydre, en Forskel, der tilsidst
spærrede for levende og aktiv Deltagelse i Livet, svarer
Kristendommens endnu mere udprægede Dualisme mellem det
Hinsidige og »denne Verden«}. Hint Rige, som forventedes
snarligt, skulde komme ad overnaturlig Vej og ophæve det
naturlige Liv og alle dets Forhold. Deraf fulgte for den
konsekvente kristelige Etik et negativt og passivt Forhold til det
naturlige Liv og den historiske Udvikling. Kun ved
Omfortolkning er der senere naaet et Kompromis, ligesom der
heller ikke manglede analoge Kompromiser indenfor den
stoiske Skole. Naar Kristendommen skal betegne et blivende
Fremskridt i Forhold til Stoicismen, maa det bero paa, om
den levende, lidenskabelige Forventning, den fulde Hengivelse
til Følelseslivets Svingninger og Tilliden til, at man
derigennem ikke behøver at miste sig selv, men netop kan lære
Livet at kende derved, --- om alt dette kan bevares og
knyttes til det virkelige Menneskelivs Udviklingsgang
gennem Historien, selv om de dogmatiske Forestillinger, der
fra først af fremkaldte disse sjælelige Tilstande, falde bort.
Og om Muligheden heraf synes det aandelige Livs Udvikling
i den nyere Tid at vidne. Selv hvor det staar Kristendommen
i dennes dogmatiske Form fjernt, har det dog optaget fra
den psykologiske og etiske Elementer af største Betydning.

Det bekendte Ord, at Verdenshistorien er
Verdensdommen, bekræfter sig især i de store aandelige Strømningers
Skæbne. De fremtræde fra først af med bestemte
Særegenheder og Ufuldkommenheder, der skyldes de Forhold, under
hvilke de bleve til. For at kunne gaa over til at faa blivende
historisk Betydning, maa de først gennemgaa en Skærsild,
i hvilken det udskilles, der strider mod, hvad Kampen med
den \emph{virkelige} Tilværelse fordrer. Baade den kristelige og
den stoiske Etik have behøvet en saadan Skærsild, og der
er værdifulde Elementer i begge, som formaa at gennemgaa
denne Prøve.

\streg{}

% --- p. 180 ---
\opage{180}
\essay[FILOSOFI SOM KUNST]{FILOSOFI SOM KUNST.}

\section*{I.}

Naar En studerer Jura, kalder vi ham Jurist, --- naar
han studerer Medicin, Mediciner, og saa fremdeles. Ganske
paa samme Maade kan man ikke kalde Den, der studerer
Filosofi, Filosof. Det vilde man kun, hvis Ordet Filosofi
historisk betegnede en rent teoretisk og videnskabelig
Bestræbelse. Men Filosofiens Historie viser, at tænkende
Bearbejdelse af visse Problemer vel er en Hovedbestemmelse
ved Begrebet Filosofi, men at der dertil ogsaa knytter sig
en anden, mere praktisk Betydning: et \emph{Liv} i og for Tanker,
eller en Stræben efter at gennemføre Tanken i Livet, --- i
større eller mindre Kredse af Livet, i hvert Tilfælde i sit
eget Liv. Tages denne Side af Begrebet Filosofi, vil Ingen
uden Fordringsfuldhed kunne kalde sig selv Filosof. Man
kunde med Kant, der oftere kommer tilbage til dette Punkt,
skelne mellem Filosofi som Viden og Filosofi som Visdom.
Kant er omhyggelig for at præcisere, at der gives en
Betydning af Ordet, i hvilken Ingen tør kalde sig selv Filosof.
I den praktiske Betydning af Ordet lægger Filosofien Beslag
paa andre aandelige Evner end de rent intellektuelle.

Den empiriske og den kritiske Filosofi skelner skarpere
mellem de to Sider ved Begrebet Filosofi end den
spekulative, som paa romantisk Vis lader de forskellige Omraader
fortone sig sammen og vil sammenfatte dem alle i én Idé,
uden at anerkende Loven om Arbejdets Deling.
Gennemførelsen af den empiriske og kritiske Metode har ført til
Sondring mellem Fag, der tidligere ikke kom til at føre
% --- p. 181 ---
\opage{181}deres særegne Tilværelse. Den filosofiske Tænken er i vore
Dage ingen samlet, enkelt Strøm, men har delt sig i mange
Strømme: Psykologi, Logik, Etik o. s. v. Det er en Deling,
der har vist sig frugtbar; men den umiddelbare Begejstring,
der let kunde vækkes, da man i en enkelt Idé troede at
sammenfatte Tankens og Tilværelsens hele Indhold, al
Erkendelse og alle Værdier, som en Perle, det var værdt at
købe for alt sit Gods, --- en saadan koncentreret
Tankebevægelse kan vore Dages filosofiske Forsken ikke frembyde.
Den har opgivet at finde de Vises Sten eller den Nøgle, der
lukker op overalt.

Hvor Arbejdets Deling finder Indpas, medfører den en
Krise. Saaledes sker det i Organismen, hvor Nødvendigheden
af Koncentration under Kampen for Tilværelsen sætter en
Grænse for Differentiationen, --- i Samfundet, idet det sociale
Problem væsentligt hænger sammen med Arbejdsdelingen, ---
og saaledes ogsaa i Videnskaben: ti naar al Tænkning
fordeler sig over de særskilte, fra hverandre forskellige
Omraader, hvorledes bliver da den aandelige Koncentration mulig,
den Helhedsopfattelse og den samlede Virken, som Livet
fordrer, naar det ikke skal føres i Blinde? --- Efter min
Opfattelse stiller dette Problem sig ikke blot for den humane
eller naturalistiske Livsanskuelse, men ogsaa for den religiøse.
Det religiøse Problem i vore Dage hænger nøje sammen
med, at religiøs Tro nu er bleven en Specialitet \emph{ved Siden
af} Videnskab, Kunst og praktisk Stræben, medens den fra
først af stod som Indbegrebet af Alt, hvad der kunde bringe
aandelig og personlig Næring\footnote{Smlgn. min Etik p. 305---307. (2 Udg. p. 419 --421).}.

De betænkelige Virkninger af Spredtheden og
Specialiseringen synes at være i Færd med at fremkalde en Reaktion.
Den, som med Opmærksomhed følger, hvad der bevæger sig
i Tiden, vil kunne fornemme, at en Reaktion i mystisk
Retning ligger i Luften. Mystisk kan man kalde denne
Bevægelse: ti mystisk er ikke blot Dunkelhed, men
Koncentration af Aandslivet, Opgaaen i en Enhed. Man trænger til
fortættet Tankenæring; man er træt af Analysen og Speciali%
% --- p. 182 ---
\opage{182}seringen og vil Intuition, Sammenfatten af Erkendelsen til en
ledende Idé, som kan fylde Personligheden og ikke blot
tilfredsstille en enkelt Side af den. Man fordrer, at Videnskaben
skal blive Kunst, eller at Kunsten skal agtes for den højeste
Videnskab.

Der er mange Træk, som Slutningen af det 19.
Aarhundrede har fælles med Slutningen af det 18. Aarhundrede.
Ogsaa da rørte sig Trangen til Koncentration og Mystik efter
den megen Analyse, og Følgen blev Reaktion over hele Linien.
Karikaturen af denne Trang ytrede sig da, ligesom nu, i
Tendentsen til Spiritisme, Okkultisme og andre
hemmelighedsfulde Retninger. Saadanne aandelige Strømningers Opstaaen
har i og for sig noget Dunkelt ved sig. I Aaret 1790 blev
Kant spurgt af en Ven om Aarsagen til det overhaandtagende
Sværmeri. Kant svarede, at det var ligesaa vanskeligt et
Spørgsmaal at besvare for Psykologerne, som det for Lægerne
var at angive Aarsagerne til den herskende Influenza. Han
mener, at man ikke skal tage Sagen for alvorligt og ikke
bruge heroiske Kure overfor den, og synes at tænke sig et
aandeligt Middel derimod i Analogi med koldt Vand, som
Lægerne brugte for Influenza\footnote{Kants Brev findes hos \emph{Borovski}: Darstellung des Lebens und
Characters Immanuel Kants. Kønigsberg 1804, p. 227.}. Karakteristisk for den
begrænsede Erfaring, man paa Kants Tid endnu havde om
aandelige Strømninger, deres Natur og Love, er det, at Kant
ikke tænkte paa, om der ikke i hele det aandelige Livs
Karakter, som det havde formet sig i den sidste Halvdel af
Aarhundredet, kunde være Noget, der naturligt maatte
hidføre en Reaktion. Vi have i Løbet af det Aarhundrede, der
siden er gaaet, havt Lejlighed til at se og studere
Bevægelserne paa det aandelige Omraade, og ere fortrolige med den
Tanke, at en stærk Strømning altid paa en eller anden Maade
er motiveret ved den foregaaende Periodes Resultater eller
Mangler. Vi have set saa mange Bølger afløse hverandre, at
vi mene at kunne slutte fra de foregaaende, hvorledes omtrent
de følgende vil være beskafne. Og vi kunne da ruste os paa
Forhaand til at tage imod, hvad der skal komme. Det
% --- p. 183 ---
\opage{183}er jo ikke sagt, at den Koldtvandskur, Kant anbefaler, er
tilstrækkelig.

Hos forskellige Forfattere fra de seneste Aar faar
Trangen til Intuition, til Tanker, i hvilke der kan leves
personsonligt % PRINTED AS IS (for personligt; person-|sonligt at the line end)
og koncentreret, til afsluttende Livsanskuelse, Luft
paa ejendommelig, til Dels meget subjektiv og barok Maade.
Af tyske Forfattere ere her Lagarde, Nietzsche og den anonyme
Forfatter til »Rembrandt als Erzieher« at nævne. I
Frankrig haves i Renan en aandfuld Repræsentant for denne
Retning. I et nyligt udkommet talentfuldt svensk Skrift af Hans
Larsson (»Intuition«) viser den samme Tendens sig. Hvad
der end kan indvendes mod disse Forfattere, saa udtrykke
de hver paa sin Maade en aandelig Fornødenhed, der ikke
lader sig tilbagetrænge. Vi kunne ikke leve paa Abstraktioner
og Specialiteter, og det Spørgsmaal kan ikke skydes tilbage,
hvilken Betydning Tænkning og Forskning da har for vor
personlige Livsførelse.

Uagtet jeg tror, og ogsaa lejlighedsvis allerede har
berørt, at Problemet lige saa vel opstaar paa det positivt religiøse
Standpunkt i vore Dage, som paa det filosofiske eller humane,
vil jeg dog her især holde mig til dette sidste Standpunkt.
Spørgsmaalet er, hvorledes Filosofien kan blive Kunst, og i
hvilken Betydning den kan blive det; \emph{at} den maa blive det, er
givet. Og som jeg begyndte med at sige, historisk ligger det i
Filosofiens Begreb ikke blot at være Viden, men ogsaa Kunst
eller Visdom. Spørgsmaalet er, hvorledes dette under vore
Forhold kan ske.

De fleste af de Forfattere, jeg før nævnte, synes ikke
at have set Problemets hele Vanskelighed. De overhugge
Knuden i Stedet for at løse den; de springe over i Mystik
og Intuition, gøre Vold paa Tænkningen, glemme, afbryde
eller afslutte den paa vilkaarlig Maade, ja se maaske endog
med romantisk Foragt ned paa den ædruelige Tænknings og
Specialundersøgelsernes Mænd. Hvad der interesserer mig
hos disse Forfattere, er derfor mere den Trang, der kommer
frem hos dem, end den Maade, hvorpaa de selv søge denne
Trang tilfredsstillet.

% --- p. 184 ---
\opage{184}Jeg skal forsøge i korte Træk at vise, af hvad Grund
Filosofien maa blive Kunst, og paa hvilken Maade denne
Kunst er at øve.

\section*{II.}

Saaledes som vor Erkendelse nu én Gang er stillet, kan
enhver afsluttende Verdensanskuelse kun være en Hypotese.
Enhver Helhedsopfattelse, enhver Fortolkning af Tilværelsen,
vi forsøge, maa godtgøre sin Gyldighed ved at sammenholdes
med nye, videre gaaende Erfaringer. Men da saadanne
indfinde sig, saa længe Livet varer i os og udenfor os, kan
Erfaringsbekræftelsen aldrig afsluttes saaledes, som det er muligt
ved mere begrænsede Opgaver. Det er altid paa
fragmentarisk Grundlag, vi bygger en Fortolkning af Tilværelsen; vi ere
stillede omtrent som Filologen, der skal tyde et brudstykkeagtigt
Manuskript, eller som Historikeren, der af spredte og
løsrevne Beretninger skal konstruere en Begivenhed eller en
Handling og dens Motiver. Den filosofiske Forsken ender i
Gisningens Kunst, ligesom Tekstkritiken og
Historieforskningen. Den Karakteristik af Tilværelsen, enhver
Verdensanskuelse tilstræber, kan da ikke komme ud over det
Hypotetiske. Grænseresultaterne, de sidste Konsekvenser, en Tænker
drager, ville stedse beholde en problematisk Karakter.

Der hører nu ikke blot Kunst til at udarbejde en
saadan sidste Hypotese, men der hører ogsaa en ejendommelig
Kunst til at arbejde med den i Praksis. Der udfordres en
Elasticitet, der stadigt fra den fuldbyrdede Konstruktion gaar
tilbage til Analysen, og formaar at gøre denne Svingning
frem og tilbage uden at miste Energien. Kunsten er at
forene Troens Begejstring med Drøftelsens Utrættelighed.
Tankearbejdets Kraft maa ikke svækkes ved Bevidstheden
om, at Resultatet stedse er problematisk. Naar den givne
Virkelighed er i Modstrid med de Ideer, vi med al vor Flid
har dannet os, er det saa Ideernes Skyld, eller er det den
begrænsede Erfarings Skyld? Her er Kunsten den, stedse
at være beredt til ny Prøvelse. Det er lettere at slaa en
% --- p. 185 ---
\opage{185}Streg over Erfaringsbekræftelsen og én Gang for alle hensynke i
Profeteren og romantisk Intuition; men det fører til Tankens
Død --- den store Fare, som al aandelig Reaktion truer med,
og som det fremfor Alt gælder at undgaa. Men den
undgaas kun ved Indøvelse i den Kunst at leve paa Ideer.
Det er mere end blot Viden; thi det stiller Krav til Villiens
Energi og Følelsens Livlighed. Det er Noget, der ikke kan
gøres af én Gang for alle. --- Om det, man kun behøver
at vide for at kunne øve det, siger en gammel tysk
Talemaade: »Das ist keine Kunst, das ist nur eine
Wissenschaft.« Overgangen fra Filosofien som Viden til Filosofien
som Kunst kræver Opbud af en personlig Kraft og
Virtuositet, som den rene Viden ikke fordrer. Det gælder om at
bevare Sammenhængen mellem vor Tænken og vor
praktiske Higen og Færden i Livets Mangfoldighed, at gøre
Overgangen saaledes, at der ingen indre Splid kommer i vort
Væsen.

Korteligt skal jeg henvise til, at en aldeles lignende
Opgave stilles den Troende. Selv om hans Tros Indhold for
ham er givet én Gang for alle og staar uforanderligt fast,
som kundgjort af en absolut Autoritet, saa forandrer Livet
sig dog stadigt, og Opgaven bliver da den at bevare
Sammenhængen mellem sin Tro og sin praktiske Deltagelse
i Livets Færd. Som hist Tænkningens Resultater, saaledes
skulle her Troens Postulater fastholdes uden at skifte Karakter
under de skiftende Forhold. Det vilde være en interessant
Opgave at undersøge, hvorledes denne Opgave er løst fra
Kirkens Side til de forskellige Tider. Kirken selv er mere
tilfreds med den Maade, paa hvilken den har løst denne
Opgave, end dens Kritikere fra højst forskellige
Standpunkter (saa forskellige som Feuerbachs og Kierkegaards) ere
det. Men her vilde jeg blot antyde, at ogsaa paa Troens
Standpunkt er et kunstnerisk Supplement nødvendigt, skønt
det naturligvis faar en anden Karakter end den filosofiske
Kunst, som ikke har Autoriteten til Forudsætning.

Et Liv i Ideer behøver ikke at forudsætte, at disse
Ideers Gyldighed staar dogmatisk fast. Det er et smukt
% --- p. 186 ---
\opage{186}Ord af Lagarde: »De sande Idealister kunne vente«\footnote{Deutsche Schriften. Gøttingen 1886 p. 481.}. Det
er kun Utaalmodigheden og Ladheden, der volder Dogmatisme.
Det gaar Mange endnu med Ideerne, som det i gamle
Dage gik Alle med Stjernerne: man var bange for, at de
skulde falde ned, naar de ikke var forsvarligt indlagte i faste
Sfærer.

Der kræves ikke blot aandelig Frihed og Elasticitet for
at leve i Ideer, uden at gøre dem til ubevægelige Dogmer;
men der hører dertil ogsaa et Sind, som ikke forarges over,
at det Store saa ofte viser sig for Forskningen at udspringe
af det Ringe og Uanseelige. Hvad der i Ideen, i
Totalanskuelsen staar saa frit og højt, har langsomt arbejdet sig
frem og skyldes Bidrag fra mange Sider, som den store Aa
opstaar af de mange smaa Bække. Analysen, den specielle
Undersøgelse fører os ligesom bag Kulisserne, viser os
Maskineriet, og der opstaar da let en Blaserethed, som udbryder:
»Er det det Hele? Er det det, der ligger bag ved?« I Stedet
for at glæde sig over den Tilværelsens Hemmelighed, der
knytter det Store og det Smaa, det Uendelige og det
Endelige sammen paa uadskillelig Maade og i stadigt Vekselspil,
finder man i Analysen Bevis for, at Livets Værdi er Illusion,
eller hvis man vil fastholde Værdien, vender man sig fra
Analyse og Specialforsken med Haan. Ogsaa i denne
Henseende behøves der kunstnerisk Virtuositet for at fastholde de
forskellige Synspunkter, der kunne synes at stride mod hinanden,
og for ikke at tabe Troen paa det Evangelium, Historien
atter og atter forkynder, at selv de største Livsmagter en
Gang maatte svøbes og lægges i en Krybbe, og at de hjælpende
og forløsende Kræfter ofte komme fra en Kant, fra hvilken
man mindst ventede dem.

Det rent videnskabelige Arbejde, den rene Tænken lægger
kun Beslag paa en enkelt Side af Personligheden. Men ved
den Opgave at arbejde med Ideerne i Livets skiftende
Mangfoldighed, i de smaa Betingelsers Virvar uden at henfalde
hverken til Dogmatisme eller Blaserethed lægges der Beslag
paa hele Personligheden. Mangen, som med Begejstring
% --- p. 187 ---
\opage{187}kastede sig over Gransken af de store Problemer for at
gennemtrænge dem med Tanken, har erfaret, at der kom et
Punkt, hvori Situationen ændredes, og hvor han fra den
mere upersonlige Søgen efter en rent objektiv Løsning af
Gaaderne stødtes tilbage i sig selv og fik nok at gøre med
at gennemarbejde de foreløbigt vundne Tanker i sit
personlige Liv, at uddanne hos sig den personlige Kunst, som hører
til for at leve i Ideerne. En begejstret Tænker fra
Renaissancetiden, Giordano Bruno, har udtrykt denne Tænkningens
Overgang til personligt Liv ved Hjælp af Myten om Aktæon,
Jægersmanden, der efter at have faaet Øje paa den
guddommelige Skønhed blev Bytte for sine egne Hunde: Hundene
ere Tankerne, der rettes mod de store Ideer; naar Øjet er
opladt for Idealet, vende Tankerne sig indad formedelst den
stærke Bevægelse i Sindet, søge ikke længere ud efter.
Problemet bliver subjektivt, personligt, ikke længere rent
objektivt. Og Bruno fremhæver, at saaledes gaar det, hvor
det ikke er specielle og endelige Genstande, Tanken rettes
imod, men det Guddommelige, Universelle\footnote{Bruno: De gl’ heroici furori. Op. ital. ed. Lagarde. Gøttingen
1888. II. p. 651 f., 723 f.}. At en Forsken,
især hvor den rettes mod de højeste Opgaver, omsættes til
en evig Stræben, er en Tanke, ved hvilken Bruno staar som
Lessings, Kants og Goethes Forgænger. Kunde Maalet naaes
én Gang for alle, var der ingen Plads her for Kunst.

Videnskaben finder selv sin Regning ved saaledes at
give Plads for en Kunst, der dog bevarer Sammenhængen
med den. Ti denne Overgang til personlig Kunst medfører,
at vi faar et nyt Erfaringsomraade. Vor Synskreds beriges.
Vi blive selv »Instanser«, selv Genstande for Eksperimenter.
Der gøres Erfaringer om Forholdet mellem Tankerne og de
personlige Kræfter, om Hines Næringsværdi og Frugtbarhed
og disses Bøjelighed og Udholdenhed. Det er menneskelige
Erfaringer i Ordets sande Betydning. Og formaa vi paa en
eller anden Maade at forme og udtrykke dem, vil de kunne
komme Andre til Gode. Der vil paa Grundlag af
saadanne Erfaringer kunne danne sig en sammenlignende Livs%
% --- p. 188 ---
\opage{188}filosofi\footnote{Smlgn. min Bog: Søren Kierkegaard som Filosof. Kbhvn. 1892. p.
52 f., 82 ff.}, der vil blive en Analogi til den sammenlignende
Religionsvidenskab. Hidtil har man kun haft nogenlunde
Respekt for Livsanskuelser, der vare fælles for Folkeslag, for
store Menneskegrupper. Jo mere Dogmatismen trækker sig
tilbage (hvad den bebudede Mystik forhaabentligt ikke vil
hindre for længe), des mere ville de individuelle Erfaringer og
den personlige Livskunst komme til deres Ret. Jeg ser heri
et Lyspunkt for Fremtiden.

Personlighedernes Forskellighed vil uundgaaeligt faa stor
Indflydelse, naar Filosofien skal gaa over til Kunst. Hvad
jeg hævder, er kun, at Kunsten fra en væsentlig Side netop
bestaar i at bevare Sammenhængen mellem Forskning og
Personlighed, eller rettere mellem den Del af Personlighedens
Væsen, som den rene Tænken lægger Beslag paa, og den,
der virker, naar Tankerne skulle fastholdes og gennemføres i
Livets Virkelighed. Det kan ikke være Andet, end at Vejene
paa dette Punkt ville skilles paa Grund af de personlige
Forskelligheder, og kun gennem en sammenlignende
Livsfilosofi kan der igen blive Mulighed for gensidig Forstaaelse.

I Livskunsten er Enhver først og fremmest henvist til
sig selv. Han har en Opgave, han skal løse paa sin Vis, og
ikke kan løse paa nogen anden Vis, naar han ikke vil ende
i Affektation eller Hykleri. --- Da Buddha blev født, ønskede
10,000 Kvinder at blive hans Amme. Og da Buddha ikke
vilde bedrøve de 9,999, mangfoldiggjorde han sig selv
saaledes, at hver Kvinde fik sit Vidunderbarn at opfostre\footnote{Billedet henter jeg fra \emph{Renan}: Feuilles détachées. 4. ed. Paris,
1892, p. 162 f., men jeg anvender det noget anderledes end han.}.
Saaledes har vi Alle hver for sig et Gudebarn at fostre. Og
saa mange Børn der end er, ægte ere de alle. Det gælder
for Enhver at udvikle \emph{sine} Kræfter, bringe sine personlige
Anlæg og Drifter sammen med de Tanker, Slægten til den
givne Tid raader over, eller med saa mange af dem, som han
kan tilegne sig, og derigennem uddanne sin Livsfilosofi, øve
den menneskelige Livskunst paa sin Vis.

% --- p. 189 ---
\opage{189}
\section*{III.}

Hidtil have vi kun betragtet den filosofiske Kunst i dens
Forhold til Ideerne, til Tænkningens sidste Hypoteser. Endnu
klarere viser denne Kunsts \emph{Nødvendighed} sig, hvor det
gælder praktisk Livsførelse. Det gælder her ikke blot om
Kontemplation og Teori, men om Handlen. Vi have ikke mere
at gøre med den teoretiske Filosofi, med Forsøget paa at
udvikle en Verdensanskuelse, men med Filosofien som Etik.
Ogsaa her viser det sig, at der er en kunstnerisk Afslutning
og Gennemførelse, som Teorien ikke direkte kan bestemme.
Den filosofiske Etik kan kun give de almindelige Principer:
den praktiske Anvendelse i det enkelte bestemte Tilfælde,
under de aldeles konkrete Forhold, kan ikke direkte udledes
af dem. Og hvad der ogsaa her vil vise sig at være den
store Kunst, er det at forene Modsætninger i sig. I det
personlige Forhold til Erkendelsens Resultater gjaldt det at
forene Troens Frejdighed med stadigt fortsat Drøftelse, Blikket
for det Store med Sansen for de smaa Betingelser. I den
praktiske Livsførelse gælder det om Foreningen af
Modsætninger som Erindring og Haab, Opgaaen i Øjeblikket og Stræben
efter omfattende Formaal, Hævdelse af den egne Individualitet
og Hengivelse til Noget, som ligger ud over denne,
uvilkaarlig Udfoldelse af Livet og vilkaarlig Indgriben. Mange flere
Modsætningsforhold vilde der her ved nærmere Undersøgelse
være at nævne, som den personlige Kunst skal forbinde til
Enhed paa en Maade, som ingen almindelig Formel kan lære,
da Betingelserne baade i Grad og Kvalitet er forskellige for
hvert Individs Vedkommende, og da Løsningen tilmed ikke
kan naas én Gang for alle, saaledes som en Sætning staar
fast, naar man blot én Gang har faaet den bevist\footnote{Smlgn. min Bog: Søren Kierkegaard som Filosof. p. 44---48.}. --- Kun
et Par af disse Modsætningsforhold skal jeg her dvæle lidt
ved, for at oplyse dem ved Eksempler.

En Fare, der truer, naar Livsførelsen ikke længere skal
bestemmes ved Overlevering og Autoritet, men ved fri
Erkendelse, er, at det Uvilkaarlige og Umiddelbare bliver kuet
% --- p. 190 ---
\opage{190}og fortrængt. Naar der er Tale om Filosofi som Kunst, vil
man maaske endog straks opfatte Sagen som om det nu
gjaldt at regulere alle Livsytringer efter bevidste Principer.
Naar Tankens Ret til at gribe bestemmende ind i Livet
anerkendes, vil der let opstaa Mistro eller Ringeagt over for
det Naturlige, Tilvante. Overleverede og Uvilkaarlige. Man
vil i Emancipationen fra dette maaske endog i og for sig
finde et Fortrin. Vilkaarligheden prises, blot fordi den giver
noget Nyt, bryder med det Vante. Bevidst Reflekteren,
Ræsonneren og Reguleren kommer derfor let til at spille for
stor en Rolle, og det gaar ud over Friskheden og den
umiddelbare Kraft. Reaktionen efter en rationalistisk
Emancipationsperiode finder maaske sin bedste Støtte i den tilsidesatte
Trang til at give sig hen i det Uvilkaarlige i Stedet for at
lade Alt bestemmes ved bevidste Forestillinger.
Kristendommens første Virkninger i den europæiske Verden beroede
for en ikke ringe Del paa, at den kom denne Trang i Møde.
Den stoiske Lære, som var den alvorligste Livsfilosofi i Tiden,
led netop af den Ensidighed, som bestaar i Mistillid til
Sindets uvilkaarlige Bevægelser. Stoikeren vilde fremfor Alt
bevare sin rationale Sindsligevægt, vilde ikke omtumles af
Stemningerne, ikke tvinges i Affekt. Han ofrede Fylden og
Dybden for at bevare Selvherredømmet. Ikke undres, ikke
sukke i sit Inderste, ikke juble overvættes, det var hans
Maksime. Kristendommen derimod gav ved sine Grundtanker
og sin Kultus Lejlighed til at gennemleve de stærkeste
Sindsbevægelser, til at lære de mest modsatte Følelser at kende
og derved ligesom udmaale, hvad Livet indeholder af Lys og
af Mørke\footnote{Se udførligere herom Afhandlingen „Hedenske Sandhedssøgere“.}. Det var det Uvilkaarliges berettigede Reaktion
mod det Vilkaarlige. Den stoiske Lære om Tankens styrende
Magt var alligevel ikke urigtig. Men Forholdet er mere
sammensat, end Stoikeren saa. Tanken er Roret, men Vind
og Vove maa drive Skibet frem. Ogsaa her er Kunsten at
bevæge sig gennem Modsætningerne, saaledes at begge
kommer til deres Ret; paa Vekselspillet af det Vilkaarlige og
det Uvilkaarlige beror Livets Klarhed, Sundhed og Dybde.
% --- p. 191 ---
\opage{191}I nyere Tid betegner Rousseaus Optræden ogsaa en
Reaktion af det Uvilkaarlige mod den bevidste Reguleren, som
den moderne Rationalismes første Periode havde sat i System.
Han opdagede en ny Verden, Hjertets, den umiddelbare
Følelses, de uvilkaarlige Instinkters Verden i Modsætning til
Forstandskulturen og den ydre, mekaniske Dannelse. At
han tog Ordet for Naturen imod Kulturen, betød, at han
igen vilde bringe det Umiddelbare, Ubevidste, Oprindelige og
Enfolde til Ære over for det bevidst Formede og Artikulerede.
Det var som en ny Kilde, der sprang ud i Ørkenen; en ny
Værdsættelse af Livet indlededes, og den Bølge, han satte i
Bevægelse i det menneskelige Aandsliv, fortsattes gennem
den romantiske Bevægelse, hvis Forløber han var, til ind i
vore Dage. Paa karakteristisk Maade fremtræder den
Rousseauske Livsfilosofi i hans Fordring om, at Opdragelsen skal
være \emph{negativ}, det vil sige, væsentligt bestaa i Bortryddelsen
af Hindringer og Udvækster, saa at Livet faar Lov at udfolde
sig uvilkaarligt. Han hævder, at kun ved dette varsomme
og indirekte Forhold kan man i det Hele lære Barnenaturen
at kende og undgaa at forkvakle den ved at gribe ind efter
forudfattede Ideer. Den vilkaarlige Indgriben forudsætter jo
altid saadanne forudfattede Ideer; og Kunsten, den vanskelige
Kunst er netop den, at gøre den fornødne Brug af de Ideer,
der grunde paa Fortidens bevidste Erfaringer, uden at krænke
eller forkvakle det, der ubevidst, i det Skjulte er spiret frem
og begynder at udfolde sig. Deri bestaar al Opdragelses,
ogsaa Selvopdragelsens Kunst. Den finder Anvendelse paa det
Gudebarn, vi hver for sig har at fostre.

Et andet Forhold, hvor den almindelige Formel, som
kan gives i den teoretiske Etik, heller ikke slaar til, saa at
den personlige Kunst maa fremarbejde Modsætningernes
Forbindelse, er det mellem Hævdelse af den egne Personlighed
og Hengivelse til Andre. Fra mange Sider lyder nu
Slagordet om Individualitetens hensynsløse Hævdelse som det
ene Fornødne. Humaniteten, siger Lagarde endog, maa vige
for Individualiteten. Men med ikke mindre Styrke rører sig
i vor Tids Bevidsthed den socialistiske Tanke. For den
Enkelte at finde det harmoniske, det vil sige det efter hans
% --- p. 192 ---
\opage{192}Natur\footnote{Om den individuelle Forholdsmæssighed af alle etiske Fordringer,
se min Afhandling „Om Forholdsloven i Etiken.“ (Etiske
Undersøgelser. Kbhvn. 1891, p. 53. 56.)} harmoniske Forhold mellem de to Modsætninger er
en stor Kunst. I stærke Ord at proklamere det ene eller
det andet Led i Modsætningsforholdet som det Eneste, det
gælder, er ingen Kunst (uden i den Forstand, at det kan
ske med Fremstillingskunst; men det er ikke den Slags
Kunst, Talen her er om). Lige saa lidt som det er Kunst,
er det Videnskab. Det er Agitation. --- I Ibsens »Bygmester
Solness« findes anskuelige Bidrag til at belyse Problemet.
Jeg skal vel vogte mig for at give nogen Fortolkning af,
hvad Stykket gaar ud paa. Det er jo ikke sagt, at det
gaar ud paa Noget. Jeg benytter kun nogle Træk af
Hovedpersonernes Karakter. Der fremtræder hos Solness en
Tendens til absolut Selvhævdelse, til en »Rovfuglemoral«, der
gør sin egen Higen til eneste Lov. (Ved Udtrykket
Rovfuglemoral maa det dog aldrig glemmes, at Rovfuglene have
deres Reder, og at den ene Ravn ikke hugger Øjnene ud
paa den anden.) Den Selvhævdelse, der først bliver mulig
gennem Hengivelsen og Opofrelsen, og den Udvidelse og
Uddybelse af alle Sindets Kræfter, som kun vindes ved
Opgaaen i Idealer, der omfatte Mere end den egne
Individualitet, --- dem kender Solness ikke. Men det viser sig dog, at
naar han kuer og bedrager Andre for at naa sit Maal, saa
er det ikke Tegn paa Magt, men netop paa Afmagt. Saa
snart Ungdomsmodet og Ungdomskraften (personificeret i
Hilde --- hvis det ikke er Forsyndelse at kalde denne skønne
Skikkelse en Personifikation) kommer frem hos ham igen,
yder han straks Retfærdighed --- nu dog for sent. Og saa
længe Bygmesterens Planer omfatter det menneskelige Livs
Interesser, lykkes de for ham; da han utilfreds dermed vil
op i de ensomme Regioner for at bygge Luftslotte, styrter
han til Jorden. Om det er Ibsens Mening, at det, der hindrer
Solness i den hensynsløse Selvhævdelse, hindrer ham \emph{med
Rette}, er umuligt at sige, og det vedkommer os heller ikke.
Men med den slaaende Kraft, som er en saa stor Psykolog
% --- p. 193 ---
\opage{193}værdig, har han vist, at Selvhævdelsen lettest gennemføres
over for dem, man \emph{ikke} kender. Ti de, vi \emph{kende} og komme
i et inderligt Forhold til, de høre til Rovfuglereden, som vi
selv, og dermed betegnes Rovfuglemoralens foreløbige Grænse
--- ikke ved en »moralsk« Lov, men ved en Naturlov.

Fru Solness er konstrueret som Modstykke. Hun gaar
helt op i Erindring og Vedhængen ved det Tidligere, i
Hensyn og Pligter. For lutter Pligter forsømmer hun sin Pligt.
Hun vil die sine Børn, fordi det er hendes Pligt, skønt
hendes Mælk i hendes sygelige Tilstand er skadelig for
Børnene. Den Venlighed, hun viser sine Omgivelser, tager
hun selv Hjertet af Livet paa ved udtrykkeligt at øve den
som Pligt. Og i det væsentligste Forhold, i hvilket hun
staar, det til hendes Mand, griber hun fejl for lutter Pligter.
Maaske har Solness Ret i, at hendes Liv er brudt, fordi
det, der var hendes egentlige Kald, Børneopfostringen, ikke
falder i hendes Lod. Men de mange Pligter gøre det i hvert
Tilfælde ikke bedre. Bedre vilde det være, om \emph{hun} gik op
i Selvhævdelsens store Pligt, fuldt ud udviklede sig til en
Personlighed, hvis Sving og hvis Varme kunde forplante sig
til hendes Nærmeste. I Stedet for at give sine smaa
udstykkede Tjenester og Hensyn --- i Grunden kun Udtryk for
en klynkende Selvhævdelse --- kunde hun da give det Højeste,
en kraftigt stræbende Personlighed. Som det nu er blevet,
lever hun ved hans Side som et aandeligt Lig --- og det
truer derfor med, at hun skal kastes ud af Reden.

Jeg tager af den store Digtning kun, hvad jeg i denne
Sammenhæng kan benytte. Det er Digterens Fortrin, at han
kan belyse de personlige Problemer ved at give os en
Fremstilling, der medtager alle individuelle Nuancer og alle specielle
Omstændigheder. (Ibsens Storhed viser sig især i Henseende
til de individuelle Nuancer, ikke saa meget i det dramatiske
Net, han vil trække sammen over sine Personer; ti det har
ofte Huller). Den filosoferende Betragtning, der maa holde
sig til almindelige Principer og typiske Eksempler, staar i en
stor Taknemmelighedsgæld til de store Digtere, hvis
Menneskekundskab former sig i individuelle Skikkelser. Men
paa den anden Side tror jeg, at man overdriver Betydningen
% --- p. 194 ---
\opage{194}af saadanne Fremstillinger, naar man mener, at Digteren
ligefrem kan behandle eller endog løse Problemer. Derfra er
han allerede afskaaren netop ved den rent individuelle Form,
i hvilken han giver sine Fremstillinger. Det Spørgsmaal,
hvor vidt det Fremstillede kan finde Anvendelse paa andre
Tilfælde, staar ubesvaret hen; det var maaske en Undtagelse,
med Hensyn til Karakterer eller Situationer, eller der kan
opstilles modstaaende Skildringer. Problemet er ikke
udtømt med den enkelte digteriske Fremstilling. Debatten
begynder først ret efter denne, idet Spørgsmaalet bliver om det
Fremstilledes Forhold til andre Muligheder, som Livet kan
frembyde. At holde sig til den individuelle Fremstilling og
anse Sagen for udtømt ved den vilde være en Dogmatisme,
som er modsat den, der ofte kan lægges den videnskabelige
Behandling af Livsspørgsmaal til Last. Livet gaar lige saa
lidt op i en enkelt Konstellation som i et abstrakt Princip,
men maa belyses fra begge Sider. Kunsten kan lige saa lidt
gøre Videnskaben overflødig, som Videnskaben kan undvære
Kunsten. Hvad det gælder om at holde fast, --- maaske
særligt i de Tider, som nu komme (hvis de Tider komme,
man bebuder) --- er, at selv, hvor Kunsten afløser Videnskaben
og supplerer den, kan Baandet dog ikke overhugges mellem
dem. Der er et Vekselforhold mellem de to Hovedformer
for Aandsvirksomhed som mellem Indaanding og Udaanding;
den ene Proces fortsætter den anden og beror paa den anden,
trods den Modsætning, der er mellem dem. Deres Forening
er nødvendig for at gøre Livet frit og sundt. Denne Forening
er vanskelig at naa; men Sokrates, Stormesteren i Filosofi
som Kunst, har jo ogsaa sagt: »Det Skønne er vanskeligt.«

\streg{}

% --- p. 195 ---
\opage{195}
\essay[FORHOLDET MELLEM TRO OG VIDEN I DETS HISTORISKE UDVIKLING]{FORHOLDET MELLEM TRO OG VIDEN I DETS HISTORISKE UDVIKLING.}

\begin{center}Et indledende Foredrag.\end{center}

\streg{}

Problemet om Tro og Viden opstaar med Nødvendighed
paa et vist Trin af Kulturudviklingen. Vi træffe det
saaledes inden for den indiske, den græske og den
muhammedanske Verden, saa snart den aandelige, særligt den
intellektuelle Udvikling har naaet et vist Punkt. Her ville vi dog
indskrænke os til at fremhæve de Hovedformer, under hvilke
det er fremtraadt inden for den evropæiske, ved
Kristendommen paavirkede Aandsudvikling.

\section*{I.}

Af Kristendommens Kildeskrifter kan ingen Belæring
hentes angaaende dette Spørgsmaal. Der eksisterer paa det
Standpunkt, ud fra hvilket de ere forfattede, ingen fra Troen
forskellig Viden. Blikket er rettet mod én eneste Genstand,
et eneste Maal, der lægger Beslag paa al Stræben og al
aandelig Kraft. Bestræbelsen for at erkende Sandheden er
her umiddelbart Et med Trangen til at vinde Sjælens Frelse
og gaar ikke sin egen Vej med sin egen Metode. Der er
ingen Forskel mellem Teori og Praksis, da Teori og
teoretisk Interesse aldeles ikke eksisterer. Naar Guddommen
selv har aabenbaaret sig synligt for Alle, --- naar det ved
Troen er muligt at naa til Forstaaelse af det Mysterium,
som indeslutter alle Visdommens og Kundskabens Skatte,
% --- p. 196 ---
\opage{196}--- saa mangler ethvert Motiv til at søge Kundskab ad andre
Veje.

Det var naturligt, at Troen fra først af maatte være
Et og Alt, saa der ikke blev Plads til andre aandelige
Fornødenheder. Den første Slægt stod endnu saa nær ved den
mægtige aandelige Bevægelse, der førte til den kristelige
Menigheds Grundlæggelse, at den fuldstændigt levede og
aandede i de store Minder og Syner. Dertil kom endnu, at
denne første Slægt levede i Troen paa Kristi snarlige
Genkomst. Da skulde Alt afsluttes, Verdensdommen fældes, og
alle menneskelige og jordiske Forhold give Plads for en
højere Tingenes Orden. Hvem havde med en saadan
Katastrofe for Øje Tid og Ro i Sindet til filosofisk Tænkning
eller videnskabelig Forsken? Videnskaben kunde lige saa
lidt som nogen anden menneskelig Bestræbelse faa nogen
Betydning. Der skulde jo slet ikke finde nogen historisk
Udvikling Sted; Maalet var naat, og Enden stod for Døren.
Man forstaar ikke det ny Testamente og dets Forhold til
Kulturbestræbelsen i dens forskellige Former, naar man ikke
fastholder dette. Det gik her med den videnskabelige
Stræben som med Stræben efter personlig Frihed og med
Indgaaelse af Ægteskab (se 1. Korinterbrev 7. Kap.).
Hvorledes kunde alt Sligt faa Betydning, eller endog blot blive
Genstand for Drift, naar alle Tanker og Kræfter maatte
anspændes for at være forberedt til den store, den absolute
Afgørelse!

\section*{II.}

Tiden gik, Afslutningen kom ikke. Kirken trøstede sig
med, at for Gud er tusind Aar som én Dag, og én Dag som
tusind Aar (2. Peters Brev 3. Kap.), og begyndte snart at
falde til Ro i Verden, at indrette sig efter Forholdene og at
tage Stilling over for de forskellige Sider af de kulturhistoriske
Bestræbelser. Over for Videnskaben blev dens Stilling den,
at den ikke længer i Troen fandt al fornøden Erkendelse,
men søgte at vise, at Troens Lærdomme satte Kronen paa al
menneskelig Erkendelse. Kirken tilegnede sig den Viden%
% --- p. 197 ---
\opage{197}skab, Grækerne havde udviklet: den græske Astronomi,
Medicin og Filosofi, og anbragte sin Teologi som Spiret
paa Videnskabens Bygning, --- paa lignende Maade som den
omdannede de græske Templer til Kirker og f. Eks. indviede
Partenon, Athenes Tempel, til den »hellige Visdom«.

Allerede i det saakaldte Johannesevangelium finde vi
efter Nogles Mening platonisk Filosofi og kristelig Tro
forbundne. Denne Bestræbelse fortsættes af Kirkefædrene og
naar sin Afslutning i Teologiens klassiske Periode hos den
middelalderlige Kirkes store Lærere. Hos \emph{Thomas fra
Aqvino} (født 1227, død 1274), den katolske Kirkes endnu i
vore Dage hædrede og anerkendte Lærer, stiller Forholdet
mellem Tro og Videnskab sig paa følgende Maade.
Videnskaben udforsker alle \emph{underordnede} og \emph{specielle} Aarsager,
medens den Troendes Tanke angaar den \emph{første}, den \emph{absolute}
Aarsag til Alt. Kun Teologien lærer os derfor ret at
forstaa Tingene; Videnskaben lærer os egentligt kun, \emph{hvad}
Tingene ere, ikke \emph{hvorfra} de stamme. Men der er ingen
Strid mellem Troen, der er Teologiens Princip, og Fornuften,
der er Videnskabens Princip. Fornuften stammer fra Gud lige
saa vel som Troen, og Troen strider ikke mod Fornuften,
skønt den gaar ud over Fornuften.

I poetisk Form finde vi dette karakteristiske
middelalderlige Standpunkt hos \emph{Dante} (født 1265, død 1321), som
var en Discipel af Thomas. I sin »Comedia« beretter Digteren,
hvorledes det, for at hans Sjæl kunde blive lutret, forundes
ham at gennemvandre den hinsidige Verden, Helvede,
Skærsilden og Paradiset. Det er hans forklarede Elskede,
Beatrice, Symbolet for den paa himmelsk Aabenbaring grundede
Teologi, som skænker ham denne Gunst og sender Vergil,
Symbolet for den naturlige Videnskab, til ham for at lede
ham gennem Underverdenen. Men Himlen tør Vergil ikke
betræde; han angiver selv, med stor Sagtmodighed, Grunden:

\begin{verse}
ti Herskeren, hvem Himlens Stad tilhører, \\
vil ej, at jeg dens Tærskel maa betræde, \\
fordi jeg mod hans Love var Oprører.
\end{verse}

Da de Vandrende derfor nærme sig Himlen, aabenbarer
Beatrice sig selv, og Vergil forsvinder.

% --- p. 198 ---
\opage{198}Saa from og ydmyg var den naturlige Fornuft bleven i
Middelalderen. Men dens Kræfter voksede, og Middelalderens
Slutning betegnes ved voldsom Splid mellem Beatrice og
Vergil, Tro og Viden. Vi træffe her Problemet stillet for
første Gang.

\section*{III.}

Den stolte Bygning, som Thomas fra Aqvino havde
fuldført, og som Dante havde besunget, vaklede snart i sin
Grundvold. Den fortsatte Grublen opdagede, at den var
bygget af Elementer, der ikke kunde bringes i Harmoni.
Allerede \emph{Vilhelm af Occam} († 1347) finder en Modstrid mellem
den naturlige Fornuft og den over enhver Lov hævede
guddommelige Almagt. Han lærer, at Noget kan være sandt i
Teologien, som ikke er det i Filosofien: dersom Bibelen og
Kirken lære Noget, som er selvmodsigende, maa det alligevel
tages for Sandhed! Teologien bygger paa overnaturlig Avtoritet
og er derfor ingen Videnskab.

Have vi da to Slags Erkendelse, to Slags Logik? Dette
Spørgsmaal, der ligger saa nær, bliver ikke fremdraget.
Occam selv var troende Katolik; han paaviste Spliden, men
forfulgte ikke dens Konsekvenser. Den katolske Kirke holdt
fast ved sin hellige Thomas og har ikke tilladt nogen anden
Opfattelse af Tro og Viden at gøre sig gældende inden for
dens Enemærker. Endnu i vort Aarhundrede ere ikke mindre
end tre katolske Filosofer blevne dømte, fordi de enten vare
mere eller mindre rationalistiske end St. Thomas. For den
katolske Kirke eksisterer der her ikke længer noget Problem.
--- Lad os da se, hvorledes Sagen stiller sig inden for
Protestantismen.

\section*{IV.}

I den nyere Tid breder Livet sig i saare mange
forskellige Retninger og har et langt mere broget og
mangesidigt Præg end i Oldtid og Middelalder. Og Kirken staar
kun som ét af de mange Omraader for Livsbevægelsen; den
% --- p. 199 ---
\opage{199}er ikke længer Et og Alt som i den første Menighed, eller
Kronen paa det Hele som i Middelalderen. Historien er ikke
\emph{først og fremmest} en Kirkehistorie, men \emph{blandt Andet} ogsaa
Kirkehistorie. Overvejende passivt modtager Kirken de nye
Livsretningers Tilbageslag og søger at finde sig til Rette
over for dem; men de udgaa ikke længer fra den.

Den kirkelige Protestantismes Forhold til Videnskaben
er omtrent det samme som Katolicismens. Luther havde vel
beraabt sig --- ej blot paa den hellige Skrift, men ogsaa paa
»klare Grunde«; men ellers anvendte han paa Fornuften
Pauli Ord, at Kvinden skal tie i Menigheden. Altsaa
Fornuften skulde have Vergils Rolle, maatte ikke betræde det
Allerhelligste. Den protestantiske Kirke vilde gerne vedblive at
spille »den guddommelige Komedie«. Men Vergil har ikke
længer vist sig medgørlig.

Den protestantiske Ortodoksi blev en ny Skolastik, men
mere ængstelig og mindre storartet end den middelalderlige.
--- Den har ikke fundet nogen Dante. De store
protestantiske Digtere gaa uden om den eller imod den. Poesiens
Forhold til Kirken danner en Parallel til Videnskabens:
Oldkirken har ingen Poesi, Katolicismen en kirkelig Poesi,
Protestantismen en uden for Kirken staaende, om ikke
antikirkelig Poesi.

For Protestantismen er \emph{det 18. Aarhundrede}, hvad det
14. var for Katolicismen. Men den Krise, som det bragte,
var langt grundigere, bundede langt dybere end
Renaissancebevægelsen. Det er \emph{den selvstændige humane Livsanskuelses
Fødselstid}. Forløberne vare naturligvis mange --- men først
nu forstodes og anerkendtes de. Først nu fandt f. Eks.
Spinoza, der i halvandet Aarhundrede havde ligget »som en
død Hund«, en Kreds, der kunde tænke hans Tanker og
tilegne sig dem. Først nu fik den teologiske Livsanskuelse
en jævnbyrdig Modstander. Jeg tænker her ikke paa noget
bestemt, afsluttet System, men paa en almindelig Tanke- og
Livsretning, som i de store Træk var fælles for Mænd som
\emph{Lessing} og \emph{Kant}, \emph{Schiller} og \emph{Goethe}. --- Som man ser, fandt
denne nye Livsanskuelse straks mere end én Dante.

Det Ejendommelige ved dette Standpunkt, hvis Hoved%
% --- p. 200 ---
\opage{200}tanker kun forbigaaende ere blevne stillede i Skygge ved
den romantiske Reaktion, er --- \emph{fra den teoretiske Side set}
---, at man over for den teologiske Tro ikke stiller en
foregiven Indsigt i de teologiske Lærdommes Umulighed, men en
Paastand om deres Ubevislighed. Kant’s store Gerning var,
fra denne Side set, at give en videnskabelig Lære om
Videnskabens Grænser. Ikke ydmyget og kuet, som en Oprører,
der er falden til Føje, ikke ude fra hæmmet, men med
selvstændig og klar Selverkendelse siger den menneskelige
Fornuft her: saa langt kan jeg naa, og ikke videre! ---
Erkendelsens Grænser følge af Erkendelsens Natur. Ved
Undersøgelse af vor Erkendelses sidste Forudsætninger og
Love afstikkes det Omraade, paa hvilket den kan bevæge
sig. Dette Omraade er uoverskueligt, men dog bestemt
afgrænset over for det, der efter sin Natur ikke kan blive
Genstand for vor Erkendelse, saaledes som denne nu én Gang
er beskaffen.

For Kant og hans Tid laa Erkendelsens Betydning ikke
i, at den skulde give os »den absolute Sandhed«, men deri,
at vi ved den stedse kunne gaa frem i Sandhed. Den evige
Søgen, den evige Stræben blev fra nu af Slagordet:

\begin{verse}
Wer immer strebend sich bemüht, \\
den können wir erlösen.
\end{verse}

Ikke som om Intet kunde findes eller opnaas: men ethvert
Fund, ethvert opnaaet Maal staar som Udgangspunkt for ny
Stræben. Det var en stor Livs- og Verdenslov, man her
gav Udtryk, --- Udviklingstanken i en af dens vigtigste
Former, en Tanke, som vor Tid har givet en realistisk
Underbygning og en nærmere kritisk Bestemmelse, men som det
er hin Slægts Ære at have udtalt for første Gang med fuld
og klar Bevidsthed.

Men endnu vigtigere var et andet Princip, der blev
slaaet fast som Hovedpunkt i den humane Livsanskuelse: \emph{den
frie Samvittigheds Princip}. Kants anden store Gerning var,
at han hævdede Muligheden og Nødvendigheden af en rent
human Etik, fastslog, at Samvittigheden lige saa lidt som
den teoretiske Forsken kan lade sig sine Grænser afstikke
% --- p. 201 ---
\opage{201}ude fra. Maalestokken for Alt, hvad vi skulle anerkende
som godt, ædelt og ophøjet, maa findes i vort eget Indre og
være bestemt ved den menneskelige Natur.

Naar de, der staa paa dette Standpunkt, ikke kunne
gaa ind paa den kirkelige Tro, saa er det, fordi den strider
ej blot mod deres Forstand, men mod deres Samvittighed.
Man skal ikke blot slutte sig til en én Gang for alle given
Formulering af Sandheden, men ogsaa bøje sig for Lærdomme,
der ikke kunne bringes til at stemme med Samvittighedens
Maalestok for Godt og Ondt.

Først her bliver Problemet skærpet, og dets egentlige
Brod træder for Dagen. Hvad der fra den ene Side gøres
til den højeste, paa en Maade den eneste Pligt, en Pligt,
hvis Ikke-Opfyldelse medfører Fortabelse, --- det afvises fra
den anden Side netop som stridende mod Samvittighedens klare
og bestemte Fordring! ---

Biskop Martensen har fra sit Standpunkt klart set, at
det vigtigste Vendepunkt i Striden om Tro og Viden
fremkom ved den etiske Humanismes Grundlæggelse i Slutningen
af forrige Aarhundrede. Han beskæftiger sig i sin
»Kristelige Etik« vidtløftigt med denne Retnings Repræsentanter og
deres Forhold til Kirkens Tro. Den sidste Forklaring af
deres Vægring ved at slutte sig til denne finder han med
Rette ikke i deres Forstand, men i deres Villie --- kun at
denne Villie efter hans Paastand er Trods mod det Hellige:
»Mod Guds Hellighed have de en hemmelig Antipati!«

Hvor underligt, eller rettere hvor forfærdeligt maa det
menneskelige Liv dog ikke tage sig ud fra det ortodokse
Standpunkt! Den redeligste Søgen og Stræben, den strengeste
Samvittighedsfuldhed forenet med den mest uforbeholdne
Anerkendelse af Ufuldkommenheden i det, der er naat, maa fra
dette Standpunkt udlægges som Antipati mod Helligheden.
--- Har man ikke før forstaaet, hvorfor Problemet om Tro og
Viden saa let sætter Lidenskaberne i Bevægelse, saa forstaar
man det her.

% --- p. 202 ---
\opage{202}
\section*{V.}

Det nye humane Standpunkt syntes at have ret god
Udsigt til at trænge igennem. Tidsalderens mest fremragende
Aander hyldede det, de ledende Statsmænd virkede i den af
det betegnede Retning, i selve Kirken var Ortodoksien
afgjort svækket, og de konfessionelle Forskelligheder traadte
tilbage for en oplyst og frisindet Tankegang, der var fælles
for Alle. Saaledes vare ved Aarhundredets Begyndelse
Forholdene ikke blot i den protestantiske, men ogsaa i den
katolske Kirke.

Men hvor langt anderledes skulde Forholdene ikke
udvikle sig i Løbet af Aarhundredet, og hvor forandrede vare
de ikke allerede efter de to første Tiaars Udløb!

Et enkelt Træk viser tydeligt Forandringen. I Aaret
1799 udgav Schleiermacher sine »Taler om Religionen til de
Dannede blandt dens Foragtere«. Han kæmpede her
ingenlunde for nogen dogmatisk Tro, men hævdede en panteistisk
Religiøsitet og forherligede Spinoza. For en saadan
Anskuelse vilde han vinde Anerkendelse hos de Dannede. Men da
han i Aaret 1821 besørgede den tredje Udgave af sin Bog,
erklærede han i Fortalen, at naar man \emph{nu} saa sig om blandt
de Dannede, vilde man snarere finde det fornødent at skrive
Taler til Heldører og Bogstavtrælle, uvidende og ukærlige
Fordømmere, Overtroiske og Hyperdogmatikere!

Fire Aar i Forvejen (1817) havde den holstenske Præst
Claus Harms til Minde om Luthers Optræden udgivet 95 nye
Teser, i hvilke han blandt Andet protesterer mod »vor Tids
Paver«: Fornuften paa det teoretiske Omraade, og
Samvittigheden paa det praktiske Omraade. --- Hvor var det dog muligt,
at en saadan Optræden kunde vinde Tilslutning og betegne en
ny Periodes Begyndelse?

Hvad var i det Hele Grunden til, at Humanismens
Gennembrud i Slutningen af det 18. Aarhundrede ikke kom
til at indlede en Udvikling, der fortsattes i lige Linie
indtil nu, men afløstes af sin skarpe og bevidste Modsætning?

Da en nærmere psykologisk-historisk Forklaring af det
vældige Omslag, der kom til at udfylde saa stor en Del af
% --- p. 203 ---
\opage{203}vort Aarhundrede, og gjorde det til »Modsætningernes Tid«,
vil kunne tjene til nærmere Belysning af vort Problems
Natur, ville vi gaa noget nærmere ind derpaa, men saaledes,
at vi især holde os til de Aarsager, der kunne findes i selve
Beskaffenheden og Ejendommeligheden af det 18. Aarhundredes
Humanisme. Der virkede naturligvis mange forskellige
Aarsager med. Saaledes fik den almindelige politiske Reaktion
ogsaa Indflydelse paa den religiøse Udvikling;
Avtoritetsprincipets Genopstandelse i Staten maatte ogsaa føre til dets
Begunstigelse paa det religiøse Omraade, og de politiske
Bagstrævere benyttede sig ofte af Religionen paa den
modbydeligste Maade. Men dersom der ikke havde været dybt
liggende indre Aarsager, vilde denne ydre Indflydelse ikke
have kunnet udrette saa Meget. ---

1. Enhver Reaktion skyldes for en stor Del en
\emph{Kontrastvirkning}. Ligesom det af Lyset trættede Øje finder Hvile
ved Mørket eller dog ved et mattere Skær, saaledes kan i
det Hele den menneskelige Aand ikke udholde en stærk og
langvarig Anspændelse af en enkelt Side i dens Natur, men
kaster sig uvilkaarligt over til den modsatte Side for at
finde en ny Hvilestilling. Med Hensyn til det Problem, vi
her tale om, er det især \emph{Modsætningen mellem den aktive og
den passive, den virkende og den modtagende Side i vor
Natur}, som bliver af Betydning. Ikke som om der i vort
aandelige Liv kan paavises Tilstande eller Fænomener, hvor
vi enten forholde os absolut virkende eller absolut modtagende:
ved enhver Form for aandeligt Liv og i enhver aandelig
Tilstand kan der tvertimod paavises baade en Aktivitet og
en Passivitet. Men disse to Momenter kunne staa i et meget
forskelligt indbyrdes Forhold; snart har det ene, snart det
andet Overvægten. Saaledes have den forskende og
grublende Tænken, den prøvende og fremadstræbende
Samvittighed og den formende og skabende Fantasi en overvejende
aktiv Karakter, og i Modsætning til dem staa i denne
Henseende Trangen til Trøst og Hengivelse, og Følelser som
Pietet og Resignation. Hvor det gælder om at frembringe
nye Former for Livet, faar den aktive Side i vor Natur
Overvægten og lægger Beslag paa hele Energien. En saa%
% --- p. 204 ---
\opage{204}dan Opgave stillede sig netop for Slægten i den sidste Del
af det 18. Aarhundrede. Det var den Tid, hvor ubegrænset
Fremadskriden var Slagordet, og hvor Træghed betragtedes
som Arvesynden. Men der maatte komme en Tid, da
Menneskenaturens andre Sider krævede deres Tilfredsstillelse, og
hvor man søgte Hvile efter det rastløse Arbejde. Og Hvilen
fandtes ikke i de nye Former. De behøvede endnu alt for
meget aktiv Medvirken, stillede stadigt Fordring om
Selvvirksomhed. Hvile kunde derimod maaske findes i de gamle
Former, som man i sin Fremstormen havde vendt Ryggen,
men om hvilke man nu opdagede, hvor vidunderligt de vare
tilpassede efter Følelseslivets Trang: i Aarhundreder eller
Aartusinder havde Slægt efter Slægt nedlagt sine Erfaringer
i dem, havde lagt dem til Rette, saa der ikke behøvedes
uafbrudt Arbejde for at bruge dem. Vor Aktivitet søger det
Nye; med vor virkende Natur leve vi i Nutiden; men den
passive og modtagende Side i os bestemmes og næres
overvejende ved den forudgaaende Tids Eftervirkninger; her
nærmer det aandelige Liv sig det Ubevidste og drager Næring
af det. Et Eksempel herpaa fra det enkelte Individs Liv
have vi i Barndomserindringernes Magt: de virke uden vor
bevidste Medvirkning og stige især op med overvældende Styrke,
naar Sindet er træt af aktivt Aandsarbejde.

Det er karakteristisk, at Schleiermacher netop ved
Overgangen til det nye Aarhundrede opstillede sin berømte
Definition af Religionen som den ubetingede Afhængighedsfølelse.
Det var som et Program for den Retning, Aandslivet nu
skulde tage. Humanismen blev afløst af Konfessionalismen,
den frie Forsken og den frie Samvittighed af
Avtoritetstroen.

Ere vi da henviste til et Sisyfusarbejde, ruller Stenen
stedse tilbage, naar vi vel have faaet den op paa Bjærget?
--- Nej, men i det Anførte ligger, at kun den Sandhed kan
blive sejrende, der kan gennemtrænge og fylde hele Livet.
Og det sker ikke med ét Ryk; det tager Tid, og først
gennem mange Svingninger skrider det fremad. Det synes
at være en psykologisk Lov, at \emph{Erkendelsen udvikler sig
hurtigere end Følelseslivet}. Kun langsomt nedfældes det aktive
% --- p. 205 ---
\opage{205}og bevidste Arbejdes Resultater i Følelsens uvilkaarlige
Liv. Før det kan ske, kommer der da naturligt en Periode,
i hvilken de nye Tanker blegne og vise sig utilfredsstillende;
Tvivl og Uro bemægtiger sig Mange, og der synes ikke at
være nogen anden Udvej end at kaste sig tilbage i
Følelseslivets gamle Former, hvis man ikke skal gaa til Grunde.
Men denne Disharmoni vil kun være forbigaaende: ere de
nye Tanker gyldige, udtrykke de virkeligt nogle af de Love,
vor Tilværelse er underlagt, saa ville de stedse arbejde sig
frem paa ny, og den ustadige Ligevægts Tid vil være forbi.
I Længden vil Følelseslivet udvikle sig i større eller mindre
Harmoni med Tankelivet: paa den Maade bleve jo selve hine
gamle Former og Dogmer til. Allerede selve
Reaktionsperioderne vise sig at bære Præg af det, de bekæmpe; i de
gamle Læderflasker vil der for en ikke ringe Del hældes ny
Vin, og kun den, der holder sig til det Udvortes, vil mene,
at Alt igen er blevet som det var før.

Vi staa her over for en aandelig Natursammenhæng, som
vi maa bøje os for, og som vi ikke kunne trodse. Den, der
har faaet ret Øje for den, vil hverken have Manges let købte
Tillid til Sejrens Nærhed eller tabe Modet ved tilsyneladende
Tilbagegang. Han véd, at ligesom de tilsyneladende
pludselige Revolutioner kun betegne, at det længe i Skjul
Forberedte nu træder frem for Dagens Lys, saaledes tyde de
pludselige Reaktioner paa, at der var Drifter og
Fornødenheder, som kun ufuldkomment bleve tilfredsstillede ved det
tilsyneladende saa glimrende Fremskridt, og som derfor maa
søge Tilfredsstillelse ved det Gamle, indtil de have lempet sig
efter det Nye, eller dette efter dem.

2. Det 18. Aarhundredes Humanisme gav ikke blot
den erkendende og virkende Del af vor Natur en afgjort
Overvægt over den følende og modtagende, men i nøje
Sammenhæng hermed havde den en afgjort
\emph{aandsaristokratisk} Karakter. Den hele Bevægelse var, og maatte efter sin
Natur være, indskrænket til en lille Kreds; det egentlige
Folk paavirkedes saare lidet af den. Til ingen Tid have
maaske de vel stillede Klasser i Samfundet vist saa stor Iver
for at gaa i Spidsen for Oplysningens og Fremskridtets Ar%
% --- p. 206 ---
\opage{206}bejde. Men Afstanden mellem de Ledende og dem, der trængte
til Ledelse, var overordentlig stor. Og Udviklingen havde
ført til et Brud med den \emph{fælles} Tro, de fælles Former for
de højeste Tanker. I Middelalderen kunde den dybest
grublende Skolastiker og den enfoldigste Træl mødes i samme
Bekendelse. Nu var denne Enhed sprængt, og der maatte,
især da samtidigt Følelsen for det Nationale og Folkelige tog
et stærkt Opsving, opstaa en Trang til ogsaa i Tro at være
Et med sit Folk. Reaktionen mod det 18. Aarhundrede fik
fra denne Side set en \emph{aandsdemokratisk} Karakter. Det var
jo --- for at tage et enkelt Eksempel --- den Overbevisning,
at »de Ulærdes Tro umuligt kan være bunden til de Lærdes
Vidnesbyrd«, der førte Grundtvig til hans »Opdagelse«.

Ligesom Spliden mellem de forskellige Sider eller Kræfter
i det aandelige Liv vil ogsaa Spliden inden for Folkets Liv
være et stadigt virkende Motiv til at vende tilbage til den
gamle Tro. Fælles Former og Udtryk for de højeste Tanker
skabes langsomt og forudsætte, at Tankerne ere gaaede over
i Kød og Blod hos Alle. Da nu de individuelle
Forskelligheder i religiøs Trang og religiøs Sans maa antages at
stige og tiltage, jo mere Religionsfriheden bliver til
Virkelighed, saa vil et saadant Fællesskab blive stedse vanskeligere
at tilvejebringe.

Et Vidnesbyrd om, hvor stor Trangen til Fællesskab kan
være, frembyder \emph{Schellings} og \emph{Hegels} Religionsfilosofi. Disse
Mænd og deres Tilhængere vare overbeviste om, at de i
deres filosofiske Ideer havde udviklet det samme Indhold, som
Kirkens Dogmer udtalte. Forskellen skulde kun angaa
Formen: hvad Teologien har som Forestilling, i Fantasiens og
Følelsens Form, det skulde Filosofien forme til abstrakt
Begreb, uden at derved det Væsentlige gik tabt. I Dogmet
om Guds Menneskeblivelse f. Eks. skulde indeholdes den
filosofiske Tanke, at selv det Højeste maa udvikle sig gennem
Kamp og Modsætninger, og at et Princip kun viser sin
Gyldighed, naar det formaar at gennemtrænge det
virkelige Liv med al dets Splid og med alle dets stridende
Forhold.

I og for sig er det fuldt berettiget at betragte Dogmerne
% --- p. 207 ---
\opage{207}som Symboler. For den Ikke-troende kunne de ikke være
Andet, naar de i det Hele skulle være Noget. Og Mange,
der ansé sig selv for meget troende og betragtes af Andre
som saadanne, staa i Virkeligheden paa dette Standpunkt.
Men det vil altid føre til mislige Følger, naar man benytter
den symbolske Opfattelse til at tilvejebringe et falskt Skin
af Fællesskab. En virkelig Enhed kan ad den Vej ikke
bringes til Veje. Paa et saadant falskt Skin beroede det
sidste Forsøg paa Forsoning af Tro og Viden, som er blevet
gjort, Hegels Religionsfilosofi, af hvilken Mange ventede sig
en »ny Æra«. Vergil skulde her ikke ydmyg underordne
sig Beatrice; de skulde følges Arm i Arm. Men denne
tilsyneladende Enighed gav kun Anledning til, at Modsætningen
mellem Tro og Viden blev udtalt i endnu langt skarpere og
mere tilspidset Form end hidtil, saaledes som det skete af
\emph{Strauss} og \emph{Ludwig Feuerbach} i Tyskland og af \emph{Søren
Kierkegaard} hos os. Disse Mænd vare enige i, at den teologiske
Tros Indhold var det Ubegribelige; Forskellen mellem dem
bestod i, at Strauss og Feuerbach sagde: Det er ubegribeligt,
og derfor kan det ikke tros! medens Kierkegaard sagde:
Netop fordi det er ubegribeligt, kan og skal det tros! ---
Strauss’s Ord: »Af falske Foreningsforsøg er der sket nok;
kun Skærpelse af Modsætningerne kan føre videre!« er
karakteristisk for Udviklingens Gang i det 19. Aarhundrede.
En berømt kirkehistorisk Forsker, \emph{F. C. Baur}, har med Rette
fundet vort Aarhundredes Ejendommelighed deri, at alle
Modsætninger --- i Videnskab, Religion og Politik --- udarbejdes
til den højeste mulige Grad. Det hænger naturligt sammen
med, at Reaktion og Revolution paa de forskellige Omraader
saa hyppigt have staaet over for hinanden; derved er der
opstaaet en klar Bevidsthed om Modsætningernes Natur og
en Bestræbelse for at udforme dem hver for sig i deres fulde
Konsekvens. Baade Nødvendigheden af at træffe et Valg
mellem uforenelige Modsætninger og Karakteren af de Veje,
man skal vælge mellem, gaar mere og mere op for den
almindelige Bevidsthed. ---

Har den blotte Trang til Fællesskab saaledes ført til
Sympati med den gamle Tro --- hvad enten denne Sympati
% --- p. 208 ---
\opage{208}førte til personlig Tilslutning til dens Bogstav eller blot til
symbolsk Udlægning af dens Udsagn, --- saa maatte dette hele
Spørgsmaal om Tro og Viden særligt stille sig paa en
alvorlig Maade for dem, der indsaa dets Sammenhæng med det
sociale Spørgsmaal. Selv om man indrømmer, at den
videnskabelige Kritik af Dogmerne fører til deres Opløsning, vil
dog den Tanke rejse sig, at disse Dogmer, saa fornuftstridige
de end maatte være, ere Manges eneste aandelige Tilhold og
Trøst, --- det, der holder deres Haab oppe under Livets tunge
Kamp, --- det Eneste, de raade over, der kan løfte deres
Blik ud over den trange Kreds, i hvilken deres Liv bevæger
sig. Fra dette Synspunkt har saaledes en af vor Tids ædleste
filosofiske Tænkere, \emph{Albert Lange}, Forfatteren til »Geschichte
des Materialismus«, udtalt sig over for Strauss og Feuerbach,
med hvem han er enig hvad den rent teoretiske Side af Sagen
angaar. Det er i det Hele karakteristisk, hvor stor en
Rolle Spørgsmaalet om Religionens Nytte og
Uundværlighed spiller i den religiøse Debat. Det gælder i første Linie
ikke: Sandhed eller ikke Sandhed? --- men: Trøst eller ikke
Trøst? --- I tidligere Tiders Debat kom denne Side af Sagen
ikke frem som Noget for sig, end sige som Hovedsagen.
Man søgte ikke efter Grunde og Motiver til at holde paa
den dogmatiske Tro trods Vanskelighederne; men det var
dem, der opgav den dogmatiske Tro, hvem det paalaa at
motivere deres Skridt.

Det er sikkert af meget stor Betydning, at det
aandelige Fællesskab ikke forsvinder, og at vor Slægt ikke mister
nogen Kraft, der kan holde den oppe i Kampen for
Tilværelsen. Det er med fuld Ret, at disse Hensyn blive
bestemmende for den Maade, paa hvilken den religiøse Kamp
føres; men selve Kampen kunne de kun for en Tid dæmpe,
da de ikke anvise nogen Vej til Problemets Løsning.

3. Endeligt havde den humane Livsanskuelse, saaledes
som det 18. Aarhundrede formaaede at fremstille den, en
ensidig \emph{abstrakt} og \emph{idealistisk} Karakter. Den udtaltes efter
sit almindelige Princip og med de vigtigste Forudsætninger
og Fordringer, den medfører, men savnede tilstrækkelig
Gennemførelse og realt Grundlag. Saa at sige ethvert Frem%
% --- p. 209 ---
\opage{209}skridt, det 19. Aarhundrede har gjort, har bidraget Sit til at
raade nogen Bod paa disse Mangler. Naturvidenskabernes
storartede Opsving har givet en fastere Basis for
Overbevisningen om den lovordnede Sammenhæng i Naturen. I
Udviklingens Begreb er der paavist en ledende Idé, om hvilken
al vor Viden paa Naturens og Historiens Omraade kan samle
sig. Den historiske Kritik og den sammenlignende
Religionsvidenskab have udvidet og berigtiget det Materiale, som er
nødvendigt til religionsfilosofiske Problemers Behandling. En
dybere gaaende psykologisk Forstaaelse af det menneskelige
Aandslivs forskellige Sider i deres indbyrdes Forhold og en
fornyet Drøftelse af de erkendelsesteoretiske Spørgsmaal giver
Problemet om Tro og Viden en noget anden Karakter end
det forhen har kunnet have. --- Alt dette er naturligvis kun
smaa Fremskridt i Forhold til Problemets Storhed. Dette vil
i lange Tider følge Slægten paa dens Vej, og den forskellige
Maade, hvorpaa det opfattes og behandles, vil kunne afgive
Maalestok for den aandelige Udviklings Grad og Retning til
de forskellige Tider. Her have vi kun tilbage at vise,
hvorledes Problemet maa stille sig i vore Dage.

\section*{VI}

Der blev i den Forhandling om Tro og Viden, som
i Tredserne fandt Sted i vor Literatur, lagt stor Vægt paa
Distinktionen mellem Tro og Teologi, og jeg selv hørte, som
Discipel af Rasmus Nielsen, til deres Tal, der knyttede ikke
ringe Forhaabninger til den. Den hjælper os aabenbart ikke,
og den er i hvert Tilfælde ikke holdbar. Saa snart Troen har
et bestemt Indhold, en bestemt Genstand, maa der ogsaa
udvikle sig en Lære om dette Indhold, som søger at paavise
dets indre Sammenhæng og dets Betydning. Ellers véd
hverken den Troende selv eller Andre, hvad det egentligt
er, han tror paa.

Nu til Dags har desuden Teologien opgivet at støtte sine
Paastande med videnskabelige Beviser. Den har opgivet alle
stolte Drømme om at faa Vergil til at gaa i Ledebaand og
% --- p. 210 ---
\opage{210}er tilfreds, naar man lader den uanfægtet fra den
videnskabelige Kritiks Side. Den erklærer sine Dogmer for \emph{ubevislige},
indrømmer maaske endog, at de frembyde Vanskeligheder for
den »naturlige« Forstand, og stiller sig væsentligt paa det
Standpunkt, at den hævder den dogmatiske Tro som en
aandelig Livsfornødenhed --- vel at mærke en Fornødenhed,
det er Pligt at have. Det skal være Tegn paa Forhærdelse
og Trods mod det Hellige, naar man ej føler den Trang, der
lader Fornuften bøje sig for det Ubegribelige. Deri ligger,
som jeg allerede har haft Lejlighed til at berøre, Problemets
egentlige Brod.

Paa den anden Side vil ingen videnskabelig Tænker
paastaa, at den absolute \emph{Umulighed} af Dogmernes Indhold
kan strengt bevises. Vi bygge i videnskabelig Forsken
paa de naturlige Aarsagers Princip. Det har hjulpet os
igennem mangen god Gang, og noget andet videnskabeligt
Princip til Forstaaelse af Tilværelsen vil neppe nogen Sinde
kunne opstilles; --- men det kan ikke bevises, at \emph{Alt} i
Tilværelsen er dette Princip underlagt.

Ja, vi maa endog gaa et Skridt videre. Samtidigt med
de vidunderlige Fremskridt, som ere gjorte i Kundskab om
Naturen og Historien, baade hvad Enkeltheder og hvad den
store Sammenhæng angaar, have vi i langt dybere Forstand
end forhen lært, at \emph{den hele Tilværelse er én stor Gaade}.
For Filosofen er det ringeste Fænomen i denne Henseende
lige saa lærerigt som det største og mest i Øjne faldende.
Det simpleste Aarsagsforhold giver os Mere end nok at tænke
paa --- indeslutter egentligt hele Verdensgaaden i sig. Alt
hvad vi forstaa og forklare, forstaa og forklare vi ved at
paavise Aarsagssammenhæng: vidste vi nu blot, hvad
Aarsagssammenhæng egentligt er, --- hvad det er, der holder
Alt i Naturen sammen og gør den til den vidunderlige
Helhed, der mere og mere afslører sig for os, --- saa vilde alle
Gaader være løste! Men til en saadan Kundskab er al
Udsigt spærret. Vore Fornemmelser, Forestillinger og Begreber
(og saaledes ogsaa Aarsagsbegrebet) ere ret gode
Hjælpemidler til at finde os til Rette i Tilværelsen, men vi have
ingen Ret til at mene, at de kunne afbilde Tilværelsens inder%
% --- p. 211 ---
\opage{211}ste Væsen for os. Vi kunne i det Højeste forme spekulative
Hypoteser om, hvorledes Tilværelsen i sig selv kunde være
beskaffen; men nogen Erfaringsbekræftelse kunne saadanne
Hypoteser i Følge Sagens Natur ikke faa.

Deraf følger nu aldeles ikke, at man bøjer sig for ethvert
teologisk Orakelsprog. Deraf følger kun, at Diskussionen om
Tro og Viden maa føres paa en noget anden Maade end den
hidtil er bleven ført. Lad os i Korthed berøre de
Hovedspørgsmaal, som her ville være at drøfte.

1. Naar den teologiske Tro vedkender sig, at dens
Paastande ere ubevislige eller maaske endog uforenelige med
nogle af Videnskabens vigtigste Forudsætninger og Resultater,
stiller den sig tilsyneladende paa et utilgængeligt og uangribeligt
Standpunkt. Der synes ikke mere at være nogen fælles Grund,
paa hvilken Kampen kan føres.

Men naar den teologiske Tro alligevel --- for at udtrykke
sit Indhold --- benytter Forestillinger og Begreber, som
ogsaa den naturlige Fornuft og den videnskabelige Tænkning
anvender, saa maa der spørges, om disse Begreber tages i
aldeles samme Betydning. Det har vist sig forbundet med
uovervindelige Vanskeligheder, selv for den dristigste
Spekulation, at gennemføre Sætninger og Domme om det, der ligger
ud over enhver mulig Erfaring. Hvad kan det nu egentligt
være, som bevirker, at Troen her kommer bedre fra det end
den videnskabelige Tænken? \emph{Hvorfor bortfalde de store
Vanskeligheder ved Dogmerne, naar man gaar fra Viden over
til Troen?} En Skabelse af Intet f. Eks., --- en første Aarsag,
--- et absolut godt Væsen, som dog er Aarsag til Alt, altsaa
ogsaa til det Onde, --- saadanne Begreber, mod hvilke den
menneskelige Tanke er tørnet atter og atter, indtil den til
sidst har opgivet Forsøgene paa at forene dem med sine
andre Begreber, --- hvorledes kunne de stille sig i et helt andet
Lys, fordi de erklæres at være Genstande for Tro, ikke for
Viden?

Svares der (som f. Eks. af Biskop Martensen i hans
Dogmatik), at den Genfødelse af Menneskenaturen, som finder
Sted hos den Troende, ogsaa omfatter Erkendelsen, saa at
% --- p. 212 ---
\opage{212}\emph{det} bliver forstaaeligt, der før var ubegribeligt eller endog
selvmodsigende, saa kunne naturligvis de Ikkegenfødte ikke
tale med om, hvorledes Tingene tage sig ud efter en saa
radikal Forandring. Men da de Troende og Genfødte dog
ikke ganske opgive Sammenhængen med den naturlige
Verden og, for saa vidt de vedblive at høre med til den, ogsaa
maa vedblive at bruge deres naturlige Fornuft, saa maatte
de dog kunne sammenligne de to Arter af Erkendelse, den
naturlige og den genfødte, og sige os, hvori Forskellen
egentligt bestaar. Her er der et Kapitel i Dogmatiken,
som Teologerne i alle de mange Aar, i hvilke Teologien
har floreret, have forsømt at skrive, nemlig: en ny eller højere
Logik og Erkendelsesteori, som kunde vise os, ved Hjælp af
hvilke nye Principer man i Troen slipper forbi de Skær, paa
hvilke man i den naturlige Tænkning bestandigt strander.
Indtil dette Kapitel er skrevet, maa vi have Ret til at
undersøge de teologiske Begreber paa samme Maade, som vi
undersøge alle andre, og fælde samme Dom om Inkonsekvenser og
uberettigede Begrebsanvendelser, som vi alle andre Steder
fælde. En saadan Undersøgelse vil føre til det Resultat, \emph{at
Dogmet i Virkeligheden staar lige saa afmægtigt over for de
store Gaader som Videnskaben}.

2. Men --- vil man da sige fra teologisk Side (især nu,
da den filosofiske Interesse i vor teologiske Verden synes at
være skrinlagt med Biskop Martensen) --- det er dog en
Misforstaaelse at mene, at Troen skal fyldestgøre en teoretisk
Trang. For Troen har Tænkning og Refleksion Intet at
betyde, og al Verdens Modgrunde falde magtesløse til Jorden
af den simple Grund, at Ingen kommer til Troen uden den,
som søger Trøst for en ængstet Samvittighed. Det er
\emph{Følelsen}, særlig den ængstede Samvittighed, ikke nogen anden
Kraft, der her virker.

Dette er ikke nogen behagelig Indvending, hvis
Meningen er, at den Talende selv ikke kan undvære den Trøst,
der ligger i Troen; thi i saa Fald maa Diskussionen høre op.
Hvem vil berøve Nogen hans Trøst? Selv den ivrigste
Kemiker vilde ikke bruge til Undersøgelse den Drik Vand, hvor%
% --- p. 213 ---
\opage{213}med et Menneske skal læske sin brændende Tørst. Enhver
Diskussion forudsætter, at Deltagerne hver for sig have saa
megen Ro og Kraft i Sindet, at de kunne se Sandheden
under Øjne, saa man ikke behøver, som over for Syge, at dølge
Noget.

Hvad der trøster den Ene, trøster desuden ikke den
Anden, eller vækker endog hans Forargelse og Indignation.
Saa snart der appelleres til Følelsen, appelleres der mere
eller mindre til noget rent Individuelt. En saadan Appel
kan aldrig være Argument; thi \emph{Følelsen i og for sig lærer
os aldeles Intet uden vor egen Tilstand}. Den mest levende
Glæde og Beroligelse, jeg føler ved en Efterretning,
garanterer Intet om denne Efterretnings Vederhæftighed. Fra min
Følelse kan jeg derfor ikke med Sikkerhed slutte til Noget,
som ligger ud over min egen Tilstand. Vore Ønsker bestemme
ikke, hvad der er Sandhed; men det kunde jo være, at H. C.
Ørsted havde Ret i, at den fulde Sandhed altid bringer sin
Trøst med sig.

Den Følelse, paa hvilken man især beraaber sig som
Støtte for den teologiske Tro, er \emph{Syndsbevidstheden}. Her
opstaar nu det store Spørgsmaal, om Syndsbevidstheden,
saaledes som Teologien opfatter den, virkeligt er en, etisk
set, sund og berettiget Følelse, --- om der ikke skulde
indeholdes Elementer i den, som ikke ere virkeligt etiske, men
som hænge sammen med barbariske Forestillinger (om en
vred Guddom, der skal forsones), og som maa udskilles,
for at Syndsbevidstheden kan blive til virkelig etisk
Angerfølelse.

Men det er netop Maalestokken for det Etiske, der er
saa forskellig paa de to Standpunkter! Og Teologien, der
før viste os fra det Dogmatiske over til det Etiske, hævder
nu omvendt, at uden Dogmetro er ingen Etik mulig. Appellen
til Syndsbevidstheden leder os derfor kun i en Kreds: thi
Syndsbevidstheden er en af bestemte dogmatiske
Forestillinger betinget Følelse, og det er da intet Under, om den
aflægger Vidnesbyrd til Fordel for Dogmet. --- Vi komme kun
ud af denne Kreds, dersom Teologien kan vise os, hvad der
hidtil ikke er lykkedes, at kun dogmatisk Tro kan være
% --- p. 214 ---
\opage{214}Grundlag for Etik, ja, at de teologiske Antagelser i det
Hele taget egne sig til at være Grundlaget for en Etik.
Dersom derimod en human Etik er mulig, og dersom en
ubevidst human Etik skulde vise sig stedse at ligge til
Grund for den teologiske Etik i dens forskellige Former
(saaledes som man kunde slutte af den Kendsgerning, at naar
Menneskene forbedre sig, forbedre ogsaa deres Guder sig),
saa vilde den skarpeste Brod være taget bort, og hvad
der saa blev tilbage af Problemet, vilde vel stedse kunne
sætte Følelsen i stærk og dyb Bevægelse, men ikke længer
spalte vor Slægt i to modsatte Lejre, der støde sammen
i uforligelig Strid om den sidste Maalestok for Godt og
Ondt.

3. Der er en Paastand, som længe utaalmodigt har
ventet paa at komme til Orde fra teologisk Side. Den
lyder omtrent saaledes: »Troen er bygget paa \emph{historiske
Kendsgerninger;} dens Indhold er først og fremmest Historie
og hverken noget udspekuleret System eller blotte
Følelsespostulater. Og historiske Kendsgerninger maa man tage,
som de ere; ved dem kan intet Fornuftræsonnement rokke.«

Naar man nu blot her tager Ordet Historie i samme
Betydning, som vi tage det i den historiske Videnskab!
Men det gaar med den teologiske Brug af dette Begreb
som med Brugen af andre Begreber (Aarsag, Etik o. s. v.):
trods Ordligheden er den virkelige Betydning af
Begreberne saare forskellig. Historisk er ikke Alt hvad der
staar i gamle Bøger eller indeholdes i gamle Traditioner.
En historisk Kendsgerning kan ofte kræve indtrængende
kritisk Forsken, før den kan gøre Fordring paa
Anerkendelse. Det store Spørgsmaal bliver da, om Teologien
vil underkaste sig den historiske Kritiks sædvanlige Love,
eller om den gør Fordring paa visse Privilegier for sine
»Kendsgerninger«. At den maa fordre saadanne
Privilegier, følger allerede af, at dens Kildeskrifter indeholde
Fortællinger om Undere, der kræve Tro for at blive
anerkendte. En historisk Forsken, der ikke anvender andre
Forudsætninger end dem, som al Videnskab maa gaa ud
% --- p. 215 ---
\opage{215}fra, fører --- som Erfaring stedse har vist --- til en helt
anden Opfattelse af Begivenhederne ved Kirkens
Grundlæggelse end den, som den teologiske Tro fastslaar. Det
gaar altsaa med Appellen til Historien som med Appellen
til Følelsen og til Etiken: den fører i en Kreds, thi
ved Historie mener man den \emph{teologisk} opfattede og vurderede
Historie. ---

\streg{}

% --- p. 216 ---
\opage{216}
\begin{center}

\essay[DET RELIGIØSE PROBLEM]{DET RELIGIØSE PROBLEM.}

Tale ved Universitetets Reformationsfest den 16. November 1897.\end{center}

\streg{}

Den Fest, vi fejre, er helliget Mindet om et af de store
Gennembrud i Historien, hvor det, der var forberedt ved
Tankens og Samvittighedens indre Arbejde, gennem kraftige
Personligheder krævede Ret til at raade ogsaa i det ydre Liv.
Reformationen staar som et Vidnesbyrd om, at der paa det
aandelige Omraade, ligesom i Naturlivet, kan finde en
Udvikling Sted, der fører til nye Livsformer. Herimod strider
det ikke, at Reformatorerne væsentligt opfattede deres Værk
som en Fornyelse af den oprindelige Kristendom. De staar
her, som Renæssancens Mænd i det Hele stod. Man vilde
paa alle Omraader en Genfødelse af Oldtiden. I Kunst og i
Literatur, i Tænkning og i Naturforskning saa vel som i
Religionen beraabte man sig paa Oldtidens Forbilleder. Den
nye Kunst inspireredes af Antiken; den nye Filosofi gik i
Skole hos Platon, Aristoteles og Lucretius; Kopernikus søgte
at skaffe sig Anerkendelse for sine nye astronomiske Teorier
ved at beraabe sig paa lignende Forestillinger hos de gamle
Pythagoræere, og Atomlæren, der skulde spille saa stor en
Rolle i den moderne Fysik og Kemi, fremstilledes først som
en Rekonstruktion af den antike Atomistik. Machiavelli har i
sine »Discorsi«, der forfattedes omtrent samtidigt med
Reformationens Udbrud, udtalt Tidens Tendens, naar han siger,
at religiøse og politiske Institutioner fra Tid til anden maa
føres tilbage til deres Udspring, bedømmes og beaandes ved
Sammenligning med dette.

% --- p. 217 ---
\opage{217}I Ordets absolute Betydning er Genfødelse af en tidligere
Tids Tanke eller Tro ikke mulig. Som Historien viser, er
Genfødelse i Virkeligheden et nyt Liv. At Fordybelsen i
Oldtiden paa Kunstens, Tænkningens og Naturforskningens
Omraade ikke bragte det Antike frem igen, men tjente til
Forberedelse for nye Former og Retninger, benegtes ikke af
Nogen, og det forholder sig heller ikke anderledes paa det
religiøse Livs Omraade. Det kunde ikke være anderledes,
hvis Livet ikke skulde gaa tilbage. Den Trang, der førte
til ny Fordybelse i Oldtiden, var i Virkeligheden en ny
Trang, en Trang, der var opstaaet ved, at nye Betingelser
og Fornødenheder havde udviklet sig paa det aandelige
Omraade.

Lige saa langt som det var fra, at den mærkelige Tid,
hvis Minde vi fejre, genfrembragte det Antike og det
Primitive, lige saa langt er det fra, at den en Gang for alle skulde
have fundet og udtalt Sandheden. »Sandhed, du erkendte
Fædreneland!« hedder det i vor Reformationskantate. Men
Sandheden er stor og omfattende, og ingen Slægt kan finde
den for alle Tider. Selve den Sandhed, Reformatorerne
erkendte, maa atter og atter prøves efter dens Gyldighed. Det
religiøse Problem staar den Dag i Dag, og i langt skarpere
og mere bevidst Form end før, i første Række blandt det
aandelige Livs Problemer. Det stilles nu paa en væsentlig
anden Maade end paa Reformationstiden. Der er kommet
nye Elementer til paa Aandslivets Omraade, Elementer,
Reformatorerne ikke kendte eller ikke agtede, og som give
selve Problemet ny Betydning og stille Vilkaarene for dets
Behandling væsentligt anderledes. Det vilde ikke være i
Reformationens Aand, om vi vilde lukke Øjnene herfor. Men
det vil være i Reformationens Aand, naar vi anstille en
Undersøgelse af de Elementer, der --- saa vidt vi hidtil kunne
se --- ere af væsentlig Indflydelse paa det religiøse Liv og
den religiøse Tro.

Der gives neppe nogen af de Videnskaber, der ere
repræsenterede her ved Universitetet, der ikke vilde kunne
give et Bidrag til en saadan Undersøgelse. Jeg skal fra mit
Standpunkt fortsætte, hvad to af mine Kolleger have indledet
% --- p. 218 ---
\opage{218}i Taler ved denne Fest. En af mine teologiske Kolleger
har talt om den historiske Kritiks Betydning ved
Fastsættelsen af den religiøse Tradition, og min astronomiske
Kollega har talt om Erfaringsvidenskabens stille, men stadigt
voksende Indflydelse paa religiøse Forestillingers Indhold og
Betydning.

Der er, saa vidt jeg ser, fire Elementer, som utvivlsomt
bestemme den Karakter, menneskelig Livsanskuelse, hvad
enten den antager Form af religiøs Tro eller ikke, antager
til en bestemt given Tid. De kunne i Korthed betegnes som
Overlevering, Iagttagelse, Tænkning og Villie.

Religion bestaar fra Slægt til Slægt som en ærværdig
Overlevering, i hvilken den Enkelte uvilkaarligt lever sig
ind. Den giver Form og Indhold for Trangen til at forstaa
Livet, dets Betydning og dets Opgaver. Et religiøst
Problem opstaar, hvad dette Punkt angaar, naar der spørges,
hvorfra Overleveringen stammer, hvor vidt man kan stole
paa dens Ægthed, og hvem der i Tvivlstilfælde skal afgøre
saadanne Spørgsmaal. I vore Dage gælder det her om den
historiske Kritiks Forhold til kirkelig Overlevering. Jeg
behøver ikke at minde om, at dette Emne netop i de seneste
Aar er kommet op ogsaa her hjemme. Et endnu mere
brændende og dybtgaaende Spørgsmaal er det, om den
Livsførelse, der nu kalder sig med Kristendommens Navn,
virkeligt er den samme, som den, der oprindeligt var knyttet til
dette Navn. Fra denne Side blev det religiøse Problem rejst
her hjemme af vort Folks største religiøse Tænker i Form af
det Spørgsmaal: Er det nye Testamentes Kristendom
virkeligt til?

Den Iagttagelse og Kundskab, den Enkelte kan vinde,
vil nødvendigvis faa Indflydelse paa hans Tro. Den vil
bestemme den Maade, paa hvilken han opfatter
Overleveringen, --- hvad han virkeligt tilegner sig, og hvad han lader
ligge, --- hvad han tager i bogstavelig Betydning, og hvad
han opfatter som blot Symbol. Heri bestaar, hvad min
astronomiske Kollega har kaldet Iagttagelsens religiøse
Betydning.

Naar jeg stiller Tænkningen som et eget Element i
% --- p. 219 ---
\opage{219}denne Sammenhæng, mener jeg ved Tænkning det
aandelige Arbejde, ved hvilket vi blive os Indholdet af vore
Forestillinger og det indbyrdes Forhold mellem vore
forskellige Forestillinger klart bevidst. Et saadant aandeligt
Arbejde foregaar uafbrudt og uvilkaarligt, med større eller
mindre Energi. Her rejser sig en hel Række Spørgsmaal.
Kunne religiøse Antagelser bevises ad Fornuftens Vej? Den
katolske Kirke har ved sin højeste Autoritet slaaet fast, at
det ikke blot er muligt at bevise den saakaldte naturlige
Religions Sandheder (særligt Guds Tilværelse), men ogsaa at
føre et Bevis for Aabenbaringens og den ufejlbare Kirkes
Nødvendighed. I den protestantiske Verden, hvor den
kritiske Filosofi (uden at man kan tilskrive den protestantiske
Kirke nogen Ære derfor) gennemgaaende har fortrængt
Skolastiken, stiller Spørgsmaalet sig anderledes. Der spørges
ikke om, hvor vidt religiøse Antagelser kunne bevises, men om
de kunne gendrives. Naar de ikke kunne gendrives, mener man
sig berettiget til at fastholde dem. Herimod rejser der sig
dog Indsigelse fra anden Side. Det er en Regel, der
anerkendes overalt i Videnskaben, at man ikke skal anvende
andre Antagelser og Formodninger end dem, der med streng
Nødvendighed fremgaar af de kritisk prøvede Erfaringer; det
er ikke nok, at en Sætning ikke strider mod, hvad der kan
bevises. Kampen staar mellem disse to Standpunkter: fra
det ene hævdes Retten til at fastholde, hvad der ikke
videnskabeligt kan gendrives; paa det andet hævdes Forpligtelsen
til ikke at antage Andet, end hvad der med streng
Nødvendighed fremgaa af virkelig Iagttagelse. Fra det sidste
Standpunkt hævdes, at den Sparsommelighedens Lov, der
gælder for videnskabelige Antagelser, ogsaa bør gøre og
efterhaanden vil gøre sig gældende paa Troens og
Livsanskuelsens Omraade, saa at man vil nærme sig til et Punkt,
hvor Erfaringsvidenskaben bliver eneste Grundlag for de
Tanker, Mennesket har Brug for i Livets Kamp. At der
foregaar en saadan Tilnærmelse, er der ingen Tvivl om.
Spørgsmaalet er, hvor langt den vil kunne gaa. En aandfuld
amerikansk Filosof, William James, har sagt, at
Gennemførelsen af Sparsommelighedens Lov paa Troens Omraade
% --- p. 220 ---
\opage{220}vilde føre til Barbari. Derom kunne vi dog Intet vide, da den
Tid, hvor der vil kunne blive Tale om Sligt, ligger fjernt
endnu, om den nogensinde kommer; skulde den komme, vil
den kunne medbringe aandelige Livsbetingelser, som vi nu
ikke kunne gøre os nogen Forestilling om. Men foreløbigt staa
vi midt i Udviklingen. Erfaringen bringer os stedse nye Ting,
ikke mindst paa de Omraader, hvor Livsspørgsmaalene høre
hjemme, det biologiske, det psykologiske og det sociale
Omraade. Al Tro gælder jo dog Livets, særligt det aandelige
Livs Bestaaen, Bevarelsen og Udviklingen af de højeste
Værdier, vi har lært at kende. Og der maa stedse kæmpes
for Livet, ikke blot for det ydre, materielle Liv, men for at
hævde det Liv, der i vore højeste og bedste Øjeblikke staar
for os som det eneste Liv, der er værdt at leve. Og her
staa vi --- trods Alt hvad Erfaringsvidenskaben har bragt ---
stadigt overfor det Gaadefulde og det Uoverskuelige. Nye
Erfaringer fortrænge gamle, nye Synsmaader vise sig
nødvendige, og den Maade, paa hvilken vi have lagt os Tingene
til Rette, sættes stedse paa nye Prøver. Med Spænding maa
netop den, der tillægger Erfaringen en afgørende Betydning
for sin Tro, følge den fortsatte Udvikling af Livet uden om
os og i os:

\begin{verse}
Neu ist immer die Erfahrung, \\
Immer ist dem Herzen bang!
\end{verse}

Denne Betragtning fører os over til det sidste, det mest
afgørende Element: den personlige Trang og Higen, der er
utilfredsstillet ved alt Givet, fordi den gaar ud paa Noget,
som Erfaringen kun ufuldkomment lader komme til
Raadighed. Gennem Værdibegrebet hænger religiøs Tro sammen
med Villien. Vor Tro bestemmes ved, hvad der for os staar
som den højeste Værdi, der giver alt Andet Værdi, og ved det
Forhold, vi mene at finde mellem Virkelighedens Verden og
denne højeste Værdi. Al Værdi beror paa et Formaal, vor
Villie med større eller mindre Bevidsthed fastholder. Saa
længe der er et Spændingsforhold mellem vor højeste Værdi
og Erfaringsvirkeligheden, saa længe er der Plads for et
religiøst Problem. Det er ganske vist først og fremmest vor
% --- p. 221 ---
\opage{221}egen Opgave at kæmpe for det højeste Værdifulde --- for
dets Sejr i os og uden for os; men det religiøse Problem
stiller sig med det Spørgsmaal: er det os alene, der kæmpe
for det, eller er der Magter og Love i Tilværelsen, der virke
i Harmoni med det og i Retning af det? --- Det er
Reformationens største Fortjeneste, at den har hævdet dette som
Grundspørgsmaalet. Reformationen lod den gamle Dogmatik
staa, men lagde hele Vægten paa Personlighedens indre
Tilslutning. I sin Catechismus major spørger Luther: Qvid est
Deus? (Hvad er Gud?) Og hans Svar er: det, hvortil du har
fuld Tro og Tillid. Det er ene Hjertets Tillid, der gør
Noget til en Gud eller en Afgud (sola cordis fiducia deum
pariter atqve idolum facit et constituit); og den rette Gud
vil du have, naar din Tro og Tillid er ret og ægte. Luther
blev den Mand, han blev, fordi han af egen Erfaring kendte
den Angst, der føles i Kampen for Livets højeste Værdier,
ikke blot for, om de ville sejre uden om os, men først og
fremmest for, om de ville sejre i os selv. Der, hvor Villiens
Grænse ligger --- indadtil og udadtil --- der hævder Villien
sig endnu i Form af Tro. Hvad Skolastiken og det romerske
Kirkesystem ikke kunde lade komme til Udfoldelse, eller hvad
de for hastigt og paa for lette Vilkaar tilfredsstillede, ---
den Enkeltes indre Trang og Krav til sig selv og til Verden,
det stillede Reformatorerne i første Række. Det er deres
store, umistelige Fortjeneste.

Den personlige Trang til Hævdelse af Livsværdierne er
den Faktor, som tilsidst bliver afgørende for den Indflydelse,
Overlevering, Iagttagelse og Tænkning faa. Den træffer sit
Valg mellem Overleveringernes og Iagttagelsernes
Mangfoldighed, betoner nogle og stiller andre i Baggrunden, og den
hæmmer eller fremmer Tænkningens Konsekvens. Det er
det mest urolige Element i det Religiøse. Det er mod den,
at Opmærksomheden mere og mere vender sig, fordi der er
opstaaet en levende Fornemmelse af, at det er her, Slaget
tilsidst skal staa. De gamle Guder, som Slægten har vendt
sig fra paa forskellige Stadier af sin Udvikling, ere aldrig
blevne gendrevne. Gudedynastier uddø, naar der ikke mere
er Trang til dem. Derfor vender Behandlingen af det reli%
% --- p. 222 ---
\opage{222}giøse Spørgsmaal sig mere og mere fra Spekulationens til
Psykologiens Omraade.

Men Menneskets Tendens til at dogmatisere følger let
med ogsaa her. Religionspsykologien er endnu meget lidet
dyrket. Først langsomt er den Overbevisning bleven
almindelig, som Luther --- i det mindste et enkelt Sted --- saa klart
har udtalt. Og efter at den var begyndt at fæstne sig, mente
man at kunne antage, at dette afgørende Element, den
personlige Trang og Higen, den inderste Vilje, maatte være den
samme hos Alle og til alle Tider. Psykologisk Iagttagelse
bekræfter ikke denne Antagelse. Der er et
Vekselvirkningsforhold mellem denne personlige Faktor og de tre andre
Faktorer: Overlevering, Iagttagelse og Tænkning bestemme
Trangens Art og Retning, lige saa vel som Trangen igen
faar Indflydelse paa deres Udvikling. Saaledes er det nu en
Gang paa det aandelige Omraade; det er et Væv, hvor en
enkelt Traad ikke kan drages ud, uden at det zitrer i det
hele. Det aandelige Livs Elementer indgaa en saa inderlig
Forening, at det er umuligt at betragte noget enkelt af dem
som uafhængigt og konstant. Det var derfor et uberettiget
Haab, naar man mente, ved at bygge paa den personlige
Trang som en ensartet og konstant Faktor at kunne opføre
en ny dogmatisk Bygning. Problemet er langt mere indviklet
end som saa. Vi staa her ved nogle af de zarteste og mest
sammensatte psykologiske Spørgsmaal.

Man kan sige, at vor Tids Filosofi bevæger sig
henimod at blive en empirisk Personlighedsfilosofi. Ved at
arbejde paa det psykologiske, erkendelsesteoretiske og etiske
Omraade, --- de Omraader, der i den sidste Menneskealder
have lagt Beslag paa næsten al filosofisk Interesse, stræber
den at vinde Indblik i det aandelige Livs Betingelser og
Former, dets Tarv og dets Trang. Man beskylder ofte den
realistiske eller positivistiske Retning i Filosofien for at
fornegte Personlighedens Ret. Jeg tror omvendt, at ingen
Retning anerkender denne Ret med mere Alvor og
Oprigtighed end Positivismen. Medens den romantiske Teologi og
Spekulation meget hurtigt fik sit Personlighedsbegreb færdigt
og straks begyndte at bygge sine ofte saa skrøbelige Tanke%
% --- p. 223 ---
\opage{223}bygninger derpaa, ser Positivismen i Personlighedsbegrebet et
Problem, der maa belyses og behandles fra saare mange Sider, ---
ad alle de mange Veje, den moderne Psykologi benytter for at
vinde Forstaaelse af det menneskelige Sjæleliv. Og medens
den romantiske Personlighedslære nærede Foragt eller Frygt
(snart mere det Ene, snart mere det Andet) for den lovmæssige
Natursammenhæng og for de bestemte Betingelser, der bære
Alt i Livet, det Højeste som det Laveste, er det den frejdige
Tro, der ligger til Grund for den moderne psykologiske Forsken,
at Aandslivets Værdi ikke forringes ved, at det er knyttet
til bestemte Betingelser og Love. I hin romantiske Fordom
rører der sig en hemmelig Vantro. I det personlige Liv er
der et uudtømmeligt Omraade for Forsken, og intetsteds er
dogmatisk Afslutning betænkeligere. Heldigvis ville Problemerne
stedse rejse sig paa ny.

Hvor vidt det nogensinde vil blive muligt at samle alle
de Elementer, vi her have omtalt, til en stor og omfattende
Verdensanskuelse, der forener Overleveringens Ærværdighed
med Tankens Konsekvens, den Rigdom og Lovmæssighed,
Iagttagelserne frembyde, med Hjertets inderste Trang, det er
et Spørgsmaal, som det er umuligt at besvare. Den eneste
Filosofi, den katolske Kirke anerkender, Skolastiken, som den
i det 13. Aarhundrede formedes af Thomas Aqvinas, og som
den af Dante benyttedes til Ramme og Tankeelement i hans
store Digtning, --- den kunde for sin Tid yde noget Saadant.
Maaske har det kun været muligt paa dette bestemte Punkt
i Historien. Maaske vil en Gang mange Traade, der for os
synes at løbe i modsat Retning, samles til et vidunderligt Hele,
om hvis Natur vi nu ikke kan have nogen Anelse. Hvilken
af disse to Muligheder vi holde os til, saa vil det foreløbigt
ikke kunne ændre Noget i vor Arbejdsmaade. Vi ville søge at
holde vor Horisont aaben, men tillige tage vore Opgaver op
der, hvor der virkeligt er et Arbejde at gøre. Der er i den
nyeste Tid en idealistisk Retning at spore i den filosofiske
Verden. Den er berettiget, for saa vidt den bekæmper den
Ignoreren af afgørende Problemer, som Reaktionen mod den
romantiske Teologi og Spekulation ofte har ført til. Sin
positive Berettigelse vil den kun kunne godtgøre, naar den
% --- p. 224 ---
\opage{224}kan paavise nye Midler til at behandle de gamle Problemer.
Men den vil bidrage til at holde Horisonten aaben, saa at de
nye Tvivl lige saa lidt formaa at lukke den for os, som de
gamle Dogmer have formaaet det. Tanken kan undvære sine
Løsninger, men den kan ikke undvære sine Problemer.

For den, hvis Livsstilling fører med sig, at han ikke blot
arbejder med Problemerne i det stille Lønkammer, men ogsaa
skal føre de Unge paa Vej under deres Stræben efter at finde
sig til Rette i Tankens Verden, for ham vil denne Betragtning
have sin særlige Betydning. Han vil søge i al sin Færd at
forbinde Sansen for den stræbende Personlighed med
Overbevisningen om, at Personligheden skal lutres og forædles gennem
Arbejdet paa at vinde Sandheden og bøje sig for dens Krav.
Han vil ønske at kunne gøre sit Arbejde i denne Aand, og
han vil ønske, at denne Aand maa raade i den Højskole, det
er hans Ære at tilhøre.

\streg{}

% --- p. 225 ---
\opage{225}
\begin{center}

\essay[KAMPEN MELLEM GAMLE OG NYE TANKER]{KAMPEN MELLEM GAMLE OG NYE TANKER.}

\emph{ET TILBAGEBLIK}\footnote{Foredrag, holdt i det norske Studentersamfund i Kristiania den 14.
April 1895.}.\end{center}

\streg{}

Homer kalder Mennesket et Væsen, der ser frem og
tilbage. De to Ting hænge nøje sammen, ti for at kunne
se frem, hvad der er det Vigtigste, maa man kunne se
tilbage, og de Tanker og Stemninger, med hvilke vi se
Fremtiden i Møde og berede os til den, ville nødvendigvis --- hvad
enten vi ere os det bevidst eller ikke --- være bestemte ved
den Maade, paa hvilken vi opfatte Fortiden. Og imod et
Aarhundredes Slutning er der særlig Anledning til et
Tilbageblik. Ganske vist er Inddelingen i Aarhundreder en
rent udvortes og borgerlig Indretning, og hverken den indre
eller den ydre Verdens Begivenheder passe ind i den Ramme,
denne Inddeling giver. Jeg vil ogsaa her forsøge at paavise,
at aandeligt talt begynder vort Aarhundrede, før man i
Almanaken begyndte at skrive 18… Maaske vil man en
Gang sige det Samme om det næste Aarhundrede; vi, der
leve midt i Tiden, kunne ikke vide det, da vi endnu ikke kunne
skelne det, der er Begyndelsen til noget Nyt, fra hvad der
hører det Gamle til; vi maa naturligvis haabe, at der findes
saadanne Begyndelser i vor Tid. Der vil dog ikke være
noget til Hinder for at benytte den ydre Anledning,
Aarhundredets borgerlige Afslutning giver, til at se tilbage paa
den aandelige Udviklings Gang siden der sidst skiftedes Aar%
% --- p. 226 ---
\opage{226}hundrede. En mere væsentlig Anledning har jeg for mit
Vedkommende faaet til et saadant Tilbageblik ved
Afslutningen af et Arbejde over Filosofiens Historie i det sidste
Aarhundrede. Det er paa Resultater af dette Arbejde, jeg for
en væsentlig Del støtter mig i de Anskuelser, jeg her vil
gøre gældende. Men før jeg, støttet til Historien, gaar over
til at gøre Rede for, hvorledes vi nu ere stillede med Hensyn
til nogle af de vigtigste Problemer paa det aandelige Livs
Omraade, maa jeg forudskikke nogle Bemærkninger om
Synspunkt og Maalestok for det aandelige Livs Udvikling.

\section*{I.}

Naar man i det Hele taget anerkender, at der foregaar
en aandelig Udvikling, ikke blot hos den Enkelte, men ogsaa
i Slægten, saa at lige saa lidt i den indre som i den ydre
Verden Alt bliver som det var, saa opstaar naturligt det
Spørgsmaal, om der ved denne Udvikling gøres et Fremskridt
eller ikke, --- om der tabes eller vindes ved de Forandringer,
der ske. I Begrebet Udvikling ligger ikke nødvendigvis,
at der bliver noget mere Værdifuldt til end der var før, ikke
heller, at de vundne Værdier bevares. Der kan tænkes en
Udvikling --- i Betydning af en lovbestemt Række af
Forandringer --- som fører til en Tilbagegang ikke blot i Værdi,
men ogsaa i Energi.

Paa den ydre Naturs Omraade har man som bekendt i
de sidste halvhundrede Aar opstillet og gennemført Loven om
Energiens Bestaaen. Fysikerne ere overbeviste om, at der i
den materielle Verden ingen Energi opstaar og ingen Energi
forgaar, men at Energien ved enhver Forandring og
Udvikling kun omsættes i nye Former. Men heller ikke her er
Energiens Bestaaen Et med Værdiens Bestaaen, ti man har
jo drøftet den Mulighed, at der i vort Verdenssystem kunde
indtræde en saadan Fordeling af Varmen, at al Bevægelse
og alt Liv vilde høre op, saa at det, der fra vort Synspunkt
(og andet Synspunkt have vi nu en Gang ikke) giver Verden
Værdi --- det, at den kan være Hjemsted for Liv og Aand
--- dermed vilde være faldet bort.

% --- p. 227 ---
\opage{227}Hvorledes ere vi nu stillede paa det aandelige Omraade?
Jeg skal her ikke komme ind paa det meget omtvistede
Spørgsmaal om det aandelige Livs Forhold til Loven om den
fysiske Energis Bestaaen. Hvorledes man end stiller sig til
dette Spørgsmaal, saa vil der, saa snart man indrømmer, at
der foregaar en Udvikling paa det aandelige Omraade, være
en naturlig Trang til at spørge, om Energien og Værdien
bevares under denne Udvikling eller ikke. Men da
Videnskaben om den aandelige Verden ikke, saaledes som
Videnskaben om den materielle Verden, er en eksakt Videnskab,
kunne vi ikke vente at komme til noget sikkert Resultat.
Allerede med Hensyn til Paavisningen af den Energis
Bevarelse, der kundgør sig i det aandelige Livs Funktioner,
ere vi meget uheldigt stillede. Vi kunne ikke paavise, at
der stadigt er samme Mulighed for aandeligt Liv tilstede i
Verden. Og selv om vi begrænse Betragtningen til en enkelt
Periode, mangle vi Midler til at overbevise os om, at der i
Løbet af Udviklingen ikke sættes Energi til eller maaske
opstaar ny Energi. Det Spørgsmaal, om der i Løbet af et eller
flere Aarhundreder er sket en Svækkelse eller en Bevarelse
eller en Forøgelse af Energien, kunne vi ikke med Sikkerhed
besvare. Selv om vi nu tillidsfuldt antage, at den Energi,
der staar til det aandelige Livs Raadighed, ikke forminskes,
saa er dermed dog ikke det Spørgsmaal besvaret, om denne
Energi finder Anvendelse paa lige saa værdifuld Maade som
tidligere. Sygelig Udvikling kan lægge Beslag paa lige saa
megen Energi som sund Udvikling, en Opløsningsproces kan
forbruge ikke mindre Energi end en Tilblivelsesproces;
Størrelserne kunne være de samme, men Fortegnene forskellige.
Selve Bevarelsen af den aandelige Energi vil foreløbigt
være af Værdi, da Livets Værdi ogsaa beror paa dets Kraft;
men det er ikke den eneste Maalestok. Foruden det
aandelige Livs Kraft spørge vi ogsaa om Rigdommen og Fylden
af dets Indhold. Jo større og rigere Nuancering, det
frembyder, jo flere Forskelligheder og Modsætninger, det formaar
at omfatte, des større Værdi tillægge vi det. Men paa den
anden Side: kun hvis Aandslivets Indholdsrigdom og
Forskelligartethed er forbunden med Evne til at virke med samlet
% --- p. 228 ---
\opage{228}Kraft, er Livet sundt. Meget let kan der blive et omvendt
Forhold mellem Kraft og Fylde. I vor Vurdering indskrænke
vi ikke Betragtningen til det enkelte Individ for sig alene.
Vi betragte kun Aandslivet som sundt, naar den enkelte
Personlighed staar i nøje Sammenhæng --- som givende og
modtagende --- til Samfundet, Slægten. Vi vurdere en Periodes
Aandsliv efter den Kraft og Selvstændighed, den enkelte
Personlighed kan besidde, uden at det aandelige Fællesskab
med andre Personligheder derfor brydes. Det er en ulykkelig
Tid, ikke blot naar der mangler Kraft eller Fylde hos det
enkelte Individ, men ogsaa naar der mangler Fællesskab,
Forstaaelse og harmonisk Vekselvirkning mellem Individerne
indbyrdes. Naar de Tanker og Følelser, der bære Livet, blive
helt forskellige hos de forskellige Individer, træder
Ligegyldighed og Foragt i Stedet for Sympati; der mangler Glæden
ved at give og Glæden ved at modtage. Netop de mest
fremragende Individer, hvis Aandsliv har Betydning for store
Kredse, lide herunder; \emph{de} føle det største Savn, fordi de
have Mulighed for at give Mere og modtage Mere end Andre.
Man slipper for mange Problemer og mange Sorger, naar
man indsnevrer sin Horizont.

Den nyere Psykologi har paa forskellig Maade paavist
Muligheden af, at den Energi, det aandelige Liv har besiddet
paa tidligere Trin, kan vedblive at bestaa under nye Former
paa senere Trin i Udviklingen. Herhen hører først og fremmest
det af Spinoza og de engelske Psykologer paaviste
Motivforskydningsfænomen. Hvad der først kun har Interesse og
vækker Stræben, fordi det er Middel til at opnaa et Formaal,
kan blive Genstand for selvstændig Interesse og efterhaanden
selv komme til at staa som Formaal. Saaledes kan en Tanke,
vi først bøje os for af Lydighed mod en Autoritet, en Tanke,
hvis Værdi for os derfor først ene beror paa vor Frygt eller
vor Tillid over for Autoriteten, senere ved sin egen Magt,
sin egen Begrundelse, staa fast for os, selv om
Autoritetsforholdet falder bort. Der kan i Samfundet udvikles
Institutioner, som skulle fyldestgøre bestemte Formaal, men som senere
hen opretholdes, efter at de oprindelige Formaal ere faldne
bort, fordi det viser sig, at de kunne tjene nye Formaal. Saa%
% --- p. 229 ---
\opage{229}danne Motivforskydninger foregaa uafbrudt hos det enkelte
Individ og i Historien, og uden dem finder man ikke
Sammenhæng i Udviklingen. En anden Form, i hvilken Kontinuiteten
gennem alle Forandringer kan bevares, er hvad man med
Benyttelse af en nærliggende Analogi har kaldt psykisk
Kemi. Sjælelige Fænomener kunne indgaa en saadan
Forbindelse, at de ikke længere gøre sig gældende hvert for
sig i deres Ejendommelighed, men lige som smelte sammen
til et nyt Fænomen, der er deres fælles Produkt, ligesom
Vandet opstaar ved Forbindelse af Ilt og Brint. Glæde og
Sorg kunne knytte sig til samme Begivenhed, og Tanken om
denne Begivenhed kan da tilsidst være forbunden med en
Grundstemning, Vemodet, der hverken lader Glæden for sig
eller Sorgen for sig træde frem, men udtrykker dem begge
paa en Gang. Her har fundet en Udligning Sted, hvorved
de tidligere sjælelige Erfaringer bevares under andre Former.
Forskydning og Udligning ere de to vigtigste Maader, paa
hvilke den aandelige Energi vil kunne bevares. Men dermed
er det ikke sagt, at Værdien er den samme i de nye Former
som i de gamle. Ikke blot det sunde, ogsaa det syge
Bevidsthedsliv frembyder Forskydninger og Udligninger.
Ogsaa i Læren om Sindssygdomme taler man om psykiske
Ækvivalenter, saaledes især, hvor Tendensen til epileptiske Anfald
faar Luft paa anden Maade, i Form af pludseligt Raseri eller
overraskende Stemningsændring. Den ene og samme Sum af
Energi, der er til Stede, kan faa Afløb paa mange Maader.
Men selv om det kan paavises, at den bevares, er det ikke
sagt, at den bevares uden Forandring af Værdi. Ja,
Forskydninger og Udligninger gaa stadigt for sig, uden Bevarelse
hverken af Kraften eller af Værdien.

Et Punkt er her især af Betydning. Paa et vist Trin af den
aandelige Udvikling, Slægtens saa vel som det enkelte Individs,
er Instinktets og Autoritetens Enehersken forbi. Synskredsen
udvides; der viser sig en større Mængde af Muligheder og
Retninger baade for Tro og for Handlen. Hidtil havde Verden
været begrænset, og der havde kun frembudt sig en eneste
Maade at opfatte Livet paa og en eneste Maade at leve det paa.
Nu slaas Dørene op til den uendelige Verden med dens mange
% --- p. 230 ---
\opage{230}Synspunkter og Livsretninger. Kan nu Kraften, Trygheden og
Inderligheden bevares, naar Indholdet bliver saa meget des
rigere? Kan det aandelige Fællesskab bevares, naar saa
mange Retninger i Tro og Liv faa deres mere eller mindre
frie Spil? --- Det var et Spørgsmaal, som først Renaissancens
og senere Oplysningens og Revolutionens Aarhundrede stillede.
Det er et Grundproblem, der under forskellige Former strækker
sig hen igennem den nyere Tids Tænkning lige fra
Renaissancens Tid til vore Dage. Renaissancen selv stillede ikke
Problemet; den var saa optagen af det Nye og Store, den
opdagede, at den ikke følte Problemets Brod. Oplysningens
Tid var ligeledes saa optagen af sin Kritik og sin Fornuft,
af at skyde det Gamle til Side og føle Glæden ved den
kritiserende og prøvende Fornufts Brug, at den ikke tænkte
paa, om der ikke ved Bruddet med Instinkt og Autoritet var
gaaet Noget tabt, de nye Former for aandeligt Liv ikke vilde
kunne erstatte. Selv om Kontinuiteten i Udviklingen kunde
paavises trods alle Renaissancer og Revolutioner (ligesom
Geologien mener at kunne paavise Kontinuiteten i
Jordudviklingen trods alle Jordrevolutioner), og selv om det kunde
paavises, at den samme Energi, der ytrede sig i Instinktets
og Autoritetstroens koncentrerede Form, nu ytrer sig i mere
spredt, fordelt og nuanceret Form, saa er det jo ikke sagt,
at disse forskellige Ytringsformer ere lige meget værd. Den
paa mange Punkter fordelte, i mange nuancerede Former
fremtrædende Kraft har maaske ikke paa noget Punkt den
Mulighed for at optræde samlet, den Koncentrationens
Sikkerhed og Vælde, som den havde under den gamle Tingenes Orden.
Bliver Livet ikke mattere under Udviklingens Gang, ikke
fordi der forsvinder Kraft, men fordi den fordeles og
nuanceres, --- ligesom Fysikere tro, at vort Verdenssystem en
Gang vil stivne, fordi Energien er fordelt, saa de store Rytmer
i Verdensprocessen ikke mere bliver mulige?

Det er Problemet. Og naar vi nu tage vor
Tidsbetragtning op, finder jeg det ejendommeligt for vort Aarhundrede,
at dette Problem er blevet følt og stillet. Og \emph{vort}
Aarhundredes Begyndelse maa sættes der, hvor Problemet først
optræder. Men saa maa vi forlade Almanakens Grænser; ti
% --- p. 231 ---
\opage{231}Problemet stilledes først af \emph{Rousseau, Lessing} og \emph{Kant}. De
vare de Første, som paa en Gang følte hele Berettigelsen af
den moderne Kritik og af Horizontens Udvidelse, og tillige
saa det som et stort Spørgsmaal, hvorledes den Revolution i
det aandelige Liv, der stod for Døren, skulde kunne foregaa
saaledes, at Livets Sundhed og Inderlighed ikke led derunder.
De hørte Oplysningens Tid til, samtidigt med at de saa dens
Grænser og indsaa Nødvendigheden af at bevare den inderste
Kærne i den gamle Tro og den gamle Ordning, selv om
Skallerne maatte kastes bort. Med disse tre Mænd begynder
det nittende Aarhundrede. De stillede dette dets egentlige
Problem.

\section*{II.}

Den filosofiske Tænknings Historie i det forløbne
Aarhundrede frembyder Træk, der godtgøre, at Problemet stadigt
er blevet følt af en hel Række af fremragende Aander, saa
at det med Rette kan kaldes vor Tids Problem.

Filosofien kan --- især fra et historisk Synspunkt ---
betragtes paa to forskellige Maader. --- Filosofiske Ideer
kunne for det Første betragtes som \emph{Symptomer} paa, hvad der
rører sig i Tiden. Hvad der hos de Fleste kun bevæger
sig som dunkel Følelse og uvilkaarlig Trang, det former sig
hos Filosofen til en Ide, ligesom det hos Digteren former
sig til et Billede. Driften til at forme i klare Tanker, hvad
der bevæger sig i Sindet, er ejendommelig for filosofiske
Naturer. Vi se derfor ofte i Filosofiens Historie Ideer komme
frem, som vedkommende Tænker ikke selv formaar at give
fuld Begrundelse af; de stikke for første Gang Hovedet ind
i Tilværelsen, dukke fra Erfaringernes Kaos og de
uvilkaarlige Tendensers dunkle Verden op i Tænkningens Dagslys, og
deres første Form er derfor ofte ufuldkommen og uden fast
Begrundelse. Først senere kan det forstaas, hvorledes de
hænge sammen med Tidens og Udviklingens hele Vilkaar. ---
Paa den anden Side er Filosofien Forsøg paa \emph{videnskabelig
Behandling} af uundgaaelige, bevidst stillede Problemer. Fra
dette Synspunkt falder Vurderingen af en Filosofi ofte helt
% --- p. 232 ---
\opage{232}anderledes ud end fra det første; men enhver
betydningsfuld filosofisk Bestræbelse vil frembyde Elementer af begge
Arter. Og der er desuden aabenbart en indre Forbindelse
mellem Filosofien som Symptom og Filosofien som
Videnskab. Fælles for begge Sider er, at Filosofien er den
menneskelige Aands Stræben efter bevidst Klarhed angaaende
sit Væsen og sine Livsvilkaar; Filosofien er \emph{videnskabelig
Selverkendelse}.

I denne Sammenhæng vil jeg dog især benytte det
første Synspunkt. Jeg finder det et mærkeligt Tidens Tegn,
at det af Rousseau, Lessing og Kant rejste Problem strækker
sig under forskellige Former gennem hele Aarhundredets
Tænkning. Atter og atter kommer den Tanke frem: Vi have
forladt Instinktets og Autoritetens trygge og faste Grund;
Kritiken og Analysen har gjort sit Værk, et Værk, der ikke
kunde holdes borte og heller ikke igen kan gøres ugjort;
men af Kritik og Analyse kunne vi ikke leve, da Livet stedse
kræver positivt Indhold og fast Grund under Fødderne; der
maa da naas frem til et nyt Stade, der forener det førstes
positive Kraft med det andets kritiske Alsidighed og
Bevægelighed! Hvis dette ikke er muligt, saa er Livet i Opløsning.
Vi kunne ganske vist endnu en Tid lang tære paa
Eftervirkningerne fra den instinktmæssige og troende Barndomstid;
men denne Kapital vil tilsidst blive opbrugt, dersom ikke
nye positive Former udvikle sig. For at bruge et Billede,
som Renan har brugt i sine »Barndomserindringer«: Efter
et Sagn fra Bretagne blev der en Gang ved en Baads
Kæntring nedsænket en Kirkeklokke i Havet. I stille Vejr ---
fortælles der --- kan man endnu høre dens Klemten fra
Havbunden. Saaledes høre vi endnu Tonerne fra den Tid,
der var forenet i fælles og kraftig Tro og i naiv Tillid til
Livet; men vil Klangen ikke en Gang dø hen? ville vi ikke
fjerne os saa langt fra dens Sted, at dens Toner ikke kunne
høres?

Det var \emph{Rousseau}, der først for Alvor rejste Problemet,
ikke blot i sine bekendte Paradokser, men i alle sine Skrifter,
i sin inderste Tankegang. Han gennemzitredes af det, blev
et Offer for det. Han blev hjemløs, fordi han hverken delte
% --- p. 233 ---
\opage{233}de Konservatives Glæde over det Bestaaende eller
Encyklopædisternes Glæde over den nye Oplysning. En lignende
Stilling indtog \emph{Lessing} over for Ortodokse og Rationalister
i Tyskland, og \emph{Kant} over for Dogmatismen og Skepticismen.
Begge ere de --- især Kant --- paavirkede af Rousseau.
Ingen af disse Mænd vilde tilbage til det Gamle, saa lidt
de vare tilfredse med det Nye. De ventede --- hver paa sin
Vis --- det tredie Rige. Kant, den største Tænker af de
tre, har tydeligst formuleret sit Forhold til de to Retninger,
fra hvilke der var at tage Afstand. Han forkastede
ligesom Skeptikerne de dogmatiske Systemer, men han gjorde
det ikke blot gennem negativ Kritik. Han søgte at forstaa,
hvorledes Systemerne kunde blive til, forstaa den aandelige
Trang og Evne, som fik Luft i dem, og som kunde blive ved
at leve, selv om den fik Luft paa anden Maade. I denne
Tilbagegaaen fra de færdige Systemer til deres første Motiver
og i denne Undersøgelse af, hvor vidt disse Motiver kunne
vedblive at virke under andre Forhold og paa andre Maader,
finder jeg Kants største Gerning. Derved blev paa en Gang
Kontinuiteten med Fortiden reddet, og Udsigten til Fremtiden
holdt aaben. Kant har derved stillet hele
Aandsvidenskabens Program. Vi lære Menneskeaanden at kende gennem
dens Værker. Men de \emph{færdige Værker}, de fuldt udførte
Trosformer og Tankesystemer blive let en død Skat uden
Forbindelse med vort eget Liv og derfor ubrugelige for
dette. Derfor maa vi trænge ind til de \emph{Kræfter, der have
frembragt Værkerne,} for at se, om de endnu staa til vor
Raadighed, og for at se, paa hvilke Vilkaar de nu maa
virke.

I det Aarhundrede, der er gaaet efter Rousseau, Lessing
og Kant, kan der paavises to Hovedretninger i den filosofiske
Tænken. Den ene kalder jeg \emph{Romantiken,} den anden
\emph{Positivismen}. Den ene søger Problemernes Løsning ad
Spekulationens og Kontemplationens Vej, den anden gennem
Konstatering af Kendsgerninger og Refleksioner over
Kendsgerninger. De staa skarpt imod hinanden, da de tage deres
Udgangspunkt ved de modsatte Grænser for Erkendelsen.
Den ene vil stige ovenfra ned, den anden nedefra op. Men
% --- p. 234 ---
\opage{234}det skulde da gaa underligt til, om de slet ikke mødtes. Og
en nærmere Undersøgelse viser da ogsaa, at begge
Retninger, hver paa sin Maade, følge Kants Program for
Aandsvidenskaben. Karakteristisk for begge Retninger er en dyb
Forstaaelse af Middelalderen og i det Hele af tidligere Tiders
Aandsliv, en Utilfredsstillethed, maaske endog en Uvillie ved
den blotte negative Kritik, og Udblik til en ny Tid, da
positiv Fylde og kritisk Refleksion, fælles Grundsyn og
individuel Frihedstrang ikke længere skulle staa fjendtligt
over for hinanden. Gennemgaaende opfatte de da ogsaa
den nærværende Periode som en Overgangstid, en
Brydningens Tid, som de ufærdige Begyndelsers Tid. Dog have
--- mærkværdigt nok --- Romantikerne større Tilbøjelighed
til at anerkende den givne Tingenes Orden end
Positivisterne.

I \emph{Fichte’s Grundzüge des gegenwärtigen Zeitalters} (1806)
finder Læren om de tre Stadier maaske sit klareste og mest
energiske Udtryk. Efter Instinktets og Autoritetens Periode
kommer en Periode, der betegnes som den tomme Friheds
Periode, hvis Karaktermærke er »Auf- und Ausklärung«.
Glad ved Oplysning og Kritik vender man sig med Stolthed
fra Fortidens Mørke. Negation og Egoisme raade. Men
Tænkeren ser hen til en ny Tid, hvor en grundigere
Videnskab og en højere Kunst vil føre til en ny Livsanskuelse,
der lader alle aandelige Fornødenheder, Hengivelsen saa vel
som Refleksionen, komme til deres Ret. Omtrent samtidigt
betegnede \emph{Saint-Simon} i Frankrig det 18. Aarhundrede som
en kritisk, revolutionær Periode, der nu skulde afløses af en
organisk Periode. Hans Disciple generaliserede
Modsætningen mellem organiske og kritiske Perioder og fandt den
komme igen paa forskellige Punkter af den historiske
Udvikling. En organisk Periode er karakteriseret ved Fasthed og
Sluttethed i Tro og Tanke, og ved aandeligt Fællesskab. I
den kritiske Periode raader derimod Tvivlen og Antagonismen
paa Troens og Associationens Bekostning. Medens Fichte er
paavirket at Kant, og gennem ham af Rousseau, synes % PRINTED AS IS: „at“
Saint-Simon nærmest at være kommen til sin Opfattelse gennem
Refleksioner over Industriens og Videnskabens Udviklings%
% --- p. 235 ---
\opage{235}gang. Paa den videre Udvikling af hans Ideer hos hans
Disciple fik Lessings Skrift om Menneskehedens Opdragelse
en ikke ringe Indflydelse.

Det vilde blive for vidtløftigt at fremdrage alle de
forskellige Former, hvorunder den samme store Tanke
fremtræder hos de romantiske og positive Filosofer. --- \emph{Hegels}
Tredelinger og hele hans dialektiske Metode ere egentligt
opstaaede ved Skematisering af de tre Stadiers Lov. Til
Grund ligge --- som det ses af Optegnelser fra Hegels
Ungdom --- Studier over Kultur- og Religionshistorie. Som
lykkelige Tider priser Hegel de Perioder, hvor der er Tro
paa Tilværelsens Kærne, og hvor fælles Forstaaelse af Livet
forener Alle. Subjektiv Kritik isolerer da ikke den Enkelte
fra det Hele og gør ham ikke splidagtig med sig selv. Men
det var Hegels Overbevisning, at det Substantielle i den
gamle Tro og i den gamle Tingenes Orden kunde bevares
gennem Refleksionens Skærsild, og i sin egen Filosofi mente
han at have paavist det. \emph{Schleiermacher} var mere kritisk,
men mente dog at kunne paavise de samme
Følelseserfaringer, af hvilke det Væsentlige i den gamle Tro var
udsprungen. Ingen af de to Tænkere vilde være ved, at de
omdannede Aandslivets gamle Former efter deres Behov. Deri
bestaar netop det Romantiske hos dem, at de umiddelbart
mene at leve Fortidens Liv om igen. --- Inden for
Positivismen er først og fremmest at nævne \emph{Auguste Comte’s}
Formulering af de tre Stadiers Lov. Den aandelige Udvikling
foregaar herefter gennem tre Stadier, det teologiske, det
metafysiske og det positive. Det metafysiske Stadium er
rent negativt og kritisk; det er en Opløsningsperiode, svarende
til hvad Fichte kaldte den tomme Friheds Periode, til hvad
Saint-Simon kaldte den kritiske Periode, og til hvad Hegel
kaldte de ulykkelige Tider, den kritiske Subjektivitets Perioder.
Der er efter Comte et aandeligt Slægtskab mellem den
positive og den teologiske Periode; de have hver sin
Helhedsopfattelse, og de forbinde den Enkelte med Slægten paa
inderlig, organisk Vis. Hos \emph{Stuart Mill} og \emph{Herbert Spencer}
finde vi lignende Ideer, ligeledes hos \emph{Schelling} og hos \emph{Eugen
Dühring}. \emph{Carlyle} er opfyldt af den Tanke, at han lever i
% --- p. 236 ---
\opage{236}en af de mest ulykkelige Tider i Historien, i Kritikens og
Mekanismens Periode. Men han er ikke mindre overbevist
om, at »de gamle evige Magter leve bestandigt«, selv om
ethvert romantisk-spekulativt Forsøg paa at rehabilitere de
gamle Former er paa Strudsens Vis at stikke Hovedet i en
Busk for ikke at se Tingene som de ere. Han er endog
overbevist om, at Fuglen Føniks ikke dør paa Baalet, for saa
paa ny at blive til, men at den stadigt er levende, selv om
vi ikke høre dens Vingeslag. Han levede paa de tre Stadiers
Lov, i Troen paa det tredie Rige.

Romantikerne ere mere tilbøjelige til at tro paa en
Identitet af de gamle og de nye Tanker, Positivisterne mere
til at opfatte Forholdet mellem det Gamle og det Nye som
Ækvivalens. Derfor ere Romantikerne mere konservative end
Positivisterne. Men fælles er Overbevisningen om
Udviklingens Kontinuitet. Medens Oplysningens og Revolutionens Tid
saa ned paa Middelalderen som lutter Mørke og Barbari,
anerkende baade Positivister og Romantikere dens store
Betydning som organisk Periode. Og medens Reaktionen i det
19. Aarhundrede saa tilbage til Oplysningens og Revolutionens
Tid som lutter Forstandshovmod, Vildskab og Negation, saa
anerkende baade Romantikere og Positivister saadanne
kritiske Perioders Nødvendighed.

\section*{III.}

Men --- kunde man sige --- denne Enighed betyder ikke
meget, naar de to Retninger ere uenige om, hvorledes det
tredie Stadium, det tredie Rige skal være. Her træder en
afgørende Modsætning frem. Og desuden: Er Læren om de
tre Stadier eller om organiske og kritiske Perioder i rytmisk
Skiften anerkendt af Alle? Er der ikke dem, som i dette
Punkt hverken ere enige med Romantikerne eller Positivisterne,
men som enten mene, at der slet intet tredje Rige \emph{behøver}
at komme, eller at der --- selv om det behøves --- intet
saadant \emph{kan} komme?

Jeg skal søge at besvare disse Indvendinger, og tager
den sidste først.

% --- p. 237 ---
\opage{237}Skønt en sammenlignende Betragtning af de to
Hovedretninger i det sidste Aarhundrede, der kunne betegnes som
Romantiken og Positivismen, finder større Overensstemmelse,
end man paa Forhaand skulde vente, med Hensyn til et saa
væsentligt Problem som det, vi her betragte, saa er der
dog et andet Træk, der ikke kan Andet end paatrænge sig
den, der betragter Aandsudviklingens Gang i vort
Aarhundrede. Det er den Maade, paa hvilken Modsætningerne
skærpes. Ikke tilfreds med den Anerkendelse og Sympati,
baade Romantik og Positivisme lægge for Dagen over for
den gamle Tro og den gamle Tingenes Orden, udfolder
Autoritetstroen sig i Løbet af Aarhundredet med stedse større
Konsekvens og i stedse mere ekstreme Former.
Mellemstandspunkter knuses. Inden for den katolske Kirke
bortfalder Gallikanismen, Pavens Ufejlbarhed proklameres, og al
anden Filosoferen end den, der følger i Thomas Aqvinas’
Spor, forkastes. Og ligesom den katolske Teologi gaar
tilbage til det 13. Aarhundrede, gaar den protestantiske
Teologi tilbage til det 17. Aarhundrede. Fra Standpunkter som
disse føler man ikke Trang til noget tredie Rige. Det ene
Fornødne antages at være traadt ind i Historien en Gang for
alle, og al Tidens Uro og Nød udledes af, at man ikke
holder sig til det. Med Medlidenhed ser man ned til
Filosoferne, der plage sig med at finde Vilkaarene for det
aandelige Livs Bestaaen, ligesom de, der staa paa Bredden, se
paa de Skibbrudnes Kamp med Bølgerne. Paa den anden
Side er der dem, der mene, at Opløsningens Værk, som
Renaissancen og Revolutionen have indledet, ikke kan stanses. De
glæde sig maaske endog ved Opløsningen og anse
Opløsningsprodukterne for Udviklingens højeste Produkt. Kritiken og
Negationen have udbrændt Evnen og Trangen hos dem til
positivt Liv, til Hengivelse og til aandeligt Fællesskab. De
føle ikke Trang til et tredie Rige, og selv om de føle den,
saa anse disse »trætte Mænd« det for umuligt, at det kan
komme; langt snarere end at hige fremad kaste de sig tilbage
i den gamle Tros Arme.

Men over for dette unægtelige Faktum, Modsætningernes
Skærpelse paa det aandelige Omraade, maa vi henvise til
% --- p. 238 ---
\opage{238}en psykologisk Lov, som allerede ovenfor blev omtalt:
Motivforskydningens Lov. Selv om gamle Former og
Traditioner endnu holdes oppe iblandt os, saa vil dog en
nærmere Betragtning vise, at der er foregaaet væsentlige
Forandringer i det subjektive Forhold, i hvilket Individerne
staa til dem. Det er i hvert Tilfælde et stort Spørgsmaal
--- et nyt Problem, eller snarere en særegen Form for det
Problem, vi her behandle --- om vor Tids Kristendom er
den samme som Urkristendommen, og hvorledes de to
forholde sig til hinanden. Tænkere, der komme fra saa
forskellige Sider som Schopenhauer, Feuerbach og
Kierkegaard, negte det. Kun Blindhed for, hvad der er det
Væsentlige i en praktisk Livsanskuelse, kan negte, at her
i hvert Tilfælde er et Problem. Og Andet vil jeg
foreløbigt ikke hævde. --- Hvad Skeptikerne angaar, gælder
det at vise, at de forveksle det gamle System med de
Kræfter, der virkede til at frembringe det. Kritiken
rammer den færdige Bygning, ikke de Hænder og den Energi,
der bragte den i Stand. Kun Erfaring kan afgøre, om ikke
disse samme Kræfter endnu ere levende, selv om Former og
Betingelser for deres Anvendelse nu ere væsentligt forandrede.
Skeptikeren bliver selv dogmatisk, naar han negter de nye
Muligheder. Der er Kræfter, man kun faar at vide, at
man har, naar man bruger dem. Og ligesom Dogmatikeren
overser den Forskydning, den Metamorfose, der er
foregaaet især i de sidste Aarhundreder, saaledes overser
Skeptikeren den Forskydning, der stadigt foregaar iblandt os,
og som gør, at ogsaa de kritiske Perioder, naar de komme
igen, faa en anden Karakter end forrige Gang. Der er
en væsentlig og en karakteristisk Forskel mellem Ludvig
Feuerbachs Forhold til Religionen og Voltaires. Saadan
Forskydning og Sammensmeltning af Motiver er Historiens store
Husraad paa den religiøse og den sociale Udviklings
Omraade.

Hvad den nærmere Beskaffenhed af det tredie Rige,
den nye organiske Periode angaar, er det intet Under, at
Afvigelserne i Forestillingerne ere store. Konstruktion af
Fremtidens Tro eller Fremtidens Samfundsorden er umulig.
% --- p. 239 ---
\opage{239}De almindelige Træk, vi kunne angive, --- Foreningen af
koncentreret Kraft og nuanceret Fylde, af bevægelig Kritik
og fast Tro, af individuel Frihed og social Organisation, ---
ere ikke tilstrækkelige til en saadan Konstruktion. Og
desuden ville de stadige Forskydninger hindre, at man paa
Forhaand skulde kunne danne sig en Forestilling om det mere
fremskredne Livs Former.

Det er efter min Opfattelse ogsaa en Misforstaaelse af
Filosofiens Betydning, naar man antager, at dens Opgave
skulde være at give en saadan Konstruktion. Det er
Filosofiens Sag at være en videnskabelig Selverkendelse, at søge
Klarhed over, hvorledes vi ere stillede i aandelig Henseende,
hvor vidt vort Livs Former svare til det Stade, vi ere naaede
til, og hvilke Kræfter der har frembragt disse Former.
Filosofen maa være tilfreds, naar han kan bidrage til
Løsningen af denne Opgave. Han sympatiserer med
Skeptikernes Uro, for saa vidt den virkeligt er en Søgen, og med
Dogmatikernes Konservatisme, for saa vidt den virkeligt
bevarer noget Værdifuldt.

Jeg for mit Vedkommende har i hvert Tilfælde kun
villet konstatere et Problem, ikke give en definitiv Løsning.
Men med Problemets Opstilling er ogsaa givet den Aand, i
hvilken, og den Metode, med hvilken vi maa arbejde. Jeg
skal fremhæve de Punkter, som jeg i denne Henseende anser
for de mest væsentlige.

For det Første ligger der i Anerkendelsen af Problemet
et Princip, man kunde kalde Ækvivalensens Princip. Hvad
der har udfyldt en væsentlig Plads inden for det aandelige
Liv, kan ikke falde bort uden Erstatning. Den blotte
Kritik kan ikke være det sidste Ord, naar Livet ikke skal
stanse. Det er Aandsvidenskabens Opgave at undersøge,
om nye Former for aandelig Energi staa rede til at
indtage deres Plads, hvis Virken aftager. Det gælder at
udfinde, hvad der giver det aandelige Liv Styrke og Værdi
til enhver given Tid. Denne Opgave er --- som allerede
berørt --- langt vanskeligere end den tilsvarende Opgave
paa det materielle Omraade, fordi vi mangle en nøjagtig
Maalestok, baade for Energien og for Værdien. Men det
% --- p. 240 ---
\opage{240}er en uafviselig Opgave, baade til teoretisk Forstaaelse af
det aandelige Liv og til praktisk Videreførelse af det. Vi
forstaa kun en Personlighed, naar vi se, hvad den drager
Næring af, og hvorledes den bruger den Næring, den
optager, ligesom vi kun forstaa, hvad der foregaar i en
Organisme, naar vi forfølge Omsætningsprocesserne i den,
lige fra Optagelsen af Elementer fra Omverdenen til
Funktionernes mangfoldige Anvendelse af de optagne Elementer.
Vi stave endnu i vor Psykologi paa en Lære om
Personligheden, om Betingelserne for, at der danner sig aandelige
Centra, og om Betingelserne for deres Sundhed. Her har
den romantiske Retning altfor ofte ladet sig nøje med at
bestemme Personlighedens Begreb in abstracto og at
opstille det som Genstand for Kontemplation, ligesom Platon
opstillede sine Ideer. Saaledes som jeg forstaar og hylder
den positivistiske Retning i Filosofien, stiler den henimod
at finde de Betingelser og Love, det personlige Liv,
ligesom alt Andet i Verden, er underlagt. Gennem
Eksperiment og Selviagttagelse, ved Benyttelse af fysiologisk og
historisk Materiale vil den finde de Love, der
sammenholder Livet indenfor det enkelte Individs Bevidsthed, saa vel
som de Love, i Følge hvilke Individets Bevidsthedsliv er
knyttet til hele Slægtens Liv. Det er saa langt fra, at en
realistisk Forsken vil føre til at fornegte Personlighedens
Idé, at den netop vil styrke denne Idés Anvendelighed
ved at føre til Indsigt i de bestemte Love, Betingelser og
Grænser for den. Den romantiske Frygt eller Foragt for
den mekaniske Sammenhæng, der nu en Gang er Bærer
og Ramme for Alt, hvad der skal kunne bestaa og gøre
sig gældende i Virkelighedens Verden, er en af de vigtigste
Hindringer for den aandelige Udvikling. Den er selv en
hemmelig Vantro, idet den hviler paa den Forudsætning,
at det forringer Aandslivets Værdi at være en lovbunden
Sammenhæng underlagt. En frejdig Tro paa
Personlighedens Idé vil baade medføre Tillid til, at den vil kunne
hævde sig under nye Former, og Iver for at undersøge den
Tingenes Orden, der kan føre det personlige Liv over fra
% --- p. 241 ---
\opage{241}gamle til nye Former. Idealet for al psykologisk og etisk
Forsken er hermed givet. Mod dette Ideal fæster
Positivisten sit Blik, naar han vil undgaa at forvirres af
Enkeltundersøgelsernes Mangfoldighed. Som ethvert Ideal har det
noget Irrationelt, det vil sige: i matematisk Forstand
Irrationelt, ved sig: ved vore virkelige Bestræbelser kunne vi kun
naa en Tilnærmelse, Bestemmelse af flere og flere Decimaler.
Men det er den samme Begrænsning, som al vor Erkendelse
er underlagt.

Under denne Bestræbelse for at tilvejebringe
Grundlaget for en Personlighedsfilosofi, der kan lære os at skelne
mellem de sande og de falske Positioner og mellem de
sande og falske Negationer, at opdage det »Organiske« og
det »Kritiske« i deres berettigede Former og sondre det fra
Surrogater og Forfalskninger, --- maa vi være klare over
en Ting: at der er saa stor Forskel mellem menneskelige
Individualiteter, at de trænge til højst forskellig Næring,
og at Værdierne ikke kunne blive ens for Alle. Den
virkelige Næringsværdi har kun det, der passer for den bestemte
individuelle Personlighed. Man vil ikke kunne faa Alle ind
under et fælles ydre Udtryk for Sandheden, naar denne
skal være personlig Sandhed for hver Enkelt. Kendetegnet
for den personlige Sandhed er, som en dansk Tænker har
sagt, at den frembringes paa ny af Individet selv: »Ingen
Livsytring har Sandhed, uden deri ligger skabende
Selvvirksomhed«. Der skal ingen Undskyldning gøres for de
individuelle Forskelligheder, som medføre, at Enhver former
Sandheden paa sin Maade. Ligesom Livet viser sin Energi
i den enkelte Personlighed ved den Fylde af Tanker og
Motiver, den formaar at koncentrere om sit ledende
Formaal, saaledes er det et Tegn paa, at Livet i Slægten er
energisk, naar der fremtræder forskelligartede, hver for sig
udprægede Personligheder, som tage Livet og Verden hver
efter sin Naturs Medfør. De »kritiske« Perioder indtræde
i Historien, naar en forudgaaende »organisk« Sammenhæng
er bleven for snever, for tilvant, --- naar Synskredsen trænger
til at udvides og nye Erfaringer og nye Synsmaader skulle
% --- p. 242 ---
\opage{242}vindes. Da kan Slægtens Fremtid kun reddes gennem de
individuelle Variationer. Hver enkelt Variation, hver enkelt
Nuance af det almindelige Mennesketema betyder Mulighed
for nye Erfaringer, nye Udblik og Opdagelser, der kunne
komme det fortsatte Liv til Gode. Der udkom i England
noget over Midten af Aarhundredet to Bøger, af to af
Englands største Forskere i dette Aarhundrede, der hver fra
sit Synspunkt talte de individuelle Variationers Sag. Det
var \emph{Stuart Mill’s} Bog om Friheden og \emph{Darwins} Bog om
Arternes Oprindelse. Stuart Mill talte den individuelle
Friheds Sag, fordi det gjaldt om, at der var saa mange
selvstændige Udgangspunkter for Udviklingen som muligt, og
Darwin paaviste de individuelle Variationer hos Dyr og
Planter som Betingelse for Udviklingen af nye Former.
Naar vi i den nyeste Tid saa ofte se Individualiteten blive
gjort gældende paa en forceret Maade, naar f. Eks. i
Literaturen sære Karakterer og sære Situationer fremstilles med
Forkærlighed, saa er det et Tegn paa en Trang og en
Sans, som er i Færd med at vaagne, og som kan føre
Adskilligt med sig, der ikke er sundt, men som dog kun ved
at faa Lov at komme ind under Livets og Erfaringens egen
Tugt kan finde sine Grænser. Ikke ved Moraliseren udefra,
men ved Tillid til det personlige Livs Evne til at finde sine
egne Regler, til at arbejde sig ud af sine Kriser til nyt
organisk Liv, --- og ved at holde de store Idealer og
Betingelser for aandelig Sundhed fast vil det lykkes at redde de
sande Værdier gennem Kritikens og den tilsyneladende
Opløsnings Skærsild.

Vor Tid frembyder paa det aandelige Omraade et Billede
af broget Forvirring. Der er dem, der holde paa det
Gamle for enhver Pris, --- der er dem, der med let Hjerte
give det til Pris, uden at føle Trang til noget Nyt, ---
der er dem, der sørge over, at de gamle Tider ere forbi,
uden at kunne haabe Noget af de nye, --- og der er
endeligt dem, der have Programmet for det Nye færdigt i alle
Rubriker og utaalmodige fordre, at det Gamle skal give
Plads for de nye Herligheder. Filosofens Opgave er det
gennem stille Forsken at finde Midler til at prøve Kraften
% --- p. 243 ---
\opage{243}og Værdien baade af det, der gaar, og det, der kommer,
og saa maaske tillige gennem den Klarhed, han maatte kunne
bringe, og den Tro til Livets uoverskuelige Muligheder, han
maatte kunne styrke, at give sit Bidrag til Tidens Arbejde.
Selv søger han ligesom Dr. Faust i sin Grublen tilbage
til den Tanke, at Aandens Verden ikke er lukket: hverken
gamle Dogmer eller nye Tvivl ville formaa at stænge den
for os!

\streg{}

\end{document}
